%========================= Classe du Document ==========================
\documentclass[a4paper,fleqn,12pt,openright,twoside]{report}
%============================= Packages ================================

\usepackage[french,english]{babel}
\usepackage[utf8]{inputenc}
\usepackage{ae,lmodern} % Vectorial font
\usepackage[T1]{fontenc}  % Zotero must export bib.tex in Western instead of UTF to avoid problem for \THORN characters

\usepackage[hidelinks, breaklinks]{hyperref} % Cick on citations, allow line breaking

\usepackage{fancyhdr} % For headers with the chapter name
\setlength{\headheight}{14.49998pt}
\fancypagestyle{plain}{%
    \fancyhead{} % get rid of headers
    \renewcommand{\headrulewidth}{0pt} % and the line
}

\fancypagestyle{fancy}{%
    \fancyhead[LO,RE]{\chaptername~\thechapter}
    \fancyhead[LE,RO]{}
    \renewcommand{\headrulewidth}{0.4pt} % and the line
}

\fancypagestyle{intro}{%
    \fancyhead[LO,RE]{Introduction}
    \fancyhead[LE,RO]{}
    \renewcommand{\headrulewidth}{0.4pt} % and the line
}

\fancypagestyle{annex}{%
    \fancyhead[LO,RE]{Annex}
    \fancyhead[LE,RO]{}
    \renewcommand{\headrulewidth}{0.4pt} % and the line
}

\fancypagestyle{conclusion}{%
    \fancyhead[LO,RE]{Conclusion}
    \fancyhead[LE,RO]{}
    \renewcommand{\headrulewidth}{0.4pt} % and the line
}

\usepackage[skip=10pt, indent=20pt]{parskip}  % Paragraph spacing
\usepackage{setspace} % Line spacing
\usepackage[hypcap=false]{caption} % When using hypcap=false, hyperref will link to the caption and not the float environment. This is because using \captionof outside a float env, it generates warnings
\usepackage[labelformat=simple]{subcaption} % Subcaptions for multiple figures
\renewcommand\thesubfigure{(\alph{subfigure})} % when referencing subfigure with subref, yields '(a)' instead of "a"
\usepackage{multicol}

\usepackage{amsmath}
\usepackage{amssymb}
\usepackage{amsfonts}
\usepackage{amsthm}
\usepackage{dsfont} % For characteristic function (\mathds{1}, similar to mathbb 1 but)
\usepackage{textcomp,gensymb} % For the \degree symbol

\usepackage{comment}  % For commenting big chuncks of text
\usepackage{multirow} % For multiple rows in tables

\usepackage[margin=2.5cm]{geometry}
\usepackage{graphicx}
\usepackage{pgfplots}
\pgfplotsset{compat=1.18}
\usepackage{tikz}

\usepackage{xcolor}
\usepackage{xspace}
\usepackage{colortbl} % For 
\usepackage{xr}  % For crossreferencing between chapters

\usepackage[tikz]{mdframed}  % For putting boxes around theorem, examples, definitions etc
\usepackage{epigraph} % Used for epigraph at the begining of the thesis 

\usepackage[acronym,toc]{glossaries}
\usepackage{breakcites} % Citations over multiple lines
%\setlength{\parindent}{0pt}  % Paragraph indentation 

\usepackage[margin,index]{fixme}

\usepackage[noabbrev,nameinlink]{cleveref} % For \cref. noabbrev: No "fig. 1", only "figure 1". nameinlink: "chapter" is clickable in "chapter 3"
\usepackage{menukeys} % For accessing characters from the keyboard, such as Alt or Ctrl

\usepackage{wrapfig} % For wrapping figures around text in Remerciements

%Import the natbib package and sets a bibliography  and citation styles
\bibliographystyle{apalike}

%=============================== Commands ===================================
\definecolor{color_census}{rgb}{0.8, 0.93, 2}
\definecolor{color_mccnn}{rgb}{1, 0.84, 0.7}

% We use colors in equations to help readability.
% Un-comment next lines if you want just black text.
%\let\oldtextcolor\textcolor
%\renewcommand{\textcolor}[2]{\oldtextcolor{black}{#2}}

\newcommand{\RPC}{\mathrm{RPC}}
\newcommand{\M}{\mathfrak{M}}
\newcommand{\X}{\mathcal{X}}
\newcommand{\low}{\underline{P}}
\newcommand{\tdt}{\times\dots\times}
\newcommand{\enum}{, ~\ldots, ~}
\newcommand{\opi}{[\![}
\newcommand{\cli}{]\!]}
\newcommand{\Bel}{\mathrm{Bel}}
\newcommand{\Pl}{\mathrm{Pl}}
\newcommand{\Nec}{\mathrm{Nec}}

\newcommand{\rowcol}{row, ~col}
\newcommand{\AD}{\mathrm{AD}}
\newcommand{\SAD}{\mathrm{SAD}}
\newcommand{\DSM}{\mathrm{DSM}}

\newcommand{\N}{\mathfrak{N}}
\newcommand{\ith}{\textsuperscript{th}~}

\DeclareMathOperator{\CDF}{cdf}
\tikzset{
    declare function={
        cdf_erf(\x,\m,\s)=1/(1 + exp(-0.07056*((\x-\m-3)/\s)^3 - 1.5976*(\x-\m-3)/\s));
    }
}
\DeclareMathOperator*{\argmin}{arg\,min} % thin space, limits underneath in displays
\DeclareMathOperator*{\argmax}{arg\,max} % thin space, limits underneath in displays
\DeclareMathOperator*{\mean}{mean}
\DeclareMathOperator*{\median}{median}

%%%%%%%%%%%%%%%%% Defining the properties of the frames for theorems etc %%%%%%%%%%%%%%%%%%%%%%%%%%%%%%%%%%%%%
\mdfsetup{skipabove=\topskip, skipbelow=\topskip}  % Define Skips between the frames and text above and bellow

%% Code taken almost directly from https://ctan.mines-albi.fr/macros/latex/contrib/mdframed/mdframed-example-default.pdf
%% Theorem environment
\newcounter{theor}
\renewcommand{\thetheor}{\arabic{theor}}
\newenvironment{theorem}[1][]{%
  \refstepcounter{theor}%
  \ifstrempty{#1}%
   {\mdfsetup{
      firstextra={\useasboundingbox (P) rectangle (P); \node[rectangle,fill=blue!30!white, anchor=west,xshift=5mm] at (P-|O) { \strut \textbf{Theorem~\thetheor}};},
      singleextra={\useasboundingbox (P) rectangle (P); \node[rectangle,fill=blue!30!white, anchor=west,xshift=5mm] at (P-|O) { \strut \textbf{Theorem~\thetheor}};}}
  }%
  {\mdfsetup{%
     firstextra={\useasboundingbox (P) rectangle (P); \node[rectangle,fill=blue!30!white,anchor=west,xshift=5mm] at (P-|O) { \strut \textbf{Theorem~\thetheor:~#1}};},
     singleextra={\useasboundingbox (P) rectangle (P); \node[rectangle,fill=blue!30!white,anchor=west,xshift=5mm] at (P-|O) { \strut \textbf{Theorem~\thetheor:~#1}};}
     }
   }%
   \mdfsetup{innertopmargin=2\topskip,linecolor=blue!70!black,%
             linewidth=2pt,topline=true,
             backgroundcolor=blue!10!white,%
             roundcorner=10pt,
             skipabove=2\topskip,}
   \begin{mdframed}[]\relax%
   }{\end{mdframed}}

%% Definition environment
\newcounter{definition}
\renewcommand{\thedefinition}{\arabic{definition}}
\newenvironment{definition}[1][]{%
 \refstepcounter{definition}%
  \ifstrempty{#1}%
  {\mdfsetup{
      firstextra={\useasboundingbox (P) rectangle (P); \node[rectangle,fill=blue!20, anchor=west,xshift=5mm] at (P-|O) { \strut \textbf{Definition~\thedefinition}};},
      singleextra={\useasboundingbox (P) rectangle (P); \node[rectangle,fill=blue!20, anchor=west,xshift=5mm] at (P-|O) { \strut \textbf{Definition~\thedefinition}};}}
  }%
  {\mdfsetup{%
     firstextra={\useasboundingbox (P) rectangle (P); \node[rectangle,fill=blue!20,anchor=west,xshift=5mm] at (P-|O) { \strut \textbf{Definition~\thedefinition:~#1}};},
     singleextra={\useasboundingbox (P) rectangle (P); \node[rectangle,fill=blue!20,anchor=west,xshift=5mm] at (P-|O) { \strut \textbf{Definition~\thedefinition:~#1}};}
     }
   }%
   \mdfsetup{innertopmargin=2\topskip,linecolor=blue!20,%
             linewidth=2pt,topline=true,%
             roundcorner=5pt,
             skipabove=2\topskip,}
   \begin{mdframed}[]\relax%
   }{\end{mdframed}}

%% Proposition environment
\newcounter{proposition}
\renewcommand{\theproposition}{\arabic{proposition}}
\newenvironment{proposition}[1][]{%
 \refstepcounter{proposition}%
 \mdfsetup{outerlinewidth=1,%
           roundcorner=10pt,%
           backgroundcolor=gray!5!white,%
           outerlinecolor=black,%
           innertopmargin=\topskip,%
           splittopskip=\topskip,%
           skipabove=2\topskip,
           }
 \ifstrempty{#1}%
 {\begin{mdframed}[]\textbf{Proposition~\theproposition:~}\relax}%
 {\begin{mdframed}[]\textbf{Proposition~\theproposition:}~\textit{#1}\par\relax}%
 }{\end{mdframed}}


%% Example Environment
%\theoremstyle{definition}  % Boldface title, Roman body
%\newmdtheoremenv[roundcorner=5pt, font=\bfseries]{example}{Example}[section]
% OR
\newcounter{example}
\renewcommand{\theexample}{\arabic{example}}
\newenvironment{example}[1][]{%
 \refstepcounter{example}%
 \mdfsetup{roundcorner=5pt,}
 \ifstrempty{#1}%
 {\begin{mdframed}[]\textbf{Example~\theexample:~}\relax}%
 {\begin{mdframed}[]\textbf{Example~\theexample:}~\textit{#1}\par\relax}%
 }{\end{mdframed}}

% Remark Environment
\newenvironment{remark}{%
   \mdfsetup{hidealllines=true,%
             leftline=true,%
             linecolor=gray,%
             middlelinewidth=.2em,%
            }
   \begin{mdframed}[]\textbf{Remark:}\relax%
   }
   {\end{mdframed}}
   
   
% Redefines the proof environment
\renewenvironment{proof}{%
   \mdfsetup{hidealllines=true,%
             leftline=true,%
             linecolor=gray!50,%
             middlelinewidth=.2em,%
            }
   \begin{mdframed}[]\textit{\textbf{Proof:}}\relax%
   }
   {\qed\end{mdframed}}
   
% Conclusion Environment
\usetikzlibrary{calc}
\tikzset{
    excursus line/.style={%
    line width=2pt,
    draw=gray!40,
    rounded corners=2ex,
    line cap=round,
    },
    excursus head/.style={
    fill=white,
    font=\bfseries,
    anchor=base west,
    },
}
\mdfdefinestyle{conclusion}{%
    singleextra={%
        \path let \p1=(P), \p2=(O) in (\x2,\y1) coordinate (Q);
        \path let \p1=(P), \p2=(O) in (\x2,\y1/2-\y2/2) coordinate (M);
        \path [excursus line] ($(O)+(4,0)$) -| (M) |- ($(P)+(-3,0)$);
        \node [excursus head] at ($(Q)+(1,-0.1)$) {Conclusion};
        },
    firstextra={%
        \path let \p1=(P), \p2=(O) in (\x2,\y1) coordinate (Q); 
        \path [excursus line] (O) |- ($(P)+(-3,0)$);
        \node [excursus head] at ($(Q)+(1,-0.1)$) {Conclusion};
    },
    secondextra={%
        \path let \p1=(P), \p2=(O) in (\x2,\y1) coordinate (Q);
        \path [excursus line] ($(O)+(4,0)$) -| (Q);
    },
    middleextra={%
        \path let \p1=(P), \p2=(O) in (\x2,\y1) coordinate (Q);
        \path [excursus line] (O) -- (Q);
    },
    middlelinewidth=2.5em,middlelinecolor=white,
    hidealllines=true,topline=true,
    innertopmargin=0.5ex,
    innerbottommargin=2.5ex,
    innerrightmargin=2pt,
    innerleftmargin=2ex,
    skipabove=4\baselineskip,
    skipbelow=2\baselineskip,
    nobreak=true,
}

\newenvironment{conclusion}{%
   \mdfsetup{style=conclusion,%
            }
   \begin{mdframed}[]\relax%
   }
{\end{mdframed}}
%%%%%%%%%%%%%%%%%%%%%%%%%%%%%%%%%%%%%%%%%%%%%%%%%%%%%%%%%%%%%%%%%%%%%%%%%%%%%%%%%%%%%%%%%%%%%%%%%%%%%%%%%%%%%%%%%%%%%
   
   
% From CVPR cvpr.sty
\makeatletter
\DeclareRobustCommand\onedot{\futurelet\@let@token\@onedot}
\def\@onedot{\ifx\@let@token.\else.\null\fi\xspace}

\def\eg{\emph{e.g}\onedot} \def\Eg{\emph{E.g}\onedot}
\def\ie{\emph{i.e}\onedot} \def\Ie{\emph{I.e}\onedot}
\def\cf{\emph{cf}\onedot} \def\Cf{\emph{Cf}\onedot}
\def\etc{\emph{etc}\onedot} \def\vs{\emph{vs}\onedot}
\def\st{\emph{s.t}\onedot}
\def\wrt{w.r.t\onedot} \def\dof{d.o.f\onedot}
\def\iid{i.i.d\onedot} \def\wolog{w.l.o.g\onedot}
\def\etal{\emph{et al}\onedot}
\def\vs{\emph{vs}\onedot}

\makeatother

%========================= Glossary ====================================

\pdfstringdefDisableCommands{%
  \def\glspl#1{<#1>}} % To silence hyperref Warning Token not allowed in a PDF string (Unicode)
%=======================================================================
%=======================================================================

% \renewcommand\thechapter{\Roman{chapter}} % Chapters are Roman numbers
\makeglossaries
\makeglossaries
\newglossaryentry{EO}
{%
    name={Earth Observation},
    description={Earth Observation involves employing remote sensing techniques for terrestrial, marine, and atmospheric monitoring}
}

\newglossaryentry{DISPARITY}
{%
    name={Disparity},
    description={The displacement of an object between two stereo images when expressed in pixels}
}

\newglossaryentry{DSM}
{%
    name=Digital Surface Model,
    description={Representation of the elevation of an area using a raster, including the vegetation, buildings \etc. Also called Digital Elevation Model. Sometimes referred to as a 2.5D model, as it is a 2D grid where each cell contains data about the elevation}
}

\newglossaryentry{DTM}
{%
    name={Digital Terrain Model},
    description={Representation of the terrain of an area using a raster, excluding the vegetation and buildings}
}

\newglossaryentry{LiDAR}
{%
    name={Light Detection And Ranging},
    description={Sometimes referred to as Laser Imaging Detection And Ranging. A LiDAR is an active sensor allowing to determine the distance between the sensor and its surroundings by measure the time between the emission of a laser beam and the detection of its reflection on the environment}
}

\newglossaryentry{Panchromatic}
{%
    name={Panchromatic image},
    description={A panchromatic image is an image representing the light in the visible spectrum, resulting in a black and white image. Satellites equipped with panchromatic sensors usually produce images of higher resolution than classical RGB sensors, although a high resolution RGB image can be obtained from a panchromatic image using a method called pansharpening}
}
\newacronym{b/h}{B/H}{Base-to-Height ratio}
\newacronym{cnes}{CNES}{Centre National d'Etudes Spatiales}
\newacronym{co3d}{CO3D}{Constellation Optique 3D}
\newacronym{dem}{DEM}{Digital Elevation Models}
\newacronym{dsm}{DSM}{Digital Surface Model}
\newacronym{dtm}{DTM}{Digital Terrain Model}
\newacronym{gis}{GIS}{Geographic Information Systems}
\newacronym{nir}{NIR}{Near Infra Red}
\newacronym{rgb}{RGB}{Red Green Blue}
\newacronym{vhr}{VHR}{Very High Resolution}
\newacronym{eo}{EO}{Earth Observation}
\newacronym{lidar}{LiDAR}{Light Detection And Ranging}
\newacronym{rpc}{RPC}{Rational Polynomial Models}

\begin{document}
    \pagestyle{empty}
    
    \pagenumbering{gobble} % No numbering for now

    \definecolor{color}{HTML}{F5F5F5}
    \definecolor{ForestGreen}{HTML}{228B22}
    %======================= Page de Garde =================================
    
    \setkeys{Gin}{draft=false}
    \begin{titlepage}
    \begin{center}
        \Huge
        \textbf{Uncertainty characterisation in stereophotogrammetry using satellite images}
            
        \vspace{0.5cm}
        \LARGE
        Caractérisation d'incertitude pour la stéréophotogrammétrie à partir d'images satellites

        \vspace{1.cm}
        \large
        Defended by\\
        \textbf{Roman MALINOWSKI}\\
        \vspace{0.8cm}
        Members of the jury\\
        
        \vspace{0.2cm}
        President: \textbf{Benjamin QUOST}\\
        Reviewers: \textbf{Isabelle BLOCH}, \textbf{Marc PIERROT DESEILLIGNY}\\
        External members: \textbf{Luc JAULIN}, \textbf{Olivier STRAUSS}\\
        Supervisors: \textbf{Sébastien DESTERCKE}, \textbf{Emmanuel DUBOIS}, \textbf{Loïc DUMAS}, \textbf{Emmanuelle SARRAZIN}\\
    
        \vspace{1cm}
        Speciality: Sciences et Technologie de l'Information et de la Communication\\~\\
        Centre National D'Etudes Spatiales - CS Group \\
        Heudiasyc, Université de Technologie de Compiègne\\
        ~\\
        19/12/2024
        \
        \begin{figure}[hb]
            \begin{subfigure}[b]{0.3\linewidth}
                \centering
                \includegraphics[height=2.5cm]{Images/General/Logo_CNES.png}
            \end{subfigure}\hfill
            \begin{subfigure}[b]{0.3\linewidth}
                \centering
                \includegraphics[height=1.5cm]{Images/General/Logo_CS.png}
            \end{subfigure}
        \end{figure}
        \begin{figure}
            \begin{subfigure}[b]{0.3\linewidth}
                \centering
                \includegraphics[height=1.3cm]{Images/General/Logo_UTC.jpg}
            \end{subfigure}\hfill
            \begin{subfigure}[b]{0.3\linewidth}
                \centering
                \includegraphics[height=1.3cm]{Images/General/Logo Région HDF - partenaire.jpg}
            \end{subfigure}
        \end{figure}
    \end{center}
\end{titlepage}

    \clearpage
    \setkeys{Gin}{draft=false}
    \chapter*{Remerciements}
Je tiens tout d'abord à remercier l'Université de technologie de Compiègne, le \acrshort{cnes} et CS Group, qui ont su s'entendre pour me permettre de réaliser cette thèse dans de bonnes conditions. J'ai eu la chance de pouvoir travailler entre Compiègne et Toulouse, et d'y rencontrer de nombreuses personnes qui m'ont épaulé durant ces trois années particulièrement enrichissantes. Je tiens ensuite à remercier les membres du jury, notamment M. Luc Jaulin pour avoir suivi mes travaux de thèse depuis nos échanges lors de la journée Imagin. Je remercie également Mme Isabelle Bloch et M. Marc Pierrot Deseilligny pour avoir accepté de rapporter ce travail et pour le temps qu’ils y ont consacré ainsi que pour leurs remarques pertinentes. Je voudrais aussi remercier chaleureusement M. Olivier Strauss, qui a manifesté son intérêt pour mes travaux depuis les premiers instants, m'a permis de découvrir son laboratoire du LIRMM, et a piqué ma curiosité avec des problèmes de traitement d'image auxquels je réfléchis encore aujourd'hui. Un grand merci à M. Benjamin Quost, qui a accepté de présider ce jury, mais aussi de partager son bureau pour que s'y déroulent mes réunions de thèse lorsque j'étais un peu perdu.

Je voudrais remercier Etienne Berthier et Simon Gascoin, pour tout l'intérêt qu'ils ont porté à mes travaux. Mais aussi pour m'avoir généreusement partagé leurs données, et mis en contact avec Liss Marie Andreassen et Brian Menounos qui ont fait de même. C'était pour moi bien plus que de simples données. J'aimerais également remercier Yoann Steux, pour tout le temps passé à résoudre mes problèmes sur \acrshort{cars}.

Bien entendu, merci à Simon, Tatha, Tom, pour avoir égayé mes soirées compiègnoises. Cassandra, Gabriela, Lou, Marion, Alice, Martin, et surtout Camille, pour tout ce qu'on a vécu au CNES, mais aussi en dehors.

Je ne dis pas merci à :
\begin{itemize}
    \item La personne qui m’a dit qu’installer \acrshort{cars} était facile
    \item Mon vélo, et surtout mes freins. Trop de temps et d’efforts passés à le réparer, et encore plus à me réparer moi-même
\end{itemize}

Un grand merci à ma soeur, mes parents et grands-parents, qui se sont intéressés à mes travaux et m’ont toujours souhaité le meilleur. Merci également à mon autre famille, chinoise, pour les déménagements et pour avoir pris soin de mes papilles gustatives. À toi, Virginie, pour m’avoir soutenu, écouté, t’être intéressée à mes travaux quand bien même ce que je racontais n’avait que peu d’intérêt. M’avoir changé les idées quand je voulais vivre autre chose que cette thèse, et avoir accepté que je n’aie parfois pas d’autre vie que celle-ci. J'espère pouvoir te rendre, un jour, tout ce que tu m'as donné.

Un grand merci à Manu, qui a fait la jonction entre ces deux mondes que sont la galaxie CNES et la communauté des probas imprécises. Sans toi, cette thèse n'aurait sûrement pas eu lieu, ou du moins pas sous cette forme. Un grand merci à Loïc, avec le regret de n'avoir pas passé encore plus de temps à travailler à tes côtés. Je t'admirerai toujours pour tes conseils et remarques qui touchent toujours juste. Mais également pour tes commentaires qui m'ont beaucoup fait rire lors de la rédaction de ce manuscrit. Sébastien, bien entendu, je te serai à jamais reconnaissant pour m'avoir guidé, conseillé et épaulé lorsque j'étais perdu dans des concepts théoriques qui me dépassaient. Pour avoir dit ce qui n'allait pas, mais aussi toujours ce qui allait bien. Merci pour m'avoir toujours accordé ton temps, même quand tu ne l'avais pas.

\begin{wrapfigure}{L}{0.2\textwidth}
  \begin{center}
    \includegraphics[width=0.2\textwidth]{tmp/balrog.png}
  \end{center}
\end{wrapfigure}
Enfin, un grand merci à Manue, qui m’a convaincu de me lancer dans cette aventure, et qui est restée jusqu’au bout — et même parfois très tard — pour enrichir mon travail de ses idées et de ses conseils. Tu m’as dit au début de ma thèse, que je la commencerai en étant sous ta supervision, et que je la terminerai en étant ton collègue. Plus qu’une collègue, j’ai terminé cette thèse avec un modèle, un guide, une amie. Tu es la flamme d’Udûn qui éclairait les recoins les plus sombres du cluster ou de l’écriture de manuscrit. Chanceux seront les prochains doctorants qui t’auront comme maître Jedi.
\begin{center}
    \includegraphics[width=0.2\textwidth]{tmp/gandalf.png}
\end{center}
\clearpage
    %==================== Glossaire, TOC etc ===============================
    \tableofcontents
    \clearpage
    
    \printglossary[type=\acronymtype,nonumberlist]
    
    %========================== Les choses sérieuses commencent ===========
    %\doublespacing
    \onehalfspace
    \setkeys{Gin}{draft=false}
    \chapter*{Résumé}\addcontentsline{toc}{chapter}{Résumé}
\newpage

\chapter*{Abstract}\addcontentsline{toc}{chapter}{Abstract}
\newpage

\chapter*{Foreword}\addcontentsline{toc}{chapter}{Foreword}
Before delving into the subject of this manuscript, we would like to give some advise on how to efficiently navigate through it. When writing this thesis, we made extensive use of the \textit{hyperref} package, so that reading it on a PDF viewer was made easier. You can thus click on citations, figure numbers, equation numbers, chapters and sections number, acronyms \etc to directly jump to the concerned part. When following a reference to a citation, a previous chapter, equations or figures located in a different part of the manuscript, it can be a laborious process to go back to the section you were reading. Depending on the OS of you computer and the app used to read the PDF document, there usually exist shortcuts to jump back to the previous view. This allows to quickly switch back and forth between chapters and sections.

For instance, imagine that you are in \Cref{chap:epistemic_uncertainty} and we make a reference to an equation form \Cref{chap:representation_of_uncertainty}. If you do not recall the equation, and quickly want to see what it is about, just click on the hyperlink to directly go to the relevant equation from \Cref{chap:representation_of_uncertainty}. Then use your system's shortcut to go back to where you were in \Cref{chap:epistemic_uncertainty}. 
\begin{itemize}
    \item Using Acrobat Reader: the short cut \keys{\Altwin + \arrowkeyleft} (left arrow key) on Windows or Linux brings you to the previous view after clicking on a hyperlink. Afterwards, you can alternate views with \keys{\Altwin + \arrowkeyright} and \keys{\Altwin + \arrowkeyleft}. On MacOS, the \keys{\Altwin} key is replaced by the \keys{\cmd} key.
    \item Using Preview on MacOS, you can add the \menu{Page History} button to the toolbar, by right-clicking on the toolbar and selecting \menu{Customize Toolbar}
    \item Using Okular on Linux, \keys{\Altwin + \shift + \arrowkeyleft} (left arrow key) brings you to the previous view after clicking on a hyperlink. Afterwards, you can alternate views with \keys{\Altwin + \shift + \arrowkeyright} and \keys{\Altwin + \shift + \arrowkeyleft}
\end{itemize}
Hopefully, this makes the reading of this thesis a more pleasant experience.

\newpage
\renewcommand{\epigraphsize}{\normalsize}
\setlength{\epigraphrule}{1pt}
\setlength{\epigraphwidth}{0.48\linewidth}
\vspace*{\fill}
\epigraph{You never talk of likelihoods on Arrakis.\\You speak only of possibilities.}{Frank Herbert, \textit{Dune}}
\vspace*{\fill}

\newpage

    \pagenumbering{arabic}  % Switch numerotation of pages to Arabic
    \pagestyle{intro}
    \setkeys{Gin}{draft=false}
    \addcontentsline{toc}{chapter}{Introduction}
\chapter*{Introduction}
\todoroman{Mettre le logo de la region Picardie sur la thèse}
\todoroman{Mettre un reading guide pour dire quel ``niveau d'expertise'' il faut en fonction des chapitres }
\pagebreak
    
    \pagestyle{fancy}
    \setkeys{Gin}{draft=false}
    \chapter{Principles of Stereophotogrammetry}\label{chap:stereophotogrammetry}

This section will present important concepts regarding photogrammetry that will be relevant throughout this thesis. \Cref{sec:dsm} will present the main 3D product we will consider, \ie \acrshort{dsm}, while \Cref{sec:co3d} will introduce the different sensors and satellites employed. \Cref{sec:classical_stero_pipeline} presents in details the different processing steps for creating \acrshort{dsm}s from stereo images. Finally, \Cref{sec:previous_work_stereo_uncertainty} addresses the modeling of uncertainty in photogrammetry.

\section{Digital Surface Models}\label{sec:dsm}
Topographical information are crucial for many applications in \acrfull{eo}, or more generally when manipulating georeferenced data. A popular and simple model containing such topographical data named \acrlong{dem}, and represents elevation data using a regular grid as in \Cref{fig:VDG_dsm}. In the literature, a distinction is usually made between two types of \acrshort{dem}: \acrlong{dsm} and \acrlong{dtm}. \acrshort{dtm}s only represent the ground surface, without man made structures (buildings, houses) or rapidly evolving volumes such as vegetation or seasonal snow. Those those finer details are instead included in \acrshort{dsm}s. \Cref{fig:DTM_DSM} illustrates the difference between the two.

\begin{figure}
    \centering
    \includegraphics[width=0.8\linewidth]{Images/Chap_1/DTM_DSM.png}
    \caption{Difference between \acrlong{dtm} and \acrlong{dsm}}
    \label{fig:DTM_DSM}
\end{figure}

\acrshort{dtm}s are used for hydrology applications, such as for estimating water reservoirs \cite{yamazaki_merit_2019}, river flow \cite{miguez-macho_incorporating_2007}, global snow resources \cite{gascoin_theia_2019} or modeling potential flood \cite{yamazaki_regional_2014}. On the other hand, the high level of details of \acrshort{dsm}s are used to monitor changes in vegetation \cite{sadeghi_canopy_2016}, melting rates of glaciers \cite{berthier_glacier_2014, rieg_pleiades_2018}, changes in volcanoes \cite{ganci_data_2022} or to precisely map snow depth on mountainous terrain \cite{marti_mapping_2016}. In urban environment, \acrshort{dsm} allow to predict the potential damage caused by floods \cite{jenkins_physics-based_2023} or to do urban planning \cite{velazco_3d_2012}. For instance, possessing \acrshort{dsm}s on the same city but for different acquisition dates allow to measure the growth rate of the city \cite{warth_dsm-based_2019}, or damages caused by an earthquake \cite{erdogan_detection_2019} in-between acquisitions. \acrshort{dsm} are also crucial for ortho-rectifying image, \ie geometrically correcting the effects of distortions between the sensor and the terrain \cite{toutin_ortho-rectification_2012}. When capturing an optical image, the satellite is not necessarily right above its target, as illustrated in \Cref{fig:ortho_nadir}. However, if the image must be used as a background for a map, it is usually required that the image must be captured vertically, referred to as nadir, in order to be free of any perspective effect. \acrshort{dsm}s allow to reproject images as if they were taken at nadir. This rectification is used for many satellite products \cite{hagolle_maja_2017}.  
\begin{figure}
    \centering
    \includegraphics[width=0.7\linewidth]{Images/Chap_1/Ortho_Nadir.png}
    \caption{Incidence angle of a satellite, and nadir direction}
    \label{fig:ortho_nadir}
\end{figure}

Many applications depends on the resolution of \acrshort{dsm}s: evaluating water reservoirs in a mountainous regions does not necessarily require the same accuracy as 3D urban modelling for instance. The precision of a \acrshort{dsm} is defined by its planimetric resolution, \ie the size of a cell, and its altimetric resolution, \ie the smallest elevation variation it can detect. The altimetric resolution is itself characterized two notions: the absolute and relative accuracy. Absolute accuracy is the precision of a \acrshort{dsm} at a pixel level, \ie how close the predicted elevation is to the true elevation at the geographic coordinates of a given cell. Relative accuracy characterizes the quality of the 3D reconstruction of a scene. For instance, a \acrshort{dsm} precisely reproducing the shape of a house and its chimney will have a high relative accuracy. However, if the house is not properly georeferenced then the \acrshort{dsm} will have a poor absolute accuracy. Relative accuracy is therefore always greater than the absolute accuracy. In this thesis, we will mainly consider relative accuracy when evaluating the quality of \acrshort{dsm}s, as our sources possess different georeferencing systems which naturally induces small planimetric errors. 

There are multiple ways of creating a \acrshort{dsm} from airborne sensors such as planes, drones or satellites. The first way is to use \acrshort{radar} interferometry, as done by Sentinel-1 satellites \cite{geudtner_sentinel-1_2014} for instance. This method presents the advantage of being able to acquire data by day and night, even in the presence of clouds. Radar sensors, such as those present in the Shuttle Radar Topography Mission (SRTM) \cite{farr_shuttle_2007} or TanDEM-X \cite{krieger_tandem-x_2007}, possess a large swath of around $30$km, which was used to produce \acrshort{dsm}s covering the majority (or totality for TanDEM-X) of emerged land. The planimetric accuracy are in the range of a dozen meters (10m for TanDEM-X or 30m for the SRTM), and relative altimetric resolution is typically of a few meters (2m for TanDEM-X and 8m for the SRTM). This is not sufficient to distinguish objects such as buildings or trees, but rather global terrain topography. \Cref{fig:etoile_radar} presents a TanDEM-X \acrshort{dsm}. SRTM and TanDEM-X \acrshort{dsm} are not meant to be produced on a regular basis, and thus do not allow for temporal analysis as such. Another method for constructing \acrshort{dsm}s is to use stereophotogrammetry \cite{tao_comprehensive_2001}. Photogrammetry is the science of recovering 3D information from optical images. For this method, images of a scene are acquired from different points of view, and depth information is recovered from the parallax effect between images, already presented in \Cref{fig:parallax}. Photogrammetry is not restricted to the production of \acrshort{dsm}s, as it has many usages in robotics, for instance on Martian rovers \cite{goldberg_stereo_2002} or autonomous driving \cite{geiger_vision_2013}. Current optical satellites have a sub-meter resolution and a large swath (20km for Pléiades for instance \cite{coeurdevey_pleiades_2012}), allowing to massively produce high-resolution \acrshort{dsm} covering the globe for a relatively low cost using photogrammetry. The altimetric resolution is typically around one meter, although it depends of the different acquisition angles of the satellites. With this resolution, buildings, trees, and even cars can be detected on the \acrshort{dsm}. A final method for producing \acrshort{dsm} is to use \acrshort{lidar} (laser sensors) \cite{khosravipour_generating_2016}. Space-borne \acrshort{lidar}s are mostly used for atmospheric or discrete measurements \cite{fouladinejad_history_2019}. For instance, NASA’s ICESat-2 \cite{jasinski_atlasicesat-2_2020} measures elevation of seas and glaciers using 6 lasers taking measurements along track, which is not adapted for reconstructing high-accuracy DSM. On the other hand, planes or drones equipped with \acrshort{lidar} sensors usually have very good accuracy, allowing to model tower cranes, lamp post or small details on buildings. In particular, \acrshort{ign} uses \acrshort{lidar} to create point clouds with a planimetric accuracy of 25 centimeters and an altimetric resolution of 5 centimeters, and then produce \acrshort{dsm}s out of it  \cite{monnet_lidarhd_2023, ign_lidar_2024}. This is part of a 5 years campaign named \acrshort{lidar} HD and started in 2021, where planes equipped with \acrshort{lidar} sensors will cover the whole French territory (except French Guiana). Those types of campaign are long and costly, and can hardly be reproduced at a global scale. Therefore, photogrammetry \acrshort{dsm}s are currently the best solution to produce \acrshort{dsm}s covering the globe, with a sub-meter resolution for relatively low costs. \Cref{fig:etoile} presents examples of \acrshort{dsm}s produced by \acrshort{radar}, photogrammetry and \acrshort{lidar}.

\begin{figure}
    \begin{subfigure}[t]{0.31\linewidth}
        \flushleft
        \includegraphics[width=\linewidth]{Images/Chap_1/etoile_low_res_dsm.png}
        \caption{\acrshort{radar} \acrshort{dsm}}
        \label{fig:etoile_radar}
    \end{subfigure}\hfill
    \begin{subfigure}[t]{0.31\linewidth}
        \centering
        \includegraphics[width=\linewidth]{Images/Chap_1/etoile_high_res_dsm.png}
        \caption{\acrshort{lidar} \acrshort{dsm}}
        \label{fig:etoile_lidar}
    \end{subfigure}\hfill
    \begin{subfigure}[t]{0.365\linewidth}
        \flushright
        \includegraphics[width=\linewidth]{Images/Chap_1/etoile_cars_dsm.png}
        \caption{Photogrammetry \acrshort{dsm}}
        \label{fig:etoile_photogrammetry}
    \end{subfigure}
    \caption{\acrshort{dsm}s over Arc de Triomphe, Paris, obtained with different sensors.}
    \label{fig:etoile}
\end{figure}

In this thesis, we use \acrshort{lidar} \acrshort{dsm} as our reference \acrshort{dsm}s due to their high quality. They will be used to evaluated the errors or uncertainties of \acrshort{dsm}s produced using stereophotogrammetry. The next section details which stereo images we will be using to produce those \acrshort{dsm}s. 

\section{Satellite Photogrammetry}\label{sec:co3d}
This section presents the different sensors, satellites and geolocation models that will be used to acquire and process stereo images.

\subsection{Different sensors}
Different types of sensors can be used to acquire optical satellite images. We detail in this section a (non-exhaustive) list of sensors of interest \cite{cnes_imagerie_2008}.

\subsubsection{Push-broom sensor}
For each observed wavelength, push-broom sensor are only composed of a single cell row, acquiring simultaneously radiometric information alongside a line perpendicular to the direction of the satellite. We use the satellite movement along track to acquire the different rows of the image, as seen in \Cref{fig:ccd_pushbroom_sensors}. As only one line of cells is needed, push-broom sensors are simple systems which can capture images continuously, while guarantying good geometrical quality along the rows of the images. A variation of those sensors are TDI sensors (Time Delay Integration). Those sensors function as a push-broom except that each row has the ability to transfer its photon charges to the next row. This allows to capture signals over a longer period of time,  thus reducing the signal-to-noise ratio. Harder to produce, TDI sensors also require a precise control of the satellite so that observed objects stay within a column of the TDI sensor. 

\subsubsection{CCD sensors}
\acrfull{ccd} are classical sensors, used in current digital cameras for instance. They consist of a grid of cells sensitive to a given wavelength configured in a checkerboard pattern. A specific type of \acrshort{ccd} sensors are Bayer matrices, where the sensor only possesses a single type of photo-sensitive cell and the different colors are captured by applying a color filter on the sensor. \acrshort{ccd} sensor possess multiple advantages, such as good geometrical quality as all pixels are acquired simultaneously, or the possibility to perform many acquisitions with various angles possible. \acrshort{ccd} sensors are relatively recent in spatial imagery as they are technologically difficult to built. Augmenting the number of pixels complicates the shutter function and occupies more space compared to push-broom sensors. It requires more radiometric calibration as each pixel results from a different photo-sensitive cell. Acquire long segments of an image is not natural as it is the case for push-broom sensors.

\begin{figure}
    \begin{subfigure}[t]{0.5\linewidth}
        \centering
        \includegraphics[height=6cm]{Images/Chap_1/pushbrrom.png}
        \caption{Push-broom sensor}
        \label{fig:pushbroom}
    \end{subfigure}
    \begin{subfigure}[t]{0.5\linewidth}
        \centering
        \includegraphics[height=6cm]{Images/Chap_1/CCD.png}
        \caption{\acrshort{ccd} sensor}
        \label{fig:ccd}
    \end{subfigure}
    \caption{Schematic representation \acrshort{ccd} and push-broom sensors}
    \label{fig:ccd_pushbroom_sensors}
\end{figure}

Those different sensors are used in most \acrfull{vhr} optical satellites as we will see in the next section.

\subsection{VHR satellites}
In this section, we present different \acrshort{vhr} optical satellites which can be used to acquire stereo images. 

Across the years, different constellations of \acrshort{vhr} satellites have been launched, either for civilian (SPOT 1-7), defense (CSO) or commercial use (Ikonos, QuickBird, Worldview 1-3). Most of them are agile satellites, using push-broom sensors that can now reach resolutions near 30cm \cite{cnes_imagerie_2008, coffer_vertical_2022}. One constellation of interest that we will consider in this thesis is the Pléiades constellation. Developed by Airbus, this constellation is composed of two identical satellites, 1A and 1B. The satellites were launched in 2011 and 2012 in an helio-synchronous orbit at $690$km, for both civilian and defense usages. They provide panchromatic images at a resolution of $70$cm (resampled at $50$cm), and \acrshort{rgb}-\acrshort{nir} images at a resolution of $2$m, with a $20$km swath (\url{https://dinamis.data-terra.org/pleiades/}) using a push-broom TDI sensor. Their high agility and revisit rate allow them to capture stereo and tri-stereo images for any location on the globe, ideal to produce \acrshort{dsm} with high accuracy. \Cref{fig:Pleiade_over_Paris} provide an example of a Pléiades stereo acquisition. A ``video'' mode is even available were dozens of image of the same scene can be acquired in the span of a few minutes \cite{lebegue_pleiades-hr_2015}. However, stereo acquisitions is not the only objective of this mission, even though the demand for those products is increasing \cite{berthier_glacier_2014, poli_radiometric_2015, rieg_pleiades_2018, loghin_potential_2020}. The acquisition of stereo images is thus provided on command, which can conflict with other usages of the satellite, and can become costly when trying to cover large areas. Pléiades is also a satellite for both defense and civilian usages, military acquisitions having the priority over civilian one. Pléiades acquisitions can also be requested using the ``Disaster Chart'' (\url{https://disasterscharter.org/fr/web/guest/home}) in order to evaluate damages caused by floods, landslide or earth quakes for instance, and to better plan and provide emergency relief to victims of those disasters.

\begin{figure}
    \centering
    \begin{subfigure}[t]{0.5\linewidth}
        \centering
        \includegraphics[height=6cm]{Images/Chap_1/Paris_Left_stereo.png}
        \caption{Left image}
        \label{fig:Pleiade_over_Paris_a}
    \end{subfigure}\hfill
    \begin{subfigure}[t]{0.5\linewidth}
        \centering
        \includegraphics[height=6cm]{Images/Chap_1/Paris_Right_stereo.png}
        \caption{Right image}
        \label{fig:Pleiade_over_Paris_b}
    \end{subfigure}
    \caption{Pansharpened Pléiades stereo images over Paris at $0.5$m of resolution. The change of point of view can be easily observed by looking at the Eiffel Tower. Pléiades \copyright \acrshort{cnes} 2023, Distribution AIRBUS DS}
    \label{fig:Pleiade_over_Paris}
\end{figure}

In order to produce a worldwide \acrshort{dsm} with $1$m resolution by 2025, the \acrfull{cnes} is launching the \acrfull{co3d} mission \cite{melet_co3d_2020}. Composed of two pairs of low-cost satellites (\Cref{fig:co3d_mission}) equipped with \acrshort{vhr} optical sensors, the mission will produce image in the \acrshort{rgb} and \acrshort{nir} spectrum at $0.5$m of resolution \cite{lebegue_co3d_2020} using a \acrshort{ccd} sensor cell, and more specifically a Bayer matrix. The pairing of satellites and \acrshort{ccd} sensor used allow for \textit{almost} simultaneous stereo image acquisition, cutting short the transient object problem (\ie objects moving/disappearing between stereo images). To be able to process the amount of data provided by the \acrshort{co3d} mission ($40\,000$ images a day, at $50$cm resolution and covering a footprint of 7km$\times$5km \cite{melet_co3d_2020, lebegue_co3d_2020}), every step of the pipeline has been developed to be highly scalable. Different acquisition schemes can also be used, such as the video mode, or even the `diamond' geometry acquisition, which acquires a quadri-stereo over a couple of days. In parallel with image quality specifications, the \acrshort{co3d} products need to abide to a height relative accuracy of $1$m on low slopes. In addition, the \acrshort{co3d} mission plans to produce a performance map supporting the output DSM. Therefore, investigating uncertainty inside a 3D stereo pipeline is essential for the implementation of this performance map.

\begin{figure}
    \centering
    \includegraphics[width=0.6\linewidth]{Images/Chap_1/CO3D_schema.png}
    \caption{Pairing of satellites for the \acrshort{co3d} mission. (Credit: \acrshort{cnes})}
    \label{fig:co3d_mission}
\end{figure}

\begin{remark}
    When an image is acquired both in panchromatic and RGB mode, it is possible to leverage the high resolution of the panchromatic image to improve that of the color image. This fusion technique is called \textit{pansharpening} \cite{loncan_hyperspectral_2015}. We use this technique for clarity in figures and other illustrations of this thesis (for instance \Cref{fig:VDG_ortho,fig:Pleiade_over_Paris}). It is important to remember that the processed images are the panchromatic images, and not the pansharpened ones which are only used for the final visualization. The \acrshort{co3d} satellites are not concerned with pansharpening, as the used Bayer matrix directly provides \acrshort{rgb} images at 50cm resolution without panchromatic band.
\end{remark} 

Although this thesis contributes to the ground segment of the \acrshort{co3d} mission, we used Pléiades stereo images in the absence of \acrshort{co3d} images (at the time this thesis is written, the \acrshort{co3d} mission has not yet been launched).

\subsection{Geolocation Models}\label{sec:sensors_rpc}
A crucial part of satellite imagery is the ability to perform georeferencement, or georegistratation, of every pixel, \ie, locate their coordinates in an Earth system of coordinates.

Different systems of coordinate exist, for instance the Geodetic Coordinates use latitude and longitude to represent the relative position of a point with regards to a reference ellipsoid, \ie a smooth and regular mathematical approximation of the Earth's shape. Another widely used reference system is the \acrfull{utm} coordinate system, which divides the globe into sixty north-south areas. Each area is $6\degree$ of longitude wide, and is approximated by a plane. The coordinates of a point in each area are Cartesian coordinates with the origin being the intersection of the equator and the meridian of the area. \acrshort{utm} coordinates are employed in GPS systems for instance. In Geodetic Coordinates or \acrshort{utm}, the elevation of a point is defined using the ellipsoid as reference. It is also possible to use the geoid, \ie the gravity equipotential surface, which is less smooth than the ellipsoid but is closer to the actual irregular shape of the Earth.

Physical models possess high geolocation accuracy, but are sensor-specific and are computationally complex. For stereo reconstruction, generalized sensor models are preferred. Specifically, we will focus on \acrfull{rpc}, provided alongside images of many satellites \cite{grodecki_ikonos_2001,devika_developement_2006}, and specifically of the \acrshort{co3d} and Pléiades satellites. Sometimes called Polynomial Mapping Functions \cite{baltsavias_metric_1992} or Rational Function Models \cite{tao_comprehensive_2001}, \acrshort{rpc} are functions allowing to transform a pixel's ground location in Geodetic Coordinates, \ie its longitude, latitude and geodetic height $(X,~Y,~Z)$, into its coordinates $(\rowcol)$ in the image space. \acrshort{rpc} encode lines of sight of the satellite, \ie the line joining the center of the sensor's cell to the ground and going through the optical center of the sensor. To improve numerical stability and minimize computation errors, the image coordinates and ground coordinates are normalized between $-1$ and $1$, using their scale factors $SF$ and mean values:
\begin{eqnarray*}
    SF_X &=& \max(X_\mathrm{max}-\overline{X},~\overline{X}-X_\mathrm{min})\\
    \Tilde{X}&=&\frac{X-\overline{X}}{SF_X}
\end{eqnarray*}
The same process is applied to $Y, ~Z,~row$ and $col$. To avoid heavy notation in this section, we will refer to every normalized coordinate using their non-normalized symbol.

\begin{figure}
    \centering
    \includegraphics[width=0.8\linewidth]{Images/Chap_1/RPC_model.png}
    \caption{Schematic representation of \acrshort{rpc} models.}
    \label{fig:RPC_models}
\end{figure}

Formally, \acrshort{rpc} are defined as rational fractions of polynomials:  
\begin{eqnarray}
    \RPC:\mathbb{R}^3 &\rightarrow&\mathbb{R}^2\nonumber\\
    (X,~Y,~Z) 	&\mapsto& \left(\frac{Num_{row}(X,Y,Z)}{Den_{row}(X,Y,Z)}, \frac{Num_{col}(X,Y,Z)}{Den_{col}(X,Y,Z)}\right)\nonumber\\
    &\mapsto&(row,col)\label{eq:rpc}
\end{eqnarray}
where $Num_{row},~Den_{row},~Num_{col}$ and $Den_{col}$ are the numerators and denominators for rows and columns respectively, expressed as polynomials with a maximum order of $3$:
\begin{eqnarray*}
    Num_{row}(X,~Y,~Z) &=& \sum_{i=0}^3~\sum_{j=0}^{3-i}~\sum_{k=0}^{3-i-j}a_{ijk}X^iY^jZ^k\\
    &=& a_{000} + a_{100} X + a_{010} Y + a_{001} Z + a_{110} XY + a_{101} XZ \\
    &&+ a_{011} YZ + a_{200} X^2 + a_{020} Y^2 + a_{002} Z^2 + a_{111} XYZ \\
    && + a_{210} X^2Y + a_{201} X^2Z + a_{120} XY^2 + a_{102} XZ^2\\
    && + a_{021} Y^2Z + a_{012} YZ^2 + a_{300} X^3 + a_{030} Y^3 + a_{003} Z^3
\end{eqnarray*}
$Den_{row},~Num_{col}$ and $Den_{col}$ respectively possess different coefficients $a_{ijk}$. The order or indexing of $a_{ijk}$ may differ in the literature. For instance, they can be numbered from $0$ to $19$ or $1$ to $20$, and do not refer to the same indeterminate. Numerical values of \acrshort{rpc} coefficient are computed using a set of reference ground control points.

It has been shown in \cite{baltsavias_metric_1992} that \acrshort{rpc} are well suited to be used for ortho-rectification and stereophotogrammetry, as they possess good accuracy and are computationally fast. For stereophotogrammetry, it is often required to use the inverse \acrshort{rpc} model. As \acrshort{rpc} encode a line of sight given an image coordinate, the use of an additional elevation coordinate is required to move from the image space to the object space (\ie to go from a 2D space to a 3D space). It can be a geoid modelling the Earth's surface, or a \acrshort{dsm} with higher resolution. Knowing the true elevation $Z$ of a pixel, we define the inverse model as:
\begin{eqnarray}
    \RPC^{-1}:\mathbb{R}^3 &\rightarrow&\mathbb{R}^3\nonumber\\
    (\rowcol,~Z) 	&\mapsto& (X,~Y,~Z) \label{eq:inverse_rpc}
\end{eqnarray}
An illustration on how \acrshort{rpc} can be used to go from the image space (\rowcol) to the object space $(X,Y,Z)$ is presented in \Cref{fig:RPC_models}. Knowing the position of pixels in the object space using \acrshort{rpc} allows to triangulate the position of an object using multiple images of the same scene. This is how photogrammetry pipelines processing satellite images are computing the elevation of the different cells of a \acrshort{dsm}. The next section will explain the inner workings of those pipelines.

\section{Structure of the Stereophotogrammetry Pipeline}\label{sec:classical_stero_pipeline}
This section dives more into details of the inner workings of stereophotogrammetry, and presents the stereo pipeline mainly considered in this thesis. Photogrammetry is the science of deducing information from photographic images \cite{kasser_photogrammetrie_2001}. A sub domain of photogrammetry is stereophotogrammetry, which specifically consists in deducing 3D information from multiple photographic images. Although multiple stereophotogrammetry setups can be achieved, for instance using structured light \cite{scharstein_high-accuracy_2003} or different wavelength \cite{geng_rainbow_1996}, we focus here on the pipelines designed for processing satellite images, which are used and studied in this thesis. 


\begin{figure}
    \begin{subfigure}[t]{0.5\linewidth}
        \flushleft
        \includegraphics[width=0.7\linewidth]{Images/Chap_1/Sterero_principle_1.png}
        \caption{}
        \label{fig:stereo_principle_1}
    \end{subfigure}
    \begin{subfigure}[t]{0.5\linewidth}
        \flushright
        \includegraphics[width=\linewidth]{Images/Chap_1/Sterero_principle_2.png}
        \caption{}
        \label{fig:stereo_principle_2}
    \end{subfigure}
    \caption{Principle of photogrammetry. The same object is captured by two different cameras in \Cref{fig:stereo_principle_1}. The position of the object in both image spaces in \Cref{fig:stereo_principle_2}, is characterized by its displacement, called \textit{disparity}.}
    \label{fig:stereo_principle}
\end{figure}

The main idea of stereophotogrammetry is to identify the parallax of objects between multiple images, and to deduce the distance between the object and the sensors from this displacement. \Cref{fig:stereo_principle} illustrates how the position of an object can be triangulated from multiple images of a scene. The two lines of sight used to triangulate the position of an object are determined by finding the location of the object in both image spaces. Given the coordinate of an object in a reference image space (for instance the left image in \Cref{fig:stereo_principle}), the corresponding line of sight in the second image is completely determined by the position of the object in the reference image space, and its displacement between the two images spaces. When expressed in pixels, this displacement is called disparity, as presented in \Cref{fig:stereo_principle_2}. Determining the disparity of every pixel of a reference image is called dense-matching, and is done by a dense-matching algorithm, also called \textit{correlator} (as it computes the best correlation between pixels of stereo images). It allows to create pairs of pixels representing the same object in their respective image space. Using the geometrical model of the camera, or the geolocation model of sensors in the case of satellite imagery, the 3D coordinates of the corresponding object of a pair of matched pixels can then be determined by intersecting lines of sight (or best approximation if they do not strictly intersect). This results in a point cloud, encoding the 3D position of every visible object in both images. In satellite photogrammetry, the point cloud can be projected onto a regular grid to obtain a \acrshort{dsm}. The following section dives more into details of satellite photogrammetry pipelines, and specifically on the \acrshort{cars} stereo pipeline.

\subsection{Different Photogrammetry Pipelines}\label{sec:different_photogrammetry_pipelines}
Several stereo pipelines processing satellite images exist in the literature. We can think of NASA's \textit{ASP} \cite{shean_automated_2016}, Centre Borelli's \textit{s2p} \cite{franchis_automatic_2014}, DLR's \textit{CATENA} \cite{kraus_fully_2013}, Ohio State University's \textit{RSP} and \textit{SETSM} \cite{qin_rpc_2016, noh_surface_2017}, \acrshort{ign}'s \textit{MicMac} \cite{rupnik_micmac_2017}. Concerning MicMac, we only consider the methods for satellite photogrammetry, as many other photogrammetry applications and methods are available in this software. For a comparison of those pipelines performances, we refer to \cite{bosch_metric_2017,qin_quality_2021}.

All those pipelines roughly possess the same structure:
\begin{itemize}
    \item Pre-processing. All those pipeline (except S2P) undergo a bundle adjustment step to refine images geometries and correct sensors misalignment.
    \item Images resampling in a convenient geometry for pixel matching. Most pipelines proposes multiple geometries for the resampling of images: the matching can directly be done in the object space, or the secondary image is compute in the geometry of the reference images, or both images are resampled into an epipolar geometry, \ie a geometry where objects move horizontally in both image spaces.
    \item Dense matching. RSP and S2P restrict the choice of the geometry to the epipolar geometry and therefore only propose to do 1D (row-wise) dense matching. Other pipeline allow for a 2D matching.
    \item Triangulation of matches to get 3D point clouds.
    \item Rasterization. When working with more than two images, pipelines all offer the possibility to fuse the obtained pair-wise \acrshort{dsm}s to obtain a global \acrshort{dsm}.
\end{itemize}
It is also noteworthy that when working with more than two images, only MicMac and ASP allow to consider all images at once, while other pipelines use multiple pairings of images. The way those pipeline handle uncertainty will be detailed in section \Cref{sec:related_work}.

We will focus on the \acrshort{cars} pipeline used in this thesis, developed by \acrshort{cnes} \cite{michel_new_2020}. A schematic of the different steps of the pipeline are presented in \Cref{fig:cars_pipeline}.

\begin{figure}
    \centering
    \includegraphics[width=\linewidth]{Images/Chap_1/CARS_pipeline_detailed.png}
    \caption{Different steps of the \acrshort{cars} pipeline}
    \label{fig:cars_pipeline}
\end{figure}

The next sections will look specifically at the \acrshort{cars} stereo pipeline that will be used throughout this thesis.

\subsection{Resampling in Epipolar Geometry}\label{sec:epipolar_geometry}

The \acrshort{cars} pipeline takes as inputs single channel images in sensor geometry alongside their geometric model (\acrshort{rpc} in our case). As the dense matching step compares a single channel in both images, the input images are single channel images. When working with Pléiades images, we use the panchromatic channel as it has the best resolution. For \acrshort{co3d} images, the green channel will be used as the Bayer matrix contains more green pixels.

In many stereo setups, cameras are aligned in such a way that objects only move horizontally between images \cite{geiger_are_2012, scharstein_high-resolution_2014, keselman_intel_2017}. This allows to restrict the search space for pixel matches to a single row instead of the whole image. Most people's eyes also present this alignment. Due to the fixed orbit of satellites, and their acquisition mode for push-broom sensors, it is not possible to maintain the alignment of cameras for satellite photogrammetry as it would be the case for classical stereo setups. Some pre-processing step are thus required to rectify images and ensuring that the displacement of an object only occurs horizontally (see \Cref{sec:epipolar_geometry}). Then matching pixels are determined by computing their disparity in a step called stereo matching presented in \Cref{sec:stereo_matching}.

The first step of the \acrshort{cars} pipeline is to resample the stereo acquisitions into a convenient geometry to carry out the dense matching. This geometry is the epipolar geometry \cite{cnes_imagerie_2008}, which is constructed so that each line of an image follow the movement of the satellite, \ie the objects move only horizontally between the reference and secondary stereo images. This greatly facilitates the matching step as the search space for a match is limited to a one dimensional space instead of a two dimensional space.

Sensors using a pinhole camera model have a perfectly defined epipolar geometry, modeled by an affine transform \cite{hartley_multiple_2004}. This not true for push-broom sensors \cite{morgan_epipolar_2004} such as those used in Pléiades satellites. In the case were no analytical model exist, approximations of the epipolar geometry must be computed \cite{oh_piecewise_2010, koh_unified_2016, michel_new_2020}. The epipolar geometry is computed using the respective geolocation models of both images $\RPC_1$ and $\RPC_2$. Indeed, given the altitude $Z$ of the object represented by pixel $(row_1, ~col_1)$ in the reference image space, we can deduce its ground location $(X, ~Y, ~Z)$ using $\RPC_1^{-1}$. The position in the secondary image space $(row_2,~col_2)$ can then be retrieved using $\RPC_2$: 
\begin{equation}
    (row_2,~col_2) = \RPC_2\circ\RPC_1^{-1}(row_1, ~col_1, ~Z)\label{eq:epipolar_transfer}
\end{equation}
The function $\RPC_2\circ\RPC_1^{-1}$ is called the co-location function $f_{1\rightarrow 2}$ and allows to switch from the reference image space to the secondary image space given an elevation $Z$. By varying the elevation in the range of considered elevations $[Z_{min},~Z_{max}]$, $f_{1\rightarrow 2}$ provides a characterization of the parallax between images. The lines described by $f_{1\rightarrow 2}$ are thus the epipolar curves in the secondary image. Similarly, $f_{2\rightarrow 1}=\RPC_1\circ\RPC_2^{-1}$ provides the epipolar curves in the reference image. The range of considered elevations $[Z_{min},~Z_{max}]$ can be determined using any elevation model. The geoid of the Earth can be used, and there also exists open models such as the NASA's SRTM \cite{farr_shuttle_2007} ($30$m and $90$m of resolution between $-56\degree$ and $60\degree$ of latitude), or ESA's Copernicus DEM ($30$m and $90$m of resolution worldwide).

In practice, determining the epipolar grids using $f_{1\rightarrow 2}$ for every pixel is computationally heavy. Instead, $f_{1\rightarrow 2}$ can be used to compute local affine approximation of the epipolar geometry (similar to what is used with pinhole cameras) as in \cite{de_franchis_stereo-rectification_2014}. When working with large images, tiling effects appear at the border between local approximations. To solve this issue, the \acrshort{cars} pipeline computes a deformation grid approximating the epipolar geometry, as the variations of epipolar lines vary slowly inside an image. The sampling rate of the grid is chosen as to be large enough to facilitate the computation while insuring to grasp variations of the epipolar geometry.

\begin{figure}
    \centering
    \includegraphics[width=\linewidth]{Images/Chap_1/epipolar_lines.png}
    \caption{Computation of an epipolar line for a pixel using small height variations}
    \label{fig:epipolar_lines}
\end{figure}

In order to determine the deformation grid, \cref{eq:epipolar_transfer} is evaluated with an elevation $Z_{coarse}$ extracted from the low resolution elevation model, and for small variations of altitude $Z_{coarse}\pm\delta$. Those three points allow to compute the direction of epipolar lines for every pixel of both images as in \Cref{fig:epipolar_lines}. This method generates the deformation grid $g_{e1}$ joining every epipolar coordinate $(row_e, ~col_e)$ to its position in the reference image $(row_1, ~col_1)$:
\begin{align}\label{eq:epipolar_grid}
    g_{e1}(row_e, ~col_e) = (row_1, ~col_1)
\end{align}

A similar grid $g_{e2}$ is determined for the secondary image, which is computed jointly with $g_{e1}$. For more details on the way epipolar grids are computed, we refer to \cite{michel_new_2020}. Using those grids, it is possible to resample the reference and secondary images in their respective epipolar geometry. 

\begin{remark}
    With this method, an object whose elevation equals that of the reference elevation model $Z_{coarse}$ (\ie the geoid, or a low resolution \acrshort{dsm} such as the Copernicus \acrshort{dem}) has a disparity of $0$ between epipolar images. After the dense matching step, a disparity of $0$ would mean that the object has the same altitude than that of the low resolution elevation model.
\end{remark}

Because the geolocation models have a limited precision, there might be a misalignment left in the epipolar grids. To correct this error, a set of \acrfull{sift} points \cite{lowe_distinctive_2004} is computed between the reference and secondary epipolar images. For every match, the difference between their rows is computed. The secondary epipolar grid is then corrected so that row differences are null on average. Once aligned epipolar grids have been obtained, stereo images can been resampled in a epipolar geometry which minimizes the errors between \acrshort{sift} points.

Along side epipolar grids can be computed the disparity to altitude ratio $r_{alt}$, which corresponds to the altimetric shift resulting from a shift of 1 in disparity. In other words, if a pixel has a disparity of $0$, then its altitude will be $Z_{coarse}$, and if it has a disparity of $1$, then its altitude will be $Z_{coarse}+r_{alt}$. $r_{alt}$ is also called the altimetric ratio, and provides an estimation of the altimetric resolution of the final \acrshort{dsm}. The altimetric ratio varies along an image, but its variations are small enough so that we can safely approximate it by a constant. We will make extensive use of the altimetric ratio $r_{alt}$ in \Cref{chap:elevation_intervals} to compare the uncertainty of \acrshort{dsm} with different altimetric resolutions.

Once the input images in sensor geometry have been resampled into epipolar geometry, the disparity of every pixel can be computed using dense matching algorithms.

\subsection{Stereo Matching}\label{sec:stereo_matching}
Dense matching refers to the pairing of every pixel between two images. It differs from sparse matching were only a restricted sets of pixels must be matched in both images. Sparse matching is used for registering the different channels of an image if they have been acquired separately, for instance using push-broom sensors. For constructing high-resolution \acrshort{dsm}, dense matching is necessary as we are interested in estimating the disparity at a pixel scale.

Dense matching can be performed once epipolar images have been computed. As stereo matching is an important problem in computer vision, multiple algorithms have been proposed to compute a dense disparity map $\mathcal{D}$ mapping erach pixel $(\rowcol)$ to its disparity $d$. Dense matching algorithms can be broadly classified into two categories: classical approaches following the steps outlined by Scharstein \etal \cite{scharstein_taxonomy_2001}, and deep-learning based methods \cite{laga_survey_2022}.
\begin{remark}
	In the domain of stereo matching, the reference image is often referred to as the \textit{left} image, and the secondary image is the \textit{right} image.
\end{remark}

Recently, deep-learning methods have greatly improved the results of stereo matching algorithms for both remote sensing applications \cite{chebbi_deepsim-nets_2023} and more general applications \cite{tosi_survey_2024} such as robotics, autonomous cars or augmented reality. Best results on famous benchmarks have been obtained using 2D an 3D convolution neural networks end-to-end deep-learning approaches \cite{guo_openstereo_2024, liu_playing_2024}. Those models usually undergo a supervised training, with many datasets being available for stereo processing, especially for autonomous cars \cite{geiger_are_2012, geiger_vision_2013}.

There are less open satellite datasets due to the costs and copyrights of satellite images, even if more datasets tends to be released \cite{bosch_semantic_2018, le_saux_data_2019, huang_urban_2022}. Obtaining ground truth data for satellite imagery is also challenging: as classical approaches such as structured light for producing ground truth disparity cannot be used. Instead, costly airborne campaigns are carried-out, and the acquired elevation data must be converted into disparity maps, which is a complex process \cite{cournet_ground_2020}. Additionally, large models often face generalization challenges, particularly when applied to images that differ from their training datasets. It is especially true in the case of satellite imagery \cite{mari_disparity_2022, jiang_rethinking_2024}, as there is a large variety of landscapes, sensors and acquisitions angles to consider. Furthermore, the radiometry can change depending on the time of the day where the acquisition occurred, and landscapes change can also vary greatly between two seasons. All those factors make the training of a performing and generalizable network for satellite imagery quite challenging. Non-supervised algorithm could avoid some of those drawbacks, but they do not present the same performances as supervised networks.

Even though deep stereo algorithms produces exciting results, we will not focus on deep end-to-end algorithms in this thesis and instead restrict ourselves to so-called \textit{classical} methods used in satellite stereo pipelines. Those methods have the advantage of relying on an extensive literature on the subject. More importantly, those methods are explainable as they do not have a ``black-box'' structure similar to that of end-to-end networks. Their interpretability is crucial when modelling and propagating their uncertainty, which is particularly relevant in the context of this thesis.

Classical approaches usually encompass the following steps\cite{scharstein_taxonomy_2001}:
\begin{itemize}
    \item matching cost computation
    \item cost aggregation
    \item disparity optimization and computation
    \item disparity refinement
\end{itemize}
\Cref{fig:stereo_matching_pipeline} illustrates those different steps. We will detail each step into more details in the following sections. We consider the case of the Pandora correlator developed at \acrshort{cnes} (\url{https://github.com/CNES/Pandora}), which is used by the \acrshort{cars} pipeline for the dense matching step.

\begin{figure}
	\centering
	\includegraphics[width=\linewidth]{Images/Chap_1/stereo-matching_pipeline.png}
	\caption{The different steps of classical dense stereo matching algorithms.}
	\label{fig:stereo_matching_pipeline}
\end{figure}

\subsubsection{Matching Cost Computation}\label{sec:cost_volume_computation}
To determine pairs of matching pixels, we first start by measuring the cost of matching two pixels together. Pixels with similar values and similar surroundings will have a low matching cost, while dissimilar pixels will have a high matching cost. Each matching cost value is evaluated using a cost function, which is a mapping $f$ from subsets of the left and right images to $\mathbb{R}$. A cost function $f$ measures the dissimilarity between the two subsets of the left and right images.. \Cref{ex:cost_functions} provides different instances of cost functions.

\begin{example}\label{ex:cost_functions}
	Simple instances of cost functions include the Sum of Absolute Differences (SAD), the Zero Normalized Crossed Correlation (ZNCC)  \cite{hannah_computer_1994}, the CENSUS transform \cite{zabih_non-parametric_1994} and MC-CNN \cite{zbontar_stereo_2016}, which will be detailed in the following paragraphs.
	
	Given two windows $W_L$ and $W_R$ from the left and right images, the SAD cost function is defined as follows:
	\begin{equation}
		f_{SAD}(W_L, W_R)  = \sum_i\sum_j | W_L(i,j) - W_R(i,j) |
	\end{equation}
	Low $f_{SAD}$ values indicate that the windows are similar, while high values indicate noticeable differences. A $f_{SAD}$ value of $0$ indicates that the windows are identical. This cost function is probably the simplest cost function one could imagine, and will be used in \Cref{chap:propagating} to didactically illustrate how uncertainty models can be propagated throughout a cost function. However, this cost function is not usually employed in practice as it is based on intensity differences, and is thus very sensitive to radiometric changes. It can however be a fast and easy way to have a first estimate of similarities between multiple patches. The SAD can also be used for motion estimation and image/video compression \cite{richardson_h264_2006}.
	
	The ZNCC cost function is defined as the correlation coefficient between both images:
	\begin{equation}
		f_{ZNCC}(W_L, W_R)  = \sum_i\sum_j \frac{(W_L(i,j)  - \tilde{W}_L) (W_R(i,j)  - \tilde{W}_R) }{\sigma_L\sigma_R}
	\end{equation}
	where $\tilde{W}$ refers to the mean value of a window, and $\sigma$ its standard deviation. Negatively correlated windows would present a ZNCC value of $-1$ and positively correlated windows present a ZNCC value of $1$. Contrary to the SAD cost function, matching windows will be indicated by a high value of the ZNCC. It is thus not \textit{strictly} a cost function but rather a similarity function. It is not a problem, as multiplying $f_{ZNCC}$ by $-1$ will transform it into a cost function. Another formulation could be to say that the $ZNCC$ is a \textit{maxitive} cost function, in the sense where potential matches are found by searching for its maximum. Conversely, the SAD is a \textit{minitive} cost function in the sense where potential matches are indicated by a minimal cost. The $ZNCC$ cost function performs well in homogeneous areas when computed over a large windows. However in an urban settings, it struggles to correctly estimate buildings boundaries.
	
	The CENSUS cost function needs to be a bit more detailed. For a squared window $W$ with a side of $2n+1$ pixels, we first compare the value of each pixel of the window with the center pixel. This gives a binary string where $1$ indicates that the value of the pixel is superior to that of the center pixel. For instance if we consider the two $3\times3$ following windows $W_L, W_R$:
	$$
    \begin{bmatrix}
        155 & 133 & 97 \\
        80 & 110 & 132 \\
        100 & 102 & 120
    \end{bmatrix}
    \qquad
    \begin{bmatrix}
        175 & 153 & 133 \\
        100 & 130 & 152 \\
        120 & 135 & 125
    \end{bmatrix}
	$$
	Then comparing each of their pixel to the center of the windows will yield the following binary strings (expressed here as matrices):
	$$
    \begin{bmatrix}
        1 & 1 & 0 \\
        0 &  & 1 \\
        0 & 0 & 1
    \end{bmatrix}
    \qquad
    \begin{bmatrix}
        1 & 1 & 1 \\
        0 &  & 1 \\
        0 & 1 & 0
    \end{bmatrix}
	$$
	The cost function $f_{CENSUS}$ is finally obtained by taking the Hamming distance (\ie the number of different bits between those two strings):
	\begin{equation*}
		f_{CENSUS}(W_L, W_R) = 3
	\end{equation*}
	
	The CENSUS cost function compares relative intensity variations, it is thus less sensitive to variations of intensities between images, such as a change of exposure for instance. Similarly to the SAD, two similar patches will tend to have a low value.
	
	The MC-CNN cost function \cite{zbontar_stereo_2016} is using a convolutional neural network architecture to measure the similarity between patches. It was train on $11\times 11$ patches from stereo images from the Kitti \cite{geiger_vision_2013, menze_object_2015} and Middlebury \cite{scharstein_taxonomy_2001,scharstein_high-accuracy_2003,hirschmuller_evaluation_2007,scharstein_learning_2007,scharstein_high-resolution_2014} datasets. They first trained a siamese neural network to compute a vector of features for each patch, then compute the dot product between both vectors to obtain a similarity measure. MC-CNN is thus a maxitive cost function. \Cref{fig:mccnn} details the architecture of the network. Although it has been trained on stereo images of autonomous cars (Kitti) and stereo images of toys (Middlebury), it generalizes well to satellite images \cite{defonte_evaluation_2021}.
	
	{\centering\includegraphics[width=0.8\linewidth]{Images/Chap_1/MCCNN.png}\captionof{figure}{MC-CNN architecture}\label{fig:mccnn}}
\end{example}

Matching costs are evaluated for every potential pairs of pixels whose corresponding disparity lies in a given disparity range. Matching cost values are stored in an array of data, called \textit{cost volume} (\Cref{fig:cost_volume}). The term volume is used as a matching cost value is determined by three coordinates in the cost volume: the row and column $(\rowcol)$ of the left pixel as well as the considered disparity $d$:
\begin{equation}\label{eq:cost_volume}
    C_V(\rowcol, ~d) = f(W_L(\rowcol),  W_R(\rowcol+d))
\end{equation}
where $W_L(\rowcol)$ is a window centered on the pixel at coordinates $(\rowcol)$ in the left image, and $W_R(\rowcol+d)$ is a window centered on the pixel at $(\rowcol+d)$ in the right image. Some padding is usually added to the images to avoid problems near borders where the column $col+d$ would not be defined. Given a pixel in the left image $(\rowcol)$, matching cost values for every considered disparity form what is called a cost curve. In theory, the correct disparity for a match is determined by finding the minimum of the cost curve (for minitive cost functions such as SAD or CENSUS). \Cref{fig:cost_volume} represents a cost volume, and one of its cost curves where potential matches have been highlighted. In the rest of this thesis, we will consider that a cost function is always minitive unless specified otherwise. In practice, directly estimating the disparities from the cost volume is not efficient as we only consider local information. The resulting disparity map is often very noisy, as can be seen in \Cref{fig:disparity_map_no_sgm}. In order to reduce this noise, one solution is to aggregate the costs of neighboring pixels in the cost volume.

\begin{figure}
	\centering
	\includegraphics[width=\linewidth]{Images/Chap_1/Cost_volume.png}
	\caption{Matching cost volume and one of its cost curve}
	\label{fig:cost_volume}
\end{figure}

\begin{figure}
    \begin{subfigure}[t]{0.42\linewidth}
        \flushleft
        \includegraphics[width=\linewidth]{Images/Chap_1/disparity_map_no_sgm.png}
        \caption{Disparity with local cost only}
        \label{fig:disparity_map_no_sgm}
    \end{subfigure}\hfill
    \begin{subfigure}[t]{0.514\linewidth}
        \flushright
        \includegraphics[width=\linewidth]{Images/Chap_1/disparity_with_sgm.png}
        \caption{Disparity with \acrshort{sgm} regularization}
        \label{fig:disparity_map_with_sgm}
    \end{subfigure}
    \caption{Disparity map without and with \acrshort{sgm} regularization.}
    \label{fig:disparity_map_with_and_without_sgm}
\end{figure}

\subsubsection{Cost Aggregation}
The usage of windows in the calculation of the cost function allows to take into consideration the surrounding of pixels to better measure their similarity. However, window based approaches, which are mostly square-shaped, also present the disadvantage of struggling to correctly identify matches near object borders \cite{hirschmuller_real-time_2002}. This is usually called an adherence effect, represented in \Cref{fig:adherence_window}. In \Cref{fig:adherence_left}, we can see that the considered pixel is at the border of an object with the left window $W_L$ spanning over both part of the border. In \Cref{fig:adherence_cost_curve}, we can see that the minimum of the cost curve (around disparity -40) is not exactly located at the true disparity. This is in part due to the fact that $W_L$ contains pixels from two objects with different disparities, which influences the matching cost values. A low matching cost thus does not necessarily mean that the center pixels constitute a match. The term ``adherence'' is used as the matching windows tend to falsely estimate the shift in disparity near objects border, as if it adhered to the object. Other work have been proposing to use a spatial weighting \cite{kuk-jin_yoon_locally_2005}, segmentation \cite{hutchison_segmentation-based_2007}, windows with different shapes \cite{ke_zhang_cross-based_2009, buades_reliable_2015} to solve this problem. 

\begin{figure}
    \begin{subfigure}[t]{0.35\linewidth}
        \flushleft
        \includegraphics[width=\linewidth]{Images/Chap_1/adherence_overview_left.png}
        \caption{Left image with $W_L$ in orange}
        \label{fig:adherence_overview_left}
    \end{subfigure}\hfill
    \begin{subfigure}[t]{0.6\linewidth}
        \flushright
        \includegraphics[width=\linewidth]{Images/Chap_1/adherence_left.png}
        \caption{Zoom over $W_L$, in orange}
        \label{fig:adherence_left}
    \end{subfigure}\\
    \begin{subfigure}[t]{\linewidth}
        \flushright
        \includegraphics[width=0.87\linewidth]{Images/Chap_1/adherence_right.png}
        \caption{Right image. Considered windows in orange, and $W_R$ for the true disparity in blue}
        \label{fig:adherence_right}
    \end{subfigure}\\
    
    \begin{subfigure}[t]{\linewidth}
        \centering
        \includegraphics[width=\linewidth]{Images/Chap_1/adherence_cost_curve.png}
        \caption{Corresponding minitive cost curve ``$-f_{ZNCC}$''}
        \label{fig:adherence_cost_curve}
    \end{subfigure}
    \caption{Example of the adherence effect. \Cref{fig:adherence_overview_left,fig:adherence_left} present the left window $W_L$ at the border of an object. \Cref{fig:adherence_right} presents all the windows from the right image in the considered disparity range, that will be compared to $W_L$. The window corresponding to the true disparity appears in blue. \Cref{fig:adherence_cost_curve} details the corresponding cost curve, where the true disparity does not correspond to the minimum of the cost curve.}
    \label{fig:adherence_window}
\end{figure}

After computing the matching cost, information from neighboring pixels can be incorporated in the cost volume. A first approach is to aggregate different parts of the cost volume. Usually, costs of pixels belonging to the same objects are aggregated using different methods for segmentation \cite{ke_zhang_cross-based_2009, ji_superpixel_2021}. This part is not always present in algorithms, as we will see we can consider global information directly when determining the disparity map.

\subsubsection{Disparity Optimization and Computation}
Computing the disparity map from the cost volume can be done in several ways. So-called local methods apply a direct \textit{winner-takes-all} strategy, where the $\argmin$ of every cost curve is kept as the selected disparity. On the other hand, global method use the information contained in the cost volume to solve an optimization problem, where the objective is to compute the disparity map $\mathcal{D}$ minimizing an energy function expressed as follows:
\begin{equation}
    E(\mathcal{D}) = E_{data}(\mathcal{D}) + \lambda E_{smooth}(\mathcal{D})
\end{equation}
where $\lambda$ is a scalar for tuning the importance of the regularization term $E_{smooth}$. Usually, the data term is directly computed from the cost volume as:
\begin{equation}\label{eq:global_methods}
    E_{data}(\mathcal{D}) = \sum_{row,~col}C_V(\rowcol,~\mathcal{D}(\rowcol))
\end{equation}
The regularization term $E_{smooth}$ can take numerous forms, usually measuring if the neighboring disparities possess similar values \cite{scharstein_taxonomy_2001}. Then a local minimum for this energy is found using various methods, such as Markov Random Fields \cite{boykov_markov_1998, sun_stereo_2003}, graph cuts \cite{kolmogorov_computing_2001} or minimum spanning trees \cite{zureiki_stereo_2008, qingxiong_yang_non-local_2012}. Those algorithms improves performances in comparison to local methods, but can be computationally expensive.

A most popular method is called \acrfull{sgm} \cite{hirschmuller_accurate_2005}: it aims at incorporating regularization constraints to the cost volume similarly to global methods, while being processed with relative low computational cost like local methods. Its name It is, in a way, in-between local and global methods. This is the method used in the \acrshort{cars} pipeline, and for experiments in this thesis. In a few words, for each pixel, the \acrshort{sgm} algorithm computes a regularized cost volume, which is based on the regular cost volume (the $E_{data}$ term), and adds cost penalties to the disparities differing from that of their neighboring pixels (the $E_{smooth}$ term). Although the formulation of the \acrshort{sgm} algorithm can be expressed in a few equations, understanding its inner workings is more complex. We will first present its mathematical formulation in \cref{eq:sgm,eq:sgm_penalties}, and use \Cref{fig:sgm} to illustrate its effect on (a portion of) a cost curve. \Cref{fig:cost_curve_with_without_sgm} presents the effect of \acrshort{sgm} regularization on different cost curves and \Cref{fig:disparity_map_with_and_without_sgm} presents the effect of \acrshort{sgm} regularization on the final disparity map. Formally, for every pixel $p=(\rowcol)$, the cost volume is explored in multiple directions as in \Cref{fig:sgm_directions}.

\begin{figure}
	\centering
	\includegraphics[width=0.5\linewidth]{Images/Chap_1/SGM_directions.png}
	\caption{\acrshort{sgm} regularization with $8$ directions}
	\label{fig:sgm_directions}
\end{figure}

The cost is regularized for each direction $r$ to take into account the best disparity $d$ along that direction. A direction can be for instance $r=(0,~1)$, meaning that we will look at same row and travel to the right of the image when browsing direction $r$. Given two positive scalars $P_1<P_2$, the regularized cost $L_r$ along direction $r$ and at disparity $d$ is expressed with the following recursive formulation:
\begin{align}\label{eq:sgm}
    L_r(p,d) = C_V(p,d) + \min_\delta \left(L_r(p-r,~\delta) + R(d, ~\delta)\right)
\end{align}
where $\delta$ is a dummy variable to explore the disparity range, and where $R(d, ~\delta)$ equals:
\begin{align}
    R(d, ~\delta) = &P_1\cdot\mathds{1}(|d-\delta|=1) + P_2\cdot\mathds{1}(|d-\delta|\geqslant 2) - \min_k L_r(p-r,~k) \label{eq:sgm_penalties}
\end{align}
Here, $\mathds{1}$ is the indicator function. The first term of \cref{eq:sgm} is the cost volume, which can be compared to $E_{data}$ in \cref{eq:global_methods}. The second term is similar to the regularisation term $E_{smooth}$. This term can be seen as the regularized cost $L_r(p-r, ~d)$ from the previous pixel in direction $r$, but shifted of $P_1$ for neighboring disparities $d\pm1$, and of $P_2$ for further disparities. The term $\min_k(L_r(p-r,~k))$ prevents $L_r(p,\cdot)$ to diverge to very high values by ensuring that the minimum of $L_r(p-r,~\delta) - \min_k(L_r(p-r,~k))$ always equals $0$. \Cref{fig:sgm} presents $L_r(p-r,~\delta) + R(d, ~\delta)$ for two different disparities $d$ and $d+1$ in the blue frame. The minimum of each curve $L_r(p-r,~\delta) + R(d, ~\delta)$ (blue arrows in the figure) are then added to the matching cost $C_V(p,d)$ to obtain the final regularized cost $L_r(p,d)$. Note that in \Cref{fig:sgm}, the minimum of $L_r(p-r,~\delta)+R(d+1, ~\delta)$ equals $0$, thus $C_V$ is unchanged for this disparity. This is because $d+1$ is the minimum of $L_r(p-r,~\cdot)$. In a way, the regularized matching cost will be a mixture between $C_V(p,\cdot)$ and $L_r(p-r,~\cdot)$. Indeed, the regularized cost $L_r(p-r,\cdot)$ indicates that disparity $d+1$ seems likely, while disparity $d$ is unlikely. On the other hand, the matching cost $C_V(p,\cdot)$ indicates that both disparities $d$ and $d+1$ are likely. The final regularized cost $L_r(p-r,\cdot)$ takes into account that both $C_V(p,\cdot)$ and $L_r(p-r,\cdot)$ agree that $d+1$ is likely, but that there is a disagreement on $d$, and thus slightly increases the regularized cost $L_r(p,d)$. Disparities that seem unlikely to both $C_V(p,\cdot)$ and $L_r(p-r,\cdot)$ lead to a larger increase of the regularized cost $L_r(p,\cdot)$.

\begin{figure}
    \begin{subfigure}[t]{0.49\linewidth}
        \flushleft
        \includegraphics[width=\linewidth]{Images/Chap_1/cost_curve_no_sgm_row_100_col_250.png}
        \caption{Cost curve without \acrshort{sgm}}
        \label{fig:cost_curve_no_sgm_row_100_col_250}
    \end{subfigure}\hfill
    \begin{subfigure}[t]{0.49\linewidth}
        \flushright
        \includegraphics[width=\linewidth]{Images/Chap_1/cost_curve_no_sgm_row_276_col_360.png}
        \caption{Cost curve without \acrshort{sgm}}
        \label{fig:cost_curve_no_sgm_row_276_col_360}
    \end{subfigure}\\
    \begin{subfigure}[t]{0.49\linewidth}
        \flushleft
        \includegraphics[width=\linewidth]{Images/Chap_1/cost_curve_sgm_row_100_col_250.png}
        \caption{Cost curve from \Cref{fig:cost_curve_no_sgm_row_100_col_250} with \acrshort{sgm}}
        \label{fig:cost_curve_sgm_row_100_col_250}
    \end{subfigure}\hfill
    \begin{subfigure}[t]{0.49\linewidth}
        \flushright
        \includegraphics[width=\linewidth]{Images/Chap_1/cost_curve_sgm_row_276_col_360.png}
        \caption{Cost curve from \Cref{fig:cost_curve_no_sgm_row_276_col_360} with \acrshort{sgm}}
        \label{fig:cost_curve_sgm_row_276_col_360}
    \end{subfigure}
    \caption{Different cost curves, with and without \acrshort{sgm} regularization}
    \label{fig:cost_curve_with_without_sgm}
\end{figure}

\begin{remark}
    For the sake of the argument, lets assume that $P_1=P_2$ for now. The value added to $C_V(p, ~d)$ in \eqref{eq:sgm} will be lower than $P_2$ only if $(L_r(p-r,~d) - \min_k L_r(p-r,~k))$ is less than $P_2$. This means that we will always add a penalty of $P_2$ to the cost $C_V$, except if $L_r(p-r,~d)$ is less than $P_2$ away from its minimum. $P_2$ thus represents the penalty that must be overcame in $L_r(p-r,~d)$ in order to not penalize $C_V(p,d)$, or slightly penalize $C_V(p,d)$.
    
    If $P_1<P_2$, then we can draw similar conclusions, except that we also accept to reduce the penalty to $C_V(p,d)$ if a neighboring disparity, \ie $d\pm1$ is less than $P_1$ away from the minimum.
\end{remark}

Formulation of \cref{eq:sgm} is recursive, and thus must be initialized. The regularization curves thus begin at the borders of the image and their value is, by convention, $0$ when undefined. For instance with $r=(0,1)$:
\begin{align*}
    L_r((row, 0),d) = C_V(row, 0 ,d)
\end{align*} 

\begin{figure}
	\centering
	\includegraphics[width=\linewidth]{Images/Chap_1/SGM.png}
	\caption{Schematic explanation of the \acrshort{sgm} algorithm in a single direction $r$. Top: regularized cost $L_r$ at $p-r$. In the blue frame $L_r(p-r,~\delta)+R(d,~\delta)$ for two consecutive disparities $d$ and $d+1$ ($L_r-\min_k L_r$ appears in dotted line for clarity). Penalty $P_1$ appears in yellow, penalty $P_2$ appears in red, the minimum of each curve is denoted by a blue arrow. Bottom left: cost volume $C_V$ at $p$. Bottom right: regularized cost $L_r$ at $p$ where the minima of \textit{all} $L_r(p-r,~\delta)+R(d,~\delta)$ have been added}
	\label{fig:sgm}
\end{figure}

When all regularized cost curves $L_r$ have been computed, they are summed to obtain the regularized cost volume $C_V^{SGM}$:
\begin{equation}
    C_V^{SGM}(p, d) = \sum_r L_r(p,d)
\end{equation}
The cost volume $C^{SGM}_V$ contains the regularized cost for every  considered direction. The number of direction used in our experiments is 8 as in \Cref{fig:sgm_directions}, but more directions can be considered in order to consider a more global coverage. Since the original paper \cite{hirschmuller_accurate_2005}, different variations of the \acrshort{sgm} algorithm have been proposed. For instance using different regularization paths \cite{facciolo_mgm_2015} or a different strategy for the aggregation of costs \cite{poggi_learning_2016}.

In order to insure smooth surfaces, \acrshort{sgm} penalties $P_1$ and $P_2$ must be high enough. This however leads to a reluctance to detect discontinuities in the disparity map. Indeed, \acrshort{sgm} penalizes disparity changes, therefore strong variations of disparity are badly reconstructed. For instance near the border of a building, the term $E_{smooth}$ tends to be predominant, leading to rounded borders and soft edges of buildings, instead of the expected sharp disparity variations. This is even reinforced with the window adherence problem presented in \ref{fig:adherence_window}. \Cref{fig:DSM_toulouse} presents a comparison of two \acrshort{dsm}: one obtained directly from \acrshort{lidar} HD data \cite{monnet_lidarhd_2023}, and the other using the \acrshort{cars} pipeline with the CENSUS cost function over a $5\times5$ window and with \acrshort{sgm} regularization. We can see on this figure that the border of the buildings obtained using \acrshort{sgm} regularization are not as sharp and precise as the ground truth provided by the \acrshort{lidar} \acrshort{dsm}. To provide a solution to this problem, one might limit the \acrshort{sgm} regularization to pixels from the same object using a segmentation \cite{dumas_improving_2022}. This method relies on the quality of the segmentation method used and can become quite costly. It has not been considered in the context of this thesis.

\begin{figure}
    \centering
    \begin{subfigure}[t]{0.5\linewidth}
        \centering
        \includegraphics[height=6cm]{Images/Chap_1/DSM_Capitole_LiDAR.png}
        \caption{\acrshort{dsm} from \acrshort{lidar} HD data,\\Place du Capitole}
        \label{fig:DSM_capitole_lidar}
    \end{subfigure}\hfill
    \begin{subfigure}[t]{0.5\linewidth}
        \centering
        \includegraphics[height=6cm]{Images/Chap_1/DSM_Capitole_CARS.png}
        \caption{\acrshort{cars} \acrshort{dsm} from Pléiades images,\\Place du Capitole}
        \label{fig:DSM_capitole_cars}
    \end{subfigure}
    \begin{subfigure}[t]{0.5\linewidth}
        \centering
        \includegraphics[height=6cm]{Images/Chap_1/DSM_Wallon_LiDAR.png}
        \caption{\acrshort{dsm} from \acrshort{lidar} HD data,\\Ernest Wallon stadium}
        \label{fig:DSM_ernest_wallon_lidar}
    \end{subfigure}\hfill
    \begin{subfigure}[t]{0.5\linewidth}
        \centering
        \includegraphics[height=6cm]{Images/Chap_1/DSM_Wallon_CARS.png}
        \caption{\acrshort{cars} \acrshort{dsm} from Pléiades images,\\Ernest Wallon stadium}
        \label{fig:DSM_ernest_wallon_cars}
    \end{subfigure}
    \caption{Different \acrshort{dsm}s over Toulouse, France. \acrshort{dsm}s were obtained by rasterizing \acrshort{lidar} HD data or by processing Pléiades stereo images with the \acrshort{cars} pipeline. The \acrshort{cars} pipeline uses a $5\times5$ CENSUS cost function and \acrshort{sgm} regularization for stereo reconstruction. \copyright \acrshort{cnes} 2017, Distribution AIRBUS DS}
    \label{fig:DSM_toulouse}
\end{figure}

Once the regularized cost volume $C_V^{SGM}$ has been computed, the classical approach in the literature is to apply a \textit{winner-takes-all} strategy to determine the disparity map $\mathcal{D}$:
\begin{align}
    \mathcal{D}(\rowcol) = \argmin_d C_V^{SGM}(row, col, d) 
\end{align}

\subsubsection{Disparity Refinement}\label{sec:postprocess_disparity}
Once the disparity map has been computed, it is usually post-processed to remove artifacts, and improve disparity resolution. It is common to add a sub-pixel refinement step, where a non-integer disparity is interpolated around the selected disparity. The main idea is to interpolate a model through the selected disparity $d=\argmin_\delta C_V(\rowcol, ~\delta)$ and its two direct neighbors. The refined disparity $d_{interp}$ is then defined as the $\argmin$ of this interpolation model. \Cref{fig:sub-pixel_refinement} presents examples of interpolated disparities from  \cite{haller_real-time_2010}, mainly a ``V''-like shape as in \ref{fig:vfit_refinement} and a parabola as in \ref{fig:parabola_refinement}. Carrying out sub-pixel refinement suggests we assume the algorithm can attain a significant level of precision, which is debatable (see \Cref{sec:uncertainty_cars}), as well as the fact that the cost function was sufficiently sampled to be correctly interpolated.

\begin{figure}
    \centering
    \begin{subfigure}[t]{0.5\linewidth}
        \centering
        \includegraphics[width=\linewidth]{Images/Chap_1/subpixel_refinment_vfit.png}
        \caption{V-fit sub-pixel refinement}
        \label{fig:vfit_refinement}
    \end{subfigure}\hfill
    \begin{subfigure}[t]{0.5\linewidth}
        \centering
        \includegraphics[width=\linewidth]{Images/Chap_1/subpixel_refinment_parabola.png}
        \caption{Parabola sub-pixel refinement}
        \label{fig:parabola_refinement}
    \end{subfigure}
    \caption{Different methods for refining the disparity $d=\argmin_\delta C_V(\rowcol, ~\delta)$}
    \label{fig:sub-pixel_refinement}
\end{figure}

The disparity map can also be filtered in order to remove local outliers. Classical strategies include the use of a mean or median filter for instance. There exists more advance filtering methods, such as bilateral filtering \cite{tomasi_bilateral_1998} which performs a weighted average, where weights depend on the proximity to other pixels both in the spatial and spectral domain.

There exist strategies allowing to detect dubious matches. For instance, cross-checking \cite{fua_combining_1991} verifies if the disparity selected would stay consistent with the one computed by reversing the roles of the images: the reference image becomes the secondary image. We first start by computing a secondary cost volume $C_V'$ by reversing the role of the images. We then obtain two disparity maps $\mathcal{D}$ and $\mathcal{D}'$, and check that they are consistent, as in theory we would that $\mathcal{D}=-\mathcal{D}'$. A pixel $(\rowcol)$ is considered consistent across disparity maps if it verifies:
\begin{equation}\label{eq:cross-checking}
    |\mathcal{D}(row, col) + \mathcal{D}'(row, col+\mathcal{D}(row, col))|\leqslant \tau
\end{equation}
where $\tau$ is a consistency threshold usually set to $1$. For errors that do not satisfy this consistency check, it has been proposed in \cite{hirschmuller_stereo_2008} to differentiate between mismatched pixels and occluded pixels. Mismatch pixels are pixels for which there exists a correct match, whereas occluded pixels are pixels that are visible in an image but masked by an object in the other. The difference is determined as follows: if there is a disparity $d\in\mathcal{D}$ such that $\mathcal{D}'(row, col+d)=-d$ then it is a mismatch. Otherwise, pixel $(\rowcol)$ is occluded. Occluded regions of an image can be filled by interpolation with the closest valid disparities or left as \textit{nodata}.  

\begin{remark}
     \Cref{eq:cross-checking} requires the computation of a second cost volume, which effectively doubles the processing time. As dense stereo matching is also the longest part of a photogrammetry pipeline, one might be reluctant to carry out a cross-checking step. However, we can make the following observation: the cost volume contains the dissimilarity between every considered pair of pixels in the disparity range, so the cost of every considered match in the first cost volume is also present in the second cost volume. In theory, for every pixel $(\rowcol)$ and for every disparity $d$ in the (reference) disparity range, it holds \cite{bebis_mutual_2008}:
     \begin{align}\label{eq:diagonal_search}
         C_V(\rowcol,~d) = C_V'(\rowcol+d, -d)
     \end{align}
     \Cref{eq:diagonal_search} holds only for cost volumes obtained after the matching cost step. However, when \acrshort{sgm} regularization or another aggregation cost method modifies cost volumes, then the there can be differences between $C_V(row,~col,~d)$ and $C_V'(row,~col+d, -d)$. In the case of \acrshort{sgm} regularization, series of tests show that the differences between cost volumes are small and marginally modify the disparity maps. Using \eqref{eq:diagonal_search} then allow to obtain both cost volumes by only computing one and re-indexing it to obtain the other.
\end{remark}

Once every step of the dense matching algorithm have been carried out, the disparity map obtained is used create pairs of lines of sight to triangulate 3D point clouds.

\subsection{Triangulation}\label{sec:triangulation}
The disparity maps allow to deduce elevation information contained in the shift of objects in both image spaces. Computing this information is carried out during the triangulation step of the stereo pipeline.

The type of optical sensor used determines the resolution of the disparity map, and therefore the density of the resulting point cloud. It will therefore influence the planimetric resolution of the final \acrshort{dsm}. We usually consider that the planimetric resolution of the \acrshort{dsm} can be similar to the resolution of the sensors. On the other hand, the altimetric resolution of the \acrshort{dsm} is determined by the altitude and positions of the satellites \cite{qin_critical_2019}. This position can be characterized by the \acrfull{b/h} ratio, as in \Cref{fig:RPC}. This ratio is computed by dividing the distance separating the stereo acquisitions by the altitude of the satellite. It indicates the angle formed between the lines of sight originating from the satellites towards an object of the scene. A high \acrshort{b/h} allows for high elevation accuracy, but possesses more occluded regions (for instance a narrow street between two high buildings), and conversely for a low \acrshort{b/h} \cite{delon_small_2007}. In natural landscapes, stereo acquisitions present high \acrshort{b/h} ratio in order to have a high altimetric accuracy, while the \acrshort{b/h} ratio is smaller in urban landscapes to limit the occlusions.  In our experiments, the \acrshort{b/h} ratio for stereo acquisitions varies between $0.1$ and $0.4$. The \acrshort{co3d} will use \acrshort{b/h} ration between $0.2$ and $0.3$. The \acrshort{b/h} will also partly determine the value of the altimetric ratio $r_{alt}$ presented in \Cref{sec:epipolar_geometry}.

The quality of the final \acrshort{dsm} is also influenced by the solar angles \cite{qin_critical_2019}, with the zenith angle being as small as possible to limit the projected shadows over the scene. Sun-synchronous satellites, such as Pléiades or \acrshort{co3d}, benefit from advantageous sun angles in order to favor good images acquisitions and good stereo reconstruction.

\begin{figure}
    \centering
    \includegraphics[width=0.8\linewidth]{Images/Chap_1/RPC.png}
    \caption{Triangulation of the position of a point using \acrshort{rpc} models. The distance $B$ between satellites as well as the global elevation $H$ can be used to determine the angle between satellites.}
    \label{fig:RPC}
\end{figure}

The computed disparity map allows to determine the position of 3D points. When working with pinhole camera models, the depth $z$ of a pixel is computed using the following formula:
\begin{equation}
	z=\frac{Bf}{d}\label{eq:z_bfd}
\end{equation}
where $B$ is the baseline between cameras, $f$ is the focal length of the camera, and $d$ is the disparity of a pixel. This formula can be found using optical geometry \cite{bolles_epipolar-plane_1987}, and illustrates the fact that pixels closer to the camera present a bigger position shift in between images. As mentioned previously, the pinhole camera model is not valid in the case of satellite imagery,  and we instead use other sensor models (see \Cref{sec:sensors_rpc}). \Cref{eq:z_bfd} thus cannot be used as such to provide accurate results. The disparity 

Instead, the disparity is used to determine pairs of lines of sight pointing to the same object. Using \acrshort{rpc} models, we can compute the 3D coordinates of every intersection of pairs of lines of sight. To do so, we first consider the \acrshort{rpc} models $\RPC_1$, $\RPC_2$ of two stereo images defined in \cref{eq:rpc}. We consider as well their epipolar grids $g_{e1}$, $g_{e2}$ from \cref{eq:epipolar_grid}. For every pixel $(row_e, ~col_e)$ from the reference epipolar image, whose disparity is $d$, the 3D point $(X, Y, Z)$ represented by $(row_e, ~col_e)$ is the point verifying the following equations:
\begin{align}
    (X, Y, Z) &= \RPC^{-1}_1(g_{e_1}(row_e,~col_e),~Z)\label{eq:triangulation_exact}\\
    (X, Y, Z) &= \RPC^{-1}_2(g_{e_2}(row_e,~col_e+d),~Z) 
\end{align}
If the lines of sight intersect, then $Z$ is found by solving the following equation:
\begin{equation}
     \RPC^{-1}_1(g_{e_1}(row_e,~col_e),~Z) = \RPC^{-1}_2(g_{e_2}(row_e,~col_e+d),~Z) 
\end{equation}
Knowing $Z$ and $g_{e_1}(row_e,~col_e)$, \cref{eq:triangulation_exact} provides the $X$ and $Y$ coordinates as well.

A known problem is that lines of sight rarely intersect. By approximating lines of sight by a starting point $P$ and a direction vector $\overrightarrow{V}$, we can instead define the coordinates $(X,Y,Z)$ as the point minimizing its distance to both lines. Approximating \acrshort{rpc}s by segments is valid on the limited range of considered altitudes. The exact formula for computing $(X,Y,Z)$ is provided in \cite{delvit_geometric_2006}:
\begin{equation}
    (X,Y,Z) = \left[ Id - V_1V_1^T + Id - V_2V_2^T \right]^{-1} \left[ (Id - V_1V_1^T)P_1 + (Id - V_2V_2^T)P_2 \right]
\end{equation}
where $Id$ is the identity matrix, and $V^T$ is the transposed vector $V$. The error due to enforcing \acrshort{rpc}s  intersection is however small (a few centimeters) in comparison to the errors occurring in the disparity map. In this thesis, we mainly focus on processing the disparity errors an therefore do not consider the approximation made during the lines of sight intersection.

Triangulating every point using the disparity map leads to a 3D point cloud. Because we know from which pixel each 3D point originates from, we can associate to every 3D point additional information such as:
\begin{itemize}
    \item The color of the reference pixel (if provided). The point cloud is thus a colored point cloud.
    \item Confidence measures computed during the dense matching step, that will presented in \Cref{sec:uncertainty_pandora}
\end{itemize}

The 3D points can be filtered to remove obvious errors. Advanced filtering methods can be implemented, such as bilateral filtering \cite{digne_bilateral_2017}, and filtering using color information or confidence measures from the disparity map (see \Cref{sec:uncertainty_pandora}) \cite{youssefi_geometrically_2024}. In the \acrshort{cars} pipeline, two different filtering steps are carried out: one for removing statistical outliers, and one for removing so-called ``small-components''.

\subsubsection{Statistical outliers filtering}
Statistical outliers are determined by considering the positions of their neighbors. For each point $P$, we compute the mean distance $\mu_P$ to its $N$ neighbors. Then we can compute for each point $P$ the mean distance $\mu$ and standard deviation $\sigma$ of its $N$ neighbors as:
\begin{align}
    \mu &= \frac{1}{N}\sum_{i=1}^N\mu_i\\
    \sigma &= \sqrt{\frac{1}{N-1}\sum_{i=1}^N(\mu_i-\mu)^2}
\end{align}

A point $P$ is considered to be a statistical outlier if the difference between its mean distance $\mu_P$ and the mean distance $\mu$ of the whole neighboring is too large:
\begin{align}\label{eq:statistical_outlier}
    \mu_P>\mu+k\sigma
\end{align}
Where $k$ is a constant usually set to $5$ by the user. We also consider $N=50$.

\subsubsection{Small-components filtering}
The other filtering method, named small-components filtering, attempts to remove small isolated clusters of points. For each point, we count the number of neighbors $N$ present within a distance $D_{max}$. If this number is inferior to a given threshold $N_{min}$, then we consider the point belongs to a small component and is removed. Formally, a point $P$ is removed if:
\begin{align}\label{eq:small_components}
    \#\{~\text{Points } Q~|~\sqrt{(P-Q)(P-Q)^T}\leqslant D_{max}~\}\leqslant N_{min}
\end{align}
where $\#$ is the cardinal of a set. We usually set $D_{max}$ to $3$m and $N_{min}$ to $50$.

\Cref{fig:point_cloud_filtering} illustrates the two filtering methods presented.

\begin{figure}
    \centering
    \includegraphics[width=\linewidth]{Images/Chap_1/Point_cloud_filtering.png}
    \caption{Example of statistical outliers and small components filtering of a 3D point cloud.}
    \label{fig:point_cloud_filtering}
\end{figure}

Filtering the point cloud results in a 3D product that can already be provided as such to users. However, point clouds, while containing more 3D information than \acrshort{dsm}s, can be hard to manipulate in conjunction with other \acrshort{gis} data. Projecting the point cloud on a regular grid to produce a \acrshort{dsm} is thus often preferred, which constitutes the final step of the stereo pipeline.

\subsection{Rasterization}\label{sec:rasterization}
Rasterization consists in projecting the 3D points onto a regular grid over the $(X,Y$) plane to produce the final \acrshort{dsm}. One of the challenges faced when projecting the point cloud is that of the low density of the point cloud relative to the \acrshort{dsm} grid. Indeed, if the density of the point cloud is high enough, multiple 3D points can be projected to the same cell, which raises the question on how to merge their 3D information. On the other hand, if the density is too enough, there may be some cells where no points are projected onto. 

The \acrshort{cars} stereo pipeline uses a Gaussian interpolation to fuse the information of point clouds. Given a cell with coordinates $(X,Y)$ of the \acrshort{dsm}, we consider every point $P_i=(X_i, Y_i, Z_i)$ in a given radius $r$ of $(X,Y)$, and note $\mathcal{N}_{XY}$ the set of all those points. The final value of the \acrshort{dsm} is then computed as the following mean with Gaussian weights:
\begin{align}\label{eq:rasterization}
    \DSM(X,Y) &= \frac{\sum_{P_i\in\mathcal{N}_{XY}}Z_i\cdot e^{-\frac{(X_i-X)^2+(Y_i-Y)^2}{2\sigma^2}}}{\sum_{P_i\in\mathcal{N}_{XY}} e^{-\frac{(X_i-X)^2+(Y_i-Y)^2}{2\sigma^2}}}
\end{align}
with $\sigma$ usually set to $0.3$m and the radius $r$ being $3$m. \Cref{fig:rasterization} illustrates the rasterization process.

Rasterizing with this method provides the advantage of smoothing the potential height variations still present in the 3D point cloud, while allocating more weight to points that are near the cell center. This method can also be found in other stereo pipelines \cite{shean_automated_2016}, while other pipelines use different weighting methods, such as Inverse Distance Weightings (IDW) \cite{rupnik_micmac_2017}, producing similar results.

\begin{figure}
    \centering
    \includegraphics[width=0.7\linewidth]{Images/Chap_1/Rasterization.png}
    \caption{Rasterization of the point clouds on a regular grid. Red points are nearer to the center $(X,Y)$ of the considered cell, while yellow points are further away.}
    \label{fig:rasterization}
\end{figure}

\begin{remark}
    If no points are projected in a given cell (or its direct neighbors), then the cell can be filled with \textit{nodata} values. Different method can be considered to fill holes, such as directly using the values of nearest valid neighboring cells, or interpolating their values. More advanced methods consists in simulating a cloth-like surface fill the holes as in \cite{lallement_bulldozer_2022}.
\end{remark}

As we seen in the previous sections, a photogrammetry pipeline consists in multiple processing steps with intermediary products. For each step of the pipeline, different methods (\eg matching, filtering \etc), parameterization and post-processes are available for each step. This broad range of solutions allow to adapt our processes to the type of images and terrain observed, but it sometimes makes it difficult to determine the configuration producing the best quality \acrshort{dsm}, or to single out a general good-working configuration. The values of the different step parameters presented previously, and that we will be using throughout this thesis, result from testing and sensitivity analyses. In this thesis, we will therefore not discuss their value and instead focus on quantifying and propagating the uncertainty of the \acrshort{cars} stereo pipeline.

\section{Uncertainty in Stereophotogrammetry}\label{sec:previous_work_stereo_uncertainty}
Producing high resolution \acrshort{dsm}s is a complex task, where many uncertainties arise. Those uncertainties can be associated with input data (noise on images, sensor model location uncertainty) or be caused by the processing of those data (resampling errors, disparity computation, rasterization method).

We will first look at attempts to characterized the uncertainty associated with \acrshort{dsm}s.

\subsection{Related Work}\label{sec:related_work}
The quality of \acrshort{dsm} can greatly vary depending on the quality/resolution of the data used, the precision of the geolocation, the performances of the processing algorithm \etc.

From a product point of view, such as TanDEM-X for instance, quality masks can be provided alongside the \acrshort{dem}, usually based on the sensor or method used for producing the \acrshort{dem}. TanDEM-X produces a height error map \cite{wessel_tandem-x_2018}, representing for each DEM pixel the standard deviation corresponding to the height error. The value is derived from the interferometric coherence, and is considered to be a stochastic error. It does not include any contributions of epistemic errors, such as erroneous orbital parameters. The altimetric and planimetric accuracy of by \acrshort{ign}'s \acrshort{lidar} HD 3D points have been measured using ground control points (25cm of planimetric relative accuracy, 5cm of altimetric relative accuracy). Other intrinsic quality parameters of the \acrshort{ign}'s \acrshort{lidar} HD numerical models have not been measured \cite{ign_lidar_2024}. However, \acrshort{ign}'s \acrshort{dtm} product, the RGE ALTI®, possesses a mask indicating the quality of the data provided. This quality mask is composed of the data source used (\acrshort{lidar}, stereo photogrammetry, \acrshort{radar}) for which the mean squared altimetric error has been determined for each source \cite{ign_rge_alti_2013}. In particular, the photogrammetry altimetric accuracy is estimated between 1.5 and 2.5m according to \cite{ign_bd-topo_2013}.

Other work try to estimate the uncertainty associated with a \acrshort{dsm} based not on the method or sensor used, but on the values of the \acrshort{dsm} itself. \cite{mesa-mingorance_accuracy_2020} reviews those different methods for assessing the accuracy of \acrshort{dsm}s. They are mostly producing confidence intervals associated with a \acrshort{dsm} \cite{oksanen_digital_2006,panagiotakis_validation_2018,deschamps-berger_apport_2021}. Those work use robust estimators to evaluate the mean, median and variance of a \acrshort{dsm} altitude distribution \cite{wang_robust_2015}, or advanced statistical methods such as variograms to measure the spatial correlation of errors \cite{hugonnet_uncertainty_2022}. Using those methods, it is possible to produce a confidence interval that globally describes the uncertainty associated to a given \acrshort{dsm} (or a single interval per slope category in \cite{hugonnet_uncertainty_2022}). Those methods estimate the uncertainty \textit{a posteriori}, meaning that they do not consider the means by which the \acrshort{dsm} was created. The uncertainty is estimated using a set of reference points, either extracted from a better resolution \acrshort{dsm}, or from GPS ground control points. In \cite{wang_robust_2015}, the confidence intervals are computed based on the residuals between the \acrshort{dsm} and a bilinear interpolation of samples from the same \acrshort{dsm}, supposing that the quality of a \acrshort{dsm} can be evaluated without using an external reference. 

Photogrammetry pipelines presented in \Cref{sec:different_photogrammetry_pipelines} also have different methods of estimating the uncertainty of the produced \acrshort{dsm}s. For instance, the NASA's ASP pipeline can take as inputs camera planimetric standard deviation of position uncertainty (expressed in meters), and propagate it during the triangulation and rasterization steps. The documentation \cite{nasa_ames_2020} details the covariance matrix propagation method used. It explicitly states that the propagated uncertainty does not represent the error between the predicted height and a hypothetical ground truth, even though they are expressed in the same unit of measure. It is rather the propagated covariance from the camera position projected in the horizontal and vertical directions, regardless of the matching errors or triangulation errors. We did not find clear documentation on how the S2P pipeline handles uncertainties, but the available code of the algorithm suggests that a confidence measure can be computed during the dense matching step of the algorithm (presented later in \Cref{sec:uncertainty_pandora}) and provided alongside the final \acrshort{dsm} as a quality mask. The confidence measure used is a measure of the consensus between the different regularized directions during the \acrshort{sgm} regularization. It is thus an dimensionless parameter, indicating if the predicted height is confident or not.  \acrshort{ign}'s MicMac software computes different errors such as the residual errors from Bundle Adjustment step \cite{ign_micmac_2022}, and provides a correlation map alongside the \acrshort{dsm}. SETSM also computes matching scores \cite{noh_surface_2017}, but it does not seem to be propagated until the final \acrshort{dsm}. DLR's CATENA code is not available, and no documentation indicates the existence of a quality map. In any case, it seems as most pipelines compute a confidence measure (usually during the dense matching step) and propagate it to create a quality mask alongside the final \acrshort{dsm}. This quality mask produced indicate the amount of confidence associated with each \acrshort{dsm} cell value, but does not indicate its potential altimetric error. 

To the best of our knowledge, there is currently no method for producing elevation confidence intervals at a pixel level, by opposition to the global or almost global intervals in \cite{oksanen_digital_2006,wang_robust_2015,hugonnet_uncertainty_2022}. Furthermore, the uncertainty arising from the different photogrammetry algorithms are not considered when producing those intervals. The main contribution of this thesis is the production of pixel-wise confidence intervals for \acrshort{dsm}s produced by photogrammetry. In the following section, we present different sources of uncertainty present in the \acrshort{cars} pipeline. 

\subsection{Uncertainty in the CARS pipeline}\label{sec:uncertainty_cars}
In this thesis, we focused on quantifying the uncertainty alongside the creation process of \acrshort{dsm}. We thus differ in this regard from previous work, which produces a single confidence intervals from the final produced \acrshort{dsm}. Multiple sources of uncertainty influence the production of a \acrshort{dsm}, which we will now detail.

A first source of uncertainty comes from the input data, \ie the radiometric uncertainty associated with stereo images and the geometric uncertainty associated with sensors models. The radiometric uncertainty is mainly due to noise from the sensor and from atmospheric effects appearing on the acquisitions. Those radiometric errors can be estimated and are relatively moderate \cite{jacobsen_radiometric_2014}. Regarding the uncertainty associated with the sensor models, satellite movements not captured by models can have a big impact on the final images, especially for push-broom sensors. For instance, vibration of the satellite can lead to biases on the geolocation of the different rows of the image. Those biases will themselves be propagated to the final \acrshort{dsm}, leading to errors of a few meters \cite{loghin_accuracy_2019}. Additionally \acrshort{rpc} models attempt to represent real lines of sight with polynomial coefficient, which is not exact and possess its own accuracy. The method used for coefficients calibration and the frequency of calibration also bring their share of uncertainty. 

Another source of uncertainty arises from the different processing steps of the stereo pipeline. First, different resamplings occur in order to convert stereo images from sensor geometry to epipolar geometry one. Those resamplings are based on \acrshort{sift} matches, presented in \Cref{sec:epipolar_geometry}. \acrshort{sift} matches are used to refine the epipolar grid depend on the performance of the \acrshort{sift} algorithm used to obtain them. In homogeneous and texture-less area, \eg glaciers, false matches can arise. This leads to wrong epipolar lines, themselves preventing correct results in the dense matching step and leading to a false 3D reconstruction.  We assume in this thesis that the refinement of epipolar grids using \acrshort{sift} points, which minimizes the mean error, provides accurate epipolar resampling grids with an epipolar error of around a tenth of a pixel \cite{franchis_automatic_2014}. However, the resampling in epipolar geometry introduce errors if the input and target resolutions do not respect Shannon criteria \cite{delon_small_2007}. In the case of Pléiades images, the acquisition resolution is usually high enough to insure correct resampling \cite{jacobsen_radiometric_2014}. 

The resampled images in epipolar geometry are then used as inputs for the dense matching step. Dense stereo matching is a complex task, for which many different algorithms exist, each potentially presenting different performances. When compared to epipolar errors, that typically are less than a pixel \cite{franchis_automatic_2014}, dense matching errors can potentially reach higher magnitudes, depending on the size of the considered disparity range (sometimes reaching hundreds of pixels). Considering those orders of magnitude, estimating and quantifying the uncertainty of the stereo matching process is therefore crucial to control the uncertainty on the output \acrshort{dsm}. Some of the sources of those errors are now presented. Because the correlator usually compares windows (in our case $5\times5$ windows using the CENSUS cost function or $11\times11$ windows using MC-CNN) and not single pixels, this can create an adherence effect near building borders, as in \Cref{fig:adherence_window}. Note that in \cite{okutomi_stereo_1994}, authors have been trying to adapt the window shape to reduce the uncertainty of stereo matching, but this method requires to iteratively compute costs on different windows, which can become quite expensive. Using a window-based correlator with \acrshort{sgm} regularization, if correctly parameterized, usually presents good performance in areas without height discontinuities. This can become a problem in urban areas, where the presence of high buildings represents an additional challenge for the correlator. As \acrshort{sgm} regularization also penalizes disparity changes, tops of buildings tend to be extended beyond their true footprint. Other processes, such as filtering or sub-pixel refinement, improve the quality of the disparity map but require specific care for handling their induced uncertainty. \Cref{sec:uncertainty_pandora} dives more into details regarding the different methods that have been developed to quantify the uncertainty of dense stereo matching.

Once the disparity has been estimated, 3D points can be triangulated by intersecting \acrshort{rpc} lines from each matched pixel. However, we saw previously that there is no guarantee that the 3D lines do intersect. If they indeed do not intersect, the 3D point is defined as the point minimizing its squared distances to both lines of sight. An alternative is to modify the geolocation line from the secondary image so that it intersects the line from the reference image. In both cases, the localization of the 3D point is not exact. This uncertainty stems from the fact that \acrshort{rpc} have a their own limited resolution, and that stereo images (resampled or not) do not necessarily point to the exact same location in the object space.

The final part of the stereo pipeline is to rasterize the point cloud onto a regular grid, thus yielding the \acrshort{dsm}. When characterizing the uncertainty on the final result, we must first agree on what the \acrshort{dsm} is assumed to represent. It is common to consider that each pixel's value should represent the average height over the cell. However, providing the maximum or minimal height might be more adapted to some scenarios: for instance if the \acrshort{dsm} is used to prepare very low altitude flights (for drones \etc), the maximum height is more relevant as one would want to avoid any foliage or power line. Those elements could disappear in the final \acrshort{dsm} if it represented the mean or median elevation. In this thesis, we will consider that the \acrshort{dsm} represents the (weighted) average height. Depending of the resolution of input images and the desired output resolution of the \acrshort{dsm}, the density of points per \acrshort{dsm} cell will vary. In our applications, the input and output resolutions are identical as we typically desire to produce a \acrshort{dsm} at $50$cm resolution from $50$cm panchromatic images. This means that there is on average one 3D point per \acrshort{dsm} cell. For occluded regions, or when we discarded stereo matches that seemed incorrect, there might be no point directly in the output cell. In this case, the value of the \acrshort{dsm} cell will be completely determined by the value of points in neighboring cells, even if their distance is high and the averaging weights are small. Interpreting the final \acrshort{dsm} as the average height on each cell is thus debatable, as the average is computed on a limited number of points, and sometimes not even belonging to the considered cell. Note that if there is no points around in a given radius, the cell will be left empty. 

Different sources of uncertainty occurring throughout the stereo pipeline were presented in the previous paragraphs. Characterizing, modelling and propagating all of those uncertainties could not be considered in the span of this thesis. We thus focus mostly on the uncertainty arising from the dense stereo matching step, as it is the source of the biggest errors in the pipeline. \Cref{chap:propagating} investigates how uncertainty from the input epipolar images can be propagated in the stereo matching step, and \Cref{chap:epistemic_uncertainty} attempts to model the processing uncertainty of stereo algorithm itself. We also propagate this uncertainty all the way to the output \acrshort{dsm} and show that it can correctly estimate the errors made during the \acrshort{dsm} production.

\subsection{Uncertainty Quantification in Dense Stereo Matching}\label{sec:uncertainty_pandora}
As stereo matching is a popular problem in computer vision, many methods for quantifying its uncertainty have been proposed in the literature. Without being exhaustive, this section presents a quick overview of the main approaches as well as the solutions currently implemented in the dense stereo correlator Pandora used in the \acrshort{cars} pipeline.

The way uncertainty is quantified in stereo matching is mainly done by producing \textit{confidence maps}, \ie a mapping for each pixel $(\rowcol)$ to a real confidence value, usually between $0$ and $1$. By convention, a value of $0$ means that we are not confident in the disparity value associated to $(\rowcol)$. Conversely, a confidence value of $1$ indicates that we are very confidence in the predicted disparity.

There are multiple sources of information that can be used to compute a confidence measure. Left and right input images and the predicted disparity map being those available to every method, but cost-based approaches can also make use of the cost volume, as it contains a quantity of useful information. Most confidence measures are handcrafted using those different information sources. For reviews on those methods, we refer to \cite{egnal_stereo_2004, hu_quantitative_2012, poggi_quantitative_2017}. With the rising use of deep-learning in stereo, many networks have also been developed to estimate the uncertainty. A review for methods using regression forests can be found in \cite{min-gyu_park_leveraging_2015}, and a more general review, including the use of 2D and 3D CNNs on the cost volume, can be found in \cite{poggi_confidence_2021}.

\begin{remark}
     Quantifying the uncertainty in stereo matching is a popular field of research. Recently, people have even been trying to evaluate the uncertainty of the confidence estimation itself, called \textit{meta-confidence} \cite{kim_meta-confidence_2022}.
\end{remark}

In \Cref{ex:confidence_measure}, we present some examples of confidence measures that use different sources of information.

\begin{example}\label{ex:confidence_measure}
    Regarding confidence measure based on the disparity map, we already presented a type of binary confidence measure in \Cref{sec:postprocess_disparity} with the cross-checking test from \cref{eq:cross-checking}. Other methods compute for instance the local variance of the disparity map, where a low variance suggests confident regions.
    
    For methods using the cost volume, a simple method for measuring the cost would be using the value of the matching score measure (MSM, \cite{egnal_stereo_2004}) for a given disparity at coordinates $(\rowcol)$:
    \begin{equation}
        MSM(\rowcol) = -C_V(row, col, \mathcal{D}(\rowcol))
    \end{equation}
    This measure can be normalized between $0$ and $1$ using the global minimum and maximum of the cost volume. The idea behind this measure is the following: a high matching cost for a selected disparity indicates that the two matched pixels are not that similar, and thus the match is not confident. Other measures using the matching cost compare the value of the first and second minimum of a cost curve, or measure the curvature of the cost curves. More advances measures make use of 3D CNNs on the entirety of the cost volume to learn an efficient confidence measure \cite{mehltretter_cnn-based_2019}.
    
    Measures based on the input images usually measure the gradient \cite{haeusler_ensemble_2013} or variance \cite{park_learning_2019} of input images. High gradient or high variances indicate highly-textured regions, which are often easier to match. The confidence is therefore higher for those pixels. 
    
    Deep-learning approaches can combine multiple sources of information (input images, cost volume, disparity map) to learn a confidence measure \cite{tosi_beyond_2018, kim_adversarial_2020}.
\end{example}


\begin{figure}
    \centering
    \begin{subfigure}[t]{0.5\linewidth}
        \centering
        \includegraphics[width=\linewidth]{Images/Chap_1/Ambiguity.png}
        \caption{}
        \label{fig:ambgiuity}
    \end{subfigure}\hfill
    \begin{subfigure}[t]{0.5\linewidth}
        \centering
        \includegraphics[width=\linewidth]{Images/Chap_1/Risk.png}
        \caption{}
        \label{fig:risk}
    \end{subfigure}\hfill
    \caption{Cost curves with different levels of ambiguity \ref{fig:ambgiuity} and risk \ref{fig:risk}}
    \label{fig:ambiguity_risk}
\end{figure}

In this thesis, we will also consider another confidence measure that can already be computed with the Pandora correlator. This measure will help us to quantify the uncertainty in the photogrammetry pipeline. This confidence measure is a confidence measure called \textit{confidence from ambiguity} \cite{sarrazin_ambiguity_2021}. This method is based on the cost volume and quantifies how easy or hard it is to single out the correct disparity in a cost curve. \Cref{fig:ambgiuity} presents two cost curves, one which is not ambiguous as there is a single well defined minimum, and an ambiguous cost curve which presents multiple values that are close to the minimum. To compute it, we first start by constructing what is called an \textit{ambiguity curve}. We do that by counting how many disparities have a cost close within a threshold $\eta$ to the minimal cost, for increasing values of $\eta$. Formally we need to define for a pixel $(\rowcol)$ the set of all disparities whose cost is within a range $\eta$ of the minimal cost, and then define the ambiguity curve as the cardinal of this set:
\begin{align}
    &\mathcal{D}_\eta = \{d ~|~ C_V(row,~col,~d) \leqslant \min_\delta C_V(\rowcol, ~\delta) + \eta\}\\
    &Amb(\rowcol, ~\eta) = \#\mathcal{D}_\eta
\end{align}
where $\#$ is the cardinal of a set. Evaluating $Amb$ for different $\eta$ gives the ambiguity curve. \Cref{fig:integral_ambiguity_1} presents a cost curve with different values of $\eta$. \Cref{fig:integral_ambiguity_2} presents the resulting ambiguity curve. On those two figures we can see for instance that for $\eta_3$ there are $6$ disparities whose cost is lower than $\min_\delta C_V(\rowcol, ~\delta) + \eta_3$. For non-ambiguous cost curves, $Amb$ will increase only for high values of $\eta$. On the contrary for ambiguous curves, $Amb$ will be high for small values of $\eta$. To obtain a scalar value from $Amb$, we compute the area under its curve, normalized by the range of $\eta$:
\begin{equation}
    \mathrm{AUC}_{Amb}(\rowcol) = \frac{1}{\max\eta-\min\eta}\int_\eta Amb(row,~col,~\eta)d\eta
\end{equation}
It results on low values for confident (non-ambiguous) cost curves, and high values for less confident (ambiguous) curves. Because confidence measures actually present high values for confident pixels, $\mathrm{AUC}_{Amb}$ is normalized and inverted in order to obtain the confidence from ambiguity $c_{Amb}$:
\begin{equation}\label{eq:confidence_from_ambiguity}
    c_{Amb}(\rowcol) = \frac{\max \mathrm{AUC}_{Amb}- \mathrm{AUC}_{Amb}(\rowcol)}{\max \mathrm{AUC}_{Amb} -\min \mathrm{AUC}_{Amb}}
\end{equation}
This way, values of $c_{Amb}$ near $0$ indicate that we are not confident in the predicted disparity. Reversely, values of $c_{Amb}$ near $1$ indicate that we are confident in the predicted disparity.

\begin{figure}
    \centering
    \begin{subfigure}[t]{0.6\linewidth}
        \centering
        \includegraphics[width=\linewidth]{Images/Chap_1/Integral_Ambiguity_1.png}
        \caption{Cost curve with different values of $\eta$. Horizontal dotted lines indicates the range of costs between $\min_\delta C_V(\rowcol, ~\delta)$ and $\min_\delta C_V(\rowcol, ~\delta)+\eta_3$}
        \label{fig:integral_ambiguity_1}
    \end{subfigure}\hfill
    \begin{subfigure}[t]{0.4\linewidth}
        \centering
        \includegraphics[width=\linewidth]{Images/Chap_1/Integral_Ambiguity_2.png}
        \caption{Associated ambiguity curve in full line. Ambiguous curve in gray dotted line}
        \label{fig:integral_ambiguity_2}
    \end{subfigure}\hfill
    \caption{Illustration of the computation of the ambiguity curve.}
    \label{fig:integral_ambiguity}
\end{figure}

\begin{comment}
    The second confidence measure that we consider is named \textit{confidence from risk}. It is designed to measure if the magnitude of a potential error is high or low. After all, if we predict a wrong disparity while the correct disparity is right next to our prediction, the impact on the final result will be smaller than if the true disparity is at the other side of the disparity range. Keeping this in mind, we apply the same methodology as for the ambiguity. We consider the set $\mathcal{D}_\eta$ of all disparities whose cost is less than $\eta$ away from the minimal cost, and define the risk as the gap between extrema of this set:
    \begin{equation}
        R(\rowcol, ~\eta) = \max_{\mathcal{D}_\eta}d - \min_{\mathcal{D}_\eta}d
    \end{equation}
    \Cref{fig:integral_risk_1} presents a cost curve with different values of $\eta$ and their associated risk $R$. \Cref{fig:integral_risk_2} presents the resulting risk for all $\eta$. For non-risky cost curves, $R$ will increase only for high values of $\eta$. On the contrary, for risky curves, $R$ will be high for small values of $\eta$. To obtain a single scalar, we compute the area under the risk curve, normalized by the range of $\eta$:
    \begin{equation}
        \mathrm{AUC}_R = \frac{1}{\max\eta-\min\eta}\int_\eta R(row,~col,~\eta)d\eta
    \end{equation}
    $\mathrm{AUC}_R$ is thus expressed in number of disparities. Contrary to the confidence from ambiguity, the risk is kept as such with its current unity. 
    
    \begin{remark}
         The risk measures the magnitude of the potential error, but not its probability. It is therefore possible to have a very risky cost curve, but which is not ambiguous at all. In other words, it means that the probability of an error occurring can be very low, but if it somehow happened, then the magnitude of this error would also be very low. The risk is thus a confidence measure that need to be completed by a measure such as ambiguity. On it own, it has less meaning than other measures.
    \end{remark}
    
    \begin{figure}
        \centering
        \begin{subfigure}[t]{0.6\linewidth}
            \centering
            \includegraphics[width=\linewidth]{Images/Chap_1/Integral_Risk_1.png}
            \caption{Cost curve with different values of $\eta$ and risk}
            \label{fig:integral_risk_1}
        \end{subfigure}\hfill
        \begin{subfigure}[t]{0.4\linewidth}
            \centering
            \includegraphics[width=\linewidth]{Images/Chap_1/Integral_Risk_2.png}
            \caption{Associated risk curve in full line. Risky curve in gray dotted line}
            \label{fig:integral_risk_2}
        \end{subfigure}\hfill
        \caption{Illustration of the computation of the risk curve.}
        \label{fig:integral_risk}
    \end{figure}
\end{comment}

The confidence from ambiguity, like other confidence measures, is designed to indicate potentially wrong matches. It however does not measure the extent of the potential error. After all, it is possible that we predict a wrong disparity, while the correct disparity is right next to our prediction. The impact of this error on the final result will be smaller than if the true disparity is at the other side of the disparity range. There is the a distinction to be made between the value of a confidence measures and the magnitude of a potential error, even if they can be correlated in some cases.

In \Cref{chap:epistemic_uncertainty,chap:elevation_intervals}, we introduce a method for computing confidence intervals, which aims to serve as a complement to confidence measures as it estimates the magnitude of the potential error. This method was motivated following discussions with different users and experts in 3D modelling from the AI4GEO consortium (\url{https://www.ai4geo.eu/}). Indeed, a common point that emerged from those discussions was the desire to produce confidence intervals alongside photogrammetry \acrshort{dsm}s. This also converged with the quality map requirement of the \acrshort{co3d} mission. We therefore developed our method with the objective of answering to this user need. This chapter presented the different concepts relative to satellite photogrammetry, and the different uncertainty associated with a photogrammetry pipeline. In the following, \Cref{chap:representation_of_uncertainty,chap:joining_credal_sets} will present relevant uncertainty models and concepts for this thesis, while \Cref{chap:propagating,chap:epistemic_uncertainty,chap:elevation_intervals} estimate the uncertainty at different steps of the photogrammetry pipeline.

\pagebreak
\blankpage

    
    \setkeys{Gin}{draft=false}
    \chapter{Mathematical Representations of Uncertainty}\label{chap:representation_of_uncertainty}

\section{Introduction}
This chapter takes some distance from the stereo vision problem presented previously, and instead describes more formally different representations of uncertainty and the tools that can be used to manipulate them. In this thesis, we consider classical probabilities (\Cref{sec:probabilities}), possibility distributions (\Cref{sec:possibilities}) and p-boxes (\Cref{sec:pboxes}) to model uncertainty. We also consider the case where we have uncertainty on multiple variables simultaneously. In this setting, the dependency between the different sources must be taken into account. We thus present a dependency model called copulas in \Cref{sec:copulas}, and some of its properties. The thorough presentation of those concepts will then be put into relations in \Cref{chap:joining_credal_sets}. 

\section{Notations}
We introduce here a few notations that will be used in this chapter.
\begin{itemize}
    \item \st means ``such that''
    \item We will often present theorems or results with $n$ different variables, spaces, or Cartesian products of $n$ elements. Notations therefore quickly become quite heavy. For that reason, we often use the notation ``$\ldots$'' to imply that we are enumerating all variables. For instance, if we apply a function $f$ to four variables $x_1, x_2, x_3, x_4$, we will write $f(x_1\enum x_4)$.
    \item When working with $n$ variables, the index $i$ will usually be used to refer to the $i$-th variable (or one of its attribute), and will be appended as a subscript when possible, otherwise as a superscript. For instance, if $x\in\mathbb{R}^n$, then $x_i$ will refer to the $i$-th component of $x$. If $m_\times\in\mathbb{R}^n$, then $m_\times^i$ will refer to $i$-th component of $m_\times$.
    \item $\opi\cdot,\cdot\cli$ refers to an interval of integers. For instance $\opi1,~4\cli=\{1,~2,~3,~4\}$.
    \item The power set of a set $\X$ is noted $2^\X$. It corresponds to the set of all sets included in $\X$. In the discrete case, if the cardinal of $\X$ is $n$, then the cardinal of the power set is $2^n$, thus the notation.
\end{itemize}

\section{Different Models to Represent Uncertainty}\label{sec:different_models_of_uncertainty}
Assessing the reliability of an engineering system requires quantifying the uncertainty of the input parameters or of the system itself, using models of uncertainty. Modeling the uncertainty can be done in various ways, depending on the type of uncertainty considered and the available measures or \textit{a priori} regarding the uncertain sources. Most common models are probability distributions, which have been studied extensively. When using those models, we know —or assume— that the information we try to estimate is of stochastic (or random) nature, and that we are able to precisely describe its structure using a probability distribution. However, there are many cases where such assumptions cannot be made: for instance when data is insufficient to determine the correct probability distribution, or when the uncertainty is not random, but epistemic. In those cases, we can use other models, such as:
\begin{itemize}
    \item fuzzy sets \cite{zadeh_fuzzy_1999} when trying to estimate the degree of truth of a statement such as ``This person is tall''
    \item intervals \cite{jaulin_applied_2001}, where no preferences are given inside a specific range of possible values
    \item \acrfull{ip} which tries to extend the concept of probabilities in order to model epistemic uncertainty \cite{augustin_introduction_2014}
\end{itemize}
This list is not exhaustive. Additionally, those different models can sometimes be equivalent. Choosing to use one or another depends on the nature of the problem and of the available data. In this chapter, we will mainly consider probability distributions (\Cref{sec:probabilities}) and \acrshort{ip} (\Cref{sec:imprecise_probabilities}). Two specific cases of \acrshort{ip} will be detailed, namely possibility distributions in \Cref{sec:possibilities} and p-boxes in \Cref{sec:pboxes}.

\begin{remark}
    Distinguishing between stochastic and epistemic uncertainty is a model accepted by many to distinguish between different situations of uncertainty. One could argue that stochastic uncertainty is, to some extent, equivalent to epistemic uncertainty. Indeed, if we had enough knowledge on initial conditions of a die throw for instance (force and torque applied on the dice, its exact shape and mass distribution \etc), as well as the exact physics model, one could predict with certainty on which side it would land. There would therefore be no aleatoric process at stake here. The question whether or not we should make the distinction between aleatoric and epistemic uncertainty is of interest regarding theoretical aspects of the nature of uncertainty. For real-life applications, however, differentiating between the two seems reasonable, as we cannot know every parameter and the exact model involved for every quantity of interest. 
\end{remark}

At the end of the day, choosing one model over another is not always straightforward. It requires being aware of the type of uncertainty faced, of the strengths and limitations of each model, of the tools available to manipulate the models \etc When it comes to less common models cited above, it fundamentally requires being aware of the existence of such models, which is not always the case for non-specialists. During this thesis, we tried to promote less common models, especially possibility distributions. We did so by presenting real life cases where they could be used while improving uncertainty modeling in the field of stereophotogrammetry (see \Cref{chap:propagating,chap:epistemic_uncertainty}).

\subsection{Probabilities}\label{sec:probabilities}
Probability measures are a classical framework to represent uncertainty. There are multiple ways of interpreting them, mainly with a \textit{frequentist} approach, or a \textit{Bayesian} approach. From a frequentist point of view, probabilities are well-fitted to represent stochastic uncertainty, \ie uncertainty regarding events that can get a different result each time we run an experiment or acquire a measure (typically, noise on a sensor). From a \textit{Bayesian} point of view, probabilities represent a state of knowledge or degree of belief, and can be updated with additional information \cite{irvine_philosophies_2009}. This leads to the notion of prior and posterior probability that will not be considered in this thesis.

We remind here basic definitions regarding probability distributions.

\begin{definition}[Probability Space]\label{def:probability_space}
    We call a probability space $(\X,~\mathcal{A},~P)$ a tuple where:
    \begin{itemize}
        \item $\X$ is the set of possible outcomes (for instance, head or tails for a toss coin)
        \item $\mathcal{A}$ is the set of all subsets of $\X$ for which a probability can be measured (for instance $\{\emptyset, \{\text{heads}\}, \{\text{tails}\}, \{\text{heads, tails}\}\}$)
        \item $P$ is the probability measure assigning a probability to each of the sets of $\mathcal{A}$. For instance, for a fair coin, $P(\emptyset)=0$, $P(\{\text{heads}\})=P(\{\text{tails}\})=0.5$ and $P(\X)=1$.
    \end{itemize}
    Note that $\mathcal{A}$ must be a $\sigma$-algebra, meaning that it is closed under complement, countable unions and countable intersections. $P$ is a probability measure if it verifies all the Kolmogorov axioms:
    \begin{itemize}
        \item $\forall A\in\mathcal{A},~P(A)\in[0,1]$
        \item $P(\X) = 1$
        \item for any countable disjoint family of sets $A_i\in\mathcal{A}$, $P(\cup_i A_i)=\sum_i~P(A_i)$
    \end{itemize}
\end{definition}

\begin{definition}[Random Variable]
    A random variable $X$ is a measurable function from $\X$ to a measurable space (which is $\mathbb{R}$ or a subset of $\mathbb{R}$). We can then measure the probability that $X$ takes its values in a measurable set $E\subseteq\mathbb{R}$:
    \begin{align*}
        P(E) = P(\{x\in\X~|~X(x)\in E\})
    \end{align*}
\end{definition}

Using a random variable allows considering the probability measure on different spaces. We then call $P$ the probability distribution of the considered random variable $X$.

Other useful concepts regarding probability distributions are \acrfull{pdf} and \acrfull{cdf}:
\begin{definition}[Cumulative Distribution Function]\label{def:cdf}
    A \acrfull{cdf} $F_X$ of a random variable $X$ with real values is the probability that $X$ will be less or equal to a number $x$. Formally, we define $F_X:\X\rightarrow[0,1]$:
    \begin{equation*}
        \forall x\in\X,~F_X(x)=P(X\leqslant x)
    \end{equation*}
\end{definition}

\begin{definition}[Density Function]\label{def:density}
    A random variable $X$ is said to possess a \acrfull{pdf} if there exists a positive integrable function $f$ over $\mathbb{R}$ \st $\forall (x_1, x_2)\in\mathbb{R}^2$:
    \begin{equation*}
        P(x_1\leqslant X \leqslant x_2) = \int_{x_1}^{x_2}f(x)dx
    \end{equation*}
    In the continuous case, the \acrshort{pdf} $f$ is also the derivative of the \acrshort{cdf} of $X$. In the discrete case, the \acrshort{pdf} is also called \textit{probability mass function}, and is defined as:
    \begin{equation*}
        f(x_1) = P(X = x_1)
    \end{equation*}
\end{definition}

In the rest of the chapter, we consider random variables $X$ on discrete spaces $\X$. If $\X=\{x_1,~\ldots,~x_n\}$, we call $\{x_i\}$ an atom of $\X$, and the probability distribution $P$ of $X$ on $\X$ is completely determined by its \acrshort{pdf}, which is the value of $P$ on atoms. For simplicity of notation, we will not always use braces around atoms when computing their probability. So we will sometimes write $P(x_1)$ instead of $P(X=\{x_1\})$.

As stated previously, probability measures are fitted to represent stochastic uncertainty. The following example illustrates why probability measures are not adapted to represent epistemic uncertainty:

\begin{example}\label{ex:proba_limitations}
    Let consider a card facing down, with a number written on its hidden side. The person who wrote the number tells you that they chose to write either $1$, $2$ or $3$ on it. We should note that because they chose to write a number, the uncertainty on its value is not random. They then ask you to evaluate your chances of guessing the correct number and its parity, \ie if it is odd or even. We first consider the random variable $X$ taking values in $\{1,~2,~3\}$. Because you have no further information and thus no preferences on the values of $X$, a common (yet arguably inadequate) decision is to associate the uniform distribution $P$ to $X$:
    \begin{equation*}
        P(X=1)=P(X=2)=P(X=3)=\frac{1}{3}
    \end{equation*}
    Concerning the parity of the number, we may consider the random variable $Y$ defined such that $Y=0$ if ``$X=1$ or $X=3$'' and $Y=1$ if ``$X=2$''. Because we have no information on the number written on the card, we also do not have any preferences on the values of $Y$. Following the same reasoning as before, one might be tempted to associate a uniform distribution to it. However, deducing the \acrshort{pdf} of $Y$ from that of $X$ yields:
    \begin{equation*}
        P(Y=0)=P(X=1)+P(X=3)=\frac{2}{3},\qquad P(Y=1)=P(X=2) = \frac{1}{3}
    \end{equation*}
    which is clearly not the \acrshort{pdf} of a uniform distribution. By assuming we have no preferences on the values of a variable, we actually deduced preferences on the values of another variable. 
    
    This example shows that one should use a uniform distribution only when they are certain that all values are equiprobable. This is due to the fact that uniform distributions are well suited for representing statements like ``all values have the same likelihood'' but not for statements like ``I have no information over the values, and thus no preference''. Indeed, a probability distribution actually contains a lot of information about a (random) variable, which is not suited to represent epistemic uncertainty.
    
    \begin{remark}
        From a Bayesian point of view, a probability can also represent a degree of belief, allowing them to represent epistemic uncertainty and not only stochastic uncertainty, in theory. However, it does not solve the expressiveness problem raised by this example. Bayesians are still reasoning with probabilities, which make no difference between ``I have no preference between these events'' and ``these events are equiprobable''. Even though they can update their prior with additional information, this would not fix the problem presented here. In the absence of additional information, basing a decision on the prior would lead to debatable conclusions.
    \end{remark}
\end{example}

\subsection{Imprecise Probabilities}\label{sec:imprecise_probabilities}
As highlighted in \Cref{ex:proba_limitations}, uncertainty cannot always be correctly modeled by probabilities, especially in a context where data is sparse. To overcome this problem, a generalization of probabilities has been introduced, called \acrfull{ip}, which provides a general framework for working with both aleatoric and epistemic uncertainty. It uses the concept of lower and upper probabilities, which are quite generic and flexible, and can be derived in more specific models. Here is a brief scope of the relevant tools it encompasses: the special case of \textit{belief functions}, themselves containing specific sub-categories such as possibility distributions (\Cref{sec:possibilities}) and probability boxes (\Cref{sec:pboxes}). It also contains probabilities presented in \Cref{sec:probabilities}, which we will call \textit{precise} probabilities in contrast to \textit{imprecise} probabilities. \Cref{fig:diagram_IP} is a non-exhaustive overview of relationships and specificities of each imprecise model.

\begin{remark}
    On a more generic level, \acrshort{ip} can be described by sets of acceptable gambles and by lower and upper expectations \cite{walley_statistical_1991,augustin_introduction_2014}. Although very interesting, we did not use them in our applications and thus do not consider them in this thesis.
\end{remark}

\begin{figure}[ht]
    {\centering
    \includegraphics[width=0.8\linewidth]{Images/Chap_2/Diagramme_IP_Bel.png}
    \caption{Diagram representing the relationship between different \acrshort{ip} models presented throughout \Cref{sec:different_models_of_uncertainty}.}
    \label{fig:diagram_IP}}
\end{figure}

As stated previously, a core concept of \acrshort{ip} is the concept of lower and upper probabilities. Similarly to precise probabilities, a lower probability $\low$ and an upper probability $\overline{P}$ are mappings from a $\sigma$-algebra $\mathcal{A}$ to $[0, 1]$. However, while a probability $P$ gives a single measure of uncertainty for every event, lower and upper probabilities provide two bounds for every event, allowing them to express more complex uncertainty structures.

\begin{remark}
    Formally, a lower probability $\low$ needs to be \textit{super-additive}, \ie to satisfy:
    \begin{equation}
        \forall A,B\in\mathcal{A}, \text{ if } A\cap B\neq\emptyset, ~\low(A\cup B)\geqslant \low(A)+\low(B) \label{eq:super_additivity}
    \end{equation}
    Conversely, an upper probability is sub-additive, meaning that it verifies the same property as \Cref{eq:super_additivity} but with the inequality reversed.
    
    Those properties are less constraining than their equivalent for precise probabilities in \Cref{def:probability_space}, so lower and upper probabilities are generally \textit{not} precise probabilities. The only case where they are precise probabilities is when $\low=\overline{P}$, because $\low$ is then additive as it is both super and sub-additive. In this case, \acrshort{ip} are actually a single precise probability. This illustrates the fact that precise probabilities are special cases of \acrshort{ip}. 
\end{remark}

Precise probabilities define a measure of uncertainty towards a random variable $X$ taking numerical values. \acrshort{ip} however, can model the uncertainty towards a random variable taking set values instead of numerical one. We then say that $X$ is a \textit{random set} instead of a random variable. That being said, \acrshort{ip} can also represent the uncertainty of random variables for which we assume a precise probability exists, but we are not able to determine it precisely. Indeed, lower and upper probabilities form the bounds of a family of precise probabilities called \textit{credal set}.
\begin{definition}[Credal set]\label{def:credal_set}
    Given a lower probability $\low$ and an upper probability $\overline{P}$, a credal set $\M$ is the set of all probabilities $P$ that are greater that $\low$ and lower that $\overline{P}$:
    \begin{equation}
        \M(\low, \overline{P}) = \{P~|~\forall A\in\mathcal{A},~\low(A)~\leqslant~P(A)~\leqslant\overline{P}(A)\}
    \end{equation}
\end{definition}
We refer to $\M(\low, \overline{P})$ as $\M$ when no confusion is possible. Credal sets allow considering multiple probabilities at once, which improves on the limited expressiveness of a single probability measure. The gap between the two bounds of a credal set reflects how imprecise is the model, in terms of epistemic uncertainty.

Conversely, we can define lower and upper bounds from a set of probabilities $\M$ as:
\begin{align*}
    \forall A\in\mathcal{A},~\low(A) =& \inf_{P\in\M}P(A)\\
    \forall A\in\mathcal{A},~\overline{P}(A) =& \sup_{P\in\M}P(A)
\end{align*}
Note that $\M$, contrary to $\M(\low, \overline{P})$, is not necessarily defined by $\low, \overline{P}$. Therefore, even though $\M$ and $\M(\low,\overline{P})$ have the same bounds on events, they are not necessarily equal and only $\M\subseteq\M(\low, \overline{P})$ holds. 

Although it is not required, we usually assume that credal sets verify additional properties expressed as follows:
\begin{definition}[Coherence and Avoiding sure loss]\label{def:coherence_sure_loss}
    A credal set $\M$ is said to avoid sure loss if it contains at least one probability measure, \ie if
    \begin{equation}
        \M \neq \emptyset\label{eq:avoid_sure_loss}
    \end{equation}
    
    In this thesis, we consider that a credal set $\M$ induces \textit{coherent} lower and upper probabilities if its lower and upper bounds on all events are attained by a probability measure, \ie if:
    \begin{equation}
        \forall A\in\mathcal{A}, \exists P, P'\in\M~\st~\low(A)=P(A) \text{ and } \overline{P}(A)=P'(A) 
    \end{equation}
\end{definition}
The bounds $\low, \overline{P}$ of a coherent credal set verify the following property:
\begin{equation}
    \forall A\in\mathcal{A}, ~\low(A) = 1 - \overline{P}(A^c)\label{eq:lower_proba_complement}
\end{equation}
which is a generalization of the classical property for computing the complement of an event with precise probabilities. This allows us to only specify the lower bound $\low$ of a credal set to completely describe it, as its upper bound is determined by $\low$ through complementation. Defining a credal set requires specifying much more constraints on the probability space than in the case of probability distributions. Indeed, a lower probability must be defined on every possible event, while a precise probability can be defined by its \acrshort{pdf} on possible values only (instead of possible events). In the case of a discrete space with $n$ elements, a probability is completely determined by its values on the $n$ atoms, while a lower bound is completely determined by its values on the $2^n$ considered events.


\begin{remark}
    Sampling from a credal set is not straightforward. Multiple methods exist, the most intuitive consisting in sampling distributions from the credal set extreme points.
\end{remark}

\begin{remark}
    With the way we defined a credal set $\M$, it is closed and convex. This is a common way of constructing credal sets, but is not the only way. We could for instance impose that all probabilities in $\M$ belong to a family of Gaussian probabilities, which would prevent $\M$ from being convex.
\end{remark}

When constructing credal sets, it is common to possess a non-convex set of probabilities $S$. In that case, we can define the convex hull of a set of probabilities $S$:
\begin{definition}[Convex Hull]\label{def:convex_hull}
    The convex hull ($CH$) of a set of probabilities $S$ is the smallest (convex) credal set containing $S$.
\end{definition}

\begin{remark}
    The convex hull is computationally heavy to evaluate. Determining the infimum and supremum of a set of probability is not trivial, as it might be required to iterate through every element of $S$. It is however very useful, especially in the case where we do not know the set bounds on every event. In that case, computing the convex hull also allows us to determine the bounds on those events.
\end{remark}

\begin{example}
    Let us consider the scenario presented in \Cref{ex:proba_limitations} and model the uncertainty with a credal set instead of a single probability distribution. Because we have no information on the value written on the card except that it is in $\{1,2,3\}$, we characterize the uncertainty by lower and upper probability $\low$, $\overline{P}$ defined as follows:
    \begin{align*}
        &\low(X=1) = \low(X=2)= \low(X=3) = 0\\
        &\overline{P}(X=1) = \overline{P}(X=2)= \overline{P}(X=3) = 1\\
        &\low(X\in\{1,2,3\}) = \overline{P}(X\in\{1,2,3\}) = 1
    \end{align*}
    We can use \Cref{eq:lower_proba_complement} for computing the bounds of remaining events.
    
    Here, the credal set is the largest credal set possible, as its bounds are always $0$ and $1$ for events that are not $\emptyset$ or $\{1,2,3\}$. We say that it is the vacuous credal set, as it does not encode any information. It however solves the problem of a contradicting probability when evaluating the value of the card and its parity presented in \Cref{ex:proba_limitations}
\end{example}

\subsection{Belief Functions}\label{sec:belief_functions}
A special case of \acrshort{ip} are belief functions, which we will detail in this section. First, we will introduce a key concept that goes along belief functions: mass distribution functions. We will then derive belief functions from it.

\begin{definition}[Mass distribution function]\label{def:mass_distribution_function}
    Let $\X$ be a set of possible outcomes, and $2^\X$ its power set. A mass distribution function (or basic probability assignment \cite{shafer_mathematical_1976}) is a function $m:2^\X\rightarrow[0,1]$ \st:
    \begin{align}
        m(\emptyset) =& 0\label{eq:mass_emptyset}\\
        \sum_{a\subseteq\X}m(a) =& 1\label{eq:mass_whole_set}
    \end{align}
\end{definition}

There are multiple ways of interpreting $m(a)$, presented as follows:
\begin{itemize}
    \item If we consider a random set $X$, then $m(a)$ encodes the probability mass that $X$ takes $a$ as its set value.
    \item If $X$ is a random variable, then $m(a)$ encodes the available evidence that the numerical value of $X$ is \textit{exactly} in $a$, without any preferences for the values within $a$. This means there could also be other evidence that $X$ is in $a'\subset a$, encoded with $m(a')$, and which could be either more or less than $m(a)$ depending on the amount of evidence available. \Cref{ex:bicycle_pressure} illustrates this with a toy scenario. 
    \item $m(a)$ can also be linked to precise probability masses defined in \Cref{def:density}. If we assume there exists an unknown underlying probability measure $P$ for $X$, $m(a)$ measures the probability mass that is assigned to $a$, but that can move freely to every point of $a$ without any preference. In other words, with more information, we could distribute $m(a)$ to every element of $a$, and doing this for all $a$ would lead to a precise probability mass distribution.
\end{itemize}

\begin{remark}
    \Cref{eq:mass_emptyset} translates the fact that there is no evidence that the uncertain variable belongs to the empty set, \ie that it is not defined. Relaxing this constraint allows accepting a certain amount of contradiction in our model.
    \Cref{eq:mass_whole_set} is a convention, which states that the total amount of evidence equals $1$. It is similar to probabilities, which cannot be more than $1$. 
\end{remark}

\begin{definition}[Focal set]\label{def:focal_set}
    Let $\X$ be a set of possible outcomes, $2^\X$ its power set, and $m:2^\X\rightarrow[0,1]$ a mass distribution function. A set $a\subseteq\X$ is called a focal set of $m$ if and only if:
    \begin{align}
        m(a)>0\label{eq:focal_set}
    \end{align}
\end{definition}
Focal sets thus represent sets of the set of possible outcomes for which we have evidence. The set of all focal sets is sometimes referred to as the \textit{core} of $m$.

\begin{definition}[Belief function, Plausibility function]\label{def:belief_plausibility}
    Let $\X$ be the set of possible outcomes, $2^\X$ its power set, and $m:2^\X\rightarrow[0,1]$ a mass distribution function.
    
    We define the belief function associated with $m$ as the function $\Bel:2^\X\rightarrow[0,1]$ who associates to all events $A$ in $2^\X$:
    \begin{align*}
        \Bel(A)=\sum_{a\subseteq A} m(a)
    \end{align*}
    
    We define the plausibility function associated with $m$ as the function $\Pl :2^\X\rightarrow[0,1]$ who associates to all events $A$ in $2^\X$:
    \begin{align*}
        \Pl(A)=\sum_{\substack{a\\a\cap A\neq\emptyset}} m(a)
    \end{align*}
    $\Bel$ and $\Pl$ are special cases of lower and upper probabilities, and thus induce a credal set $\M(\Bel, \Pl)$ defined as:
    \begin{align*}
        \M(\Bel, \Pl) =\{P~|~\forall A\subseteq\X, \Bel(A) \leq P(A) \leq \Pl(A) \}
    \end{align*}
\end{definition}
We can interpret $\Bel(A)$ as the amount of evidence that fully support $A$, and $\Pl(A)$ as the amount of evidence that is consistent (or does not contradict) with $A$.

$\Bel$ and $\Pl$ are special cases of lower and upper probabilities, and verify \eqref{eq:lower_proba_complement} as for all events $A$ it holds:
\begin{align*}
    \Bel(A)&=\sum_{a\subseteq A} m(a)\\
    &= \sum_{a\in \X} m(a) - \sum_{a\not\subseteq A} m(a)\\
    &= 1 - \sum_{\substack{a\\a\cap A^c\neq\emptyset}} m(a)\\
    &=1-\Pl(A^c)
\end{align*}
Because $\Pl$ is completely determined by $\Bel$ through complementation, we will refer to $\M(\Bel,\Pl)$ simply as $\M(\Bel)$.

Belief functions possess interesting properties, making the credal set they induce both coherent and avoiding sure loss, motivating their extensive usage.

\begin{example}[Defining mass and belief functions]\label{ex:bicycle_pressure}
    Let us imagine an experiment where you try to estimate the pressure of the tires of your bike, but your bicycle pump only have graduation every $1$ bar. You are able to do three measurements:
    \begin{itemize}
        \item The first measure, the needle seems to be between $4$ and $5$ bar.
        \item The second measure, the needle seems to be between $4.5$ and $6$ bar.
        \item The third measure, the needle seems to be between $4.5$ and $5.5$ bar.
    \end{itemize}
    Let us say that you trust your last measurement the most, because you got used to the movement of the needle. It is now possible to model the uncertainty of your tire pressure using available evidence, encoded in $m$:\newline
    \begin{minipage}{0.3\textwidth}
        \begin{align*}
            m([4,~5]) ~=~ 0.3
        \end{align*}
    \end{minipage}\hfill
    \begin{minipage}{0.3\textwidth}
        \begin{align*}
            m([4.5,~6]) ~=~ 0.3
        \end{align*}
    \end{minipage}\hfill
    \begin{minipage}{0.3\textwidth}
        \begin{align*}
             m([4.5,5.5]) ~=~ 0.4
        \end{align*}
    \end{minipage}\bigskip\newline
    Based on this, we are now able to express our degree of belief and of plausibility for all events. For instance:
    \begin{itemize}
        \item Our degree of belief that the pressure lies in $[4, 5]$ is $\Bel([4, 5])=0.3$, that it lies in $[4, 5.5]$ is $\Bel([4, 5.5])=0.7$ and that it lies in $[4, 6]$ is $\Bel([4, 6])=1$
        \item The degree of plausibility that the pressure equals $5$ is $\Pl (5)=1$ (totally plausible), that it equals $5.5$ is $\Pl (5.5)=0.7$ and that it equals $4$ is $\Pl (4)=0.3$.
    \end{itemize}
\end{example}

\subsection{Possibility Distributions}\label{sec:possibilities}
Other convenient models of uncertainty are possibility distributions, which are a specific case of imprecise probabilities, as presented in \Cref{fig:diagram_IP}. We will see that they induce a particular type of belief functions, and will be used in our applications.

\begin{definition}[Possibility distribution]\label{def:possibility}
    Let $\X$ be the set of possible outcomes. A possibility distribution is a function $\pi: \X \rightarrow [0,1]$ satisfying:
    \begin{equation}
    	\exists x \in \X, \pi(x) = 1 \label{eq:possibility}
    \end{equation}
    The value $\pi(x)$ represents the degree of possibility of $x$, with $\pi(x) = 1$ indicating full possibility, and $\pi(x) = 0$ indicating impossibility.
\end{definition}

Another notion closely related to possibility distributions is that of \(\alpha\)-cuts:
\begin{definition}[\(\alpha\)-cut]\label{def:alpha_cut}
    Let $\pi: \X \rightarrow [0,1]$ be a possibility distribution. Given any \(\alpha\in[0,1]\), we define the \(\alpha\)-cut of \(\pi\) as:
    \begin{align}
        \alpha_\pi=\{x~|~\pi(x)\geqslant\alpha\} \label{eq:alpha_cut}   
    \end{align}
    An \(\alpha\)-cut is thus the set of all elements of $\X$ whose possibility level is more than $\alpha$.
\end{definition}

\begin{definition}[Necessity and Possibility measures]\label{def:possibility_credal_set}
    It has been proven in \cite{dubois_when_1992} that a possibility distribution defines a specific type of plausibility functions called \textit{possibility} function and noted \(\Pi\). It also defines a specific belief function by duality called \textit{necessity} function and noted \(\Nec\), as well as a credal set $\M(\pi)$. They are defined as:
    \begin{align}
        \Pi(A) &= \sup_{x\in A}\pi(x)\\
        \Nec(A)& =1-\sup_{x\in A^c}\pi(x)\label{eq:bel_pl}\\
        \M(\pi)& =\{P~|~\forall A,~P(A)\leqslant \sup_{x\in A}\pi(x)\}\\
        &=\{P~|~\forall A,~P(A)\leqslant \Pi(A)\} \label{eq:credal_set_possibility}\\
        &= \{P~|~\forall A,~P(A)\geqslant \Nec(A)\}
\end{align}
\end{definition}

\begin{example}[Defining a possibility distribution]\label{ex:bicycle_pressure_possibility}
    Let us imagine the same setting as \Cref{ex:bicycle_pressure}, where we try to estimate the pressure of the tires of our bike, but our bicycle pump only has graduations every $1$ bar. We are able to do the following measurements:
    \begin{itemize}
        \item During the first measurement, the needle seems to be around $4$ bar.
        \item During the second measurement, the needle seems to be around $5$ bar.
        \item During the third measurement, the needle seems to be around $4.5$ bar.
    \end{itemize}
    Let us say that we trust our measurement with a precision of $\pm0.5$ bars. For simplicity, we also only consider pressure values that are integers or half integers. Taking into considerations all the measurements, we can define a possibility distribution $\pi$ as:
    \begin{itemize}
        \item The value of $4.5$ being the most possible, $\pi(4.5)=1$.
        \item Values between $4$ and $5$ bar being highly possible, we can for instance fix $\pi(4)=\pi(5)=0.8$
        \item Values $3.5$ and $5.5$ are unlikely but not impossible, we can say that $\pi(3.5)=\pi(5.5)=0.3$
        \item Other values are impossible, thus they have a possibility of $0$.
    \end{itemize}
    The possibility distribution $\pi$ is represented in \Cref{fig:possibility_distribution}. 
    Based on this, we are now able to express the degree of necessity and possibility for all events. For instance:
    \begin{itemize}
        \item The degree of necessity that the pressure lies between $4$ and $5$ bar is $\Nec([4, 5])=1-\sup_{\rho\not\in[4,5]}\pi(\rho)=0.7$
        \item The degree of necessity that the pressure lies between $3.5$ and $5.5$ bar is: 
        \begin{align*}
            \Nec([3.5,~ 5.5]) = 1-\sup_{\rho\not\in[3.5,5.5]}\pi(\rho) = 1
        \end{align*}
        meaning that the pressure is necessarily in this range.
        \item The degree of possibility that the pressure is either $3.5$ or  $5$ is $\Pi (\{3.5,~5\})=\sup_{\rho\in\{3.5,~5\}}\pi(\rho)=0.9$ (mostly possible).
        \item \etc for every possible event.
    \end{itemize}
\end{example}

\begin{figure}[!ht]
    \centering
    \begin{tikzpicture}[scale=1]
        \begin{axis}[%
          xlabel=Pressure $\rho$ (in bar),
          ylabel=$\pi(\rho)$,
          xmin=2, xmax=7,
          ymin=-0.1, ymax=1.1,
          ytick={0, 0.3, 0.8, 1},
          yticklabels={$0$, $\alpha=0.3$, 0.8, $1$},
          ],
          \node (a) at (2, 0) {$\times$};
          \node (b) at (2.5, 0) {$\times$};
          \node (c) at (3, 0) {$\times$};
          \node (d) at (3.5, 0.3) {$\times$};
          \node (e) at (4, 0.8) {$\times$};
          \node (f) at (4.5, 1) {$\times$};
          \node (g) at (5, 0.8) {$\times$};
          \node (h) at (5.5, 0.3) {$\times$};
          \node (i) at (6, 0) {$\times$};
          \node (j) at (6.5, 0) {$\times$};
          \node (k) at (7, 0) {$\times$};
          
          \draw [thick, dashed, gray] (a.center) -- (b.center) -- (c.center) -- (d.center) -- (e.center) -- (f.center) -- (g.center) -- (h.center) -- (i.center) -- (j.center) -- (k.center) ;
          \draw [<->, >=latex, ultra thick, blue] (d.east) -- (h.west) node [pos=0.5, above] {$\alpha_\pi$};
        \end{axis}
    \end{tikzpicture}
    \caption{Possibility distribution of \Cref{ex:bicycle_pressure_possibility} and one of its $\alpha$-cut in blue}
    \label{fig:possibility_distribution}
\end{figure}

\begin{remark}
    If you are familiar with fuzzy sets, you may have noticed that possibility distributions are similar to membership functions of a fuzzy set. Links between fuzzy sets and possibility measures have been explored in \cite{zadeh_fuzzy_1999}.
\end{remark}

Focal sets of necessity functions can be determined directly from the possibility distribution by looking at their \( \alpha \)-cuts. It has been proven that the core $\mathcal{C}$ (the set containing all focal sets) of a necessity function is (\cite{couso_necessity_2001}):
\begin{align*}
    \mathcal{C} =& \{\alpha_\pi~|~\alpha\in[0,1]\}\\
    =& \{ ~\{x\in\X~|~\pi(x)\geqslant\alpha\}~|~\alpha\in[0,1]~\}
\end{align*}
With the way focal sets are defined, they form a nested family of sets with regard to inclusion. Indeed, if an element of \(\X\) belongs to an \(\alpha\)-cut, then its possibility is greater than \(\alpha\) and therefore belongs to any other \(\alpha'\)-cut with a lower \(\alpha'\). For simplicity, we will assume that the focal sets \(a_1\enum a_n\) are already numbered using the inclusion order, \ie \( a_1\subset\ldots\subset a_n\). In this case, we will refer to the inclusion order as the ``natural'' order for possibility distributions.

\begin{remark}
    The fact that focal sets form a nested family of sets in the case of possibility distributions also implies that if $\X$ is finite and contains $n$ elements, then there can be \textit{at most} $n$ focal sets. For comparison, belief functions can have a maximum of $2^n-1$ focal sets (as the empty set cannot be a focal set). This means that necessity functions have fewer degrees of freedom than (some) belief functions, and thus can express fewer uncertainty structures. This drawback comes with the advantage of being more straightforward to construct, as we only need to specify the mass of $n$ focal sets (or the possibility of the $n$ elements of $\X$) instead of $2^n-1$. Indeed, when we think of a random variable like the outcome of a die, it can seem more natural for someone to specify degrees of possibility for each side separately than it is to specify degrees of plausibility for different sets of outcomes. 
    
    As such, possibility distribution have been used to model experts' opinion in domain such as water contamination \cite{bardossy_l-_1995}, soil contamination and radioactive risk assessment \cite{baudrit_comparing_2005,baudrit_representation_2005,baudrit_joint_2007} or weather forecasting \cite{le_carrer_beyond_2021}. Following the same philosophy, we will use possibility distributions in \Cref{chap:epistemic_uncertainty} to model the uncertainty of a measure of similarity between two image patches.
\end{remark}

Specifying a probability distribution often comes down to specifying the probability mass function over all atoms of the set of possible outcomes. In a way, possibility distributions are constructed the same way, as we specify the possibility (or the upper bounds) of every atom. One important difference is that the condition ``the sum of all masses must be equal to $1$'' is relaxed into a less constraining condition ``the possibility distribution must be equal to $1$ at least once''. In that respect, it is easier to construct a well-defined possibility distribution than it is to construct a well-defined probability distribution. However, the comparison stops there, as the two models does not represent the same type of uncertainty.

\begin{remark}
    Any probability distribution $P$ is a belief function $\Bel$, for which focal sets are only composed of singletons (atoms) and the mass distribution function of $\Bel$ equals the probability mass function of $P$ on atoms. However, a possibility distribution cannot model a probability distribution. Indeed, this would impose that its necessity $\Nec$ and plausibility $\Pi$ functions verify:
    \begin{align*}
        \forall A, &~\Nec(A) = \Pi(A)\\
        \Leftrightarrow&~1-\sup_{x\in A^c}\pi(x) = \sup_{x\in A}\pi(x)\\
        \Leftrightarrow&~ \sup_{x\in A}\pi(x) + \sup_{x\in A^c}\pi(x) = 1
    \end{align*}
    Consider this equation for any $x'$ verifying $\pi(x')=1$. This leads to the conclusion that any $x\neq x'$ has a possibility of $0$, implying that it is impossible for a random set or random variable to take any other value than $x'$, rendering it not-random. 
\end{remark}

\subsection{P-boxes}\label{sec:pboxes}
Another special type of belief function that is commonly used is that of \textit{probability boxes}, more commonly called p-boxes. Formally, a p-box is a pair of precise cumulative distribution functions $[\underline{F}, ~\overline{F}]$ defining lower and upper bounds on all cumulative events:

\begin{definition}[P-box]\label{def:p-box}
    Let $\X$ be the set of possible outcomes. A p-box is a pair of \acrshort{cdf}s $[\underline{F}, ~\overline{F}]$ from $\X$ to $[0, ~1]$ such that:
    \begin{equation}
    	\forall x \in \X,~\underline{F}(x) \leqslant \overline{F}(x) \label{eq:p-box}
    \end{equation}
    If $\X$ is not a subset of $\mathbb{R}$, then there must exist a total order on $\X$ to define a generalized p-box \cite{destercke_unifying_2008}. 
\end{definition}

\begin{remark}
    A probability distribution can be both determined by specifying its values on every atom or by specifying its values on cumulative events. We saw in \Cref{sec:possibilities} that possibility distributions are defined by bounds on atoms. Because p-boxes are defined by bounds on cumulative events, we could say that a p-box is the ``imprecise way'' of defining a probability using cumulative events, and a possibility distribution is the ``imprecise way'' of defining a probability using atoms.
\end{remark}

The credal set $\M$ induced by a p-box $[\underline{F}, ~\overline{F}]$ is:
\begin{align}
    \M([\underline{F}, ~\overline{F}]) = \{~ F ~|~ \forall x\in\X, ~\underline{F}(x)\leqslant F(x)\leqslant \overline{F}(x) ~\}
\end{align}

\begin{definition}[Focal sets of p-boxes]
    P-boxes are special cases of belief functions. It has been proven in \cite{destercke_unifying_2008} that focal sets of p-boxes have a specific form. Although focal sets shapes are reminiscent of possibilities' $\alpha$-cuts (see \Cref{fig:p-box}), they are a bit more complex to express formally. If $\X=\{x_1\enum x_n\}$ with $x_1\leqslant\ldots\leqslant x_n$, then focal sets $\alpha_{[\underline{F}, ~\overline{F}]}$ of $[\underline{F}, ~\overline{F}]$ are given for every $\alpha\in[0,~1]$ by the following expression:
    \begin{align}
        \alpha_{[\underline{F}, ~\overline{F}]}= \opi \overline{F}^{-1}(\alpha),~\underline{F}^{-1}(\alpha)\cli\label{eq:pbox_focal_set}
    \end{align}
    where $\opi\cdot,~\cdot\cli$ are intervals of integers, and $\overline{F}^{-1}$, $\underline{F}^{-1}$ are the respective pseudo-inverse of $\overline{F}$ and $\underline{F}$ defined for every $\alpha\in[0,1]$ by:
    \begin{align*}
        \overline{F}^{-1}(\alpha) =& \min \{x_i ~\st ~\overline{F}(x_i)\geqslant \alpha\}\\
        \underline{F}^{-1}(\alpha) =& \min \{x_i ~\st ~\underline{F}(x_i)\geqslant \alpha\}
    \end{align*}
    
    Still in \cite{destercke_unifying_2008}), it has been shown that the mass of each focal set $\opi x_i, ~x_j\cli$ equals to :
    \begin{align}
        m(\opi x_i, ~x_j\cli) = \min(\overline{F}(x_i), \underline{F}(x_j)) - \max(\overline{F}(x_{i-1}), \underline{F}(x_{j-1}))
    \end{align}
    With the convention that $\max(\overline{F}(x_{i-1}), \underline{F}(x_{j-1}))=0$ if $x_{i}$ is the first element and thus $x_{i-1}$ is ill-defined. 
\end{definition}

Because of their shape, focal sets of p-boxes can be totally ordered through the so-called lattice ordering on $\X$. It can be easily observed by looking at \Cref{fig:p-box}. If $a$ and $b$ are two focal sets of the same p-box $[\underline{F}, ~\overline{F}]$, then they are ordered as follows:
\begin{align}
    a\leqslant b \Leftrightarrow \min(a)\leqslant\min(b) \text{ and } \max(a)\leqslant \max(b)
\end{align}
As $\underline{F}\leqslant\overline{F}$, we are assured that there cannot be any case where $\min(a)<\min(b)$ and $\max(a)> \max(b)$. We can also define the order of focal sets using the definition of \Cref{eq:pbox_focal_set} for every $\alpha,~\beta\in[0,1]^2$:
\begin{align*}
    \opi\overline{F}^{-1}(\alpha), ~\underline{F}^{-1}(\alpha)\cli \leqslant \opi\overline{F}^{-1}(\beta), ~\underline{F}^{-1}(\beta)\cli \Leftrightarrow \alpha\leqslant\beta
\end{align*}

Given the shape of focal sets, there can be at most $2n-1$ of them in the set of $n$ possible outcomes. This is more degree of freedom than possibility distributions, but less than general belief functions.
\begin{remark}
    Contrary to possibility distributions that cannot equal to a single probability distribution, if a p-box $[\underline{F}, ~\overline{F}]$ verifies $\underline{F}=\overline{F}$, then its credal set is composed of a single probability distribution whose \acrshort{cdf} $F$ equals $\underline{F}$ and $\overline{F}$. P-boxes are thus generalizations of precise probability distributions.
\end{remark}

\begin{figure}[!ht]
    \centering
    \begin{tikzpicture}[scale=1]
        \begin{axis}[%
          xlabel=$x$,
          ylabel=$\CDF(x)$,
          %grid=major,
          domain=0:15,
          %legend entries={$\underline{F}$, $\overline{F}$},
          legend pos=south east,%
          ],
            \addplot[dotted, black, samples=16, mark=triangle, mark options={solid}, mark size=3pt] {cdf_erf(x, 1.5, 1.5)} node [pos=0.4, above left] {$\overline{F}$};
            \addlegendentry{$\overline{F}$}
            
            \addplot[dotted, black, samples=16, mark=square, mark options={solid}, mark size=3pt] {cdf_erf(x, 7.5, 1.5)} node [pos=0.6, below right] {$\underline{F}$};
            \addlegendentry{$\underline{F}$}
            
            \addplot[dotted, gray, samples=16, mark=+, mark options={solid, color=gray}, mark size=3pt] {cdf_erf(x, 4.5, 2)};
            \addlegendentry{$F$}
            
            \node (a) at (4, {cdf_erf(4, 1.5, 1.5)}) {};
            \node (b) at (10, {cdf_erf(10, 7.5, 1.5)}) {};
            
            \draw [<->, >=latex, ultra thick, blue] (a.east) -- (b.west) node [pos=0.5, above] {$\alpha_{[\underline{F}, ~\overline{F}]}$};
        \end{axis}
        \end{tikzpicture}
    \caption{A p-box $[\underline{F}, ~\overline{F}]$, a precise \acrshort{cdf} $F$ in its credal set, and one of its focal elements $\alpha_{[\underline{F}, ~\overline{F}]}$ in blue}
    \label{fig:p-box}
\end{figure}


\section{Dependency Models: Copulas}\label{sec:copulas}
During previous sections, we presented different models of uncertainty that will be considered throughout this thesis. When we will be aggregating and propagating uncertainty over multiple variables in \Cref{chap:joining_credal_sets,chap:propagating}, we will need to take into account the dependencies between our uncertain variables. In this section, we will present dependency models known as copulas, which are mathematical tools used to represent the dependency between multiple random variables. Copulas can represent many types of dependency, ranging from complete monotonicity to complete counter-monotonicity, including independence between variables. \Cref{sec:copula_def} will present the mathematical definition of a copula, as well as practical families of copulas and how they can model different dependencies. Finally, we present how to generate multivariate samples from a copula, which will be used in \Cref{sec:montecarlo}.

\subsection{Core Definitions and Examples}\label{sec:copula_def}
In the following, let $n\in\mathbb{N}^*$ be the number of uncertainty variables considered (either represented by random variables or random sets). We first introduce copulas, which are mapping from $[0,1]^n\rightarrow [0,1]$ verifying a number of properties, and that can model dependencies when considering \hyperref[theorem:sklar]{Sklar's Theorem}, which will be presented in this section.

\begin{definition}
     A copula is a multivariate cumulative distribution function $C:[0,1]^{n}\rightarrow [0,1]$ whose marginals follow uniform distributions on $[0,1]$. It can be interpreted as a joint cumulative distribution of $n$ random variables. For all $i\in\opi 1,n\cli $, we will refer to $u_i\in[0,1]$ as its $i$-th variable (or marginal). A copula satisfies a number of properties:
\begin{align}
    &\text{if }\exists j\in\opi 1,n\cli  \text{ \st }~u_j=0, \text{ then }C(u_1\enum u_j\enum u_n)=0\label{eq:zero_copula}\\
    &\forall i\in\opi 1,n\cli ,~C(1,1\enum1,u_i,1\enum1)=u_i\label{eq:copula_ones}\\
    &\forall (v_1\enum v_n)\in[0,1]^n \text{ \st }~\forall i\in\opi 1,n\cli ,~v_i\geqslant u_i\nonumber\\
    &\sum_{(w_1\enum w_n)\in\prod_{i=1}^n\{u_i, v_i\}}(-1)^{|\{i~|~w_i=u_i\}|}C(w_1\enum w_n)\geqslant 0\label{eq:cop_hvolume}
\end{align}
where $\prod_{i=1}^n$ is the Cartesian product of $n$ elements, meaning that $(w_1\enum w_n)\in\prod_{i=1}^n\{u_i, v_i\}$ is a tuple of $n$ elements, where each element is either $u_i$ or $v_i$. Additionally, $|\{i~|~w_i=u_i\}|$ refers to the cardinal of the set $\{i~|~w_i=u_i\}$. An interpretation of the value of \Cref{eq:cop_hvolume} is presented in the next remark.
\end{definition}

The first term in \Cref{eq:cop_hvolume} is also called H-volume or hyper-volume. It is used to compute joint probability mass assignments in the precise case (and also in the imprecise case, see \Cref{sec:joint_mass}). We will now use the following notation to refer to the $H$-volume:
\begin{eqnarray}\label{eq:hvolume}
    &&\forall i\in\opi 1,n\cli ,~\forall~0\leqslant u_i \leqslant v_i \leqslant 1,\nonumber\\
    &&H^{v_1\enum v_n}_{u_1\enum u_n}=\sum_{(w_1\enum w_n)\in\prod_{i=1}^n\{u_i, v_i\}}(-1)^{|\{i~|~w_i=u_i\}|}C(w_1\enum w_n)
\end{eqnarray}

\begin{remark}
    The formula of the H-volume actually represents the probability that $n$-uniform random variables are in the hyper rectangle $[u_1,v_1]\tdt[u_n,v_n]$. However, it is difficult to see this interpretation in the general case just by looking at the formula. For simplicity, consider the two-dimensional case. Using the interpretation of a copula $C$ as a \acrshort{cdf}, we can image two random uniform variables $U_1$ and $U_2$ on $[0,1]$ for which $C$ is their \acrshort{cdf}s. We thus have for all $(u_1,u_2)\in[0,1]^2$:
    \begin{equation*}
        P(U_1\leqslant u_1, U_2\leqslant u_2) = C(u_1, u_2)
    \end{equation*}
    Let $(u_1,u_2)\in[0,1]^2$ and $(v_1,v_2)\in[0,1]^2$ \st $u_1\leqslant v_1$ and $u_2\leqslant v_2$. 
    Computing the H-volume of $C$ between $(v_1,v_2)$ and $(u_1,u_2)$ yields:
    \begin{align}
        H^{v_1,v_2}_{u_1,u_2} =& ~C(v_1, v_2) - C(v_1, u_2) - C(u_1, v_2) + C(u_1, u_2)\nonumber\\
        =&~ P(U_1\leqslant v_1, ~U_2\leqslant v_2) - P(U_1\leqslant v_1, ~U_2\leqslant u_2)\nonumber\\
        &- P(U_1\leqslant u_1, ~U_2\leqslant v_2) + P(U_1\leqslant u_1, ~U_2\leqslant u_2)\nonumber\\
        =&~ P(U_1\leqslant v_1, ~u_2 < U_2 \leqslant v_2) - P(U_1\leqslant u_1, ~u_2 < U_2\leqslant v_2)\nonumber\\
        =&~ P(u_1 < U_1\leqslant v_1, ~u_2 < U_2\leqslant v_2)\label{eq:hvol_link_with_proba}
    \end{align}
    
    {\centering
    \begin{tikzpicture}[trim axis left, trim axis right, scale=1.1]
        \begin{axis}[%
            xmin=0, xmax=1,
            ymin=0, ymax=1,
            xtick={0.3, 0.8, 1},
            xticklabels={$u_1$, $v_1$, $1$},
            ytick={0, 0.4, 0.8, 1},
            yticklabels={$0$, $u_2$, $v_2$, $1$},
            xtick pos=left,
            ytick pos=left],
            
            \addplot[domain=0:1,samples=40,opacity=0.2]({x},{0.1/x});
            \addplot[domain=0:1,samples=40,opacity=0.2]({x},{0.3/x});
            \addplot[domain=0:1,samples=40,opacity=0.2]({x},{0.5/x});
            \addplot[domain=0:1,samples=40,opacity=0.2]({x},{0.7/x});
            \addplot[domain=0:1,samples=40,opacity=0.2]({x},{0.9/x});
            
            \node (a) at (0.8, 0.85) {};
            \filldraw[black] (a.center) circle (2pt) node[anchor=south,yshift=2pt]{$C(v_1,v_2)$};
            
            \node (b) at (0.3, 0.85) {};
            \filldraw[black] (b.center) circle (2pt) node[anchor=south,yshift=2pt]{$C(u_1,v_2)$};
            
            \node (c) at (0.3, 0.4) {};
            \filldraw[black] (c.center) circle (2pt) node[anchor=east,rotate=45]{$C(u_1,u_2)$};
            
            \node (d) at (0.8, 0.4) {};
            \filldraw[black] (d.center) circle (2pt) node[anchor=west, rotate=45]{$C(v_1,u_2)$};
            
            \draw[ultra thick, black, fill=gray, opacity=0.3] (c.center) rectangle (a.center);
            \draw[ultra thick, black] (c.center) rectangle (a.center) node[pos=0.5,yshift=7pt]{$P(u_1<U_1\leqslant v_1,$} node[pos=0.5, yshift=-7pt, xshift=7pt]{$u_2<U_2\leqslant v_2)$};
            
            \node (e) at (0, 0.85) {};
            \node (f) at (0, 0.4) {};
            \node (g) at (0.8, 0) {};
            \node (h) at (0.3, 0) {};
            \draw[thick, dashed, black] (e.center) -- (b.center);
            \draw[thick, dashed, black] (f.center) -- (c.center);
            \draw[thick, dashed, black] (g.center) -- (d.center);
            \draw[thick, dashed, black] (h.center) -- (c.center);
            
            \node (m) at (0.3, 0.05) {};
            \node (n) at (0.8, 0.05) {};
            \node (o) at (0.05, 0.85) {};
            \node (p) at (0.05, 0.4) {};
            
            \draw [stealth-stealth, thick, black] (m.center) -- (n.center) node[pos=0.5, above]{$u_1<U_1\leqslant v_1$};
            \draw [stealth-stealth, thick, black] (o.center) -- (p.center) node[pos=0.5, anchor=north, rotate=-90, yshift=20pt]{$u_2<U_2\leqslant v_2$};
        \end{axis}
    \end{tikzpicture}\captionof{figure}{Schematic representation of the H-volume}\par}
    
    \Cref{eq:hvol_link_with_proba} means that the H-volume represent the probability of the event
    \begin{equation*}
        u_1 < U_1\leqslant v_1, ~u_2 < U_2\leqslant v_2
    \end{equation*}
    or in other words, the probability that $(U_1, U_2)$ is in the hyper rectangle $[u_1,v_1]\times[u_2,v_2]$ (the intervals can be open or closed in the continuous case, the probability remains the same). Verifying this result for the $n$-dimensional case can be done similarly. \Cref{ex:hvolume} illustrates how the H-volume can be used to compute the discrete joint mass distribution function in the two-dimensional case. 
\end{remark}

A central theorem regarding copulas is \hyperref[theorem:sklar]{Sklar's Theorem} \cite{sklar_fonctions_1959}:
\begin{theorem}[Sklar's Theorem]\label{theorem:sklar}
    Let $F:\X_1\tdt\X_n\rightarrow[0,1]$ be a multivariate cumulative distribution function, where $\X_i\subseteq\overline{\mathbb{R}}$. The marginals $F_i$ of $F$ are defined as $\forall i\in\opi 1,n\cli , \forall x\in\X_i, F_i(x) = F( +\infty\enum  +\infty, x,  +\infty\enum +\infty)$ where $x$ is the $i$-th component of $F$. If all $F_i$ are continuous, then a unique copula $C$ exists:
    \begin{align}
        \forall (x_1\enum x_n)\in \overline{\mathbb{R}}^n, F(x_1\enum x_n)=C(F_1(x_1)\enum F_n(x_n))\label{eq:sklar_equality}
    \end{align}
    If some $F_i$ are not continuous, then $C$ is unique on the product of the ranges of all $F_i$.
    
    The reverse is also true: any copula applied to univariate cumulative distribution functions yields a multivariate cumulative distribution function whose marginals are the univariate \acrshort{cdf}s.
\end{theorem}

\hyperref[theorem:sklar]{Sklar's Theorem} thus allow to express any multivariate \acrshort{cdf} by means of its marginal \acrshort{cdf}s. Conversely, we can join multiple \acrshort{cdf}s with a copula to create a multivariate \acrshort{cdf}. A copula thus expresses the dependency between a multivariate \acrshort{cdf} and its marginals.
 
\begin{remark}
    For marginals $F_i$ that are not continuous, then there can exist multiple copula $C$ verifying \Cref{eq:sklar_equality}. However, if we note $F_i(\X_i)$ the image of $X_i$ through $F_i$, then there exists a unique copula $C$ on the ranges of images $F_1(\X_1)\tdt F_n(\X_n)$. The restriction of a copula to a subset of $I^n$ containing $0$ and $1$ is called a sub-copula. Because we work in discrete spaces, we will mostly work with sub-copulas, but the difference will be mostly transparent.
\end{remark}

Copulas are very useful to represent the dependencies between multiple uncertain variables. As such, they can play a key role in uncertainty propagation problems, explored in \Cref{chap:propagating}.

\begin{example}[Usefulness of the H-Volume]\label{ex:hvolume}
    This example will illustrate how the H-volume of a copula is used to compute the probability mass function of a multivariate probability. Let us imagine a game where a dealer throws two coins, and we are interested in the joint result of the throws. We consider the two random variables $X_1$ and $X_2$ indicating the results of each throw:
    \begin{align*}
        &X_1=0\text{ if the first coin is heads, otherwise }X_1=1\\
        &X_2=0\text{ if the second coin is heads, otherwise }X_2=1
    \end{align*}
    Let $P_1$ and $P_2$ be the probability distributions of $X_1$ and $X_2$ respectively, and assume that both heads and tails are possible outcomes for both coins. We now consider the joint probability distribution $P$ associated with the random variable $(X_1,~X_2)$, and we want to compute the probability mass distribution of $P$. We denote $F_1$, $F_2$ and $F$ the respective \acrshort{cdf}s of $P_1$, $P_2$ and $P$. With this definition, $F_1$ and $F_2$ are the marginals of $F$, and Sklar's theorem states that there exists a copula $C$ such that $F=C(F_1,~F_2)$. 
    
    The probability mass distribution of $P$ can be computed by using the H-volume. For instance, let us start by computing $P(1,1)$. By noticing the fact that:
    \begin{equation*}
        \{x\in\X~|~X_1(x)=1\}=\{x\in\X~|~F_1(X_1(x))=F_1(1)\}
    \end{equation*}
    we can write that:
    \begin{align*}
        P(X_1 = 1, X_2=1) =& ~P(F_1(X_1)= 1, F_2(X_2) = 1)\\
        =& ~P(F_1(0) < F_1(X_1)\leqslant F_1(1), ~F_2(0) < F_2(X_2)\leqslant F_2(1))
    \end{align*}
    A common result in statistics states that random variables of the form $U_1=F_1(X_1)$ and $U_2=F_2(X_2)$ are uniform on $[0,1]$. We can thus apply the result from \Cref{eq:hvol_link_with_proba}, which yields
    \begin{align*}
        P(X_1 = 1, X_2=1) =& ~P(F_1(0) < U_1 \leqslant F_1(1), ~F_2(0) < U_2 \leqslant F_2(1))\\
        =& ~H^{F_1(1),~F_2(1)}_{F_1(0),~F_2(0)}
    \end{align*}
    The probability of the atom $(1,1)$ is therefore equal to the H-volume computed between \acrshort{cdf}s $(F_1, F_2)$ at $(1, 1)$ and at $(0,0)$. Following a similar reasoning, we can compute the probability of every atom and express them as H-volumes:
    \begin{align*}
        P(X_1=0, X_2=1) =& ~H_{0, F_2(0)}^{F_1(0), F_2(1)}\\
        P(X_1=1, X_2=0) =& ~H_{F_1(0), 0}^{F_1(1), F_2(0)}\\
        P(X_1=0, X_2=0) =& ~H_{0,0}^{F_1(0), F_2(0)}
    \end{align*}
    It is possible to generalize our observation: the probability of an atom $(x_1,~\ldots,~x_n)$ is the H-volume computed between marginals \acrshort{cdf}s $(F_1,~\ldots,~F_n)$ at $(x_1,~\ldots,~x_n)$ and at the marginal atoms that precedes them. If $x_1$ is the smallest number, then the marginal \acrshort{cdf} $F_1$ before it equals $0$, \etc.
    \Cref{ex:copulas} presents numerical applications of this example with different copulas.
\end{example}

It follows from \eqref{eq:cop_hvolume} that a copula is a component-wise increasing mapping. All copulas are actually dominating and dominated by two bounds (called lower and upper Fréchet–Hoeffding bounds):
\begin{align}
    &\forall u_i \in [0,1]^n,\nonumber\\
    &\max(0, 1-n+\sum_{i=1}^n u_i) \leqslant C(u_1\enum u_n) \leqslant \min(u_1\enum u_n)
\end{align}
The upper bound is a copula, usually called the Minimum copula $C_M$. It is used to model co-monotonic variables, \ie variables for which high values occur at the same time (or similarly, where low values tend to occur simultaneously). Co-monotonicity implies a maximal covariance between variables.

The lower bound is a copula only in the case $n=2$, called the \L ukasiewicz copula $C_L(u_1,u_2)=\max(0,u_1+u_2-1)$. It is used to model counter-monotonic variables, \ie variables with a perfect negative dependence between them. This explains why the lower bound is not a copula in dimensions higher than $2$. Indeed, if $X$ has a perfect negative dependence with $Y$ and $Z$, then $Y$ and $Z$ cannot share a perfect negative dependence. However, for every $u_1\enum u_n$, there always exists a copula $C$ attaining the lower bound (which can differ for different $u_1\enum u_n$):
\begin{eqnarray*}
    \forall (u_1\enum u_n)\in[0,1]^n,~\exists C \text{ \st }~C(u_1\enum u_n) = \max(0, 1-n+\sum_{i=1}^n u_i)
\end{eqnarray*}
Independence between variables is modeled by the product copula $C_\Pi$:
\begin{equation*}
    C_\Pi(u_1\enum u_n)=\prod_{i=1}^n u_i = u_1\cdot\ldots\cdot u_n
\end{equation*}
where ``$\cdot$'' refers to the product between two scalars, as the symbol $\times$ is already used for the Cartesian product. The product copula will be used later in \Cref{subsection:product_copula}. Graphical representations of the \L ukasiewicz, product and Min copulas are displayed in \Cref{fig:copulas}. \Cref{ex:copulas} presents a setting where the Product, Minimum and \L ukasiewicz copulas are used to model dependency between random variables.

\begin{example}[Different copulas for different dependencies]\label{ex:copulas}
    Let us try to illustrate how different copulas can represent different dependencies. 
    Consider the same setting as \Cref{ex:hvolume} with two coins being thrown. For the purpose of the example, assume that the dealer throws the coins in a separate room, and comes back to tell the result. We thus never see if he is cheating or not. He only provides us this piece of information: coins seems fair when looked at separately. We therefore have the following marginals:
    \begin{eqnarray*}
    \begin{cases}
        P_1(\text{heads}) = P_1(0) = 0.5\\
        P_1(\text{tails}) = P_1(1) = 0.5\\
    \end{cases}
    \qquad\text{ and }\qquad
    \begin{cases}
        P_2(\text{heads}) = P_2(0) = 0.5\\
        P_2(\text{tails}) = P_2(1) = 0.5
    \end{cases}
    \end{eqnarray*}
    
    \begin{itemize}
        \item Assume that the dealer is not cheating and that the two coin throws are independent. In that case, the product copula $C_\Pi(u,v)=u\cdot v$ must be used to represent the independence between variables.
        Using results from the previous example, it holds that:
        \begin{align*}
            P(1,1) =& H_{F_1(1), F_2(1)}^{F_1(0), F_2(0)} = 1\cdot1 - 0.5\cdot1 - 1\cdot0.5 + 0.5\cdot0.5\\
            =& 0.25\\
            P(1,0) =& H_{F_1(0), 0}^{F_1(1), F_2(0)} = 0.25 \\
            P(0,1) =& H_{0, F_2(0)}^{F_1(0), F_2(1)} = 0.25 \\
            P(0,0) =& H_{0, 0}^{F_1(0), F_2(0)} = 0.25
        \end{align*}
        Remark that we indeed find the same results as if we directly multiplied the marginal probability mass distributions: $P(1,1) = P_1(1)\cdot P_2(1)$, \etc. We thus observe the famous result: if $P_1$ and $P_2$ are independent, then $P=P_1\cdot P_2$.
        
        \item Imagine now that the dealer is not being fair, and actually forces the second throw to land on the same side as the first one (the coins will still seem fair when looked at separately). This kind of dependency is modeled by the Minimum copula $C_M(u,v)=\min(u,v)$. In this case, the joint probability is computed as follows: 
		\begin{align*}
            P(1,1) =& H_{F_1(1), F_2(1)}^{F_1(0), F_2(0)} = \min(1,1) - \min(0.5,1) - \min(1,0.5) + \min(0.5,0.5)\\
            =& 0.5\\
            P(1,0) =& H_{F_1(0), 0}^{F_1(1), F_2(0)} =\min(1,0.5) - \min(1, 0) - \min(0.5,0.5) + \min(0,0.5)\\
            =& 0\\
            P(0,1) =& H_{0, F_2(0)}^{F_1(0), F_2(1)} = 0 \\
            P(0,0) =& H_{0, 0}^{F_1(0), F_2(0)} = \min(0.5,0.5)=0.5
        \end{align*}
        The values taken by the joint probability are now completely different from the independence case. We see that values $(0,1)$ and $(1,0)$ are indeed impossible to obtain, while $(1,1)$ and $(0,0)$ are equiprobable.
        
		\item Imagine now that the dealer is still not being fair, but this time forces the second coin to land on the first coin's opposite side. In other words, if the first coin lands on heads, then the dealer puts the second coin on tails, and inversely. Looking at marginal distributions separately will still suggest that the coins are fair. However, they appear fully counter-monotone when looked at jointly. In this case, the dependency is modeled by the \L ukasiewicz copula $C_L(u,v)=\max(0,u+v-1)$, and the joint probability equals:
		\begin{align*}
            P(1,1) =& H_{F_1(1), F_2(1)}^{F_1(0), F_2(0)} = \max(0, 1+1-1) - \max(0, 0.5+1-1)\\
            &- \max(0, 1+0.5-1) + \max(0, 0.5+0.5-1)\\
            =& 0\\
            P(1,0) =& H_{F_1(0), 0}^{F_1(1), F_2(0)} =\max(0, 1+0.5-1) - \max(0, 1+0-1)\\
            &- \max(0, 0.5+0.5-1) + \max(0, 0+0.5-1)\\
            =& 0.5\\
            P(0,1) =& H_{0, F_2(0)}^{F_1(0), F_2(1)} = 0.5 \\
            P(0,0) =& H_{0, 0}^{F_1(0), F_2(0)} = \max(0, 0.5+0.5-1)=0
        \end{align*}
        The values taken by the joint probability is now completely different than in the other cases. We see that values $(1,1)$ and $(0,0)$ are indeed impossible to obtain, while $(0,1)$ and $(1,0)$ are equiprobable.
	\end{itemize}
	Those three cases indicate how copulas can represent very different dependency structures from the same marginals. It also makes it intuitive that the \L ukasiewicz copula only allows values that are ``opposite'', whereas the Minimum copula only allows values that are similar. In those examples, the dependency is so important that knowing the result of one coin throw determines the result of the second, which is therefore modeled by extreme copulas, \ie the upper and lower Fréchet-Hoeffding bounds.  We will see in the following that there are other families of copula that allow less ``extreme'' dependencies.
\end{example}

To complete this overview of copulas, let us present other copulas that can be generated using a single parameter $\theta$ in the case $n=2$. Some famous families of copulas are presented in \Cref{tab:family_of_copula}. Those families are quite common in the literature, but this list is not exhaustive.

Another important family of copulas is the family of Gaussian copulas. Each Gaussian copula is generated with a correlation matrix $R\in[-1,1]^{(n,n)}$:
\begin{align}
    C_R(u_1\enum u_n)=\Phi_R(\Phi^{-1}(u_1)\enum \Phi^{-1}(u_n))
\end{align} where $\Phi_R$ is the joint multivariate cumulative distribution function of a Gaussian variable with correlation matrix $R$, and $\Phi^{-1}$ is the inverse cumulative distribution function of a univariate Gaussian variable. We do not know an exact form for $\Phi_R$, but we can compute it by integrating its associated Gaussian \acrshort{pdf}:
\begin{align}
   \Phi_R&(\Phi^{-1}(u_1)\enum \Phi^{-1}(u_n)) =\nonumber\\ &\int_{-\infty}^{\Phi^{-1}(u_1)}\ldots\int_{-\infty}^{\Phi^{-1}(u_n)}\frac{1}{\sqrt{(2\pi)^{n}|R|}}\exp\left(-\frac{1}{2}\begin{bmatrix}x_1 & \ldots & x_{n}\end{bmatrix}R^{-1}\begin{bmatrix}x_1 \\ \ldots \\ x_{n}\end{bmatrix}\right)dx_1\ldots dx_n
\end{align}
Where $|R|$ is the determinant of $R$. This family of copulas will be used in \Cref{sec:sources_of_uncertainty} to model the dependency between the random intensities of pixels of stereo images, for instance. 
\begin{remark}
    The product copula is actually a Gaussian copula with the identity matrix $\mathbb{I}_n$ as its covariance matrix:
    \begin{align*}
        \Phi_R&(\Phi^{-1}(u_1)\enum \Phi^{-1}(u_n)) =\\ &\int_{-\infty}^{\Phi^{-1}(u_1)}\ldots\int_{-\infty}^{\Phi^{-1}(u_n)}\frac{1}{\sqrt{(2\pi)^{n}|\mathbb{I}_n|}}\exp\left(-\frac{1}{2}\begin{bmatrix}x_1 & \ldots & x_{n}\end{bmatrix}\mathbb{I}_n^{-1}\begin{bmatrix}x_1 \\ \ldots \\ x_{n}\end{bmatrix}\right)dx_1\ldots dx_n\\
        =& \int_{-\infty}^{\Phi^{-1}(u_1)}\ldots\int_{-\infty}^{\Phi^{-1}(u_n)}\frac{1}{\sqrt{(2\pi)^{n}}}\exp(-\frac{1}{2}\sum_{i=1}^nx_i^2)dx_1\ldots dx_n\\
        =& \int_{-\infty}^{\Phi^{-1}(u_1)}\frac{1}{\sqrt{(2\pi)}}\exp(-\frac{1}{2}x_1^2)dx_1\ldots\int_{-\infty}^{\Phi^{-1}(u_n)}\frac{1}{\sqrt{(2\pi)}}\exp(-\frac{1}{2}x_n^2)dx_n\\
        =& \int_{-\infty}^{\Phi^{-1}(u_1)}\Phi'(x_1)dx_1\ldots\int_{-\infty}^{\Phi^{-1}(u_n)}\Phi'(x_n)dx_n\\
        =& ~u_1\ldots u_n
    \end{align*}
\end{remark}

\begin{figure}
    \centering
        \begin{tikzpicture}[scale=0.6]
        \begin{axis}[
            xmin=0,xmax=1,
            ymin=0,ymax=1,
            xtick={0, 0.5, 1},
		    xticklabels={$0$, $u_1$, $1$},
		    ytick={0.5, 1},
		    yticklabels={$u_2$, $1$},
		    xtick style={draw=none},
		    ytick style={draw=none},
            title={\L ukasiewicz copula $C_L$},
            ]
            \addplot [domain=0:1,samples=40,style=dashed,color=gray]({x},{1-x});
            \addplot [domain=0:1,samples=40,style=dashed,color=gray]({x},{1.2-x}); 
            \addplot [domain=0:1,samples=40,style=dashed,color=gray]({x},{1.4-x}); 
            \addplot [domain=0:1,samples=40,style=dashed,color=gray]({x},{1.6-x}); 
            \addplot [domain=0:1,samples=40,style=dashed,color=gray]({x},{1.8-x});

            \node[rotate=-45, fill=white, rounded corners=2pt, inner sep=1pt] (x) at (0.5, 0.5) {0};
            \node[rotate=-45, fill=white, rounded corners=2pt, inner sep=1pt] (x) at (0.6, 0.6) {0.2};
            \node[rotate=-45, fill=white, rounded corners=2pt, inner sep=1pt] (x) at (0.7, 0.7) {0.4};
            \node[rotate=-45, fill=white, rounded corners=2pt, inner sep=1pt] (x) at (0.8, 0.8) {0.6};
            \node[rotate=-45, fill=white, rounded corners=2pt, inner sep=1pt] (x) at (0.9, 0.9) {0.8};
        \end{axis}
        \end{tikzpicture}\quad\begin{tikzpicture}[scale=0.6]
        \begin{axis}[
            xmin=0,xmax=1,
            ymin=0,ymax=1,
            xtick={0, 0.5, 1},
		    xticklabels={$0$, $u_1$, $1$},
		    ytick={0.5, 1},
		    yticklabels={$u_2$, $1$},
		    xtick style={draw=none},
		    ytick style={draw=none},
            title={Product copula $C_\Pi$},
            ]
            \addplot[domain=0:1,samples=40,style=dashed,color=gray]({x},{0.1/x});
            \addplot[domain=0:1,samples=40,style=dashed,color=gray]({x},{0.3/x}); 
            \addplot[domain=0:1,samples=40,style=dashed,color=gray]({x},{0.5/x}); 
            \addplot[domain=0:1,samples=40,style=dashed,color=gray]({x},{0.7/x}); 
            \addplot[domain=0:1,samples=40,style=dashed,color=gray]({x},{0.9/x});

            \node[rotate=-45, fill=white, rounded corners=2pt, inner sep=1pt] (x) at (0.32, 0.32) {0.1};
            \node[rotate=-45, fill=white, rounded corners=2pt, inner sep=1pt] (x) at (0.55, 0.55) {0.3};
            \node[rotate=-45, fill=white, rounded corners=2pt, inner sep=1pt] (x) at (0.71, 0.71) {0.5};
            \node[rotate=-45, fill=white, rounded corners=2pt, inner sep=1pt] (x) at (0.84, 0.84) {0.7};
            \node[rotate=-45, fill=white, rounded corners=2pt, inner sep=1pt] (x) at (0.95, 0.95) {0.9};
            
        \end{axis}
        \end{tikzpicture}\quad\begin{tikzpicture}[scale=0.6]
        \begin{axis}[
            xmin=0,xmax=1,
            ymin=0,ymax=1,
            xtick={0, 0.5, 1},
		    xticklabels={$0$, $u_1$, $1$},
		    ytick={0.5, 1},
		    yticklabels={$u_2$, $1$},
		    xtick style={draw=none},
		    ytick style={draw=none},
            title={Minimum copula $C_M$},
            ]
            \addplot [domain=0:1,samples=40,style=dashed,color=gray]({max(x, 0.1)},{max(1-x/0.1,0.1)});
            \addplot [domain=0:1,samples=40,style=dashed,color=gray]({max(x, 0.3)},{max(1-x/0.3,0.3)});
            \addplot [domain=0:1,samples=40,style=dashed,color=gray]({max(x, 0.5)},{max(1-x/0.5,0.5)});
            \addplot [domain=0:1,samples=40,style=dashed,color=gray]({max(x, 0.7)},{max(1-x/0.7,0.7)});
            \addplot [domain=0:1,samples=40,style=dashed,color=gray]({max(x, 0.9)},{max(1-x/0.9,0.9)});

            \node[fill=white, rounded corners=2pt, inner sep=1pt] (x) at (0.55, 0.1) {0.1};
            \node[fill=white, rounded corners=2pt, inner sep=1pt] (x) at (0.65, 0.3) {0.3};
            \node[fill=white, rounded corners=2pt, inner sep=1pt] (x) at (0.75, 0.5) {0.5};
            \node[fill=white, rounded corners=2pt, inner sep=1pt] (x) at (0.85, 0.7) {0.7};
            \node[fill=white, rounded corners=2pt, inner sep=1pt] (x) at (0.95, 0.9) {0.9};
        \end{axis}
        \end{tikzpicture}
    \caption{Bird view of the \L ukasiewicz, product and Min copulas for $n=2$. Dashed gray lines represent isolines of the copulas.}
    \label{fig:copulas}
\end{figure}


{\renewcommand{\arraystretch}{2}%
\begin{table}[!ht]
    \centering
    \begin{tabular}{|c|c|c|c|c|}
        \hline
        Family & $C(u_1,u_2)$ & $\theta \in $ & D-convex & D-concave \\
        \hline\hline
        Ali-Mikhail-Haq & $\frac{u_1u_2}{1-\theta(1-u_1)(1-u_2)}$ & $[-1,1)$ & $\theta\leqslant0$ & $\theta\geqslant 0$ \\
        \hline
        Clayton & $\left[\max(u_1^{-\theta}+u_2^{-\theta}-1,0)\right]^{-1/\theta}$ & $[-1,\infty)\backslash\{0\}$ & $\theta<0$ & $\theta>0$ \\
        \hline
        Frank & $-\frac{1}{\theta}\ln(1+\frac{(e^{-\theta u_1}-1)(e^{-\theta u_2}-1)}{e^{-\theta}-1})$ & $\mathbb{R}\backslash \{0\}$ & $\theta<0$ & $\theta>0$ \\
        \hline
        Gumbel & $u_1u_2\exp(-\theta \ln{u_1}\ln{u_2})$ & $(0,1]$ & $\theta\in(0,1]$  & Never \\
        \hline
    \end{tabular}
    \caption{Examples of families of copulas in the case $n=2$ which can be generated using a parameter \( \theta \). D-convexity/concavity is detailed in the Annex \Cref{sec:dconvexity}}
    \label{tab:family_of_copula}
\end{table}}

\begin{remark}
	As a copula is also a multivariate \acrshort{cdf}, one can imagine an ``imprecise'' copula \cite{montes_sklars_2015} similarly to what can be done with univariate probability distribution (see \Cref{sec:imprecise_probabilities}). Imprecise copulas allow modeling partially known dependencies, but can be hard to manipulate at times, as for instance the lower and upper bounds of an imprecise copula are not necessarily copulas themselves. In this thesis, we will not consider imprecise copulas.
\end{remark}

\subsection{Sampling from a Copula}\label{sec:sampling_copula}
As a copula represents the \acrshort{cdf} of a multivariate random variable, it is possible to sample from it. This section details a method for sampling from copulas in general, and a special method for sampling from Gaussian copulas. \Cref{chap:propagating} will use those methods for Monte Carlo sampling. For simplicity, let us first present a method for sampling in the case $n=2$. Given a copula $C$, and two \acrshort{cdf}s $F_X$ and $F_Y$, a method to generate a pair of observations $(x, y)$ from a joint \acrshort{cdf} $C(F_X, F_Y)$ is the following:

\begin{itemize}
    \item Sample two independent samples $u_1, u_2$ from a uniform distribution on [0,1]
    \item Set $v=\partial C^{-1}(u_2)$ where $\partial C^{-1}$ is the quasi-inverse of the partial derivative of $C$ with respect to its first variable (which exists almost everywhere and is invertible).
    \item We now have a sample $(u_1, v)$ from a multivariate random variable. Its marginals follow a uniform distribution on $[0,1]$, and its associated copula is $C$.
    \item The desired pair is $(x,y) = (F^{-1}_X(u_1), F^{-1}_Y(v))$, with $F_X^{-1}$, $F_Y^{-1}$ being the quasi-inverses of the marginals \acrshort{cdf}s.
\end{itemize}

We do not present the $n$-dimensional general case here as it is a bit more complex, but it can be found in \cite{cherubini_copula_2004}. However, drawing samples from a Gaussian $n$-copula with a correlation matrix $R$ are simpler to obtain:
\begin{itemize}
    \item Compute the Cholesky decomposition $A$ of the correlation matrix $R$
    \item Draw $n$ independent random samples $u=(u_1\enum u_n)'$ from $\mathcal{N}(0,1)$, where $\mathcal{N}$ is the normal distribution.
    \item Set $v=Au$
    \item Set $w_k=\Phi(v_k)$ where $\Phi$ is the univariate normal cumulative distribution function
    \item The desired draw is $(x_1\enum x_n)=(F^{-1}_1(w_1)\enum F^{-1}_n(w_n))$ with $F_i^{-1}$, being the quasi-inverse of the $i$-th marginal \acrshort{cdf}.
\end{itemize}

\begin{conclusion}
    In this chapter, we presented different uncertainty models, whose usage depends on the type of uncertainty and evidence available.  Choosing between the different models is often a trade-off between expressivity, compatibility and ``realism'' of the hypotheses of the considered problem. We also presented a way to combine probabilities by taking into account the dependency between random variables, using copulas. Copulas are well suited for probabilities, but their usage with \acrshort{ip} is more complex. \Cref{chap:joining_credal_sets} will investigate how \acrshort{ip} models can be combined using a copula, and \Cref{chap:propagating} will use copulas and imprecise models in a stereo matching problem.
\end{conclusion}

\clearpage
    
    \setkeys{Gin}{draft=false}
    \chapter{Using Copulas for Joining Credal Sets}\label{chap:joining_credal_sets}
In the presence of multiple sources of uncertainty, the question of how to join those sources arises. Using probability distributions, the usual method for joining probability distributions is to use copulas (section \ref{sec:copulas}). However, there are multiple ways of joining imprecise models using a copula. Sections \ref{sec:robust_method} , \ref{sec:joint_mass} and \ref{sec:aggregation_method} present three different ways of joining credal sets using a copula. Those three methods are not equivalent, we thus explore their similarities and differences in section \ref{sec:inclusions_between_methods}. The methods used to join possibility distributions will serve to propagate the uncertainty in chapter \ref{chap:propagating}.

\section{Methods for Joining Credal Sets with Copulas}\label{sec:methods_for_joining_credal_sets}\commanue{ça manque une peu de liant entre les éléments. Tu enchaines par fois les remarques et c'est au lecteur de trouver tout seul el fil conducteur. Le début de section est bien. Le fait d'avoir des encarts pour les remarques, les proposition etc ça facilite la lecture mais il te faut qq phrase de lien entre.}
\todoroman{Inclure les schema du premier poster qui vulgarisent les méthodes.}
When working with multiple sources of uncertain information, one must take into account the dependency between the different sources in order to correctly determine the joint model. Section \ref{sec:copulas} introduced copulas, a type of dependency model that we will be using for joining the uncertainty model on images intensities.\commanue{ça fait redite avec au dessus. A mon avis tu peux remonter et fusionner} 

When working with precise probabilities, using a copula $C$ to join univariate Cumulative Distribution Functions (CDF), called marginals, into a mutlivariate CDF is straightforward. When working with marginals modeled by Imprecise Probabilities, joining the models with a copula on the product space is not as trivial. We present here three different approaches to create a multivariate credal set using a copula and imprecise marginals. The first approach in section \ref{sec:robust_method} maintain the classical interpretation found of a copula found in Sklar's theorem, but is hard to compute. The second approach in section \ref{sec:joint_mass} takes some distances with the classical interpretation of a copula, but is easier to compute. Finally a third approach is presented in section \ref{sec:aggregation_method}, which completely let go\commanue{alors je ne sais pas si c'est mon côté anti reine des neiges mais let go je sais pas: put aside} of Sklar's theorem interpretation, but is very easy to compute. Inclusion relationships between those three approaches are explored in section \ref{sec:inclusions_between_methods}.

We consider here $n\in\mathbb{N}^*$ uncertain variables $(X_i)_{1\leqslant i\leqslant n}$ taking values respectively in an totally ordered finite space $\X_i$. Index $i$ will usually refer to the $i$-th random variable (or random set). We note $\M_i$ the credal set representing the uncertainty of $X_i$, and $C$ a $n$-copula. Focal sets of a belief function will be noted $a^i_k$, where $k$ refers to the $k$-th focal if they are numbered. We also note $\bigsqcup$ the union of disjoint elements. Finally, we must introduce the concept of cylindrical sets which will come in handy for specifying definition domains of various mappings. 
\begin{definition}[Cylindrical sets]
    Let $\X_1,~\dots,~\X_n$ be $n$ sets and let $\X$\commanue{tu n'as pas un autre symbole car là je vais me mélanger. A moins que tu voulais mettre $\times$} be the Cartesian product of $\X_1,~\dots,~\X_n$. We call a cylindrical (or cylinder) set $X$ of $\X$ a set which can be written as a Cartesian product of elements of $\X_1,~\dots,~\X_n$, \ie:
    \begin{equation}
        X\subseteq\X \text{ is cylindrical} ~\Leftrightarrow~ \exists (X_1,~\dots,~X_n),~\st~X=X_1\tdt X_n \label{eq:cylindrical_sets}
    \end{equation}
\end{definition}

\subsection{Point-wise Aggregation}\label{sec:robust_method}
A first way of creating a joint credal set is to sample a precise marginal for each marginal credal sets $\M_i$ and to use Sklar's theorem with the copula $C$ to create a precise multivariate CDF. The set of all resulting CDFs is thus\commanue{alors j'ai du mal avec la notation, je pense que mon problème vient de l'expression sample a precise marginal for each marginal credal sets. Je dois pas comprendre ce qu'il faut ou disons que je vois pas comment cela se traduit en notation}:
\begin{eqnarray}\label{eq:robust_ancestor}
    \mathcal{S}(C,\M_i) = \{F~|~\forall x_i \in \X_i,~ F(x_1,\dots,x_n) = C(F_1(x_1), \dots, F_n(x_n)),~F_i\in\M_i \}
\end{eqnarray}

This set is not is not guaranteed to be convex \cite{schmelzer_random_2023}. We thus define the joint credal set as the convex hull $\M_{robust}$ of $\mathcal{S}$

\begin{definition}[Robust Credal Set]\label{def:robust_credal_set}
    We define the robust credal set obtained by joining $n$ marginal credal sets $\M_i$ and a copula $C$ as:
    \begin{eqnarray}\label{eq:robust_set}
        \M_{robust}(C,\M_i) = CH(\{C(F_1, \dots, F_n),~F_i\in\M_i \})
    \end{eqnarray}
    where $CH$ is the convex hull presented in definition \ref{def:convex_hull}.
\end{definition}

We refer to this joint credal set as $\M_{robust}(C,\M_i)$ as it contains every element of the marginal credal sets with copula $C$ as their dependency model. We will omit ``$(C,\M_i)$'' when there are no confusion possible to avoid using heavy notations. As we take the convex hull of $\mathcal{S}$, it is interesting to notice that it can contain additional multivariate CDFs that do not possess the copula $C$ as the dependency model of its marginals. This credal set is usually hard to compute for events that are not Cartesian products of marginal cumulative events.

\begin{remark}
    Due to the convex hull, we can both consider $\M_{robust}$ as only containing probabilities on the set $\X=\X_1\tdt\X_n$ and containing probabilities on $2^\X$. This will also be the case for the other credal sets. In general, we will consider the latter, although the former has been considered in the literature (for instance \cite{schmelzer_characterizing_2012, schmelzer_random_2023}).
\end{remark}

Figure \ref{fig:schema_m_robust} represents a schematic of how the lower probability $P_{robust}$\commanue{c'est pas $\low_{robust}$ ?} of $\M_{robust}$ is obtained. First, two samples $F_1$ and $F_2$ (represented by ``$+$'' signs) are sampled from the marginals credal sets defined by $\Bel_1$ and $\Bel_2$. They are then joined using the copula $C$ to create a multivariate $CDF$ represented by a ``x'' sign. This process is then reproduced for every combination of samples in the marginal credal sets, and $\low_{robust}$ is the lower bound of those samples on every event. 
\begin{figure}[!ht]
    \centering
    \begin{tikzpicture}[scale=1]
        \draw [red, line width=0.4mm] plot [smooth cycle, tension=1.1] coordinates {(5,1) (7,2) (5,3) (3,2)};
        \draw [cyan, line width=0.4mm] plot [smooth cycle, tension=2] coordinates {(0,4) (1,3) (2,4) (1,5)};
        
        \draw (7.5, 2) node (Belx) {\color{red}$\Bel_1$};
        \draw (1, 5.5) node (Bely) {\color{cyan}$\Bel_2$};
        
        \draw (4.2, 2.2) node (a) {\textbf{+}};
        \draw (4.5, 1.9) node (A) {$F_1$};
        \draw [gray] (5, 1.2) node (+) {+};
        \draw [gray] (6.8, 2) node (+) {+};
        \draw [gray] (5, 2.8) node (+) {+};
        \draw [gray] (3.2, 2) node (+) {+};
        
        
        \draw (1.3, 4) node (b) {\textbf{+}};
        \draw (0.8, 4) node (B) {$F_2$};
        \draw [gray] (0.2, 3.2) node (+) {+};
        \draw [gray] (0.2, 4.8) node (+) {+};
        \draw [gray] (1.8, 3.2) node (+) {+};
        \draw [gray] (1.8, 4.8) node (+) {+};
        
        \draw (4, 4.5) node (c) {x};
        \draw (5.2, 4.5) node (C) {$C(F_1,~F_2)$};
        \draw[->] (a) to [bend right] node[midway, left, inner sep=0.5cm] {$C$}(c);
        \draw[->] (b) to [bend left] (c);
        
        \draw [gray] (3.8, 4.8) node (c) {x};
        \draw [gray] (6.8, 4.7) node (c) {x};
        \draw [gray] (5.7, 5.7) node (c) {x};
        \draw [gray] (6.3, 3.5) node (c) {x};
        \draw [gray] (5.5, 4) node (c) {x};
        
        \draw [violet, dashed, line width=0.4mm] plot [smooth cycle,tension=1.2] coordinates {(3.7, 4.2) (5, 4) (6.5, 3.5) (6.5, 5.5) (4.3, 5.4)};
        
        \draw (8, 5) node (robust) {\color{violet}{$\underline{P}_{robust}$}};
    \end{tikzpicture}
    \caption{Schematic representation of the lower probability $P_{robust}$ of $\M_{robust}$}\label{fig:schema_m_robust}
\end{figure}

\subsection{Copula Applied to Cumulated Mass Functions}\label{sec:joint_mass}
In this section, we will present another way of creating a joint credal set from multiple marginal ones. Consider the same copula $C$ as before and the same marginal credal sets $\M_i$. Each credal set $\M_i$ is fully determined by a mass distribution function $m_i$, which is strictly positive over its $N_i$ focal sets $a^i_1, \dots, a^i_{N_i}$. As described in \cite{ferson_dependence_2004}, it is possible to use the cumulative mass distribution functions as marginals of the copula to create a joint mass distribution function, granted that there is a complete order defined on the focal sets. Links between copulas and belief functions have been investigated in the continuous case in \cite{schmelzer_joint_2015, schmelzer_multivariate_2019}, the special case of necessity functions in \cite{schmelzer_sklars_2015} and of p-boxes in \cite{schmelzer_random_2023}.

Let suppose, without loss of generality, that the marginal focal sets are numbered according to the order $\preceq_i$: $a^i_1\preceq_ia^i_2\preceq_i\dots\preceq_i a^i_{N_i}$. The idea behind this method is to replace the precise marginal CDFs by cumulative masses, to conserve the philosophy behind Sklar's theorem. We thus first define the joint mass $m_\times$ on the product space of focal sets as follows:

\begin{definition}[Joint Mass]\label{def:joint_mass}
    Let $m_1$, \dots, $m_n$ be mass distribution functions over their respective power sets of $\X_1$, \dots, $\X_n$. Suppose that focal sets in each $\X_i$ are ordered and that $a_{k_i}^i$ is the $k_i$-th focal set of $m_i$ according to the chosen order. We define $m_\times$ as the H-volume of copula $C$ computed over the cumulative marginal masses:
    \begin{eqnarray}\label{eq:joint_mass}
        m_\times(a^1_{k_1}\tdt a^n_{k_n}) = H_{\sum_{k=0}^{k_1-1}m_1(a^1_k), \dots, \sum_{k=0}^{k_n-1}m_n(a^n_k)}^{\sum_{k=0}^{k_1}m_1(a^1_k), \dots, \sum_{k=0}^{k_n}m_n(a^n_k)}
    \end{eqnarray}
    with the convention that $\forall i,\, a^i_0=\emptyset$. It is not \textit{strictly} a focal set but allows to deal with the case $k_i=1$ as $m_i(a^i_0)=0$. For sets that are not of the form $a^1_{k_1}\tdt a^n_{k_n}$, the mass $m_\times$ is null.
\end{definition}

\begin{proposition}
    The function $m_\times$ defined in equation \eqref{eq:joint_mass} is a correctly defined mass distribution function over $\X$. 
\end{proposition}
\begin{proof}
    To be a mass distribution function over $\X$, $m_\times$ must verify the 3 properties of definition \ref{def:mass_distribution_function}.
    
    By construction, it holds that $m_\times(\emptyset)=0$, and the properties of the H-volume impose that $m_\times\in[0,1]$.
    
    There are multiple way of proving that $\sum_{A\subseteq\X}m_\times(A)=1$. A direct proof can be done in the case $n=2$, but the notations become quite heavy for any $n>2$. Instead let us use the interpretation of a copula as a multivariate CDF. This method will be also be used in future proofs.
    
    For all \(i\in\opi 0,\, n\cli\) let $F_i$ be a CDF over $[0,\, N_i]$, with, $F_i(j)=\sum_{k=0}^j m_i(a_k^i)$. By Sklar's theorem, $F=C(F_1,\dots, F_n)$ is a multivariate CDF over $[0,\, N_1]\tdt[0,\, N_n]$. Thus it holds that $F([0,\, N_1]\tdt[0,\, N_n])=1$ and\commanue{Alors là j'ai du mal avec les notation. Après c'est une preuve donc c'est pas grave}:
    \begin{eqnarray*}
        F([0, N_1]\tdt[0, N_n]) &=& F([0, ~\dots, ~0])+\\
        &&F\left(\bigsqcup_{k_1=0}^{N_1-1}\dots\bigsqcup_{k_n=0}^{N_n-1}\left(]k_1, k_1+1]\tdt]k_n, k_n+1]\right)\right)\\
        &=& 0 + \sum_{k_1=0}^{N_1-1}\dots\sum_{k_n=0}^{N_n-1} F(]k_1, k_1+1]\tdt]k_n, k_n+1])\\
        && \text{(CDF of an union of disjoint elements)}\\
        &=& \sum_{k_1=0}^{N_1-1}\dots\sum_{k_n=0}^{N_n-1} H_{F_1(k_1), \dots, F_n(k_n)}^{F_1(k_1+1), \dots, F_n(k_n+1)}\\
        &=& \sum_{k_1=1}^{N_1}\dots\sum_{k_n=1}^{N_n} H_{\sum_{k=0}^{k_1-1}m_1(a^1_k), \dots, \sum_{k=0}^{k_n}m_n(a^n_k)}^{\sum_{k=0}^{k_1}m_1(a^1_k), \dots, \sum_{k=0}^{k_n}m_n(a^n_k)}\\
        &=& \sum_{(a^1_{k_1}\tdt a^n_{k_n})\subseteq\X}m_\times(a^1_{k_1}\tdt a^n_{k_n})
    \end{eqnarray*}
Therefore it holds that:
\begin{eqnarray}
    \sum_{A\subseteq\X}m_\times(A)=\sum_{(a^1_{k_1}\tdt a^n_{k_n})\subseteq\X}m_\times(a^1_{k_1}\tdt a^n_{k_n})=1
\end{eqnarray}
which proves that $m_\times$ is a mass distribution function.
\end{proof}

Having defined a mass distribution function on the product space $\X$, we thus define the joint credal set $\M_{mass}$ and its belief function $\Bel_\times$ as:
\begin{definition}[Joint Belief Function and its Credal Set]\label{def:cumulative_masses_credal_set}
    Let $m_\times$ be the mass distribution function from definition \ref{def:joint_mass}. We note $\Bel_\times$ its joint belief function over the power set of $\X=\X_1\tdt \X_n$, \ie:
    \begin{align}
        \forall A\subseteq\X, \Bel_\times(A)=\sum_{a\subseteq A}m_\times(a)\label{eq:joint_belief}
    \end{align}
    Because $\Bel_\times$ is a belief function, it defines a credal set $\M_{mass}$\commanue{ça me gêne que tu ne mettes qu'une borne à ton credal set tu ne devrais pas aussi mettre la borne sup}:
    \begin{eqnarray}
        \M_{mass}(C,\M_i) = \{~P:2^\X\rightarrow [0,1] ~|~ \forall A\subseteq\X, P(A)\geqslant \Bel_\times(A)~\}\label{eq:credal_set_mass}
    \end{eqnarray}
\end{definition}

Figure \ref{fig:schema_credal_masss} presents a schematic of $\Bel_\times$ similarly to what was presented with $\M_{robust}$ in figure \ref{fig:schema_m_robust}\commanue{j'ai plus de mal à comprendre le schéma dans ce cas, notamment les flêches. Et il faudrait peut-être pas utiliser le même code couleur pour la fonction Belx}. In this case, the copula is applied to the lower envelopes instead of being applied point-wise to every CDF from marginal sets. This can lead to differences between $\M_{robust}$ and $\M_{mass}$ \commanue{Peut-être superposer les deux résultats sur la figure}.
\begin{figure}
    \centering
    \begin{tikzpicture}[scale=1]
        \draw (7.5, 2) node (Belx) {\color{red}$\Bel_1$};
        \draw (1, 5.5) node (Bely) {\color{cyan}$\Bel_2$};
        
        \draw (5.1, 2.9) node (a) {};
        \draw (5, 4.05) node (b) {};
        \draw (1.88, 4.4) node (c) {};
        \draw (3.6, 4.8) node (d) {};
        
        \draw[->] (a) to [bend right] node[midway, left, inner sep=0.5cm] {} (b.south);
        \draw[->] (c) to [bend left] node[midway, left, inner sep=0.5cm] {} (d.south);
        \draw (3, 4) node (e) {$C$};
        
        \draw [red, line width=0.4mm] plot [smooth cycle, tension=1.1] coordinates {(5,1) (7,2) (5,3) (3,2)};
        \draw [cyan, line width=0.4mm] plot [smooth cycle, tension=2] coordinates {(0,4) (1,3) (2,4) (1,5)};
        \draw [violet, dashed, line width=0.4mm] plot [smooth cycle,tension=1.2] coordinates {(3.7, 3.5) (5, 4) (6.5, 4) (6.5, 5.5) (4, 6) };
        
        \draw (7.5, 4.2) node (mass) {\color{violet}$\Bel_\times$};
        
        \draw [gray] (4.2, 2.2) node (+) {+};
        \draw [gray] (5, 1.2) node (+) {+};
        \draw [gray] (6.8, 2) node (+) {+};
        \draw [gray] (5, 2.8) node (+) {+};
        \draw [gray] (3.2, 2) node (+) {+};
        
        \draw [gray] (1.3, 4) node (+) {+};
        \draw [gray] (0.2, 3.2) node (+) {+};
        \draw [gray] (0.2, 4.8) node (+) {+};
        \draw [gray] (1.8, 3.2) node (+) {+};
        \draw [gray] (1.8, 4.8) node (+) {+};
        
        \draw [gray] (4, 4.5) node (x) {x};
        \draw [gray] (3.8, 4.8) node (x) {x};
        \draw [gray] (6.8, 4.7) node (x) {x};
        \draw [gray] (5.7, 5.7) node (x) {x};
        \draw [gray] (6.3, 3.5) node (x) {x};
        \draw [gray] (5.5, 4) node (x) {x};
    \end{tikzpicture}
    \caption{Schematic representation of the lower probability $\Bel_\times$ of $\M_{mass}$} \label{fig:schema_credal_masss}
\end{figure}

With this way of defining the multivariate mass, the choice of arbitrary orders $\preceq_i$ can have a significant impact on the value of the multivariate mass function\commanue{Alors je te crois mais tu saurais l'expliquer avec les mains. Je ne vois pas où l'ordre inflût. Je complète ma remarque, tu donnes une explication dans l'exemple qui suit donc peut-être que tu pourrais a minima pointer vers l'exemple ou là remarque qui suit}. Those order will specifically be ``natural'' orders in sections \ref{subsec:necessity_functions}, \ref{subsec:pboxes} and \ref{subsec:multiple_models}, in the sense where there exists an intuitive total order (inspired by the order on reals). When no natural order exists, the arbitrary choice of the order can greatly impact the output mass or belief functions, as illustrated by the following example.
\todoroman{Quand j'ai $n$ variables j'utilise des notations du style $a_i^n$ (pas trop le choix). Par contre quand j'ai que 2 variables, ça vaudrait peut être le coup de repasser sur des $a_i$ et $b_i$ pour plus de lisibilité ?}\commanue{alors je sais pas je suis habitué à a pour les focal sets si tu décides de les appeler b ça va me perturber}
\begin{example}\label{ex:joint_mass}
    Let us present an example in the case $n=2$. Let $m_1$ be a mass distribution function over the power set $2^{\X_1}$ of $\X_1$ with two focal sets $a_1,\, a_2$. Similarly, let $m_2$ be a mass distribution function over the power set $2^\X_2$\commanue{je pense que tu as un pb de notation car c'est pas pareil que pour le power set 1} with two focal sets $a_1^2,~a_2^2$. Throughout this example we will assume that $a_1^2\preceq_2a_2^2$. We consider the minimum copula $C(u,v)=\min(u,v)$ and that all masses are equal $m_1(a_1^1)=m_1(a_2^1)=m_2(a_1^2)=m_2(a_2^2)=0.5$.
    If there is no natural order on the focal sets of $m_1$, we have to chose an arbitrary one:\par
    $\bullet$ If "$a_1^1\preceq_1a_2^1$" is the arbitrary order (see figure \ref{fig:joint_distrib_arb}), then
    \begin{eqnarray*}
        m_\times(a_1^1,~a_1^2) &=& C(m_1(a_1^1),~m_2(a_1^2))\\
        &=& 0.5\\
        m_\times(a_2^1,~a_1^2) &=& C(1,~m_2(a_1^2)) - C(m_1(a_1^1),~m_2(a_1^2))\\
        &=& m_2(a_1^2) - C(m_1(a_1^1),~m_2(a_1^2))\\
        &=& 0
    \end{eqnarray*}\par
    $\bullet$ If "$a_2^1\preceq_1a_1^1$" is the arbitrary order, then
    \begin{eqnarray*}
        m_\times(a_1^1,~a_1^2) &=& C(1,~m_2(a_1^2)) - C(m_1(a_2^1),~m_2(a_1^2))\\
        &=& m_2(a_1^2) - C(m_1(a_2^1),~m_2(a_1^2))\\
        &=& 0\\
        m_\times(a_2^1,~a_1^2) &=& C(m_1(a_2^1),~m_2(a_1^2))\\
        &=& 0.5
    \end{eqnarray*}
        
    This illustrates that different orders lead to different masses and thus to different credal sets.
    
    \begin{center}
    \includegraphics[width=0.8\linewidth]{Images/Chap_3/M_mass_h_volume.png}\captionof{figure}{Joint mass diagram for different orders.}\label{fig:joint_distrib_arb}
    \end{center}
\end{example}

\begin{remark}
    The reason why $\M_{robust}$ is usually different from $\M_{mass}$ is mainly because of the different orders that can exist on focal sets. Consider for instance the minimum copula already presented in example \ref{ex:copulas}:
    \begin{itemize}
        \item In the precise setting, the minimum copula associates the highest probabilities to events with similar values (high-high or low-low), and the lowest probabilities to events with opposite values (low-high).
        \item In the imprecise setting, the concept of high or low values for focal sets does not usually exist. We thus replace it by an order $\preceq$ on focal sets, determining which set is considered ``low'' and which is ``high'' (regardless of the real values actually contained in the set). The minimum copula then associates the highest mass to joint focal sets with similar ``values'' in the sense of the order $\preceq$, and the lowest mass to sets with opposite values in the sense of the order $\preceq$. For instance in the bivariate case, with marginal focal sets $a_1\preceq_1 a_2$ and $b_1\preceq_2 b_2$, using the minimum copula will assign high masses to $a_1\times b_1$ and $a_2\times b_2$ and low masses to $a_1\times b_2$ and $a_2\times b_1$.
    \end{itemize}
    Assigning a high mass to sets containing both low and high values at the same time is something that would not occur in the precise case but is possbile in the imprecise case. This explains the differences between credal sets $\M_{mass}$ and $\M_{robust}$. 
\end{remark}

We saw that $\M_{robust}$ and $\M_{mass}$ can be distinct credal sets\commanue{Alors c'est peut-être mon niveau d'anglais mais distinct c'est pour ils sont différents ou pour intersection nulle.}. However, because $\M_{robust}$ is difficult to compute, it would be interesting to still use $\M_{mass}$ to approximate it, \ie to verify that $\M_{robust}\subseteq\M_{mass}$ (outer approximation) or $\M_{mass}\subseteq\M_{robust}$ (inner approximation). In example \ref{ex:mass_values}, we show that there is in general no reason for such a relation to exist. Furthermore, if we found an order allowing this relationship, then this order is copula dependent, as changing the copula might break the inclusion. 
\begin{example}\label{ex:mass_values}
    Consider the following setting:
    \begin{itemize}
        \item We consider two spaces $\X_1=\X_2=\{1,~2,~3\}$
        \item We consider two (identical) mass distribution functions $m_1$, $m_2$, each respectively possessing two focal sets $\{2\}$ and $\{1,3\}$.
        \item $m_1(\{2\}) = m_1(\{1,~3\}) = m_2(\{2\})= m_2(\{1,~3\}) = 0.5$
        \item We want to join the credal sets induced by the mass functions using the minimum copula.
    \end{itemize}
    We will compute the bounds of $\M_{mass}$ and $\M_{robust}$ to compare them. The marginals masses being identical and the copula being symmetrical, many results can be obtained by symmetry.
    
    Let us first compute the bounds of $\M_{robust}$. Marginal mass $m_1$ imposes that each marginal probability $P$ will verify:
    \begin{align*}
        \Bel_1(\{2\}) \leqslant &P(\{2\}) \leqslant \Pl_1(\{2\})\\
        0.5 = \sum_{a\subseteq\{2\}}m_1(a) \leqslant &P(\{2\}) \leqslant \sum_{a\cap\{2\}\neq\emptyset}m_1(a) = 0.5
    \end{align*}
    The same result holds for $\{1,~3\}$. We therefore have:
    \begin{align*}
        &P(2)=0.5\\
        &P(1)+P(3)=0.5
    \end{align*}
    And we can compute the same for $m_2$. Looking at those equations, we can deduce that every $P\in\M_{robust}$ verifies\commanue{je suis perdue pourquoi c'est pas $\{2\}$ et $\{1,3\}$, j'ai raté qqch... et après j'avoue que c'est pas évident de se rappeler des formules donc je sais pas comment tu arrives à ce résultat}:
    \begin{align*}
        0 \leqslant ~&P(\{1\},~ \{1\}) ~\leqslant 0.5\\
        0.5 \leqslant ~&P(\{1, ~2\},~ \{1, ~2\}) ~\leqslant 1
    \end{align*}
    
    Let us now compute the bounds of $\M_{mass}$. Choosing an order between $\{1, ~3\}$ and $\{2\}$ is not intuitive. Suppose that there is a reason which encourages us to chose different orders for the focal sets of $m_1$ and for those of $m_2$, so that $\{1,~3\} \preceq_1 \{2\}$ and $\{2\} \preceq_2 \{1,~3\}$. In this case, it holds that\commanue{Alors je sais que c'est censé être l'exemple qui permet de comprendre mais le truc c'est que j'ai du mal à me rappeler des formules donc soit tu renvoies aux équations soit tu rajoutes une ligne pour rappeler la formule}:
    \begin{align*}
        m_\times(\{1, ~3\}, \{2\}) =& ~\min(0.5, ~0.5)\\
        =& ~0.5\\
        m_\times(\{2\}, \{1, ~3\}) =& ~\min(1, ~1) - \min(0.5, ~1)\\
        & - \min(1, ~0.5) + \min(0.5, ~0.5)\\
        =& ~0.5\\
        m_\times(\{2\}, \{2\}) =& ~0\\
        m_\times(\{1, ~3\}, \{1, ~3\}) =& ~0
    \end{align*}
    Thus every probability $P\in\M_{mass}$ will verify\commanue{Même question c'est pas les même focal sets}:
    \begin{align*}
        0 \leqslant ~&P(\{1\},~ \{1\})~ \leqslant 0\\
        0 \leqslant ~&P(\{1, ~2\},~ \{1, ~2\})~ \leqslant 1
    \end{align*}
    
    Looking at the bounds of $\M_{mass}$ and $\M_{robust}$ on cumulative event $\{1\}\times\{1\}$, we can see that $\M_{robust}\not\subseteq \M_{mass}$. Looking at the bounds on $\{1,2\}\times\{1,2\}$, we can see that $\M_{mass}\not\subseteq \M_{robust}$.
    
    \begin{remark}
        When computing the join mass $m_\times$, we observed that sets $\{2\}\times\{2\}$ and $\{1, 3\}\times\{1, 3\}$ have a null mass, while sets $\{2\}\times\{1, 3\}$ and $\{1, 3\}\times\{2\}$ have a mass of $0.5$. It goes against the expected behaviour of the minimal copula in the precise case, which associates high probabilities to similar values (such as high-high or low-low values). This behaviour stems from the fact that we purposefully chosen different orders $\preceq_1$ and $\preceq_2$ on focal sets, while we cannot reverse the orders on reals in the precise case.
        
        It is legitimate to wonder what would happen if we used the same orders. In that case, it would seem that $\M_{mass}\subset\M_{robust}$. However using the \L ukasiewicz copula, it is possible to keep the same order (where $\{1,~3\} \preceq_1 \{2\}$ and $\{1,~3\} \preceq_2 \{2\}$) while still observing that $\M_{robust}\not\subseteq \M_{mass}$ and $\M_{mass}\not\subseteq \M_{robust}$. Indeed the following bounds for every $P$ in $\M_{robust}$ hold and are reached\commanue{donc là perdue à nouveau, c'est pas le résultat (là d'accord si tu veux) mais c'est pas les focal sets énoncés}:
        \begin{align*}
            0 \leqslant ~&P(\{1\}\times\{3\})~ \leqslant 0.5\\
            0 \leqslant ~&P(\{1\}\times\{2\}\cup\{2\}\times\{1\})~ \leqslant 0.5
        \end{align*}
        While $P$ in $\M_{mass}$ verifies:
        \begin{align*}
            0 \leqslant ~&P(\{1\}\times\{3\})~ \leqslant 0\\
            0 \leqslant ~&P(\{1\}\times\{2\}\cup\{2\}\times\{1\})~ \leqslant 1
        \end{align*}
        Finding the values of those bounds follow the same procedure as before but is a much more tedious process, which is thus not detailed here. \comroman{Je n'ai pas inclu les calculs car ils sont un peu longs et fastidieux (il faut expliciter chaque évenement en une somme d'évennements cumulés disjoints, et ensuite simplifier l'expression en utilisant la formulation de la copule). Est-ce que je devrais inclure les calculs en annexe ?}
        
        This means that if we found an order allowing to verify $\M_{robust}\subseteq \M_{mass}$, changing the copula while still keeping the same order might lead to $\M_{robust}\not\subseteq \M_{mass}$.
    \end{remark}
\end{example}

We considered until now that an order has to be chosen arbitrarily. However, special cases of belief functions exhibit a natural order on their focal sets, for instance p-boxes and possibilities. Those special cases will be explored in section \ref{sec:inclusions_between_methods}.

\subsection{Copulas Applied to Belief Functions}\label{sec:aggregation_method}
Another way of joining credal sets with a copula is by directly applying the copula to their lower envelope $\low_i$ for every event:
\begin{definition}[Aggregated Credal Set]\label{def:aggregation_credal_set}
    Given a copula $C$ and $n$ marginal credal sets whose lower probabilities are $\low_1,~\dots,~\low_n$, we define the credal set $\M_{agg}$ over the power set of $\X=\X_1\tdt X_n$ as\commanue{encore une fois tu ne mets qu'une borne alors que dans l'autre définition il y a deux}:
    \begin{eqnarray}
        \M_{agg} = CH(\{P ~|~\forall A_i\subseteq\X_i, P(A_1,\dots, A_n)\geqslant C(\low_1(A_1), \dots, \low_n(A_n)\})\label{eq:copula_on_lower_proba}
    \end{eqnarray}
    where $CH$ is the convex hull from definition \ref{def:convex_hull}.
\end{definition}

Constraints on this set only occur on cylindrical sets of product space $\X_1\tdt \X_n$ (from equation \ref{eq:cylindrical_sets})\commanue{je comprends pas ce que tu veux dire par constraints}, contrary to $\M_{mass}$ or $\M_{robust}$, we thus take the convex hull for extending its definition to probabilities on $\X$\commanue{en fait je ne comprends pas la phrase}. We note $\M_{agg}$ this credal set as it uses the copula as an aggregation operator only, without conserving the meaning associated to copulas by Sklar's theorem. In this regard, $\M_{agg}$ has less meaning than $\M_{robust}$ or $\M_{mass}$, but presents the advantage of being easier to compute on cylindrical sets. Figure \ref{fig:meaning_computation} sums up the performances of the different methods in terms of computation cost and meaningfulness.

\begin{figure}[!hb]
    \centering
    \begin{tikzpicture}[scale=1]
        \draw (0, 0) node (a) {};
        \draw (10, 0) node (b) {};
        \draw (0, 1.5) node (c) {};
        \draw (10, 1.5) node (d) {};
        
        \draw[->, ultra thick, darkgray] (a) to node[midway, left, inner sep=0.5cm] {} (b);
        \draw[->, ultra thick, darkgray] (d) to node[midway, left, inner sep=0.5cm] {} (c);
        
        \draw [darkgray] (9, -0.5) node (z) {\small Easy to compute};
        \draw [darkgray] (0, 2) node (z) {\small Meaningful};
        
        \draw (9, 0.75) node (z) {$\M_{agg}$};
        \draw (5.65, 0.75) node (z) {$\M_{mass}$};
        \draw (2, 0.75) node (z) {$\M_{robust}$};
    \end{tikzpicture}
    \caption{Comparing different methods of joining credal sets with a copula.}
    \label{fig:meaning_computation}
\end{figure}

In general, applying the copula directly to the lower probabilities as in \eqref{eq:copula_on_lower_proba} does not produce a coherent lower probability avoiding sure loss (see definition \ref{def:coherence_sure_loss}). For instance, let us consider $\X_1 = \X_2 = \{1, 2\}$, two lower previsions $\low_1$ and $\low_2$ such that $\low_1(\{1\}) = \low_1(\{2\}) = \low_2(\{1\}) = \low_2(\{2\}) = 0.5$. Joining those two lower probabilities using the minimum copula $C(u,v)=\min(u,v)$ gives a mapping $\low$ which is neither coherent nor avoid sure loss, as presented in Table \ref{tab:non_coherent_lower}\commanue{je ne sais pas si c'est parce que je fatigue mais je vois pas pourquoi is neither coherent nor avoid sure loss}.

\begin{table}[!ht]
    \centering
    \begin{tabular}{|c||c|c|}
        \hline
        \hspace{0.2cm} $\low$ \hspace{0.2cm} & \hspace{0.2cm} $\{1\}$ \hspace{0.2cm} & \hspace{0.2cm} $\{2\}$ \hspace{0.2cm} \\\hline\hline
        $\{1\}$ & $0.5$ & $0.5$ \\\hline
        $\{2\}$ & $0.5$ & $0.5$\\
        \hline
        \end{tabular}
        \caption{$\low = \min(\low_1, \low_2)$}
        \label{tab:non_coherent_lower}
\end{table}

\begin{proposition}
    In the special case of the product copula $C_\Pi$, the joint lower probability $\low_{C_\Pi}$ induced by \eqref{eq:copula_on_lower_proba} avoids sure loss if its marginals also avoid sure loss. It follows that for all copulas $C$ dominated by the product copula (\ie $C_\Pi\geqslant C$), and for every $\M_i$ avoiding sure loss, $\M_{agg}(C, \M_i)$ is a non empty credal set.
\end{proposition}

\begin{proof}
    For $i\in\{1,\dots,n\}$, let $\low_i$ be a lower probability avoiding sure lost, \ie whose credal set $\M_i$ contains at least one probability distribution $P_i$. Let us define a multivariate probability $P$ on every $(A_1\tdt A_n)\subseteq\X$ as:
    \begin{eqnarray*}
        P(A_1\tdt A_n) = P_1(A_1)\tdt P_n(A_n)
    \end{eqnarray*}
    Defining $P$ on $(A_1\tdt A_n)\subseteq\X$ is sufficient as those sets contain every atom of $\X$.
    Because $\forall i, P_i\in\M_i,~P_i\geqslant \low_i$, then:
    \begin{eqnarray*}
        P(A_1\tdt A_n) \geqslant \low_1(A_1)\tdt \low_n(A_n) &=& C_\Pi(\low_1(A_1), \dots, \low_n(A_n))\\
        &=& \low_{C_\Pi}(A_1\tdt A_n)
    \end{eqnarray*}
which means $P\in\M_{agg}(C_\Pi, \M_i)$. Therefore $\M_{agg}(C_\Pi, \M_i)$ avoids sure loss if every $\M_i$ avoids sure loss.

Let $C$ be a copula dominated by $C_\Pi$ (\ie $C_\Pi\geqslant C$), and $\low_C$ the lower probability associated with $\M_{agg}(C, \M_i)$. Then it holds that for all $(A_1\tdt A_n)\subseteq\X$:
\begin{eqnarray*}
\low_{C_\Pi}(A_1\tdt A_n)&=&C_\Pi(\low_1(A_1), \dots, \low_n(A_n)) \\
    &\geqslant& C(\low_1(A_1),\dots, \low_n(A_n)) = \low_C(A_1\tdt A_n)
\end{eqnarray*}
which implies that $\M_{agg}(C_\Pi, \M_i)\subseteq \M_{agg}(C, \M_i)$.
\end{proof}

\begin{proposition}
    Conversely, no lower probability $\low_C$ obtained using \eqref{eq:copula_on_lower_proba} with a copula $C$ strictly superior to the product copula by\commanue{ce by me gêne} is guaranteed to avoid sure loss. It all depends on the marginal credal sets $\low_i$.
\end{proposition}

\begin{proof}
    Let $C$ be a copula strictly superior to the product. Then there exists $(u_1,\dots,u_n)\in[0,1]^2$ such that:
    \begin{eqnarray*}
        C(u_1,\dots, u_n)~>~u_1\dots u_n
    \end{eqnarray*}
    Let $\M_i$ be marginals credal sets such that $\low_i$ are \textit{precise} probabilities, and that:
    \begin{eqnarray*}
        \forall i,~\exists A_i\in\X_i,~\low_i(A_i)=u_i
    \end{eqnarray*}
    Suppose that $\M_{agg}(C, \M_i)$ avoids sure loss, \ie there is a probability $P$ such that $P\geqslant\low_C$. Let $S$ be a collection of disjoint cylindrical sets of $\X$ (defined in equation \eqref{eq:cylindrical_sets}) covering the complementary event $(A_1\tdt A_n)^c$ of $(A_1\tdt A_n)$. $S$ is defined so that $(A_1\tdt A_n)^c=\bigsqcup_{s\in S}s$\.
    Then\commanue{je ne suis pas sure de comprendre pourquoi tu as besoin de passer par S},
    \begin{eqnarray*}
        P(\X) &=& P\left(\left(A_1\tdt A_n\right)\bigsqcup\left(A_1\tdt A_n\right)^c\right)\\
        &=& P\left(A_1\tdt A_n\right)+P\left((A_1\tdt A_n)^c\right)\\
        &=& P\left(A_1\tdt A_n\right)+\sum_{(s_1\tdt s_n)\in S}P(s_1\tdt s_n)\\
        &\geqslant& \low_C\left(A_1\tdt A_n\right)+\sum_{(s_1\tdt s_n)\in S}\low_C(s_1\tdt s_n)\\
        &>& \low_{C_\Pi}\left(A_1\tdt A_n\right)+\sum_{(s_1\tdt s_n)\in S}\low_{C_\Pi}(s_1\tdt s_n)
    \end{eqnarray*}
    Because we chose $\low_i$ so that they are precise probabilities, their product is also a precise probability. Using the fact that summing probabilities of disjoint events is equal to the probability of their union:
    \begin{eqnarray*}
        \low_{C_\Pi}\left(A_1\tdt A_n\right)+\sum_{(s_1\tdt s_n)\in S}\low_{C_\Pi}(s_1\tdt s_n) &=& \low_{C_\Pi}\left(A_1\tdt A_n\right)+\\
        &&\low_{C_\Pi}\left((A_1\tdt A_n)^c\right)\\
        &=& 1
    \end{eqnarray*}
    This means that $P(\X)>1$ which is impossible. Thus $\M_{agg}(C, \M_i)=\emptyset$ and $\M_{agg}(C, \M_i)$ does not avoid sure loss.
\end{proof}

\section{Inclusions Between Joint Credal Sets}\label{sec:inclusions_between_methods}\commanue{Possible inclusion relationships between joint credal sets}
Section \ref{sec:methods_for_joining_credal_sets} presented three methods for joining marginal credal sets using a copula. In the general case, there is no reason for the three methods to lead to the same multivariate credal sets. However, for some specific cases on the copulas or on the marginal credal sets, it is possible to find inclusion relationships between the methods. This section explores some of those specific cases.\commanue{Peut-être ajouter pourquoi c'est pratique et intéressant de mettre en lumière cette relation, toujours dans un esprit de guider le pauvre lecteur physicien vers le grand savoir mathématique. Ce serait bien d'avoir comment on peut ensuite ce servir de ces relations dans une section. Il me semble que cela correspond à une de tes publis tu devrais aussi le mentionner dans l'intro. Remarque générale que je mets ici car sinon ce serait pas drôle mais on va pas forcément ce que tu apportes via ta thèse au chapitre 3. Autre remarque ces cas spécifiques sont des cas idéaux que tu prends ou ce sont des cas qui peuvent plus ou moins fréquemment arriver.} 

\subsection{Using the Product Copula}\label{subsection:product_copula}
In this section, we will consider the case of the product copula $C_\Pi$, representing independence between variables. Using this copula in the robust approach defined by equation \eqref{eq:robust_set} is referred as the strong product in \cite{kacprzyk_factorisation_2010}. Let us denote $\low_{robust}$ the infimum of $\M_{robust}(C_\Pi, \M_i)$ and $\mathcal{S}$ the set from which $\M_{robust}$ is computed (eq. \eqref{eq:robust_ancestor}).
For cylindrical sets $(A_1, \dots, A_n)$ of $\X$, it holds that:
\begin{eqnarray*}
    \low_{robust}(A_1\tdt A_n) &=& \inf\{P(A_1\tdt A_n)~|~P\in\mathcal{S}\}\\
    &=&\inf\{\low_1(A_1)\dots\low_n(A_n)~|~P_i\in\M_i\}\\
    &=&\inf\{\low_1(A_1)~|~P_1\in\M_1\}\dots\inf\{\low_n(A_n)~|~P_n\in\M_n\}\\
    &=&\low_1(A_1)\dots\low_n(A_n)
\end{eqnarray*}
We can split the infimum of a product as a product of infima because we consider mappings with positive values. As this is equivalent of applying the copula directly to the marginals, $\M_{robust}$ and $\M_{agg}$ have the same bounds on cylindrical events. On other events, the lower probabilities are defined as the infimum of the credal sets, thus all bounds are the same and it holds that\commanue{je comprends pas la conclusion aux autres événements. Comment tu passes de c'est égal, à c'est inclus}:
\begin{eqnarray}
    \M_{robust}(C_\Pi, \M_i)\subseteq\M_{agg}(C_\Pi, \M_i)
\end{eqnarray}
This result can also be found in \cite{couso_survey_2000}.
\begin{proposition}\commanue{je trouve que la proposition poppe une peu toute seule. C'est bien d'avoir les encarts pour la lisibilité mais il faudrait les introduire}
    In the case of the product copula $C_\Pi$, the arbitrary orders on the marginal focal sets have no impact on the value of the joint mass $m_\times$ defined in \eqref{eq:joint_mass}. If $a^1_{k_1}, ~\dots, ~a^n_{k_n}$ is a focal set of $m_1, ~\dots, ~m_n$, then $m_\times$ is given by:
    \begin{eqnarray}\label{eq:joint_mass_product}
        m_\times(a^1_{k_1}\tdt a^n_{k_n}) = m_1(a^1_{k_1})\dots m_n(a^n_{k_n})
    \end{eqnarray}
\end{proposition}

\begin{proof}
    Equation \eqref{eq:joint_mass_product} shows that the order on marginal focal sets does not matter in the case of the product copula. We thus show\commanue{have to show / need to show} that equation \eqref{eq:joint_mass_product} holds.
    For simplicity and coherence with the notations of equation \eqref{eq:hvolume}, we will note for all $i\in[0,n]$, $u^i=\sum_{k=0}^{k_i-1}m_i(a_k^i)$, $v^i_k=\sum_{k=0}^{k_i}m_i(a_k^i)$\commanue{il y a des trucs en indice et en exposant mais c'ets pas raccord avec la suite}. $\Pi_{i=1}^n\{u_i, v_i\}$ will refer to the Cartesian product $\{u_1, v_1\}\times\{u_2, v_2\}\tdt\{u_n, v_n\}$ and we will note $C_\Pi$ and $H$ as the product copula and its H-volume regardless of their number of marginals. Those notations established, it holds that:
    \begin{eqnarray*}
        m_\times(a^1_{k_1}\tdt a^{n}_{k_{n}}) &=& H_{u_1, \dots, u_{n}}^{v_1, \dots, v_{n}}\\
        &=&\sum_{\substack{(w_1, \dots, w_{n})\in\\\Pi_{i=1}^{n}\{u_i, v_i\}}}(-1)^{|\{k~|~w_k=u_k\}|}C_\Pi(w_1, \dots, w_{n})\\
        &=&\sum_{\substack{(w_1, \dots, w_{n-1})\in\\\Pi_{i=1}^{n-1}\{u_i, v_i\}}}(-1)^{|\{k~|~w_k=\sum_{k=0}^{k_i-1}m_i(a^i_k)\}|}(w_1 \dots w_{n-1}v_{n})\\
        && + \sum_{\substack{(w_1, \dots, w_{n-1})\in\\\Pi_{i=1}^{n-1}\{u_i, v_i\}}}(-1)^{|\{k~|~w_k=\sum_{k=0}^{k_i-1}m_i(a^i_k)\}|+1}(w_1 \dots w_{n-1}u_n)\commanue{je ne suis pas sure de comprendre cette ligne}\\
        &=& v_{n}H_{u_1, \dots, u_{n-1}}^{v_1, \dots, v_{n-1}} - u_{n}H_{u_1, \dots, u_{n-1}}^{v_1, \dots, v_{n-1}}\\
        &=& m_{n}(a^{n}_{k_{n}})H_{u_1, \dots, u_{n-1}}^{v_1, \dots, v_{n-1}}\commanue{pareil pas compris}
    \end{eqnarray*}
    Doing the same procedure for every variable leads to:
    \begin{eqnarray*}
        m_\times(a^1_{k_1}\tdt a^n_{k_n}) = m_1(a^1_{k_1})\dots m_n(a^n_{k_n})
    \end{eqnarray*}
    which concludes the proof.
\end{proof}

The mass $m_\times$ corresponds to the notion of random set independence presented in \cite{dempster_upper_1967, couso_survey_2000}. Let $\Bel_\times$ be the belief function associated to $m_\times$, and $\forall i\in[1,n], \Bel_i$ the mass function associated to $m_i$. Then for cylindrical sets $(A_1, \dots, A_n)$ of $\X$, it holds that:
\begin{eqnarray}
    \Bel_\times(A_1\tdt A_n) &=& \sum_{(a^1\tdt a^n)\subseteq (A_1, \dots, A_n)}m_\times(a^1\tdt a^n)\nonumber\\
    &=& \sum_{(a^1\tdt a^n)\subseteq (A_1, \dots, A_n)}m_1(a_1)\dots m_n(a^n)\nonumber\\
    &=& (\sum_{a^1\subseteq A_1}m_1(a^1))\dots(\sum_{a^n\subseteq A_n}m_n(a^n))\nonumber\\
    &=& \Bel_1(A_1)\dots \Bel_n(A_n)
\end{eqnarray}

This means that in the case of the product copula $C_\Pi$ with marginals being belief functions, $\M_{robust},~\M_{mass}$ and $\M_{agg}$ all coincide on cylindrical sets. Thus\commanue{je comprends pas le thus. ça coïncide donc inclusion. Si ça coïncide, sur ma planète c'est une égalité. C'ets le même problème que précédemment, je comprends pas la généralisation},
\begin{eqnarray*}
    \M_{mass}\subseteq\M_{agg}
\end{eqnarray*}

\subsection{Using the Natural Ordering of Necessity Functions}\label{subsec:necessity_functions}
Necessity functions $\Nec$, also called minitive belief functions, are a special type of belief function verify $\Nec(A\cap B) =\min(\Nec(A), \Nec(B))$ for all events $A, B$ in their domain of definition. In the multivariate case, when a belief\commanue{function?} is only defined on cylindrical sets, then the minitive property becomes:
\begin{eqnarray}
    &&\forall (A_1\tdt A_n)\subseteq\X, (B_1\tdt B_n)\subseteq\X,\nonumber\\
    &&\Nec(A_1\cap B_1,\dots, A_n\cap B_n) = \min_{(S_1,\dots,S_n)\in\Pi_i\{A_i,B_i\}}(\Nec(S_1, \dots, S_n))
\end{eqnarray}
In \cite{schmelzer_joint_2015}, the author showed that in order to describe the relation between a multivariate belief function and its marginals, in the bivariate case, it is necessary to use a family of sub-copulas: one copula for each tuple of increasing family of events\commanue{je ne suis pas sure de voir ce que cela représente exactement et c'est dommage car tu te sers bcp du terma par la suite}. We remind that a sub-copula is a restriction of a copula to a subset of the unit hyper-cube $[0,1]$ as presented in section \ref{sec:copula_def}.

\begin{theorem}[Sklar's Theorem for Belief Functions \cite{schmelzer_joint_2015}]\label{theorem:sklar_belief}
    Let $\Bel :2^{\X_1}\times2^{\X_2}\rightarrow[0,1]$ be a bivariate belief function and let $\Bel_1$ and $\Bel_2$ denote its marginals over $2^{\X_1}$ and $2^{\X_2}$ respectively. Furthermore, let $\mathcal{I}_1$ and $\mathcal{I}_2$ denote increasing families of subsets of $\X_1$ and $\X_2$. Then there exists a unique sub-copula $C^{\mathcal{I}_1,\mathcal{I}_2}$ on  $\Bel_1(\mathcal{I}_1)\times \Bel_2(\mathcal{I}_2)$ such that:
    \begin{eqnarray}
        \Bel (L_1, L_2) = C^{\mathcal{I}_1,\mathcal{I}_2}(\Bel_1(L_1), \Bel_2(L_2))
    \end{eqnarray}
    for all $L_1\in\mathcal{I}_1,L_2\in\mathcal{I}_2$.
\end{theorem}
For the reverse to be true, it is necessary that $\X_1\in\mathcal{I}_1, \X_2\in\mathcal{I}_2$. Example 1 of \cite{schmelzer_joint_2015} illustrate the need of a copula for each increasing family of events.

Necessity functions are completely determined by their focal sets which form an increasing family of events\commanue{C'est une propriété qui découle de quoi.C'est par définition de necessity function ?}. Thus by applying Sklar's theorem for belief functions (theorem \ref{theorem:sklar_belief}), it holds that joining two necessity functions with a copula $C$ as in \eqref{eq:copula_on_lower_proba} yields a bivariate belief function (which is not necessarily a necessity function):
\begin{equation}
    \Bel = C(\Nec_1, \Nec_2)\label{eq:sklar_on_necessity}
\end{equation}
where $\Nec_1$ and $\Nec_2$ are the marginal necessity functions. The proof of those results where shown in \cite{schmelzer_joint_2015,schmelzer_sklars_2015}. In the following, we will consider that the focal sets $a^i$ of a necessity functions $\Nec_i$ are already ranked using the natural ordering $\preceq_i$\commanue{De ce que j'ai compris d'avant c'était pas facile d'avoir un ordre donc c'est une contrainte forte ou pas que tu imposes.}:
\begin{eqnarray}
    \forall (k,j)\in[1, N_i]^2,~k\leqslant j ~\Leftrightarrow ~ a^i_k \preceq_i a^i_j ~\Leftrightarrow ~ a^i_k\subseteq a^i_j
\end{eqnarray}

\begin{proposition}\label{prop:sklar_necessity}
    Joining two marginal necessity functions $\Nec_1, \Nec_2$ with a copula $C$ as in \eqref{eq:sklar_on_necessity} or using the bivariate mass function as in \eqref{eq:joint_mass} with the natural inclusion ordering yields the same bivariate belief function.
\end{proposition}

\begin{proof}
    If we denote by $\Bel_\times$ the belief function defined in \eqref{eq:joint_mass} where the ordering is the inclusion ordering $\preceq_i$ for $i\in[1,2]$. For convenience and with respects to the notations of equation \eqref{eq:hvolume}, we note: $u^i_k=\sum_{j=0}^{k}m_i(a_j^i)$ and consider that $a^i_0=\emptyset$.
    For all focal elements $a_k^1$ of $\Nec_1$ and $a^2_j$ of $\Nec_2$, it holds that:
    \begin{eqnarray*}
        \Bel_\times(a^1_k, a^2_j) &=& \sum_{a^1_p\subseteq a_k^1}\sum_{a^2_q\subseteq a_j^2}m_\times(a^1_p, a^2_q) = \sum_{p=1}^k\sum_{q=1}^j m_\times(a^1_p, a^2_q)\\
        &=&\sum_{p=1}^k\sum_{q=1}^j (~C(u^1_p, u^2_q) + C(u^1_{p-1}, u^2_{q-1}) \\
        &&- C(u^1_{p-1}, u^2_{q}) - C(u^1_{p}, u^2_{q-1})~)\\
        &=&\sum_{p=1}^k\sum_{q=1}^jC(u^1_p, u^2_q) + \sum_{p=0}^{k-1}\sum_{q=0}^{j-1}C(u^1_p, u^2_q) \\
        &&- \sum_{p=0}^{k-1}\sum_{q=1}^jC(u^1_p, u^2_q) - \sum_{p=1}^k\sum_{q=0}^{j-1}C(u^1_p, u^2_q)\\
        &=& C(u^1_k, u^2_j) = C\left(\Nec_1(a_k^1), \Nec_2(a_j^2)\right)\\
    \end{eqnarray*}
    This is the case for necessity function but not in general where\commanue{Peut-être tourner la phrase différemment en disant que dans le cas général la démonstartion n'est pas valide car la première égalité est fausse. Je trouve que la phrase actuelke porte à confusion.}:
    \begin{align*}
        \sum_{a^1_p\subseteq a_k^1}\sum_{a^2_q\subseteq a_j^2}m_\times(a^1_p, a^2_q) \neq \sum_{p=1}^k\sum_{q=1}^j m_\times(a^1_p, a^2_q)
    \end{align*}
\end{proof}

Proposition (\ref{prop:sklar_necessity}) considers two marginals. However, it still holds for $n$ marginals, not covered in \cite{schmelzer_sklars_2015}.
\begin{proposition}
   Joining $n$ marginal necessity functions $\Nec_1,~\dots,~\Nec_n$ with a n-copula $C$ as in \eqref{eq:sklar_on_necessity} or using the multivariate variate mass function as in \eqref{eq:joint_mass} with the natural inclusion ordering yields the same multivariate belief function. In other words, for every cylindrical set $(A_1, \dots, A_n)\subseteq\X$, it holds that:
\begin{eqnarray}
    \Bel_\times(A_1\tdt A_n) = C\left(\Nec_1(A_1), \dots, \Nec_n(A_n)\right)
\end{eqnarray}
\end{proposition}

\begin{proof}\commanue{Bon là j'avoue c'est le retour des notations où j'ai du mal. Donc à voir si il faut mettre la preuve complète en annexe pour les neuneus comme moi. }
    The proof is similar to the one of equation \eqref{eq:joint_mass}, but this time computing the mass of $(a_{k_1}^1\tdt a_{k_n}^n)$ using $F([0,k_1]\tdt[0,k_n])$ and noticing that \begin{eqnarray*}
        \sum_{(a^1_{p_1}\tdt a^n_{p_n})\subseteq(a^1_{k_1}\tdt a^n_{k_n})}m_\times(a^1_{p_1}\tdt a^n_{p_n})=\sum_{p_1=1}^{k_1}\dots\sum_{p_n=1}^{k_n} m_\times(a^1_{p_1}\tdt a^n_{p_n})
    \end{eqnarray*} because all marginals are necessity functions, and the natural inclusion ordered is used for ranking their focal sets.
\end{proof}

\begin{proposition}\commanue{On se croirait à Dungeon and Dragons: la classe matheux peut faire popper des propositions à l'infini. Encore une fois, je pense qu'on gagnerait en clarté si tu mettais une phrase de transition avant les propositions.}
    Let $m_\times$ be a joint mass obtained using \eqref{eq:joint_mass}. If $m_\times$ verifies
    \begin{align*}
        \sum_{a^1_p\subseteq a_k^1}\sum_{a^2_q\subseteq a_j^2}m_\times(a^1_p, a^2_q) = \sum_{p=1}^k\sum_{q=1}^j m_\times(a^1_p, a^2_q)
    \end{align*}
    for all marginal focal sets $(a^1_k)$ ,$(a^2_j)$ then its marginals masses correspond to necessity functions. The reverse implication is immediate if we use the natural inclusion order.
\end{proposition}
\begin{proof}
    Let $m_\times$ be a joint mass obtained using \eqref{eq:joint_mass}, with marginal focal sets $(a^1_k)_{1\leqslant k\leqslant N_1}$, $(a^2_j)_{1\leqslant j\leqslant N_2}$ and marginal masses $m_1,m_2$. 
    Let $\Bel_\times$ be its associated belief function verifying\commanue{je mettrais bien $\Bel_\times$ dans la ligne en dessous afin que la manière dont est défini $\Bel_\times$ soit clair}
    \begin{align*}
        \sum_{a^1_p\subseteq a^1_k}\sum_{a^2_q\subseteq a_j^2}m_\times(a^1_p, a^2_q) = \sum_{p=1}^k\sum_{q=1}^j m_\times(a^1_p, a^2_q)
    \end{align*}
    for all marginal focal sets $a^1_k,a^2_j$. It is easy to check that\commanue{C'est bizarrre de mettre $m_1(a^1_p)$ au milieu, je le mettrais plutôt au début}
    \begin{align*}
        \sum_{a^2_q\subseteq \X_2}m_\times(a^1_p, a^2_q)=m_1(a^1_p)=\sum_{q=1}^{N_2}m_\times(a^1_p, a^2_q)    
    \end{align*}
    (as we sum the H-volume over a complete partition of $[0,1]$). Thus it holds that:
    \begin{align*}
        \sum_{p=1}^k m_1(a^1_p) = \Bel_\times(a^1_k,\X_2)=\sum_{a^1_p\subseteq a_k^1}m_1(a^1_p)
    \end{align*}
    This result is not sufficient to prove the inclusion of focal sets (there could be a set $a^1_p\subseteq a^1_k,~p>k$ with the same mass value than another set $a^1_{p'}\not\subseteq a^1_k,~p' < k$). Let us show by induction that for all $(k, p)\in\opi1,N_1\cli^2$, $a_1\subset\dots\subset a_k$ and $a_p\not\subseteq a_k$ if $p>k$.
    For the case $k=1$, it holds that:
    \begin{align*}
        \sum_{p=1}^1m_1(a^1_p) &= \sum_{a_p\subseteq a_1}m_1(a^1_p)\\
        \Leftrightarrow m_1(a^1_1) &= m_1(a^1_1) + \sum_{a^1_p\subset a^1_1}m_1(a^1_p)\\
        \Leftrightarrow 0 &=\sum_{a_p\subset a_1}m_1(a^1_p)
    \end{align*}
    which means that no focal set is a strict subset of $a_1$.
    Suppose that there is $k\in\opi1,N_1\cli$ such that: $a_1\subset\dots\subset a_k$ and $\forall p>k,~a_p\not\subseteq a_k$. In particular, $a_{k+1}\not\subseteq a_k$. It holds that:
    \begin{align*}
        \sum_{p=1}^{k+1}m_1(a^1_p) &= \sum_{a^1_p\subseteq a^1_{k+1}}m_1(a^1_p)\\
        \Leftrightarrow m_1(a^1_{k+1}) + \sum_{p=1}^{k} m_1(a^1_p) &= \sum_{a^1_p\subseteq a^1_{k+1}}m_1(a^1_p)\\
        \Leftrightarrow m_1(a^1_{k+1}) + \sum_{a^1_p\subseteq a^1_k} m_1(a^1_p) &= \sum_{a^1_p\subseteq a^1_{k+1}}m_1(a^1_p)\\
        \implies m_1(a^1_{k+1}) &= \sum_{\substack{a^1_p\subseteq a^1_{k+1} \\ a^1_p\not\subseteq a^1_k}} m_1(a^1_p)\\
        \implies m_1(a^1_{k+1}) &= m_1(a^1_{k+1}) + \sum_{\substack{a^1_p\subset a^1_{k+1} \\ a^1_p\not\subseteq a^1_k}} m_1(a^1_p)\\
        \Leftrightarrow 0 &= \sum_{\substack{a^1_p\subset a^1_{k+1} \\ a^1_p\not\subseteq a^1_k}} m_1(a^1_p)
    \end{align*}
    Which means that either there is no focal set that is a strict subset of $a^1_{k+1}$, or that they are all included in $a^1_{k}$. The first case is discarded as:
    \begin{align*}
        \sum_{a^1_p\subseteq a^1_{k+1}}m_1(a^1_p) &= \sum_{p=1}^{k+1}m_1(a^1_p)\\
        \Leftrightarrow \sum_{a^1_p\subset a^1_{k+1}}m_1(a^1_p) &= \sum_{p=1}^{k}m_1(a^1_p)>0
    \end{align*}
    thus $a_1\subset\dots\subset a_{k+1}$.
    Finally, it also holds that:
    \begin{align*}
        \sum_{a^1_p\subseteq a^1_{k+1}}m_1(a^1_p) &= \sum_{p=1}^{k+1}m_1(a^1_p)\\
        \implies \sum_{\substack{a^1_p\subseteq a^1_{k+1} \\ p>k+1}} m_1(a^1_p) + \sum_{\substack{a^1_p\subseteq a^1_{k+1} \\ p\leqslant k+1}} m_1(a^1_p) &= \sum_{p=1}^{k+1}m_1(a^1_p)\\
        \Leftrightarrow \sum_{\substack{a^1_p\subseteq a^1_{k+1} \\ p>k+1}} m_1(a^1_p) &= 0\\
    \end{align*}
    meaning that for all $p>k+1$, $a^1_p\not\subseteq a^1_{k+1}$ which ends the proof by induction. Because all focal sets form a nested family of sets, $\Bel_1$ is a necessity function. The proof for $\Bel $ is identical\commanue{C'est pas $\Bel_2$ et ensuite on utilise le théorème du début}. 
\end{proof}

As the lower probabilities of $\M_{agg}$ and $\M_{mass}$ coincide on cylindrical sets \commanue{c'ets un truc que tu as démontré avant, je le retrouve pas. Ce chapitre 3 est pire que le manuel des joueurs de Dungeons and Dragons.}, the following set equality holds for credal sets $\M_i$ whose lower probabilities are necessity functions\commanue{Bon là c'est surement évident mais pas pour moi}:
\begin{eqnarray}\label{eq:inclusion_necessity}
    \M_{mass}(C, \M_i) \subseteq \M_{agg}(C, \M_i)
\end{eqnarray}

Without further assumptions, there is no inclusion relations between $\M_{robust}$ and $\M_{agg}$ or $\M_{mass}$. The following examples present cases where $\inf\M_{robust}<\inf\M_{agg}$ or $\inf\M_{mass}<\inf\M_{robust}$, proving that it is not always possible to get an (inner or outer) approximation of $\M_{robust}$ using $\M_{mass}$ or $\M_{agg}$.

\begin{example}\label{ex:necessity}
    Let $n=2$. Consider $\X_1=\{x^1_1, x^1_2\}$ and $\X_2=\{x^2_1, x^2_2\}$. Let us define two possibility distribution $\pi_1$ and $\pi_2$ over $\X_1$ and $\X_2$ respectively, such that:

    \begin{eqnarray*}
    \begin{cases}
        \pi_1(x^1_1) = 0.1\\
        \pi_1(x^1_2) = 1
    \end{cases}
    \qquad\text{ and }\qquad
    \begin{cases}
        \pi_2(x^2_1)=1\\
        \pi_2(x^2_2)=0.1
    \end{cases}
    \end{eqnarray*}
    
    For $i\in\{1,2\}$, $\pi_i$ generates a necessity measure $\Nec_i$, a possibility measure $\Pi_i$ and a credal set $\M_i$. Let $P_1$ and $P_2$ be two probabilities respectively included in $\M_1$ and $\M_2$,  whose values are indicated in Table \ref{tab:proba_distrib_1}. 
    
    \begin{center}
    \begin{tabular}{|c|c|c|}
        \hline
        $\X_1$ & $x^1_1$ & $x^1_2$\\
        \hline\hline
        $\Nec_1$ & 0 & 0.9\\
        \hline
        $P_1$  & 0.1 & 0.9\\
        \hline
        $\Pi_1$ & 0.1 & 1\\
        \hline
    \end{tabular}
    \quad
    \begin{tabular}{|c|c|c|}
        \hline
        $\X_2$ & $x^2_1$ & $x^2_2$\\
        \hline\hline
        $\Nec_2$ & 0.9 & 0\\
        \hline
        $P_2$  & 0.9 & 0.1\\
        \hline
        $\Pi_2$ & 1 & 0.1\\
        \hline
    \end{tabular}
    \captionof{table}{Probability distributions over $\X_1$ and $\X_2$}
    \label{tab:proba_distrib_1}
    \end{center}

    
    We first consider here the Minimum copula $C_M(u,v)=\min(u,v)$. After joining $P_1$ and $P_2$ with $C_M$ to obtain a joint probability $P_\times$, let us compare its value with the value of the bivariate necessity function $C_M(\Nec_1, \Nec_2)$ on the same event $\{x_2\}\times\{y_1\}$\commanue{je suis perdue, il vient d'où cet event ?}.
    \begin{eqnarray*}
        \Bel_\times(\{x^1_2\}\times\{x^2_1\}) &=& C_M\left(\Nec_1(\{x^1_2\}), \Nec_2(\{x^2_1\})\right)\\
        &=& \min(0.9,~0.9) = 0.9\\
        P_\times(\{x^1_2\}\times\{x^2_1\}) &=& C_M(P_1(\X_1), P_2(\{x^2_1\})) - C_M(P_1(\{x^1_1\}),P_2(x^2_1))\commanue{je vois pas de quelle précédente formule sort le calcul}\\
        &=& \min(1,~0.9) - \min(0.1,~0.9) = 0.8
    \end{eqnarray*}
    
    We have $P_\times\in\M_{robust}, ~P_\times\not\in\M_{mass}$, which proves that $\M_{robust}\not\subseteq\M_{mass}$\commanue{donc là pareil perdue, je vois pas comment tu obtiens les appertenances de $P_\times$}.
    
    Let us now compare the lower bound of $\underline{P}$\commanue{c'est bizarre la formulation $\underline{P}$, c'est pas déjà un lower bound} of $\M_{robust}$ with that of $\M_{agg}$ but considering the \L ukasiewicz copula $C_L(u,v)=\max(u+v-1,0)$:
    \begin{eqnarray*}
        \Bel_\times(\{x^1_2\}\times\{x^2_1\}) &=& C_L(\Nec_1(\{x^1_2\}), \Nec_2(\{x^2_1\}))\\
        &=& \max(0,~0.9 + 0.9 - 1) = 0.8\\
        \low(\{x^1_2\}\times\{x^2_1\}) &=& \inf_{P_1\in\M_1, P_2\in\M_2}\{C_L(P_1(\{\X_1\}),~P_2(\{x^2_1\})) - C_L(P_1(\{x^1_1\}),\\
        && P_2(\{x^2_1\}))\}\commanue{une petite aide pour retrouver la formule utilisée}\\
        &=& \inf_{P_1\in\M_1, P_2\in\M_2}\{\max(0,~P_1(\X_1) + P_2(\{x^2_1\}) - 1)\\
        && - \max(0,~P_1(\{x^1_1\}) + P_2(\{x^2_1\}) -1)\}\\
        &=& \inf_{P_1\in\M_1, P_2\in\M_2}\{P_2(\{x^2_1\}) - \max(0,~P_1(\{x^1_1\}) + P_2(\{x^2_1\}) - 1 )\}\\
        &=& \inf_{P_1\in\M_1, P_2\in\M_2}\max(P_2(\{x^2_1\}),\\
        && P_2(\{x^2_1\})-P_1(\{x^1_1\}) - P_2(\{x^2_1\}) + 1 )\\
        &\geqslant& \min(\inf_{P_2\in\M_2}\{P_2(\{x^2_1\})\},~\inf_{P_1\in\M_1}\{1-P_1(\{x^1_1\})\}) = 0.9
    \end{eqnarray*}
    
    On this event $\low>Nel_\times$\commanue{c'est pas $\Bel_\times$} and because the belief function is coherent then $\M_{mass}\not\subseteq\M_{robust}$ and $\M_{agg}\not\subseteq\M_{robust}$\commanue{alors là il me manque des étapes pour comprendre le raisonnement}.
\end{example}

\subsection{Using the Natural Ordering of P-boxes}\label{subsec:pboxes}
P-boxes are special cases of belief functions that resemble the most classical CDFs. They are defined with two CDF $\underline{F},~\overline{F}$ such that $\underline{F}\leqslant\overline{F}$, their focal sets $a_\alpha$ are of the form $a_\alpha=[\overline{F}^{-1}(\alpha), \underline{F}^{-1}(\alpha)]$ with $\alpha\in[0,1]$ \cite{destercke_unifying_2008}, where $F^{-1}$ is the inverse of a CDF (or pseudo-inverse if not properly defined). It it thus possible to define a natural order on the focal sets. Let $a_\alpha$ and $a_\beta$ be two focal sets of a p-box $[\underline{F},~\overline{F}]$, with $(\alpha,\beta)\in[0,1]^2$. The natural order $\preceq$ on the focal sets is defined as follows:
\begin{align}
    a_\alpha\preceq a_\beta ~\Leftrightarrow~ \overline{F}^{-1}(\alpha)\leqslant\overline{F}^{-1}(\beta) \text{ and } \underline{F}^{-1}(\alpha)\leqslant\underline{F}^{-1}(\beta)~\Leftrightarrow~ \alpha\leqslant\beta\label{eq:order_pbox}
\end{align}

As it is the case for necessity functions, it is possible to find cases where the sets $\M_{robust}\not\subseteq \M_{mass}$ or $\M_{agg}\not\subseteq\M_{mass}$\commanue{c'est pas robust qui devrait apparaître ?}.
\begin{example}\label{ex:pbox}\commanue{Une petite phrase avant pour éviter les pops d'exemple et de proposition.}
    Consider the following p-boxes:
    \begin{center}
    \begin{tabular}{|c|c|c|}
        \hline
        $\X_1$ & $x^1_1$ & $x^1_2$\\
        \hline\hline
        $\underline{F}_1$ & 0 & 1\\
        \hline
        $\overline{F}_1$ & 0.1 & 1\\
        \hline
    \end{tabular}
    \qquad
    \begin{tabular}{|c|c|c|}
        \hline
        $\X_2$ & $x^2_1$ & $x^2_2$\\
        \hline\hline
        $\underline{F}_2$ & 0.9 & 1\\
        \hline
        $\overline{F}_2$ & 1 & 1\\
        \hline
    \end{tabular}
    \captionof{table}{P-boxes over $\X_1$ and $\X_2$}\label{tab:example_pbox}
    \end{center}
    The p-boxes from Table \ref{tab:example_pbox} lead to the same belief functions as in example \ref{ex:necessity}, thus leading to the same conclusions \ie $\M_{robust}\not\subseteq\M_{mass}$, $\M_{mass}\not\subseteq\M_{robust}$ and $\M_{agg}\not\subseteq\M_{robust}$\commanue{C'est pas les mêmes que dans la phrase juste avant l'example}.
\end{example}

As stated previously, p-boxes are very closely related to CDF which can motivate one to apply Sklar's theorem (Theorem \ref{theorem:sklar}) to the lower CDF and upper CDF respectively. Given $n$ p-boxes $[\underline{F}_1,~\overline{F}_1],~\dots,~[\underline{F}_n,~\overline{F}_n]$ defined over $\X_1,\dots,\X_n$ and a copula $C$, we can define the lower and upper bounds of a $n$ variate CDF as:
\begin{align*}
    \underline{F}_\times&=C(\underline{F}_1,\dots, \underline{F}_n)\\
    \overline{F}_\times&=C(\overline{F}_1,\dots, \overline{F}_n)
\end{align*}
Sklar's theorem states that $\underline{F}_\times$ and $\overline{F}_\times$ are both CDF, which means that $[\underline{F}_\times, \overline{F}_\times]$ is a multivariate p-box \cite{pelessoni_bivariate_2016, montes_sklars_2015} over cylindrical sets, defining a credal set $\M$. Clearly, the bounds of $\M_{robust}$ on cumulative events are the same as those of the credal set $\M$ induced by the multivariate p-box $[C(\underline{F}_1,\dots, \underline{F}_n),~C(\overline{F}_1,\dots, \overline{F}_n)]$\commanue{Clearly pas pour mo mais bon c'est comme ça. Après c'est quoi cumulative event ? Après ce qui est parfois perturbant c'est que temps en temsp c'est sur des cylindrical sets et des fois non. C'ets pas toujours simple de suivre la succession de remarques}.

\begin{proposition}\label{prop:convexity_pbox}
    When joining masses with a copula $C$ using the natural order on p-boxes \eqref{eq:order_pbox}, it holds that:
    \begin{itemize}
        \item if $C$ is D-convex, then $\M_{mass}\subseteq \M_{agg}$.
        \item if $C$ is D-concave then $\M_{agg}\subseteq\M_{mass}$ 
    \end{itemize}
\end{proposition}

\begin{proof}
    Consider $n$ p-boxes $[\underline{F}_1,~\overline{F}_1],~\dots,~[\underline{F}_n,~\overline{F}_n]$, a D-convex copula $C$ and the natural order on focal sets $(a^i_k)_{1\leqslant k \leqslant N_i}$ of each marginal p-box $[\underline{F}_i,~\overline{F}_i]$. We will denote $m_\times$ the joint mass functions obtained using equations \eqref{eq:joint_mass} and $\Bel_\times$ its associated belief function. We will also refer to $\underline{P}$ as the lower probability associated to $\M_{agg}$ from equation \eqref{eq:copula_on_lower_proba}.

    When considering the natural order on focal sets of a p-box \eqref{eq:order_pbox}, it holds that for every focal set $a^i_p$, the set $\{k~|~a^i_k\subseteq a^i_p\}$ is composed of consecutive integers. In the following, we denote by $\underline{p}$ and $\overline{p}$ the lowest and highest indices of the focal sets included in $a^i_p$. This means that $\{a^i_{\underline{p}}, \dots, a^i_{\overline{p}}\}$ is the set of all focal sets included in $a^i_p$. 

    Let $a^1_{p_1},\dots,a^n_{p_n}$ be focal sets of $m_1,\dots,m_n$. We note $u^i_p=\sum_{k=1}^p m_i(a^i_k)$. It then holds that:
    \begin{align*}
        \Bel_\times(a^1_{p_1}, \dots, a^n_{p_n}) &= \sum_{a^1_{k_1}\subseteq a^1_{p_1}}\dots\sum_{a^n_{k_n}\subseteq a^n_{p_n}}m_\times(a^1_{k_1},\dots, a^n_{k_n})\\
        &= \sum_{k_1=\underline{p}_1}^{\overline{p}_1}\dots\sum_{k_n=\underline{p}_n}^{\overline{p}_n}m_\times(a^1_{k_1},\dots, a^n_{k_n})\\
        &= \sum_{k_1=\underline{p}_1}^{\overline{p}_1}\dots\sum_{k_n=\underline{p}_n}^{\overline{p}_n}H_{u^1_{k_1-1},\dots,u^n_{k_n-1}}^{u^1_{k_1},\dots,u^n_{k_n}}
    \end{align*}
    As the H-volume is computed over a partitioning of $[u^1_{\underline{p}_1-1}, u^1_{\overline{p}_1}]\tdt[u^n_{\underline{p}_n-1} , u^n_{\overline{p}_n}]$, it is possible to greatly simplify the sums. The proof is the same as the proof of \eqref{eq:joint_mass} except that the CDF is computed over $[u^1_{\underline{p}_1-1}, u^1_{\overline{p}_1}]\tdt[u^n_{\underline{p}_n-1} , u^n_{\overline{p}_n}]$ and not $[0,1]^n$. This yields\commanue{Bon là je te crois, je laiss qqn d'autre vérifier les indices}:
    \begin{align*}
        \Bel_\times(a^1_{p_1}, \dots, a^n_{p_n}) = H_{u^1_{\underline{p}_1-1},\dots,u^n_{\underline{p}_n-1}}^{u^1_{\overline{p}_1},\dots,u^n_{\overline{p}_n}}
    \end{align*}
    On the other hand, it holds that:
    \begin{align*}
        \underline{P}(a^1_{p_1}, \dots, a^n_{p_n}) &= C(\Bel_1(a^1_{p_1}),\dots,\Bel_n(a^n_{p_n}))\\
        &= C(\sum_{k_1=\underline{p}_1}^{\overline{p}_1} m_1(a^1_{k_1}), \dots, \sum_{k_n=\underline{p}_n}^{\overline{p}_n} m_1(a^n_{k_n}))\\
        &= C(u^1_{\overline{p}_1} - u^1_{\underline{p}_1-1}, \dots, u^n_{\overline{p}_n} - u^n_{\underline{p}_n-1})
    \end{align*}
    Using equation \eqref{eq:convex_diff_hvol} yields:
    \begin{align*}
        \Bel_\times(a^1_{p_1}, \dots, a^n_{p_n}) \geqslant \underline{P}(a^1_{p_1}, \dots, a^n_{p_n})
    \end{align*}
    The inequality is reversed if $C$ is D-concave, which concludes the proof.
\end{proof}

Figure \ref{fig:copula_convex} illustrates the difference between $\Bel_\times$ and $\underline{P}$ in the case $n=2$\commanue{Peut-être commenter un plus donc on retrouve bien la relation démontrer plus haut.}.

\begin{figure}[!ht]
    \centering
    \begin{tikzpicture}
        \begin{axis}[%
            xlabel=$u_1$,
            ylabel=$u_2$,
            xmin=0, xmax=1,
            ymin=0, ymax=1,
            xtick={0.35, 0.55,  0.9, 1},
            xticklabels={$u^1_{\overline{p}_1} - u^1_{\underline{p}_1-1}~~$, $u^1_{\underline{p}_1-1}$, $u^1_{\overline{p}_1}$, $1$},
            ytick={0, 0.3, 0.5,  0.8, 1},
            yticklabels={$0$, $u^2_{\overline{p}_2} - u^2_{\underline{p}_2-1}$, $u^2_{\underline{p}_2-1}$, $u^2_{\overline{p}_2}$, $1$},
            xticklabel style={rotate=30},
            xtick pos=left,
            ytick pos=left,],
            
            \addplot [domain=0:1,samples=40,draw opacity=0.3,color=gray]({x},{1-x});
            \addplot [domain=0:1,samples=40,draw opacity=0.3,color=gray]({x},{1.2-x}); 
            \addplot [domain=0:1,samples=40,draw opacity=0.3,color=gray]({x},{1.4-x}); 
            \addplot [domain=0:1,samples=40,draw opacity=0.3,color=gray]({x},{1.6-x}); 
            \addplot [domain=0:1,samples=40,draw opacity=0.3,color=gray]({x},{1.8-x});

            \node[rotate=-45, fill=white, rounded corners=2pt, inner sep=1pt] (x) at (0.1, 0.9) {\color{gray}0};
            \node[rotate=-45, fill=white, rounded corners=2pt, inner sep=1pt] (x) at (0.3, 0.9) {\color{gray}0.2};
            \node[rotate=-45, fill=white, rounded corners=2pt, inner sep=1pt] (x) at (0.5, 0.9) {\color{gray}0.4};
            \node[rotate=-45, fill=white, rounded corners=2pt, inner sep=1pt] (x) at (0.7, 0.9) {\color{gray}0.6};
            \node[rotate=-45, fill=white, rounded corners=2pt, inner sep=1pt] (x) at (0.9, 0.9) {\color{gray}0.8};
            
            
            \node (a) at (0.9, 0.8) {};
            \node (b) at (0.55, 0.8) {};
            \node (c) at (0.55, 0.5) {};
            \node (d) at (0.9, 0.5) {};
            \draw[ultra thick, black] (c.center) rectangle (a.center) node[pos=0.5]{$\Bel_\times$};
            
            \node (e) at (0, 0.8) {};
            \node (f) at (0, 0.5) {};
            \node (g) at (0.9, 0) {};
            \node (h) at (0.55, 0) {};
            \draw[thick, dashed, black] (e.center) -- (b.center);
            \draw[thick, dashed, black] (f.center) -- (c.center);
            \draw[thick, dashed, black] (g.center) -- (d.center);
            \draw[thick, dashed, black] (h.center) -- (c.center);
            
            \node (i) at (0.35, 0.3) {};
            \node (j) at (0, 0) {};
            \draw[ultra thick, black] (i.center) rectangle (j.center) node[pos=0.5]{$\underline{P}$};
            
            \node (m) at (0.55, 0.45) {};
            \node (n) at (0.9, 0.45) {};
            \node (o) at (0.5, 0.8) {};
            \node (p) at (0.5, 0.5) {};
            \draw [stealth-stealth, thick, black] (m.center) -- (n.center) node[pos=0.5, below]{$u^1_{\overline{p}_1} - u^1_{\underline{p}_1-1}$};
            \draw [stealth-stealth, thick, black] (o.center) -- (p.center) node[pos=0.5, left]{$u^2_{\overline{p}_2} - u^2_{\underline{p}_2-1}$};
        
        \end{axis}
    \end{tikzpicture}
    \caption{Bird view of the \L ukaciewicz 2-copula $C_L$, where the gray lines are the isolines of the copula. $\Bel_\times$ and $\underline{P}$ are represented in the case where the marginals are p-boxes. The thick rectangles represent the bounds on which to compute the H-volume. Numbers $u^i_k$ use the notation of the proof of proposition \ref{prop:convexity_pbox}.}
    \label{fig:copula_convex}
\end{figure}

\subsection{Joining Different Types of Models}\label{subsec:multiple_models}\commanue{Je ne sais pas si on peut pas améliorer le titre, j'ai compris de ce dont on allait parler grâce à l'intro qui suit}
Different models\commanue{juste le terme models tout seul ça fait un peu vague} can be used for\commanue{representing / characterizing ?} multiple sources of uncertainty. When some marginals are modeled using possibility distributions and other by p-boxes, it is still possible to derive results similar as those of sections \ref{subsec:necessity_functions} and \ref{subsec:pboxes} if we use natural orders on the marginal focal sets\commanue{On ne peut mixer que possiblity et p-boxes, ou on peut faire la cuisine que l'on veut si on a un ordre naturel. Je dirais que l'intro permet de comprendre ce dont tu vas parler donc c'est un point positif. Mais je trouve que tu pourrais détailler un peu plus, c'est un peu rapide.}. 
\begin{proposition}\commanue{C'est le retour du sort Apparition de proposition de la classe matheux, c'est un sort qui donne 2d6 de maux de tête si on rate sa sauvegarde aux tests de maths. C'est vraiment pas cool comme sort ! }
    When joining marginal credal sets induced by a mix of possibility distributions and p-boxes using a copula $C$ the following inclusion hold:
    \begin{itemize}
        \item if $C$ is D-convex, then $\M_{mass}\subseteq \M_{agg}$.
        \item if $C$ is D-concave then $\M_{agg}\subseteq\M_{mass}$ 
    \end{itemize}
\end{proposition}

\begin{proof}
    The proof is similar to the proof of proposition \ref{prop:convexity_pbox} using the fact that for every focal set $a^i_p$ of a possibility distribution, we can still define $\underline{p}_i$ and $\overline{p}_i$ as $\underline{p}_i=1$ and $\overline{p}_i=p$. The rest of the proof is identical.
\end{proof}

\subsection{Using Other Orders}\commanue{donc là si j'ai bien compris c'est pas l'ordre naturel mais même modèles pour représenter l'incertitude. A voir si il faut pas adapetr le titre ou au moins l'intro car on comprend les conditions dans lesquelles tu te places avec l'exemple, c'est un peu tard.}
Instead of considering the natural ordering, or when such an order does not exists\commanue{Je me dis que tu gagnerais en clarté si tu donnais des exemples. Ou si tu donnais des éléments au lecteur comme est-ce que tel cas arrive souvent.}, one could consider an arbitrary order between focal sets when defining $\M_{mass}$ as in \eqref{eq:joint_mass}. A few questions arise: is there always an arbitrary order allowing $\M_{robust}\subseteq\M_{mass}$? If such an order exists, is it possible to explicit it in advance, \ie without computing the lower bounds of the credal sets\commanue{là je suis perdue sur cette remarque. Je comprends la phrase précédente car $\M_{mass}$ est plus simple à calculer. En plus il y a des s dans ta phrase qui me laissent perplexe.}\commanue{On peut pas montrer non plus d'autres realtions ? Ces travaux viennent de ton papier. Donc en gros conclusion si pas d'ordre naturel on est mal donc ?}

It appears that an order allowing $\M_{robust}\subseteq\M_{mass}$ does not always exist. To prove it, let us find an example where no order allows either inclusion.
\begin{example}
Consider the Clayton copula for $\theta=2$ and $n=2$:
\begin{eqnarray*}
    \forall (u_1,u_2)\in\mathbb{R}^2\backslash(0,0), ~C(u_1,u_2)=\frac{u_1u_2}{\sqrt{{u_1}^2+{u_2}^2-{u_2}^2{u_2}^2}}
\end{eqnarray*}
and $C(0,0)=0$ by continuity (we simplified the expression of the copula given in table \ref{tab:familiy_of_copula}). Let us consider $\X_1=\{x^1_1,\, x^1_2,\, x^1_3\}$, $\X_2=\{x^2_1,\, x^2_2,\, x^2_3\}$ and two possibility distributions $\pi_1$, $\pi_2$ over $\X_1$ and $\X_2$ respectively:
\begin{eqnarray*}
    \pi_1(x^1_1)=\pi_2(x^2_1)=0.2 \qquad \pi_1(x^1_2)=\pi_2(x^2_2)=1 \qquad \pi_1(x^1_3)=\pi_2(x^2_3)=0.7 
\end{eqnarray*}
and the marginal credal sets $\M(\pi_1)$, $\M(\pi_2)$ they induce. We note the focal sets as $a^1_1=\{x^1_2\},~a^1_2=\{x^1_2,\, x^1_3\},~a^1_3=\{x^1_1,\,x^1_2,\,x^1_3\}$ and $a^2_1=\{x^2_2\},~a^2_2=\{x^2_2,\, x^2_3\},~a^2_3=\{x^2_1,\, x^2_2,\, x^2_3\}$. By joining $\M(\pi_1)$ and $\M(\pi_2)$ using $C$, we can obtain the lower probability $\low$ of $\M_{robust}$ using equation \eqref{eq:robust_set}, and a total of 6 belief functions $\Bel^{\preceq_1,\preceq_2}_\times$ using equation \eqref{eq:joint_mass} depending on the orders $\preceq_1, \preceq_2$ used to join the marginal masses. For instance if $\preceq_1$ is such that $a^1_3\preceq_1 a^1_1 \preceq_1 a^1_2$ and $\preceq_2$ such that $a^2_1 \preceq_2 a^2_2 \preceq_2 a^2_3$, then the bivariate mass $m_\times$ would be defined as follows:
\begin{eqnarray*}
    m_\times(a^1_3\times a^2_1) &=& C(m_1(a^1_3), m_2(a^2_1))\\
    m_\times(a^1_1\times a^2_1) &=& C(m_1(a^1_3)+m_1(a^1_1), m_2(a^2_1)) - C(m_1(a^1_3), m_2(a^2_1))\\
    m_\times(a^1_3\times a^2_2) &=& C(m_1(a^1_3), m_2(a^2_1) + m_2(a^2_2)) - C(m_2(a^1_3), m_2(a^2_1))\\
    &\dots&
\end{eqnarray*}

If we consider the two events $E_1=\{x^1_2\}\times\{x^2_2,~x^2_3\}$ and $E_2=\{x^1_2,~x^1_3\}\times\{x^2_2\}$, it is possible to show that $\low(E_1)=\low(E_2)\approx0.131$ (by symmetry of the problem) which is obtained for:
\begin{align*}
    &P_1(x^1_1)=0, \qquad&P_1(x^1_2)=0.3, \qquad&P_1(x^1_3)=0.7,\\
    &P_2(x^2_1)=0.2, \qquad&P_2(x^2_2)=0.3, \qquad&P_2(x^2_3)=0.5
\end{align*}
for $E_1$, and similarly with $P_1$ and $P_2$ reversed for $E_2$. Those values were estimated by running simulations, but their exact value can be computed by solving an optimization problem as $C$ is differentiable (although it is a bit tedious to compute). Then for all orders $\preceq_1,~\preceq_2$ on the focal sets of $\pi_1$ and $\pi_2$ it holds:

 \begin{align*}
    \begin{cases}
        \begin{split}
            \Bel_\times^{\preceq_1,\preceq_2}(E_1) \leqslant& \low(E_1)\\
            \Bel_\times^{\preceq_1,\preceq_2}(E_2) >& \low(E_2)
        \end{split}
    \end{cases}
    \qquad\text{ or }\qquad
    \begin{cases}
        \begin{split}
            \Bel_\times^{\preceq_1,\preceq_2}(E_1) >& \low(E_1)\\
            \Bel_\times^{\preceq_1,\preceq_2}(E_2) \leqslant& \low(E_2)
        \end{split}
    \end{cases}
\end{align*}

The table \ref{tab:beliefs_orders} shows a rounded value of $\Bel^{\preceq_1,\preceq_2}_\times$ for all orders for events $E_1$, $E_2$.

\begin{center}
\begin{tabular}{|c||c|c|c|}
\hline
$\Bel^{\preceq_1,\preceq_2}_\times(\mathbf{E_1})$ & $a^2_3\preceq_2a^2_2\preceq_2a^2_1$ & $a^2_3\preceq_2a^2_1\preceq_2a^2_2$ & $a^2_2\preceq_2a^2_3\preceq_2a^2_1$ \\ \hline\hline
$a^1_3\preceq_1a^1_2\preceq_1a^1_1$ & 0.296 & 0.296 & 0.224 \\ \hline
$a^1_3\preceq_1a^1_1\preceq_1a^1_2$ & 0.254 & 0.254 & 0.240 \\ \hline
$a^1_2\preceq_1a^1_3\preceq_1a^1_1$ & 0.296 & 0.296 & 0.224 \\ \hline
$a^1_1\preceq_1a^1_3\preceq_1a^1_2$ & \textbf{0.131} & \textbf{0.131} & 0.279 \\ \hline
$a^1_2\preceq_1a^1_1\preceq_1a^1_3$ & 0.291 & 0.291 & 0.216 \\ \hline
$a^1_1\preceq_1a^1_2\preceq_1a^1_3$ & \textbf{0.131} & \textbf{0.131} & 0.279 \\ \hline
\end{tabular}

\vspace{0.5cm}

\begin{tabular}{|c||c|c|c|}
\hline
$\Bel^{\preceq_1,\preceq_2}_\times(\mathbf{E_1})$ & $a^2_1\preceq_2a^2_3\preceq_2a^2_2$ & $a^2_2\preceq_2a^2_1\preceq_2a^2_3$ & $a^2_1\preceq_2a^2_2\preceq_2a^2_3$ \\ \hline\hline
$a^1_3\preceq_1a^1_2\preceq_1a^1_1$ & 0.259 & 0.180 & 0.180 \\ \hline
$a^1_3\preceq_1a^1_1\preceq_1a^1_2$ & 0.208 & 0.270 & 0.270 \\ \hline
$a^1_2\preceq_1a^1_3\preceq_1a^1_1$ & 0.259 & 0.180 & 0.180 \\ \hline
$a^1_1\preceq_1a^1_3\preceq_1a^1_2$ & 0.251 & 0.293 & 0.293 \\ \hline
$a^1_2\preceq_1a^1_1\preceq_1a^1_3$ & 0.236 & 0.218 & 0.218 \\ \hline
$a^1_1\preceq_1a^1_2\preceq_1a^1_3$ & 0.251 & 0.293 & 0.293 \\ \hline
\end{tabular}

\vspace{1cm}

\begin{tabular}{|c||c|c|c|}
\hline
$\Bel^{\preceq_1,\preceq_2}_\times(\mathbf{E_2})$ & $a^2_3\preceq_2a^2_2\preceq_2a^2_1$ & $a^2_3\preceq_2a^2_1\preceq_2a^2_2$ & $a^2_2\preceq_2a^2_3\preceq_2a^2_1$ \\ \hline\hline
$a^1_3\preceq_2a^1_2\preceq_2a^1_1$ & 0.296 & 0.254 & 0.296 \\ \hline
$a^1_3\preceq_1a^1_1\preceq_1a^1_2$ & 0.296 & 0.254 & 0.296 \\ \hline
$a^1_2\preceq_1a^1_3\preceq_1a^1_1$ & 0.224 & 0.240 & 0.224 \\ \hline
$a^1_1\preceq_1a^1_3\preceq_1a^1_2$ & 0.259 & 0.208 & 0.259 \\ \hline
$a^1_2\preceq_1a^1_1\preceq_1a^1_3$ & 0.180 & 0.270 & 0.180 \\ \hline
$a^1_1\preceq_1a^1_2\preceq_1a^1_3$ & 0.180 & 0.270 & 0.180 \\ \hline
\end{tabular}

\vspace{0.5cm}

\begin{tabular}{|c||c|c|c|}
\hline
$\Bel^{\preceq_1,\preceq_2}_\times(\mathbf{E_2})$ & $a^2_1\preceq_2a^2_3\preceq_2a^2_2$ & $a^2_2\preceq_2a^2_1\preceq_2a^2_3$ & $a^2_1\preceq_2a^2_2\preceq_2a^2_3$ \\ \hline\hline
$a^1_3\preceq_2a^1_2\preceq_2a^1_1$ & \textbf{0.131} & 0.291 & \textbf{0.131} \\ \hline
$a^1_3\preceq_1a^1_1\preceq_1a^1_2$ & \textbf{0.131} & 0.291 & \textbf{0.131} \\ \hline
$a^1_2\preceq_1a^1_3\preceq_1a^1_1$ & 0.279 & 0.216 & 0.279 \\ \hline
$a^1_1\preceq_1a^1_3\preceq_1a^1_2$ & 0.251 & 0.236 & 0.251 \\ \hline
$a^1_2\preceq_1a^1_1\preceq_1a^1_3$ & 0.293 & 0.218 & 0.293 \\ \hline
$a^1_1\preceq_1a^1_2\preceq_1a^1_3$ & 0.293 & 0.218 & 0.293 \\ \hline
\end{tabular}
\captionof{table}{Value of $\Bel^{\preceq_1,\preceq_2}_\times$ for $E_1$ and $E_2$ depending on the arbitrary orders $\preceq_1$, $\preceq_2$. Values in bald font represent the minimal value attained by the different belief functions, where $\Bel^{\preceq_1,\preceq_2}_\times(E)=\low(E)$.}
\label{tab:beliefs_orders}
\end{center}

Answering the question ``if an order allowing $\M_{robust}\subseteq\M_{mass}$ exists, is it possible to explicit it in advance?'' is not as trivial. Indeed the order will be dependent of the copula. We ran simulations where we created different possibility distributions (for $card(\X_1)=card(\X_2)\in\{3,4\}$), sampled the marginal credal sets and joined them with a copula $C$ to get $\inf\M_{robust}$ and compared its values with the one of $\Bel^{\preceq_1,\preceq_2}$.
\begin{itemize}
	\item When running the simulation with the lower Fréchet–Hoeffding copula $C(u,v)=\max(0,u+v-1)$, we observe that a specific couple of order always works: $\preceq_1:~a^1_3\preceq_1a^1_2\preceq_1a^1_1$ and $\preceq_2~a^2_3\preceq_2a^2_2\preceq_2a^2_1$, while other orders do not permit the inclusion $\M_{robust}\subseteq\M_{mass}$ all the time.
\end{itemize}
The order is thus copula dependent.

We ran simulation for the lower and upper Fréchet-Hoeffding copulas, as well as with the copulas presented in table \ref{tab:familiy_of_copula}. We were only able to find examples where no orders were working for the Ali-Mikail-Haq copula with $\theta>0$, for the Gumbel copula and for Clayton copulas. 
\end{example}
\commanue{Alors quoi ? pas de conclusion ? C'est quoi cette fin de chapitre digne d'un vieux téléfilm sur france 3. Il faut résumer les éléments clés du chapitre et donner envie à ton lecteur qui vient de subir une épreuve de continuer.}
\pagebreak
\blankpage
    
    \setkeys{Gin}{draft=false}
    \chapter{Propagating the Uncertainty from Stereo Images to the Cost Volume}\label{chap:propagating}
The previous chapter presented different methods for joining credal sets using a copula\commanue{Alors moi je tournerais la phrase au passif In the previous chapter, different methods for joining credal sets using a copula have been presented}. In this chapter, we will see how those results can be applied to propagate the uncertainty in a stereo matching example. More specifically, we will use possibility distributions from \Cref{chap:representation_of_uncertainty} as uncertainty models on the intensity values of epipolar images used in stereo matching. We will then propagate the uncertainty from those images to the cost volume using results from \Cref{sec:necessity_functions}. We will also show that propagating the uncertainty has the potential to improve the disparity map derived from the cost volume. This chapter takes up work and data already published \cite{malinowski_copulas_2022, malinowski_uncertainty_2023}.

It is important to note that in this disparity estimation problem, we only account for the uncertainty in our input image intensities, without considering the uncertainty in our cost function's ability to correctly identify the true disparity as its minimum. In other words, we do not account for the uncertainty arising from the difference between ``two patches are very similar'' and ``the pixels at the center of the patches are homologous''. We refer to \Cref{fig:adherence_window} and the discussion in \Cref{sec:stereo_matching} regarding the adherence problem for more details. We focus in this chapter on the propagation of uncertainty. We will thus consider a stereo matching pipeline with a simple cost function and no \acrshort{sgm} regularization. Indeed, propagating the uncertainty through a cost volume optimization is too complex and computationally expensive to be solved with this chapter's method. We will consider the different problem of uncertainty modelling with \acrshort{sgm} regularization in the following chapter, \Cref{chap:epistemic_uncertainty}. 

\section{Context and Hypotheses for Uncertainty Propagation }\label{sec:sources_of_uncertainty}
\subsection{Considered Stereo Matching Pipeline}
We consider the Sum of Absolute Differences (SAD) as our cost function, presented in \Cref{sec:cost_volume_computation}, and reminded here. Given patches $W_L\subset I_L$ and $W_R\subset I_R$ of the same shape with $n$ pixels (usually squares):
\begin{align}
    \SAD(W_L, W_R) = \sum_{(p_i, q_i)\in (W_L, W_R)}|I_L(p_i) - I_R(q_i)|\label{eq:SAD}
\end{align}
where $p_i$ and $q_i$ are pixels at the same position $i$ in their patch. For convenience purposes, we will refer to the Absolute Difference between two pixels as AD (when there is no sum involved). An illustration of windows from stereo images to compare is presented in \Cref{fig:Cones}, and an illustration of the AD between pixels and $\SAD$ cost function can be found in \Cref{fig:SAD}.

\begin{figure}[ht]
  \centering
  \includegraphics[width=0.8\linewidth]{Images/Chap_4/Cones.png}
  \caption{Homologous pixels in a pair of images. From \cite{malinowski_uncertainty_2024}}\label{fig:Cones}
\end{figure}

\begin{remark}
    Although the $\SAD$ is not the best performing cost function for dense matching, it is both fast and easily parallelizable. It is often use for comparison in cost-based stereo algorithms \cite{hirschmuller_evaluation_2007, zbontar_stereo_2016}, or in other applications. Considering a simple cost function such as $\SAD$ is relevant for different reasons:
    \begin{itemize}
        \item Simple stereo matching algorithms are often used as a quick and easy method for estimating the disparity.
        \item Simple cost functions (such as $\SAD$, ZNCC, \etc) considered here are still used in other problems, such as video compression for instance \cite{richardson_h264_2006}.
        \item As both stereo matching and uncertainty propagation can be complicated problems, considering a simple stereo algorithm allow us to remain (relatively) simple and didactic in our explanations. 
    \end{itemize}
\end{remark}

As stated previously, we will not consider in this chapter the use of Semi Global Matching (or similar regularization methods), as it necessitates many operations due to its recursive formulation and thus greatly increases complexity for the uncertainty propagation. We diverge from state-of-the-art methods to keep explanations as simple as possible, however the modeling of uncertainty for more advanced cost functions and \acrshort{sgm} methods will be considered in \Cref{chap:epistemic_uncertainty}.

\begin{figure}
    \centering
    \includegraphics[width=0.5\linewidth]{Images/Chap_4/SAD.png}
    \caption{Diagram representing the $SAD$ cost function between two $3\times3$ patches. From \cite{malinowski_uncertainty_2024}}
    \label{fig:SAD}
\end{figure}

In our experiments, we considered the ``Cones'' images from the 2003 Middlebury stereo dataset (\url{https://vision.middlebury.edu/stereo/data/scenes2003/}), as presented in \Cref{fig:color_cones_image}. The two images have a of size $375 \times 450$, and the range of considered disparities is $[-60, 0]$.

\subsection{Uncertainty Model for Epipolar Images Intensities}
To maintain simplicity in this section, we will not consider panchromatic images, such as Pléiades products, encoding the reflectance values as positive integer, usually contained in $[0, 5000]$. Instead, we consider grayscale images that have intensity levels quantified within the range $[0, 255]$, which will represent our measurable space $\X$.

\begin{remark}
    This hypothesis is not constraining as we can easily normalize reflectances in order to encode them using 8-bit integers value, although doing so reduces the precision of the initial images. Moreover, this normalization step is often required to fit a given format, for instance if the images must be processed by a CNN trained on 8-bit integers.
\end{remark}

We hypothesize that a pixel's intensity value can deviate around its observed value with a range of $\pm i_\sigma$, with the observed value being the most likely. This specific hypothesis remains simple and relatively plausible with regards to the processing leading to the epipolar images. We then suppose that the uncertainty from the noise of the sensor capturing the image, from pre-processing steps such as atmospheric correction\commanue{alors cette étape je ne sais, après tu as tout de même des étapes de corrections radiométriques et géométriques} or epipolar resampling (see \Cref{sec:uncertainty_cars}) or from the quantification of observed radiometric values into integers\commanue{je ne suis pas sure de comprendre}, are not exactly known, but can be model by a possibility distribution. Consequently, we model the uncertainty of each pixel $p\in I_L,I_R$ intensity with a possibility distribution $\pi$, centered around the observed intensity $i_p\in[0,255]$:
\begin{equation}\label{eq:pixel_possibility}
    \pi(i_p)=1,\quad \pi(i_p\pm i_\sigma)=\alpha\,,
\end{equation}
with $\alpha \in [0,1]$. To remain simple, we chose $i_\sigma=1$ in the following, which is relatively narrow but allows simplification without impacting our reasoning. Similarly, we do not consider multiple $\alpha$ values in order to limit the number of focal sets considered. Choosing thus simple possibility distribution is convenient, but results obtained in this chapter can easily be extended to more complex possibility distributions. In our simulation, $\alpha = 0.3$ for pixels in the left image and $\alpha = 0.4$ for pixels in the right image\commanue{Je ne sais pas si tu synthétises les valeurs quelque part mais ce serait cool, ça peut être un tableau en annexe, sinon ça oblige à relire la section pour retrouver les valeurs}. We use different values of $\alpha$ for the left and right images because the uncertainty model may vary between images due to differences in exposure, noise levels, or camera calibration. The values on themselves are chosen arbitrarily for the purpose of this example. From a credal set interpretation, this model effectively states that we accept any probability distribution supported within $[i_p - 1, i_p + 1]$ where the probability measure $P$ satisfies $\{P(A) \leq \sup_{i \in A} \pi(i)\}$ as an acceptable model for our uncertainty. The mass distribution function $m_p$ associated to this credal set possesses two focal sets $a^p$:
\begin{eqnarray}
    &m_p(a^p_1=\opi i_p, i_p\cli)=1-\alpha\,\nonumber\\
    &m_p(a^p_2=\opi i_p-1, i_p + 1\cli)=\alpha\,\label{eq:pixel_mass}
\end{eqnarray}
with $\opi\cdot, \cdot\cli$ referring to integer intervals. In particular, $\opi i_p, i_p\cli$ correspond to the singleton $\{i_p\}$.

\begin{remark}
    The hypothesis of modelling the uncertainty on image intensities by possibility distributions does not consider uncertainty from potentially bigger sources of errors, such as satellite vibrations during the acquisition, or errors in the computations of epipolar lines. Those type of errors have been encountered on some Pléiades acquisitions, and lead to significant biases and errors on the final DSM, that our simple model does not account for. However, the CO3D mission will acquire images using a CCD matrix sensor rather than push-broom sensors used in Pléiades, so we can safely assume that those problems should not be encountered on images from the CO3D mission\commanue{Bon disons plutôt que la correction géométrique des images a pour objectifs de tenir compte de ces erreurs éventuelles et de les minimiser.}.
\end{remark}


\subsection{Dependency Model between Epipolar Images}
We presented the uncertainty models for pixels in both images, but we also need to define the dependency model between every pixel of both images. Indeed, as some pixels between images represent the light reflected by the same object, it seems natural that their (uncertain) values are correlated. In our case, we propose to model their dependency with the product copula if the pixels are not from the same physical object, meaning that the value of their intensities are independent. For pixels belonging to the same object, we model their dependency by a Gaussian copula with a covariance matrix $R$. Those copulas were presented in \cref{eq:gaussian_copula} from \Cref{chap:representation_of_uncertainty}. We remind here the formulation of a Gaussian $n$-copula $C_R$:
\begin{align}
    C_R(u_1 \enum u_n)=\Phi_R(\Phi^{-1}(u_1) \enum \Phi^{-1}(u_n))\label{eq:gaussian_copula}
\end{align} where $\Phi_R$ is the joint multivariate CDF of a Gaussian variable with correlation matrix $R$, and $\Phi^{-1}$ is the inverse CDF of a univariate Gaussian variable.

The correlation values inside the covariance matrix are based on a segmentation $S:(I_L\cup I_R)\rightarrow\opi1,K\cli$, $K\in\mathbb{N}$, of the images. This segmentation is the result of a \textit{k-means} clustering performed on the ground truth disparity map. Example of such a clustering is presented in \Cref{fig:clustering_example}, with $K=8$. Given the segmentation $S$ and two pixels $(p, q)\in(I_L\cup I_R)^2$, their covariance is determined by
\begin{equation}\label{eq:correlation}
    \sigma(p, q) =
    \begin{cases}
        1 &\text{ if }p=q,\\
        \rho_k, &\text{ if } p\ne q\text{ and }S(p)=S(q)\,, \\
        0 & \text{otherwise}\,.
    \end{cases}
\end{equation}
where $0<\rho_k<1$ is the correlation of pixels belong to segment $k\in\opi 1, N\cli$. Given a set of pixels $\{p_1 \enum p_n\}\subseteq(I_L\cup I_R)^2$, their covariance matrix $R$ is therefore:
\begin{align}
    R = \begin{bmatrix}
        1 & \sigma(p_{1}, p_{2}) & \dots & \sigma(p_{1}, p_{n-1}) & \sigma(p_{1}, p_{n})\\
        \sigma(p_{2}, p_{1}) & 1 & \dots & \sigma(p_{2}, p_{n-1}) & \sigma(p_{2}, p_{n})\\
        \dots & \dots & \dots & \dots & \dots\\
        \sigma(p_{n-1}, p_{1}) & \sigma(p_{n-1}, p_{2}) & \dots & 1 & \sigma(p_{n-1}, p_{n})\\
        \sigma(p_{n}, p_{1}) & \sigma(p_{n}, p_{2}) & \dots & \sigma(p_{n}, p_{n-1}) & 1
    \end{bmatrix}
\end{align}

In practice, we will only compute the correlation matrix between the two windows from the reference and secondary images that are compared. The windows have a $3\times 3$ shape, we thus consider Gaussian $18-$copulas to model the dependency between pixels. This copula will be used for joining marginal masses in the uncertainty propagation step using \cref{eq:joint_mass} from the previous chapter. It will also be used to sample for Monte Carlo sampling, used for evaluating the correct uncertainty propagation in \Cref{sec:montecarlo}\commanue{Il y a beaucoup de used dans les phrases, je trouve que cette pĥrase complique les choses. Peut-être juste mentionner qu'on utilise la copule gaussienne quand on arrive à ces sujets}.

\begin{remark}\commanue{remarque globale avec les encarts remark, fais attention à ne pas en abuser car on a tendance à perdre le fil. Cette remarque tu pourrais la mettre au début de la section quand tu dis que tu choisis des copules gaussiennes.}
    Gaussian copulas are popular and simple copulas used to represent dependencies between more than $2$ variables. By comparison, the different $2$-copulas presented in \Cref{sec:copulas} cannot always be defined in more than $2$ dimensions or possess a quite complex formulation. Another method for modelling the dependency for more than $2$ variables is to express a $n$-copula as a combination of $2$-copulas, which is called a vine copula (\cite{czado_vine_2022}). However this is a complex subject that is not adapted to our type of dependency, and is therefore not explored in this thesis.
\end{remark}

\begin{figure}
    \centering
    \begin{subfigure}[t]{0.5\linewidth}
        \centering
        \includegraphics[width=0.9\linewidth]{Images/Chap_4/im2.png}
        \caption{Colored left image}
        \label{fig:color_cones_image}
    \end{subfigure}\hfill
    \begin{subfigure}[t]{0.5\linewidth}
        \centering
        \includegraphics[width=0.9\linewidth]{Images/Chap_4/cluster.jpg}
        \caption{Proposed clustering of the image (\textit{k-means} with $K=8$)}
        \label{fig:cluster}
    \end{subfigure}
    \caption{Middlebury 2003 Cones left image, and a clustering computed from the disparity ground truth.}
    \label{fig:clustering_example}
\end{figure}

During our simulations, the segmentation contains $K=8$ different clusters. For every $k$ in $\opi1,K\cli$, $\rho_k$ is assigned an value between $0.9$ and $1$ in order to really emphasize on their correlation.\commanue{à mon avis tu peux regrouper cette phrase dans le milieu paragraphe précédent où tu décris ton modèle. Là elle se retrouve un peu isolée}

\begin{remark}
    The segmentation is based on the ground truth disparity map. This means that two objects with similar disparities located at opposite sides of the image will be considered as belonging to the same object and thus correlated. In practice, those pixels are never compared as we only measure the dissimilarity between small windows in a restricted disparity range. The clustering is thus only used at a \textit{local} scale. 
    
    Secondly, the segmentation allows to compute the Gaussian copula, which will be used to propagate the uncertainty models. We will validate the uncertainty propagation in \Cref{sec:montecarlo} with Monte Carlo samples using the same copula. As long as the copulas used are the same between the propagation and validation, it does not really matter which correlation matrix is used. We still tried to create a realistic but simple dependency model for the sake of the example, but it is not required from a theoretical point of view\commanue{là trop de remarques tue la remarques. Tu n'es pas obligé de faire des teasers sur tous les points que tu vas abordés. Je me limiterais à donner les infos qui correspondent au modèle de dépendence entre les images qui est le thème de la section. Si par la suite tu as besoin de réutiliser les copules, bah tu le dis.}.
\end{remark}

\section{Propagating the Uncertainty with Belief Functions and a Copula}
Having defined both marginal models for pixel intensities using possibility distributions as well as dependency models using copulas, we can now join them all to construct a multivariate uncertainty model, as seen in \Cref{sec:methods_for_joining_credal_sets}. We will see in this section how the multivariate models can then be used to compute the uncertainty regarding the cost curve. 

\subsection{From Multivariate Uncertainty Models to the Propagated Model}
We first detail how multivariate models are used to propagate uncertainty in the precise case. We will then do the same in the imprecise setting by analogy.

Consider a mapping $f:\X_1\times\X_2\rightarrow\mathcal{Z}$ from a product space $\X_1\times\X_2$ to a space $\mathcal{Z}$, which propagates multiple random variables $X_1$, $X_2$ to a new random variable $Z=f(X_1, ~X_2)$. In our case, $f$ will be the $\SAD$ cost function propagating the intensities to the matching cost. When considering precise probabilities, the probability of $Z$ on atoms $z$ is obtained by summing the probabilities of all events $X_1=x_1$, $X_2=x_2$ whose image by $f$ is $z$: 
\begin{align}\label{eq:precise_propagate_proba}
    \forall z\in\mathcal{Z}, P_Z(z)=\sum_{\substack{x_1,x_2\\z=f(x_1,x_2)}}P(x_1,x_2).
\end{align}
where $P(x_1,x_2)$ is the joint probability, which is linked to its marginals by a copula $C$. $P_Z$ is completely determined by evaluated every combination of atoms of $\X_1$ and $\X_2$.

\begin{example}\label{ex:propagated_probability}
    Consider the same setting as example \ref{ex:copulas}, where a dealer throws two coins in a separate room. The coins seem fair when looked at independently. In this example, a coin landing on heads rewards you with $1$€ (or any currency of your choice), and a coin landing on tails rewards you with $0$€. So you earn $2$€ if both coins land on heads, $1$€ if only one coin land on heads, and $0$€ if both coins land on tails. We are interesting in the uncertainty regarding you earnings, noted $Z$.
    
    In example \ref{ex:copulas}, we consider $3$ different cases, each leading to a different copula $C$ modelling the dependency between the probability $P_1$ of first coin and $P_2$ of the second coin. For each copula, we then computed the joint probability $P$.
    
    In the first case, where coin throws were independent, we saw that the joint probability $P$ was:\\
    \begin{minipage}[b]{0.5\linewidth}
    \begin{align*}
        & P(\text{heads}, ~\text{heads}) = 0.25\\
        & P(\text{heads}, ~\text{tails}) = 0.25
    \end{align*}
    \end{minipage}
    \begin{minipage}[b]{0.5\linewidth}
    \begin{align*}
        & P(\text{tails}, ~\text{tails}) = 0.25 \\
        & P(\text{tails}, ~\text{heads}) = 0.25
    \end{align*}
    \end{minipage}
    In that case, it holds that the probability $P_Z$ of our earnings are:
    \begin{align*}
        & P_Z(Z=2) = P(\text{heads}, ~\text{heads}) = 0.25\\
        & P_Z(Z=1) = P(\text{heads}, ~\text{tails}) + P(\text{tails}, ~\text{heads}) = 0.5 \\
        & P_Z(Z=0) = P(\text{tails}, ~\text{tails}) = 0.25
    \end{align*}
    
    In the second case, where coin throws were rigged to land on the same side, we saw that the joint probability $P$ was:\\
    \noindent
    \begin{minipage}[b]{0.5\linewidth}
    \begin{align*}
        & P(\text{heads}, ~\text{heads}) = 0.5\\
        & P(\text{heads}, ~\text{tails}) = 0
    \end{align*}
    \end{minipage}
    \begin{minipage}[b]{0.5\linewidth}
    \begin{align*}
        & P(\text{tails}, ~\text{tails}) = 0.5\\
        & P(\text{tails}, ~\text{heads}) = 0 
    \end{align*}
    \end{minipage}
    Following the same methodology, we have:\\
    \noindent\begin{minipage}[b]{0.33\linewidth}
    \begin{align*}
        & P_Z(Z=2) = 0.5
    \end{align*}
    \end{minipage}
    \begin{minipage}[b]{0.33\linewidth}
    \begin{align*}
        & P_Z(Z=1) = 0
    \end{align*}
    \end{minipage}
    \begin{minipage}[b]{0.33\linewidth}
    \begin{align*}
        & P_Z(Z=0) = 0.5
    \end{align*}
    \end{minipage}
    
    In the third case, where coin throws were rigged to land on opposite sides, we saw that the joint probability $P$ was:\\
    \noindent\begin{minipage}[b]{0.5\linewidth}
    \begin{align*}
        & P(\text{heads}, ~\text{heads}) = 0\\
        & P(\text{heads}, ~\text{tails}) = 0.5
    \end{align*}
    \end{minipage}
    \begin{minipage}[b]{0.5\linewidth}
    \begin{align*}
        & P(\text{tails}, ~\text{tails}) = 0\\
        & P(\text{tails}, ~\text{heads}) = 0.5
    \end{align*}
    \end{minipage}
    Therefore:\\
    \begin{minipage}[b]{0.33\linewidth}
    \begin{align*}
        & P_Z(Z=2) = 0
    \end{align*}
    \end{minipage}
    \begin{minipage}[b]{0.33\linewidth}
    \begin{align*}
        & P_Z(Z=1) = 1
    \end{align*}
    \end{minipage}
    \begin{minipage}[b]{0.33\linewidth}
    \begin{align*}
        & P_Z(Z=0) = 0
    \end{align*}
    \end{minipage}
\end{example}

Determining every $(x_1, x_2)$, whose image by $f$ equals $z$, is not always trivial. This becomes even more complex when considering $n>2$ marginal variables. Similarly, the joint probability $P(x_1,x_2)$ is computed using a H-volume, which is the sum of $2^n$ terms, thus also increasing exponentially with the dimension.

There are multiple ways of extending \cref{eq:precise_propagate_proba} to the imprecise setting, as there are multiple ways of aggregating imprecise models using a copula. We presented\commanue{tu emploies bcp cette formule, donc de temps en temps described/} in \Cref{chap:joining_credal_sets} three methods for joining marginal credal sets, creating three different multivariate credal sets $\M_{robust}$, $\M_{mass}$ and $\M_{agg}$.

The robust approach of extending \cref{eq:precise_propagate_proba} is based on the robust approach from \Cref{sec:robust_method}. Given $n$ marginal credal sets $\M_i$, we can join them using a copula $C$ into a credal set $\M_{robust}$ as in \eqref{eq:robust_set}. The propagated uncertain model $\M^Z_{robust}$ is then defined as:
\begin{align}
    \M^Z_{robust} = \{P_Z~|~\forall z\in Z, ~P_Z(z)=\sum_{\substack{x_1 \enum x_n\\z=f(x_1 \enum x_n)}}P(x_1 \enum x_n), ~P\in\M_{robust}\}
\end{align}
Practically, this set is computed by sampling every probability $P_i$ from each marginal credal set $\M_i$ and joining them using \hyperref[theorem:sklar]{Sklar's Theorem} into a multivariate probability $P$. Then for each $z=P(x_1 \enum x_n)$, we can compute $P_Z$ from $P$ using \eqref{eq:precise_propagate_proba}. In short, each sample $(P_1 \enum P_n)$ leads to a new $P$, itself leading to a new $P_Z$. Sampling through every $(P_1 \enum P_n)$ thus leads to the estimation of the uncertain model of $Z$. This method is complicated to compute, but correctly propagates the uncertainty.

We saw that it was easier to compute $\M_{mass}$ than $\M_{robust}$. Another way\commanue{mot de liaison entre les deux phrase. So another way or Thus another way} of approximating the uncertainty model of $Z$ is to compute it using $\M_{mass}$. This is done by replacing the probability on atoms from \cref{eq:precise_propagate_proba} with the joint mass $m_\times$ from \cref{eq:joint_mass} \cite{gray_dependent_2021}. Consider $n$ uncertain variables $X_i$, each modeled by a mass distribution function whose $j$-th focal set is noted $a^i_j$. Given their multivariate mass $m_\times$, it is possible to compute the mass distribution function $m_Z$ of a random set $Z=f(X_1 \enum X_n)$ as:
\begin{align}
    \forall a^Z\subseteq\mathcal{Z}, m_Z(a^Z) = \sum_{\substack{a^1_i \enum a^n_j\\a^Z=f(a^1_i \enum  a^n_j)}}m_\times(a^1_i \enum a^n_j)\label{eq:mass_propagated}
\end{align}
leading to a credal set $\M_{mass}^Z$:
\begin{align}
    \M_{mass}^Z &= \{~P_Z~|~\forall A\subseteq\mathcal{Z},~P_Z(A)\geqslant \sum_{a^Z\subseteq A}m_Z(a^Z)\}\\
    &=\{~P_Z~|~\forall A\subseteq\mathcal{Z},~P_Z(A)\geqslant \sum_{\substack{a^1\enum a^n \\ f(a_1\enum a_n)\subseteq A}}m_\times(a^1\enum a^n)~\}
\end{align}

In order to compute the propagated mass $m_Z$ (and its associated belief function) from \cref{eq:mass_propagated}, two difficulties arise. The first one is to determine what the focal sets $a_Z$ of $m_Z$ will be, which corresponds to the subscript $a^Z=f(a^n_1 \enum  a^n_j)$ of the previous sum. Computing the image of $f$ for every combination of focal sets $(a^1_i \enum a^n_j)$ is even more difficult than in the precise case, as we are computing images of sets instead of real numbers. The second difficulty is to compute the joint mass $m_\times$, as in the case of the $\SAD$ it requires to compute the H-volume of an $18$-copula. Those difficulties will be addressed in the following sections \ref{sec:propagated_focal_sets} and \ref{sec:propagated_masses}

We saw in \Cref{chap:joining_credal_sets} that in the situation where marginals are possibility distributions, $\M_{mass}$ and $\M_{agg}$ have the same bounds on Cartesian products of events (see \cref{eq:inclusion_necessity}). We will thus only compute the lower bounds $Bel_\times$ of $\M_{mass}$ as it directly provides the bounds of $\M_{agg}$ on those events\commanue{tu n'as pas inversé $\M_{mass}$  et $\M_{agg}$ dans ta phrase ? $\M_{agg}$ c'ets pas celui qui est plus facile à calculer}.

\subsection{Determining the Propagated Focal Sets}\label{sec:propagated_focal_sets}
In this section, we will detail how we compute the $\SAD$ image of marginal focal sets. Computing the image of sets needs to be treated with caution in the general case. However in our case, because we chose marginal focal sets with a simple expression, and because we are using a relatively regular cost function, computing the image is significantly easier.

Given a pixel $p$, we consider the mass distribution $m_p$ of \cref{eq:pixel_mass} and its two focal sets $a_1^p$ and $a_2^p$ from \cref{eq:pixel_mass}. For every pair of pixels $p\in I_L, q\in I_R$, we note $\AD_{pq}=|i_p - i_q|$, where $i$ refers to a pixel's intensity. Given $m_p$, There exist $3$ focal sets related to the absolute difference:
\begin{itemize}
    \item $a^{\AD}_1$ is the image of the AD of $a^p_1$ and $a^q_1$
    \item $a^{\AD}_2$ is the image of the AD of $a^p_2$ and $a^q_1$ or $a^p_1$ and $a^q_2$
    \item  $a^{\AD}_3$ is the image of the AD of $a^p_2$ and $a^q_2$
\end{itemize}
The non-monotonicity of the absolute value around $0$ needs to be taken into account to compute their exact image through the AD. Indeed, if a value $x$ is in $[-1,1]$, then its absolute value will be in $[0,1]$. Applying this remark to the AD yields the following focal sets:
\begin{align*}
    a^{\AD}_1&=\opi\AD_{pq},~\AD_{pq}\cli\,,\\
    a^{\AD}_2&=\opi\max(0, \AD_{pq} - 1),~\AD_{pq} + 1\cli\,,\\
    a^{\AD}_3&=\opi\max(0, \AD_{pq} - 2),~\AD_{pq} + 2\cli\,,
\end{align*}

Focal sets of the final $\SAD$ are then computed by simply summing the bounds of the $9$ AD focal sets (as in \Cref{fig:SAD}), for every combination $(a^{\AD_1}_{k_1}\enum a^{\AD_9}_{k_9})$ of those focal sets:
\begin{align}
    a^\SAD=\sum_{i=1}^9a^{\AD_i}_{k_i}
\end{align}

\begin{remark}
    In many cases, different combinations of $\AD$ focal sets will lead to the same $\SAD$ focal set. Actually, if every $\AD$ is greater than $2$\commanue{alors je ne suis pas sure de comprendre} (so that each $a^\AD_i$ is symmetric with regards to its $\AD$) there will only be $19$ focal sets $a^\SAD$ for the $\SAD$, and they are of the following form:
    \begin{align}\label{eq:nested_SAD}
        a^\SAD = \opi\SAD-t,\SAD+t\cli,\text{ with }t\in\opi0, ~18\cli
    \end{align}
    In comparison, there are $3^9=19683$ different combinations of $\AD$ focal sets.
\end{remark}

\begin{remark}\commanue{pourquoi tu mets ce qui concerne SAD dans des remarques ?}
    In the previous remark, we established that when every $\AD$ is greater than $2$, it holds that:
    \begin{align*}
        a^\SAD = \opi\SAD-t,\SAD+t\cli,\text{ with }t\in\opi0, ~18\cli
    \end{align*}
    This equation translates the fact that focal sets of the $\SAD$ form an increasing family of events. This also means that $\Bel_\SAD$ is actually a necessity function. We can thus compute a possibility distribution $\pi_\SAD$ to represent the uncertainty of the $\SAD$. This can be useful for graphically representing the uncertainty of the $\SAD$, if we later need to build joint uncertainty models using the $\SAD$, or if we want to propagate the $\SAD$ uncertainty even further.
    
    We proved in \cite{malinowski_uncertainty_2024} that in order to propagate marginal possibility distributions $\pi_i$ into a possibility distribution $\pi_Z$ using a copula and a propagating function $f$, a sufficient condition was that\commanue{tu n'as pas deux conditions dans ce qui suit}:
    \begin{itemize}
        \item $f$ is a monotone function applied to a linear combination $\alpha_1X_1+\dots+\alpha_nX_n+\beta$ of marginal variables $X_i$
        \item Each $\pi_i$ is symmetrical and uni-modal (meaning that all values $x_i$ such that $\pi_i(x_i)=1$ are adjacent). For instance all triangular possibility distributions verify this condition
    \end{itemize}
    Here is an example of an uni-modal symmetrical possibility
    \centering{\begin{tikzpicture}[scale=1]
        \begin{axis}[%
          xlabel=$x$,
          ylabel=$\pi(x)$,
          xmin=2, xmax=8,
          ymin=-0.05, ymax=0.55,
          ytick={0, 0.15, 0.4, 0.5},
          yticklabels={$0$, $0.3$, $0.8$, $1$},
          scale only axis,
          width=7cm, height=3.5cm,
          ],
          \node (a) at (2, 0) {$\times$};
          \node (b) at (2.5, 0) {$\times$};
          \node (c) at (3, 0) {$\times$};
          \node (d) at (3.5, 0.15) {$\times$};
          \node (e) at (4, 0.4) {$\times$};
          \node (f) at (4.5, 0.5) {$\times$};
          \node (g) at (5, 0.5) {$\times$};
          \node (h) at (5.5, 0.5) {$\times$};
          \node (i) at (6, 0.4) {$\times$};
          \node (j) at (6.5, 0.15) {$\times$};
          \node (k) at (7, 0) {$\times$};
          \node (l) at (7.5, 0) {$\times$};
          \node (m) at (8, 0) {$\times$};
          
          \draw [thick, dashed, gray] (a.center) -- (b.center) -- (c.center) -- (d.center) -- (e.center) -- (f.center) -- (g.center) -- (h.center) -- (i.center) -- (j.center) -- (k.center) -- (l.center) -- (m.center) ;
          \draw [decorate,decoration={brace,amplitude=5pt,mirror,raise=1ex},gray] (f.south) -- (h.south) node[midway,yshift=-2em]{\color{gray}{mode}};
        \end{axis}
    \end{tikzpicture}}
\end{remark}

\subsection{Computing the Mass of Propagated Focal Sets}\label{sec:propagated_masses}
%\section{Leveraging Specificities to Accelerate Computations}
Now that the bounds of the $\SAD$ have been computed\commanue{alors moi j'ai compris qu'on avait calculé les focal sets. J'ai pas bien vu le calcul des bornes}, we need to compute the associated mass. Computing the joint mass over two $3 \times 3$ windows is significantly more complex. For each combination of marginal focal sets, the joint mass $m_{\times}$ is computed using the H-volume of an $18$-copula, involving a sum of $2^{18}$ terms. Given that the uncertainty of each of the 18 pixels is represented by $2$ focal sets, we need to evaluate $2^{18}$ combinations of these focal sets in total\commanue{c'est juste moi qui comprends pas la subtilité }. This computation can thus become quite costly in memory and computation time, especially when computing it over a whole image.

In the case of the family of Gaussian copulas, their expression given by \cref{eq:gaussian_copula} show that we need to compute the multivariate CDF. An $18$-variate Gaussian CDF does not possess a known analytic formula, it is thus computed by integrating its density (so integrating a $18$-variate function), as expressed below:
\begin{align}\label{eq:gaussian_cdf}
    F(x_1 &\enum x_{18}) =\nonumber\\ &\int_{-\infty}^{x_1}\dots\int_{-\infty}^{x_{18}}\frac{1}{\sqrt{(2\pi)^{18}|R|}}\exp(-\frac{1}{2}\begin{bmatrix}x_1 & \dots & x_{18}\end{bmatrix}R^{-1}\begin{bmatrix}x_1 \\ \dots \\ x_{18}\end{bmatrix})dx_1\dots dx_{18}
\end{align}
Where $|R|$ is the determinant of $R$. Computing this CDF can quickly become time-consuming ($40ms$ on average\footnote{With an AMD EPYC7713  64-Core Processor at 2GHz, using Python and the SciPy library}). For each pixel and each disparity, we saw that we need to compute the joint mass of $2^{18}$ focal sets, each time necessitating $2^{18}$ evaluation of its copula. Because there are around $170\,000$ pixels and $60$ disparities to be evaluated, the processing time is too large to be computed as such. We will instead see that we can leverage specificities of our problem to drastically reduce the computation time.
 
The first idea is to notice than if we can divide our variables into multiple mutually independent sets of variables, then the evaluation of the $18$-copula can be separated into the evaluation of multiple lower dimension copulas. To verify this statement consider the following independent sets of variables $\{X_1 \enum X_k\}$ and $\{X_{k+1} \enum X_n\}$ with $k\in\opi1,n-1\cli$. Let $F_1 \enum  F_n$ be their marginals CDF. And let $F_{(1 \enum n)}$ be the joint CDF of all variables, $F_{(1 \enum k)}$ the joint CDF of the first set of variables and $F_{(k+1 \enum n)}$ the joint CDF of the second set. The two sets are independent means that for all $(x_1 \enum x_n)\in\X_1\tdt\X_n$ it holds that:
\begin{align*}
    F_{(1 \enum n)}(x_1 \enum x_n) = F_{(1 \enum k)}(x_1 \enum x_k)\cdot F_{k+1 \enum n}(x_{k+1} \enum x_n)
\end{align*}
Using Sklar's theorem, there exist a $n$-copula $C$, a $k$-copula $C'$ and a $n-k$-copula $C''$ respectively linking $F_{(1 \enum n)}$, $F_{(1 \enum k)}$, and $F_{(k+1 \enum n)}$ to their marginals:
\begin{align}
    C(F_1(x_1) \enum F_n(x_n)) = C'(F_1(x_1) \enum F_k(x_k))\cdot C''(F_{k+1}(x_{k+1}) \enum F_n(x_n))\label{eq:copula_factorizing}
\end{align}

\begin{remark}
    We stated earlier that we did not considered vine copulas \cite{czado_vine_2022}, which are a way of constructing multivariate copulas by composition of bivariate (conditional) copulas. The decomposition $C$ into $C'$ and $C''$ actually follows the same idea of decomposing a copula into smaller ones. So even though we are not using vine copulas, we use a similar philosophy in our computations.
\end{remark}

Establishing \cref{eq:copula_factorizing} becomes interesting once we put it in relation with the following property:
\begin{proposition}[H-Volume factorizing]\label{prop:hvol_factorizing}
    Let $1<k<n$. If a $n$-copula $C$ can be expressed as the product of a $k$-copula $C'$ and a $(n-k)$-copula $C''$, then the H-volume of $C$ is the product of the H-volume $H'$ of $C'$ and the H-volume $H''$ of $C''$. This means that for all $(u_1 \enum u_n)\in[0,1]^n$ and for all $(v_1 \enum v_n)\in[0,1]^n$ such that $u_i\leqslant v_i$, it holds that:
    \begin{align}
        H_{u_1 \enum u_n}^{v_1 \enum v_n} = {H'}_{u_1 \enum u_k}^{v_1 \enum v_k}\cdot{H''}_{u_{k+1} \enum u_n}^{v_{k+1} \enum v_n}
    \end{align}
\end{proposition}
\begin{proof}
    Let $1<k<n$, $C$ a $n$-copula, $C'$ a $k$-copula and $C''$ a $n-k$ copula. $H$, $H'$, $H''$ are the respective H-volume of $C$, $C'$, $C''$. then
    \begin{align*}
        %{H'}_{u_1 \enum  u_k}^{v_1 \enum  v_k}\times {H''}&_{u_{k+1} \enum  u_n}^{v_{k+1} \enum  v_n} = \left(\sum_{w_i\in\Pi_{i=1}^{k}\{u_i,v_i\}}(-1)^{|\{w_i~|~w_i=u_i\}|}C'(w_1 \enum w_k)\right)\\
        {H'}
        \begin{matrix}
            v_1 \enum  v_k\\
            u_1 \enum  u_k
        \end{matrix}
        \cdot {H''}&
        \begin{matrix}
            v_{k+1} \enum  v_n\\
            u_{k+1} \enum  u_n
        \end{matrix}= \left(\sum_{w_i\in\Pi_{i=1}^{k}\{u_i,v_i\}}(-1)^{|\{w_i~|~w_i=u_i\}|}C'(w_1 \enum w_k)\right)\\
        &\times\left(\sum_{w_j\in\Pi_{j=k+1}^{n}\{u_j,v_j\}}(-1)^{|\{w_j~|~w_j=u_j\}|}C''(w_{k+1} \enum w_n)\right)\\
        =& \sum_{w_i\in\Pi_{i=1}^{k}\{u_i,v_i\}}\times\sum_{w_j\in\Pi_{j=k+1}^{n}\{u_j,v_j\}}(-1)^{|\{w_i~|~w_i=u_i,~i\leqslant k\}|}\\
        &\times(-1)^{|\{w_j~|~w_j=u_j,~ j>k\}|}C'(w_1 \enum w_k)C''(w_{k+1} \enum w_n)\\
        =&\sum_{w_i\in\Pi_{i=1}^{n}\{u_i,v_i\}}(-1)^{|\{w_i~|~w_i=u_i\}|}C'(w_1 \enum w_k)\\
        &\times C''(w_{k+1} \enum w_n)\\
        =&H
        \begin{matrix}
            v_1 \enum  v_n\\
            u_1 \enum  u_n
        \end{matrix}
        %=&H_{u_1 \enum  u_n}^{v_1 \enum  v_n}
    \end{align*}
\end{proof}
Using the result of proposition \ref{prop:hvol_factorizing}, we can now compute the joint mass as a product of two lower dimension copulas. It is easier to compute as we only integrate a $k$-dimensional function and a $n-k$-dimensional function, instead of a $n$-dimensional one. Similarly the $H$ volume is not the sum of $2^{18}$ terms anymore, but the sum of $2^{k}$ and $2^{n-k}$ terms.

For comparison, consider that we can split the aforementioned Gaussian $18$-copula into two Gaussian $9$-copula. Computing the value of single mass used to take around $10\,500s$, but now takes around $6s$, so around $1\,700$ times faster. This demonstrates the substantial time savings achieved by decomposing the problem into smaller, independent parts. This improvements apply directly to our application. Indeed \cref{eq:correlation} yields the following correlation between two pixels $p$, $q$ given the segmentation $S$:
\begin{equation*}
    \sigma(p, q) =
    \begin{cases}
        1 &\text{ if }p=q,\\
        \rho_k, &\text{ if } p\ne q\text{ and }S(p)=S(q)\,, \\
        0 & \text{otherwise}\,.
    \end{cases}
\end{equation*}
From this, we can split the set $\{p_1 \enum p_{18}\}$ of pixels in at most $K=8$ mutually independent sets $S_k = \{p_i~|~S(p_i)=k\}$, each set containing pixels from the same cluster. We can thus compute the H-volume of each cluster independently.

Another additional way of reducing the computation time is to avoid computing the same $H$-volume multiple times. Consider a set $S_k$ with $k_L$ pixels from the left image and $k_R$ pixels from the right image. The correlation matrix $R_k$ for this cluster is:
\begin{equation}
    R_k=\begin{pmatrix}
        1 &  & \rho_k & \dots & \rho_k\\
         & & & & \vdots\\
        \rho_k &  & 1 & & \rho_k\\
        \vdots &  &  & & \\
        \rho_k & \dots & \rho_k &  & 1
    \end{pmatrix}\label{eq:corr_matrix_sym}
\end{equation}
which implies that the Gaussian $(k_L+k_R)$-copula for this set is symmetrical. The joint mass $m^{S_k}_\times$ of this set is then also symmetrical. Pixels from the left image will share the same mass $m_L$ for their two focal sets and pixel from the right image will also share the same mass $m_R$, therefore computing $m^{S_k}_\times$ on every possible combination of cumulative masses is redundant. 
\begin{example}
    Let us imagine a set $S_k$ with $k_L=2$ pixels $p_1$, $p_2$ from the left image and $k_R=2$ pixels $p_3$, $p_4$ from the right image. Each pixel $p_i$ has two focal sets $a^i_1$ and $a^i_2$.
    We know that:
    \begin{align*}
        m_L(a_1^1) &= m_L(a_1^2) \qquad m_L(a_2^1) = m_L(a_2^2)\\
        m_R(a_1^3) &= m_R(a_1^4) \qquad m_R(a_2^3) = m_R(a_2^4)
    \end{align*}
    Because of the symmetry of the Gaussian $4$-copula of the set $S_k$, the joint mass $m^{S_k}_\times$ computed as the H-volume on cumulative masses (\cref{eq:joint_mass}) verifies:
    \begin{align*}
        m^{S_k}_\times(a_1^1, a_2^2, a_1^3, a_2^4) &= H
        \begin{matrix}
            m_L(a_1^1), & m_L(a_1^1)+m_L(a_2^2), & m_R(a_1^3), & m_R(a_1^4)+m_R(a_2^4)\\
            0, & m_L(a^2_1), & 0, & m_R(a_1^4)
        \end{matrix}\\
        &=H
        \begin{matrix}
            m_L(a_1^1)+m_L(a_2^1), & m_L(a_1^2), & m_R(a_1^3)+m_R(a_2^3), & m_R(a_1^4)\\
            m_L(a^1_1), & 0, & m_R(a_1^3), & 0
        \end{matrix}\\
        &\text{(by symmetry of the copula)}\\
        &= m^{S_k}_\times(a_2^1, a_1^2, a_2^3, a_1^4)
    \end{align*}
    We can see that we do not need to compute $m^{S_k}_\times$ on every possible combination of focal sets as many combination have the same joint mass $m^{S_k}_\times$.
\end{example}
For each set $S_k$, we only have to compute $(k_L+1)(k_R+1)$ values of $m^{S_k}_\times$ rather than $2^{k_L+k_R}$.

We saw that it was complex to compute the joint mass of the $\SAD$. However, by taking advantage of the potential factorization of the copula into smaller copulas, and by leveraging some symmetries of the problem, we can greatly reduce the time and number of computations required.

\section{Results and discussions}
In the previous section, we presented how we can use $\M_{robust}$ and $\M_{mass}$ to propagate the uncertainty from the input images into the uncertainty of the $\SAD$. Evaluating this uncertainty for every pixel and every considered disparity results in the uncertainty models of every value of the cost volume (see \cref{eq:cost_volume} from \Cref{sec:stereo_matching} for more details on the cost volume). We saw in \Cref{sec:necessity_functions} that $\M_{mass}$ and $\M_{robust}$ are different credal sets, and that neither set is guaranteed to be included in the other. This section will then compare the different models and see if $\M^Z_{mass}$ can be used to approximate $\M^Z_{robust}$. \Cref{sec:envelopes_plausibility} will present visualisations of  $\M^Z_{mass}$ using plausibility envelopes. \Cref{sec:montecarlo} will present visualisations of $\M^Z_{robust}$ using Monte Carlo samples and compare them to the plausibility envelopes. Finally, \Cref{sec:sad_improvements} will estimate the potential improvements unlocked by estimating the uncertainty.

\subsection{Envelopes Defined by Plausibility Levels}\label{sec:envelopes_plausibility}
Having defined efficient ways of computing the $\SAD$ focal sets in \Cref{sec:propagated_masses} and the joint mass in \ref{sec:propagated_focal_sets}, we can determine the belief function $\Bel_\SAD$ associated to every estimation of the $\SAD$ between two $3\times3$ windows. $\Bel_\SAD$ is deducted from the mass $m^\SAD_\times$ computed in \cref{eq:mass_propagated}.

For each pixel $p_L=(row, ~col)$ in the left image and for each disparity $d$, we calculated the $\SAD$ cost between the window centered on $p_L$ in the left image and the window centered on $p_R=(row, ~col+d)$ in the right image. We are usually interested in going through each considered disparity $d$ to obtain the cost curve for pixel $p_L$. From this cost curve, a \textit{winner-takes-all} strategy is applied to find the correct disparity. Due to its interest, we want to visualise the uncertainty of the entirety of the cost curve, not just a single value.

It can be hard to graphically represent focal sets, or similarly belief functions, especially when there is a large number of sets to consider. We are usually more keen to represent uncertainty on singletons, as we do with probability densities, possibility distributions or, to a certain extent, p-boxes (although in that case singletons represent cumulative events). Following that logic, we will consider the plausibility of singletons as a way of representing the uncertainty graphically.
\begin{remark}
    Considering the belief on singletons instead of the plausibility does not make much sense as the belief of singletons is null. The only exception is the precise $\SAD$ value, because it is the only singleton that is also a focal set.
\end{remark}
We will first explain how we  graphically represent the uncertainty on the special case of possibility distributions, and then extend this representation to all $\SAD$ belief functions. 

As shown in \Cref{sec:propagated_focal_sets}, focal sets representing the $\SAD$ uncertainty are defined as intervals containing the ``precise'' $\SAD$ value. We saw in a remark from \Cref{sec:propagated_focal_sets} that the uncertainty of the $\SAD$ could be represented by a possibility distribution $\pi_\SAD$ in the special case where all $\AD$ leading to the $\SAD$ are greater that $2$. In this case, we can easily plot this possibility distribution, at least for a few degrees of possibility. Given a degree of possibility $\gamma$, we plot the bounds of the largest interval $I_\gamma=[\underline{I}_\gamma, \overline{I}_\gamma]$ whose possibility is greater than $\gamma$. The possibility measure is computed as $\Pi_\SAD(A)=\sup_{z\in A}\pi_\SAD(z)$ or $\Pi_\SAD(A)=\sum_{a\cap A\neq\emptyset}m^\SAD_\times(a)$ using \cref{eq:bel_pl}. This way, given $\gamma\in[0,1]$, the bounds $\underline{I}_\gamma, \overline{I}_\gamma$ to plot are:
\begin{align}\label{eq:possibility_bounds}
    \underline{I}_\gamma = \argmin_{z}\{\pi_\SAD(z)\geqslant\gamma\}, \qquad \overline{I}_\gamma = \argmax_{z}\{\pi_\SAD(z)\geqslant\gamma\}
\end{align}
Because focal sets of $\pi_\SAD$ are increasing intervals, then every $\SAD$ value between $\underline{I}_\gamma$ and $\overline{I}_\gamma$ will have a possibility superior to $\gamma$. The bounds thus define the envelope of values with a degree of possibility greater than $\gamma$. Plotting envelopes for different values of $\gamma$ allows to visualize a representation of the $\SAD$ uncertainty. We will define the envelopes in the following, but readers can already have a broad idea of what those envelopes look like by looking at \Cref{fig:belief_curves}. 

In \cref{eq:possibility_bounds}, $\underline{I}_\gamma$ and $\overline{I}_\gamma$ are the upper and lower bounds of the same focal set. If we relax this constraint, then we can extend this definition of $\underline{I}_\gamma, \overline{I}_\gamma$ to any type of $\SAD$ plausibility (or equivalently belief) function $\Pl_\SAD$. For all focal sets $a=[\underline{a}, \overline{a}]$ of the $\SAD$ plausibility function, the bounds $\underline{I}_\gamma, \overline{I}_\gamma$ are defined as:
\begin{align}\label{eq:belief_bounds}
    \underline{I}_\gamma = \min\{\underline{a} ~|~ \Pl_\SAD([\underline{a}, \overline{a}])\geqslant\gamma\}, \qquad \overline{I}_\gamma = \max\{\overline{a} ~|~ \Pl_\SAD([\underline{a}, \overline{a}])\geqslant\gamma\}
\end{align}

\begin{remark}
    \Cref{eq:belief_bounds} details how we chose to represent a plausibility function (or equivalently a credal set) by some envelopes. One can wonder if we can reconstruct the plausibility function from the envelopes. It is straightforward to see that it is the case if and only if the plausibility function is a possibility measure, because it is then fully determined by its values on singletons.
    
    However if the plausibility is not a possibility measure, then the possibility distribution $\pi'$ defined by envelopes induces the smallest credal set $\M(\pi')$ containing $\M(\Pl_\SAD)$. In that regard, we chose to represent a plausibility by its best outer approximating possibility.
\end{remark}
Now that the we decided on how to plot our uncertainty models, we can present some results. We arbitrarily considered different values for plausibility levels $\gamma$:
\begin{itemize}
    \item The first value is $\gamma=1$. It corresponds to the $\SAD$ value with the highest degree of possibility. This value is unique in our case, and corresponds to the $\SAD$ value that would have been computed without considering the uncertainty.
    \item We then consider values $\gamma=0.9$ and $\gamma=0.85$. They allow to give an estimation of the dispersion of envelopes with high plausibility. With them, we can estimate if highly plausible values are near the $\SAD$ values or not.
    \item $\gamma=0.5$ gives a moderate plausibility estimation. From a credal set point of view, the probability that the $\SAD$ value lies in this set is at least $0.5$.
    \item The last value is $\gamma=0$. For this value, the inequalities in \cref{eq:belief_bounds} are actually strict (otherwise the inequality is verified by all imaginable values). These bounds represent the support of the plausibility function, or in other words, the range of values covered by focal sets.
\end{itemize}

\begin{figure}
    \centering
    \begin{subfigure}[t]{0.48\linewidth}
        \centering
        \includegraphics[width=\linewidth]{Images/Chap_4/bel_independence_100_120.png}
        \caption{$\SAD$ envelopes using the product Copula}
        \label{fig:belief_independence}
    \end{subfigure}\hfill
    \begin{subfigure}[t]{0.48\linewidth}
        \centering
        \includegraphics[width=\linewidth]{Images/Chap_4/bel_independence_100_120_zoom.png}
        \caption{Detailed view of the rectangular section of \ref{fig:belief_independence}}
        \label{fig:belief_independence_zoom}
    \end{subfigure}
    \begin{subfigure}[t]{0.48\linewidth}
        \centering
        \includegraphics[width=\linewidth]{Images/Chap_4/bel_100_120.png}
        \caption{$\SAD$ envelopes using the Gaussian Copula}
        \label{fig:belief_gaussian}
    \end{subfigure}\hfill
    \begin{subfigure}[t]{0.48\linewidth}
        \centering
        \includegraphics[width=\linewidth]{Images/Chap_4/bel_100_120_zoom.png}
        \caption{Detailed view of the rectangular section of \ref{fig:belief_gaussian}}
        \label{fig:belief_gaussian_zoom}
    \end{subfigure}
    \caption{Plausibility levels of a cost curve for the product copula $C_\Pi$ and the Gaussian copula $C_R$. The green vertical line represents the true disparity. Rectangular sections from the left figures are detailed on the right. From \cite{malinowski_uncertainty_2024}.}\commanue{j'ai pas l'impression que tu mentionnes à quel pixel ces courbes correspondent}
    \label{fig:belief_curves}
\end{figure}

A visualization of different plausibility levels of the $\SAD$ cost curve are presented in \Cref{fig:belief_curves}, computed with the product copula $C_\Pi$ and a Gaussian copula $C_R$. Both copula have the same support $\Pl>0$, and the same value for $\Pl=1$. Envelopes of the other plausibility levels are however different. \Cref{fig:belief_independence_zoom} and \Cref{fig:belief_gaussian_zoom} display the fact that the values covered by the plausibility levels vary with the copula used. Plausibility levels $0.85$ and $0.5$ are more concentrated around plausibility level $1$ in the case of the product copula than in the case of the Gaussian copula. Conversely, plausibility level $0.9$ is closer to plausibility level $1$ in the case of the Gaussian copula. This is due to the fact that the Gaussian copula $C_R$ is more co-monotone than the product copula given the correlation matrix $R$ described in \eqref{eq:correlation}\commanue{avec cette macro on ne sait pas que c'est une équation et ensuite là tu pointes sigma et pas R a priori c'est pas le bon label}.
\begin{remark}
    Both the Gaussian and the product copula have the same support. This is because when computing the joint mass of marginal focal sets, the copula is regular enough to assign a non-null mass to every focal set. This would not have been the case if we took a copula close to the lower or upper Fréchet-Hoeffding bound (the lower bound is not a copula for $n>2$). Indeed, those copulas are less regular, and can assign a null mass to joint events, as we saw in the case of probabilities in example \ref{ex:copulas} and more recently in example \ref{ex:propagated_probability}.
\end{remark}

\subsection{Estimating Propagated Credal Sets Using Monte Carlo Sampling}\label{sec:montecarlo}
In \Cref{chap:joining_credal_sets}, we defined $3$ methods for creating multivariate credal sets: $\M_{robust}$, $\M_{mass}$ and $\M_{agg}$. In the case where marginals were possibility distributions, we saw that the bounds of $\M_{agg}$ and $\M_{mass}$ were the same on Cartesian products, so we only considered $\M_{mass}$ in our application. In the previous sections, we propagated the uncertainty by using the approach from $\M_{mass}$. In this section, we will try to estimate the propagated uncertainty using $\M_{robust}$, and compare it to previous results so as to evaluate whether or not we can use $\M_{mass}$ to outer/inner approximate $\M_{robust}$.

We remind here some definitions and results regarding $\M_{robust}$. $\M_{robust}$ is the convex hull $CH$ of the set of every CDF from $n$ marginal credal sets $\M_i$ (in our case $n=18$, one per pixel involved in the $\SAD$ computation) joined with a copula $C$ (\cref{eq:robust_set}):
\begin{align*}
    \M_{robust} = CH(\{F=C(F_1 \enum F_n),~F_i\in\M_i \})
\end{align*}
We can estimate $\M_{robust}$ using Monte Carlo samplings: we first generate marginal probability distributions $F_1\enum F_n$ belonging in their respective marginal credal sets \cite{troffaes_note_2017}, then sample from the joint CDF $F$\commanue{Donc là c'est la version résumé et ensuite tu détailles donc il faudra mettre un petit mot de tarnsition. Je sais pas To do so ou IN detail ou un autre choix}. For each of the $18$ considered pixels, we sample a marginal CDF $F_i$ belonging to the credal set defined in \cref{eq:pixel_possibility}. The sampling from credal sets is not random: we generate probability distributions in such a way that for each marginal event $A$, the probability range $[\Nec(A), \Pi(A)]$ is sampled uniformly. With this method, we get a good coverage of each marginal credal set. We are also insuring that lower and upper bounds on events are all reached at least once, in order to include ``extreme'' distributions in our simulations. Once we have a CDF $F_i$ from each marginal CDF $F_i$, we can sample from the joint CDF using the method detailed in \Cref{sec:sampling_copula}. This yields a noised version of both left and right images, where the noised dependency is modeled by the provided copula, and its distribution is coherent with the possibility distributions from \cref{eq:pixel_possibility}. We can then compute a noised version of the cost volume from those images. For each cost curve, this provides a Monte Carlo samples of the $\SAD$ cost curves.

We compute $10\,000$ samples CDFs from each marginal credal set to construct the multivariate CDF. Each multivariate CDF is then itself sampled $10\,000$ times. The Central Limit Theorem states that $N$ Monte Carlo samples gives you an approximation with a precision of around $\frac{1}{\sqrt{n}}$\commanue{pourquoi le n est petit dans la formule ?} (it also depends on the standard deviation of the density you are estimating, but it is a good enough approximation), so $N=10\,000$ provides a satisfying precision for our application.  

In practice, we do not simulate a full noised pair of images at once, as it would require to sample from a copula of very large dimension (the number of pixels in both images). This is not realistically feasible, even thought it would ensure that each time a pixel is considered in the cost volume, the same noise samples are used. We instead generate noise samples for each row separately: noised values of pixels will not change during the evaluation of a cost curve for different disparities, or between the cost curves of pixels of the same row. However, their values might change between cost curves of pixels belonging to different rows. For instance, let's consider a pixel $p=(row,~col)\in I_L$ for which we computed a noised intensity $i_p$. We will use the same noised value $i_p$ of intensity when computing the $\SAD$ of every pixel $q=(row,~col')\in I_L$. But when computing the $\SAD$ of every pixel $q=(row-1,~col')\in I_L$, we will use a different Monte Carlo draw $i_p'$ for its intensity, which will stay the same for every pixel of row $row-1$. Because we draw a high number of Monte Carlo draws for each row, proceeding as such should not be noticeable. 

\begin{remark}
    As we construct the correlation matrices $R$ based on the segmentation of the left and right images, the cost curves presented in \Cref{fig:montecarlo_gauss_100_120,fig:montecarlo_gauss_200_150} mostly use a different copula for each disparity. Providing the values of those matrices would necessitate to represent around $60$ $18\times18$ correlation matrices, and is thus not provided here.
\end{remark}

\begin{figure}
    \centering
    \begin{subfigure}[t]{1\linewidth}
        \centering
        \includegraphics[width=\linewidth]{Images/Chap_4/cost_curve_100_120.png}
        \caption{Plausibility levels and Monte Carlo sampling using a Gaussian copula}
        \label{fig:montecarlo_gauss_100_120_large}
    \end{subfigure}\\
    \begin{subfigure}[t]{0.45\linewidth}
        \centering
        \includegraphics[width=\linewidth]{Images/Chap_4/cost_curve_100_120_zoom1.png}
        \caption{Zoom over the first rectangle}
        \label{fig:montecarlo_gauss_100_120_zoom1}
    \end{subfigure}\hfill
    \begin{subfigure}[t]{0.45\linewidth}
        \centering
        \includegraphics[width=\linewidth]{Images/Chap_4/cost_curve_100_120_zoom2.png}
        \caption{Zoom over the second rectangle}
        \label{fig:montecarlo_gauss_100_120_zoom2}
    \end{subfigure}
    \caption{Plausibility levels and Monte Carlo sampling for a pixel of the left image. From \cite{malinowski_uncertainty_2024}.}
    \label{fig:montecarlo_gauss_100_120}
\end{figure}

\begin{figure}
    \centering
    \begin{subfigure}[t]{\linewidth}
        \centering
        \includegraphics[width=\linewidth]{Images/Chap_4/cost_curve_200_150.png}
        \caption{Plausibility levels and Monte Carlo sampling using a Gaussian copula}
        \label{fig:montecarlo_gauss_200_150_large}
    \end{subfigure}\\
    \begin{subfigure}[t]{0.45\linewidth}
        \centering
        \includegraphics[width=\linewidth]{Images/Chap_4/cost_curve_200_150_zoom1.png}
        \caption{Zoom over the first rectangle}
        \label{fig:montecarlo_gauss_200_150_zoom1}
    \end{subfigure}
    \hfill
    \begin{subfigure}[t]{0.45\linewidth}
        \centering
        \includegraphics[width=\linewidth]{Images/Chap_4/cost_curve_200_150_zoom2.png}
        \caption{Zoom over the second rectangle}
        \label{fig:montecarlo_gauss_200_150_zoom2}
    \end{subfigure}
    \caption{Plausibility levels and Monte Carlo sampling for a pixel of the left image.  From \cite{malinowski_uncertainty_2024}.}
    \label{fig:montecarlo_gauss_200_150}
\end{figure}

Monte Carlo draws using the Gaussian copula and marginals credal sets of \cref{eq:pixel_possibility} are plotted in \Cref{fig:montecarlo_gauss_100_120,fig:montecarlo_gauss_200_150}. They respectively correspond to the cost curves of pixel located at positions $(100,~120)$ and $(200,~150)$ on the left image\commanue{alors les deux figures ont actuellement la même légende tu pourrais préciser la localisation du pixel dans la légende et cerise sur le gateau tu pourrais afficher l'image des cônes ou dire à quoi correspondent les deux pixels choisis car perso moi je ne sais pas à quel endroit ça correspond.}. Plausibility levels from \Cref{fig:belief_curves}\commanue{est-ce qu'il ne faudrait pas dire que tu affiches les plausibility levels de ces pixels calculés comme décrit pécédemment car dans la figure 4.4 tu ne considères qu'un seul pixel} also appear for comparison. The support envelopes from plausibility levels $\Pl>0$ correctly contain all Monte Carlo samplings for all considered copulas. We can observe in \Cref{fig:montecarlo_gauss_200_150_zoom2} that plausibility levels sometimes fail to correctly grasp the fluctuations of the dispersion of the samples, even though they correctly contain Monte Carlo samples. More specifically, Monte Carlo draws are first dense around disparity $-37$, then seem to spread around $-35$, and finally regather around disparity $-32$. The plausibility envelopes are more regular in this disparity range. This illustrates the fact the ``true'' point-wise credal $\M_{robust}$ set described in \Cref{sec:robust_method} is different from the joint credal set $\M_{mass}$ from \Cref{sec:joint_mass}. 

Although some differences persist between those sets, \Cref{fig:montecarlo_gauss_100_120_zoom2,fig:montecarlo_gauss_200_150_zoom2} suggest that the point-wise credal set $\M_{robust}$ can be outer approximated by $\M_{mass}$ in our applications. To quantify this observation, we can compute the proportion of Monte Carlo samples contained inside the plausibility envelopes for different plausibility levels $\gamma$:
\begin{align}
    \textrm{coverage}_\gamma = \frac{\#\{\textrm{ Monte Carlo samples }\in[\underline{I}_\gamma,~\overline{I}_\gamma]\}}{\#\{\textrm{ Monte Carlo samples }\}}
\end{align}
The coverage for the considered values of $\gamma$ are presented in \Cref{tab:Coverage}, where the first two rows represent the coverage of \Cref{fig:montecarlo_gauss_100_120,fig:montecarlo_gauss_200_150}. The global coverage over the whole left image is presented in the last row of the table. The coverage is always $100\%$ for $\gamma=0$, which shows that every sample is contained inside the support envelopes and that $\M_{robust}$ seems to be a subset of $\M_{mass}$. On the other hand, plausibility level $\gamma=0.5$ contain most Monte Carlo samples which means that the bounds could be significantly reduced while still being a good estimation of $\M_{robust}$. Finally, the variation of the coverage for plausibility levels $0.85$ and $0.9$ translates the previous observation that $\M_{mass}$ does not capture the variations of $\M_{robust}$ bounds and that those two set can substantially differ. Having estimated the uncertainty of the cost volume, we will see in the following section how we can leverage the upper and lower bounds for different plausibility levels to improve the disparity map derived from the cost volume.  

\begin{table}[ht]
\centering
\begin{tabular}{|c|c|c|c|c|}
\hline
\rowcolor[HTML]{C0C0C0}
$p=(row,col)$ & $\gamma=0.9$  & $\gamma=0.85$ & $\gamma=0.5$  & $\gamma=0$   \\ \hline
\cellcolor[HTML]{C0C0C0}$(100,~120)$     & $64,5\%$ & $94,5\%$ & $99,0\%$ & $100\%$ \\ \hline
\cellcolor[HTML]{C0C0C0}$(200,~150)$     & $30,0\%$ & $82,6\%$ & $95,2\%$ & $100\%$ \\ \hline
\cellcolor[HTML]{C0C0C0}Global        & $41,1\%$ & $87,6\%$ & $96,8\%$ & $100\%$ \\ \hline
\end{tabular}
\caption{Average coverage for various plausibility levels $\gamma$ and for different pixels $p$ of the left image.}\label{tab:Coverage}
\end{table}

\subsection{Leveraging Confidence Envelopes for Potential Improvements}\label{sec:sad_improvements}

Knowing the uncertainty, represented as confidence envelopes in our case, can provide valuable insights into potential matches. This section outlines different observations suggesting that incorporating this information can enhance the performance of stereo-matching algorithms.

From the cost volume $C_V$, we usually apply a \textit{winner-takes-all} strategy to compute a disparity prediction. Given a pixel $(row, ~col)$, the predicted disparity $\tilde{d}$ is defined as:
\begin{align}\label{eq:winner_takes_all}
    \tilde{d}(row, ~col) = \argmin_{d}C_V(row, ~col, ~d)
\end{align}
A common metric for evaluating stereo algorithm performance is the proportion of pixels for which the absolute difference between the true disparity $d_{true}$ and the predicted disparity $\tilde{d}$ is less than one pixel. The score \( s \) is defined as:
\begin{equation}\label{eq:score_d1}
    s = \frac{\#\{(row, col) \text{ such that } |d_{\mathrm{true}}(row, col) - \tilde{d}(row, col)| < 1\}}{\#\{(row, col)\}}\,.
\end{equation}
This metric is used as we usually only consider integer disparities (otherwise we need to upsample the stereo images\commanue{oui alors là c'est un peu un raccourci dans le remarque, vu ton pipeline (pas de raffinement) tu es sur des disparités entières. Mais la métrique elle est valable sur ce qu'on veut.}) but the ground truth disparity can be any real number in the disparity range\commanue{bah tu peux utiliser cette métrique avec n'importe quelle valeur de seuil. Là question est de savoir quelle est la précision de ta VT. Bref je comprends pas trop la phrase.}. We thus consider a disparity to be ``correct'' if it is less than one pixel away from the true disparity.

Having computed envelopes $\underline{I}_\gamma$, $\overline{I}_\gamma$ on the cost volume for different plausibility levels $\gamma$, the cost curve is not unique anymore, and we can instead consider all cost curves $C_{\underline{I}_\gamma,\overline{I}_\gamma}$ contained within the plausibility envelopes. Instead of a single predicted disparity, we can now compute the set of all potential disparities, defined as the set of predicted disparities from every cost curve contained in the plausibility envelopes. Given a pixel $(row, col)$, a confidence level $\gamma \in [0, 1]$, and plausibility envelopes $\underline{I}_\gamma(row, ~col, ~d)$, $\overline{I}_\gamma(row, ~col, ~d)$, the set of potential disparities $D_\gamma^{row, col}$ is defined as\commanue{Il y a un signe $\leqslant$ au milieu d'un intervalle}:
\begin{align}
    D_\gamma^{row, col} = \{d \mid& d=\argmin_\delta C_{\underline{I}_\gamma,\overline{I}_\gamma}(row,~col,~\delta),\\
    &\forall C_{\underline{I}_\gamma,\overline{I}_\gamma}(row, ~col, ~\delta)\in[\underline{I}_\gamma(row, ~col, ~\delta)\leqslant \overline{I}_\gamma(row, ~col, ~\delta)]\}\nonumber
\end{align}
There is actually a simpler way of defining and computing $D_\gamma$, which is to notice that a disparity can be the minimum of a cost curve $C_{\underline{I}_\gamma,\overline{I}_\gamma}$ if and only if the lower bound $\underline{I}_\gamma$ for this disparity is less than every upper envelope $\overline{I}_\gamma$\commanue{je comprends pas le terme every upper enveloppe. Tu fixes gamma et donc tu regardes qu'une seule enveloppe et tu t'assures que tu sois en dessous du min de la borne sup, ou alors j'ai rien compris}:
\begin{equation}
    D_\gamma^{row, col} = \{d \mid \underline{I}_\gamma(row, ~col, ~d) \leq \min_\delta\overline{I}_\gamma(row, col, \delta)\}
\end{equation}
\Cref{fig:potential_disparities} provides a schematic example of $D_\gamma(row, col)$.
\begin{figure}[ht]
    \centering
    \includegraphics[width=\linewidth]{Images/Chap_4/Potentiel_disparities.png}
    \caption{Example of a set of potential disparities $D_\gamma$. The minimum of the upper envelope $\min\overline{I}_\gamma$ in dashed gray line. The vertical green line represents the true disparity.}
    \label{fig:potential_disparities}
\end{figure}

When the cost volume is considered without its uncertainty, many disparities are not correctly estimated by \cref{eq:winner_takes_all}. We thus want to quantify the potential score improvements contained in the set of potential disparities $D_\gamma$. To do so, we consider that $D_\gamma$ is a potential improvements if it contains the true disparity $d_{true}$. From this we can compute the proportion of potential improvements $s_\gamma^{\text{opt}}$ over the whole image, and compare it to the score $s$ computed without uncertainty from \cref{eq:score_d1}. The proportion of potential improvements, or optimal score, \( s_\gamma^{\text{opt}} \) achievable using the set of possible disparities is computed as follows\commanue{là ça fait redite donc introduis et définis ton score optimal et ensuite tu dis que tu le compares. Là tu introduis, dis que tu vas le comparer et ensuite tu le définis. C'est le tiercé dans le désordre.}:
\begin{equation}\label{eq:optimal_score}
    s_\gamma^{\text{opt}} = \frac{\#\{(row, col) \mid \min_{d \in D_\gamma^{row, col}} |d_{\mathrm{true}}(row, col) - d| < 1\}}{\#\{(row, col)\}}
\end{equation}
In \Cref{fig:potential_disparities}, we can see that the predicted disparity $\tilde{d}$ is far away from the true disparity $d_{true}$, but the set of possible disparity $D_\gamma$ does indeed contain, the true disparity. This would have been counted as a potential improvement.

\begin{remark}
    \Cref{eq:optimal_score} defines the optimal score that \textit{could} have been obtained if we used an ideal cost volume contained in the envelopes. We do not provide a method for determining this ideal cost volume. In reality, even with a good strategy to leverage uncertainty information to obtain a better disparity map, the new score $s$ from \cref{eq:score_d1} would be lower than $s_\gamma^{\text{opt}}$.
\end{remark}

We define the potential gain as \( \Delta s_\gamma = s_\gamma^{\text{opt}} - s \). $\Delta s_\gamma$ measures the proportion of pixel that benefit from the method, \ie pixels $(row, ~col)$ verifying:
\begin{equation}
    |d_{\mathrm{true}}(row, ~col) - \tilde{d}(row, ~col)| \geq 1 \quad \text{and} \quad \min_{d \in D_\gamma^{row, col}} |d_{\mathrm{true}}(row, ~col) - d| < 1
\end{equation}

Examples of optimal scores and potential gains for various $\gamma$ values are provided in \Cref{tab:optimal_score}. We can see that while the potential gain for $\gamma = 0.9$ is low, it increases significantly for lower values of $\gamma$. The base score is around $53\%$, which is a relatively low score for dense matching, but that was expected as we are using the $\SAD$ without \acrshort{sgm} regularization or additional post processing. With $\gamma=0.85$, we could reach an optimal score of $67\%$, and using $\gamma=0.5$ or $\gamma=0$, the score could at best lie between $74\%$ and $82\%$. This is a quite significant improvement for a cost volume computed with the $\SAD$ cost function. For comparison, a $\SAD$ cost volume regularized with \acrshort{sgm} algorithm would have a score $s$ of $68\%$.

\Cref{fig:improvements} displays the spatial distribution of pixels that can benefit from this method. Pixels in occluded regions, \ie pixels only appearing in one of the two images, are highlighted in orange. As occluded pixels cannot be improved as no true disparity exists, we do not consider them when computing the different statistics introduced in this section. Pixels where potential improvement can occur appear in blue.

We can see that pixels with potential improvements are not randomly distributed. These pixels are typically found in homogeneous areas where multiple disparities have low matching costs, where the cost curves would be similar to the one represented in \Cref{fig:montecarlo_gauss_200_150_large}.
\comroman{Avec le recul ça aurait aussi été intéressant de voir si le nombre de disparité ou la dispersion de $D_\gamma$ pouvait être une mesure de confiance, calculer la courbe d'erreur comme dans le papier sur l'amb}\commanu{à mettre en perspective donc...}. 

\begin{table}[ht]
\centering
\begin{tabular}{|c|c|c|c|c|}
\hline
\rowcolor[HTML]{C0C0C0} 
$s=52.87\%$                                & $\gamma=0.9$ & $\gamma=0.85$ & $\gamma=0.5$ & $\gamma=0$ \\ \hline
\cellcolor[HTML]{C0C0C0}$s_\gamma^{opt}$   & $56.92\%$    & $66.99\%$     & $74.28\%$    & $81.75\%$  \\ \hline
\cellcolor[HTML]{C0C0C0}$\Delta s_\gamma$ & $4.05\%$     & $16.11\%$     & $21.41\%$    & $28.87\%$  \\ \hline
\end{tabular}
\caption{Optimal score and potential gain for different plausibility $\gamma$. The potential gain is computed with regards to the usual score $s=52,87\%$.}\label{tab:optimal_score}
\end{table}

\begin{figure}[ht]
    \centering
    \begin{subfigure}[t]{0.48\linewidth}
        \centering
        \includegraphics[width=\linewidth]{Images/Chap_4/Improvements_Pl=0.9.png}
        \caption{$\gamma=0.9,~\Delta s_\gamma=4.05\%$}
        \label{fig:improvements_a}
    \end{subfigure}
    \begin{subfigure}[t]{0.48\linewidth}
        \centering
        \includegraphics[width=\linewidth]{Images/Chap_4/Improvements_Pl=0.85.png}
        \caption{$\gamma=0.85,~\Delta s_\gamma=16.11\%$}
        \label{fig:improvements_b}
    \end{subfigure}\\
    \begin{subfigure}[t]{0.48\linewidth}
        \centering
        \includegraphics[width=\linewidth]{Images/Chap_4/Improvements_Pl=0.5.png}
        \caption{$\gamma=0.5,~\Delta s_\gamma=21.41\%$}
        \label{fig:improvements_c}
    \end{subfigure}
    \begin{subfigure}[t]{0.48\linewidth}
        \centering
        \includegraphics[width=\linewidth]{Images/Chap_4/Improvements_Pl=0.0.png}
        \caption{$\gamma=0,~\Delta s_\gamma=28.87\%$}
        \label{fig:improvements_d}
    \end{subfigure}
    \caption{Spatial disposition of potential improvements for different values of $\gamma$. Pixels with potential improvements appear in blue. Occluded pixels, for which a correct disparity does not exist, appear in orange. Grayscale left image is displayed on the background. From \cite{malinowski_uncertainty_2024}.}
    \label{fig:improvements}
\end{figure}

In this Chapter\commanue{oh mon dieu une conclusion, j'ai pas l'habitude, peut-être faire une section dédiée, c'est toi qui voit}, we modeled the uncertainty on stereo images and propagated it until the cost volume. We compared different methods from \Cref{chap:joining_credal_sets} for propagating the uncertainty, and evaluated the potential improvements unlocked by this uncertainty estimation. However, the stereo algorithm considered for computing the cost volume was intentionally simple to reduce the complexity of the problem. For the \acrshort{co3d} mission, different cost functions (CENSUS, MC-CNN) and \acrshort{sgm} regularization would instead be used. Furthermore, although the uncertainty on input images have an influence on the final uncertainty, the major part of the disparity map uncertainty comes from the processing uncertainty of the stereo algorithm itself. Quantifying and propagating this processing uncertainty will be thus be the subject of \Cref{chap:epistemic_uncertainty}.

\pagebreak
\blankpage
    
    \setkeys{Gin}{draft=false}
    \chapter{Propagating the Epistemic Uncertainty of the Cost Volume to the DSM}\label{chap:epistemic_uncertainty}
\todoroman{Citer la publication d'origine pour les figure (license IEEE ``Copyright 2024 IEEE. Published in the 2024 IEEE International Geoscience and Remote Sensing Symposium (IGARSS 2024), scheduled for 7 - 12 July, 2024 in Athens, Greece. Personal use of this material is permitted. However, permission to reprint/republish this material for advertising or promotional purposes or for creating new collective works for resale or redistribution to servers or lists, or to reuse any copyrighted component of this work in other works, must be obtained from the IEEE. Contact: Manager, Copyrights and Permissions / IEEE Service Center / 445 Hoes Lane / P.O. Box 1331 / Piscataway, NJ 08855-1331, USA. Telephone: + Intl. 908-562-3966.'')}

The previous chapter detailed the propagation of  uncertainty in a stereo matching problem, using the $\SAD$ function to compute the cost volume. However, real cases of stereo matching usually do not use such simple cost functions, but rather more complex ones. In particular, the CARS pipeline that will process the \acrshort{co3d} data will use the CENSUS and MC-CNN cost functions  \cite{zabih_non-parametric_1994,zbontar_stereo_2016} to compute the cost volume, followed by a \acrshort{sgm} regularization. Those methods produce far better results and are thus favored for stereo pipelines. Unfortunately, we are not able to propagate the uncertainty to those cost functions and regularization as we did in the previous chapter as it would be too complex and computationally heavy. The same observation can be made for any deep leaning method \cite{laga_survey_2022} as the data is processed through many consecutive layers. Furthermore, the previous chapter was restricted to the propagation of the uncertainty from input images to the cost volume, but did not attempt to quantify the uncertainty of the stereo matching process itself, \ie the algorithm ability to correctly\commanue{properly? c'est pour éviter la répétition avec correct, à moins que ce soit voulu} identify the correct disparity. A bad performing stereo matching algorithm can produce errors of great magnitude regardless of the uncertainty on input images, while more advanced algorithm may even present good performance despite noised input images. Using the semantics of this thesis, we will refer to the uncertainty of the stereo algorithm itself as its epistemic uncertainty. Indeed, it does not result form any aleatoric process, but rather a lack of knowledge on how to correctly and automatically identify the correct disparity\commanue{pareil tu as une répétition entre correctly et correct dans la phrase}.

This aforementioned epistemic uncertainty\commanue{alors moi perso je mettrais This aforementioned epistemic uncertainty modeling in stereo matching mais c'est pas une obligation} has been the subject of many works in the literature \cite{hu_quantitative_2012,  poggi_confidence_2021,wang_uncertainty_2022}, designed for so-called ``classical methods'' using cost volumes obtained from cost functions, or for learning-based methods. This uncertainty is quantified using ``confidence measures'', associating a value between $0$ and $1$ to each predicted disparity, $0$ meaning that the prediction should be questioned,  and $1$ meaning that the prediction is most certainly correct. We refer to \Cref{sec:uncertainty_pandora} for more details on confidence measures in stereo matching pipelines. In this chapter, we will study how possibility distributions are able to model the epistemic uncertainty associated with a cost volume, and then deduce disparity confidence intervals from the possibility distribution. This approach is complementary to classical confidence estimations as it is not meant to indicate whether or not we trust a prediction but rather to provide information on \textit{where} the correct disparity should be. We then propagate those disparity confidence intervals in the rest of the stereo pipeline to obtain height confidence intervals, and evaluate their performances\commanue{Est-ce que dans l'intro de ce chapitre il ne faudrait pas aussi que tu rappelles que les utilisateurs veulent un intervalle à 90\% et donc blabla comme tu es gentil, tu vas chercher à répondre à ce besoin et donc construction à fournir un intervalle de disparité à 90\%}.

\todoroman{Vérifier que dans le chapitre 1 j'insiste bien sur ce que c'est une cost curve}
\section{Producing Confidence Intervals}
\subsection{Possibility Distributions as Uncertain Models for Cost Curves}
In this section, we will detail how possibility distributions can be used to model the epistemic uncertainty associated with cost curves. We first present a quick reminder of concepts and notations regarding cost volumes presented in \Cref{sec:stereo_matching}, as we will base our model on them. Cost volume based approaches, considered here, compare every pixel from the left image $I_L$ to pixels from the same row in the right image $I_R$, in a given disparity range. The comparison is done using a cost function $f$, measuring the dissimilarity between two windows centered around pixel $p$ and $q$. All evaluations using this cost function are stored in a cost volume $C_V$:
\begin{align}
	C_V(row, ~col, ~d) = f(I_L(row, ~col), I_R(row, ~col+d))
\end{align}
where $d$ is the considered disparity. In this chapter, we consider that the cost volume then\commanue{je mettrais pas le then} undergoes a \acrshort{sgm} regularization step, which modifies the cost values\commanue{its values} to take into account more global information. Based on the observation that the disparity map is usually piece-wise regular in a scene, \acrshort{sgm} regularization has been designed to increase the cost of disparities for which no consensus exist among neighbouring pixels. This way, only disparities that seem plausible and relatively regular compared to neighbouring disparities are favored in the cost volume.

For every pixel $(row, ~col)$ from the left image, we refer to its cost curve as the cost volume at coordinates $(row,~col)$ for every considered disparity $d$, \ie the cost volume for which we fixed the first two variables. Cost curves are of much importance as we estimate the disparity of a pixel solely based on its cost curve. Indeed, we define the predicted disparity $\tilde{d}$ of a pixel $(row, ~col)$ as:
\begin{align}
	\tilde{d} = \argmin_d C_V(row, ~col, ~d)
\end{align} 
\Cref{fig:tuto_dense_matching} presents a cost curve and its true disparity. We can see that the true disparity can be correctly estimated by looking at the minimum of the cost curve.

\begin{figure}
    \centering
    \begin{subfigure}[t]{0.073\linewidth}
        \centering
        \includegraphics[width=\linewidth]{Images/Chap_5/tuto_left_patch.png}
        \caption{}
        \label{fig:tuto_a}
    \end{subfigure}
    \hfill\begin{subfigure}[t]{\linewidth}
        \flushright
        \includegraphics[width=0.927\linewidth]{Images/Chap_5/tuto_right_patch.png}
        \caption{}
        \label{fig:tuto_b}
    \end{subfigure}
    \begin{subfigure}[t]{\linewidth}
        \centering
        \includegraphics[width=\linewidth]{Images/Chap_5/tuto_cost_curve.png}
        \caption{Cost curve}
        \label{fig:tuto_c}
    \end{subfigure}
    \caption{Cost curve (\Cref{fig:tuto_c}) obtained from comparing a patch from the left image (\Cref{fig:tuto_a}) to patches from the right image (\Cref{fig:tuto_b}). The matching patch and its corresponding true disparity are indicated using dashed green lines.}
    \label{fig:tuto_dense_matching}
\end{figure}

In the following, we propose to consider possibility distributions to model the uncertainty associated with the choice of the predicted disparity from a cost curve. The values taken by possibility distributions will be based on the available information \ie the values of the cost volume. Before getting into details, let us justify this model. Possibility distributions are relatively simple model to use in comparison with imprecise probabilities or belief functions for instance, as we only need to specify a constraint on the atoms, and not on every event. As such, they have been used to model the uncertainty associated with an expert's opinion in applications such as groundwater contamination for instance in \cite{bardossy_l-_1995} and \cite{baudrit_joint_2007}. Since cost curves result in both:
\begin{itemize}
	\item dissimilarity measures between patches 
	\item a semi-global fusion of the information contained in the cost volume due to \acrshort{sgm} regularization
\end{itemize}it does not seem far-stretched to consider them equivalent to an expert stating his opinion on how likely two pixels should be matched. For this reason, possibility distributions seem appropriate to model the epistemic uncertainty of the cost volume.

In order to use possibility distributions, we first need to transform the available information, in our case the values contained in the cost curves, into degrees of possibility. \Cref{eq:possibility}\commanue{Plutôt que de faire référence à une équation située 100 pages avant, tu pourrais juste dire The definition of possibility distribution ou un truc équivalent} imposes that the values must lie between $0$ and $1$, and that the value $1$ must be attained at least once. We therefore propose to normalize each cost curve so that its minimal dissimilarity value equals a possibility degree of $1$, and that greater dissimilarity values are closer to $0$ in possibility. However, simply normalizing the values of each cost curve between $0$ and $1$ would artificially stretch the cost curve as seen in \Cref{fig:cost_curves_b}. It is especially blatant in the case for the orange dashed curve from \Cref{fig:cost_curves_a} as the range of its values is quite narrow compared to the blue curve, but they are both stretched to $[0,~1]$ in \Cref{fig:cost_curves_b}. In order to avoid this effect, we instead normalize every cost curve using the global minimum and global maximum of the cost volume, as:
\begin{align}
	C_V^{norm}(row, ~col, ~d) = \frac{C_V(row, ~col, ~d) - \max_{r,c,\delta}C_V(r,~c,~\delta)}{\min_{r,c,\delta}C_V(r,~c,~\delta) - \max_{r,c,\delta}C_V(r,~c,~\delta)}\label{eq:normalized_cost_curve}
\end{align}
Minima of the cost curve become maxima with this normalization. One problem remains, it is that unless the global maximum of the cost volume is attained in a cost curve, the normalized cost curve will never reach $1$. Therefore, it will not be a possibility distribution. This problem can be observed in \Cref{fig:cost_curves_c}. We thus add a constant to the normalized cost curve to obtain a possibility distribution $\pi_{row,~col}(d)$:
\begin{align}
	\pi_{row,~col}(d) = C_V^{norm}(row, ~col, ~d) + 1 - \max_\delta C_V^{norm}(row, ~col, ~\delta)\label{eq:possibility_cost_curve}
\end{align}
\Cref{fig:cost_curves_d} displays the possibility distributions obtained from the cost curves of \Cref{fig:cost_curves_a}.

\begin{figure}
    \centering
    \begin{subfigure}[t]{0.47\linewidth}
        \centering
        \includegraphics[width=\linewidth]{Images/Chap_5/cost_curve_not_normalized.png}
        \caption{Two cost curves}
        \label{fig:cost_curves_a}
    \end{subfigure}\hfill
    \begin{subfigure}[t]{0.47\linewidth}
        \centering
        \includegraphics[width=\linewidth]{Images/Chap_5/cost_curve_bad_normalized.png}
        \caption{Normalized cost curves with local extrema}
        \label{fig:cost_curves_b}
    \end{subfigure}
    \begin{subfigure}[t]{0.47\linewidth}
        \centering
        \includegraphics[width=\linewidth]{Images/Chap_5/cost_curve_normalized.png}
        \caption{Normalized cost curves $C_V^{norm}$ with global extrema}
        \label{fig:cost_curves_c}
    \end{subfigure}\hfill
    \begin{subfigure}[t]{0.47\linewidth}
        \centering
        \includegraphics[width=\linewidth]{Images/Chap_5/cost_curve_possibility_distribution.png}
        \caption{Possibility distributions resulting from the cost curves}
        \label{fig:cost_curves_d}
    \end{subfigure}\hfill
    \caption{Transformation of cost curves (CENSUS + \acrshort{sgm} on Middlebury Cones) into possibility distributions. \Cref{fig:cost_curves_a} represents two cost curves that are normalized differently in \Cref{fig:cost_curves_b,fig:cost_curves_c}. \Cref{fig:cost_curves_d} uses the normalization of \Cref{fig:cost_curves_c} to create possibility distributions.}
    \label{fig:cost_curves_to_possibility}
\end{figure}

As stated previously, global extrema in \cref{eq:normalized_cost_curve} are employed to minimize the stretching effect when converting cost curves into possibility distributions. However, we also could have used the theoretical extrema of a cost curve instead. For instance, the CENSUS cost function on a $5\times5$ window provides values between $0$ and $C_{max}=24$. Adding \acrshort{sgm} regularization with penalty $P_2$ on $8$ directions yields cost volumes values between $0$ and $8\times(C_{max}+P_2)$ \cite{hirschmuller_accurate_2005}. However, this maximal cost is rarely attained in real case scenarios, and\commanue{je remplacerais and par therefore it} is too pessimistic and tends to over-compress the normalized cost curves. It is instead preferred to use global extrema of the cost volume for the normalization, as we suppose the best and worst match should have similar cost values across different scenes. This hypothesis is not restrictive for the images we consider in our stereo matching problem, as the size and diversity inside each scene lead to similar extrema. 

\subsection{From Possibilities to Disparity Confidence Intervals}
With the possibility distributions defined, our next objective is to establish a set of most possible disparities. We decided to aim for sets containing the true disparity $90\%$ of the time. This \( 90\%\) value was chosen following discussions with users and other experts working at CNES, IGN and more generally in the AI4GEO consortium\commanue{moi j'aurais introduit cette info avant dans l'intro et là tu peux le rappeler. Mais là c'est toi qui voit}. To define this set of most possible disparities, we compute the $\alpha$-cut (\cref{eq:alpha_cut}\commanue{tu fais un renvoi à l'équation, moi je renverrais à la définition. Bon c'est au même endroit de toute façon}) of the possibilities, or in other words, the set of all disparities $D_\alpha$ whose possibility is superior than $\alpha$:
\begin{align}
    D_\alpha=\{~ d ~|~ \pi_{row,col}(d)\geqslant\alpha\}\label{eq:set_of_possible_disparities}
\end{align}
By looking at possibility distributions obtained from different cost curves for which we know the true disparity, we first fixed the value $\alpha$ at $0.9$. In depth study of this parameter will be tested later, in order to see if it depends on the cost function, the type of scene considered, and to provide general guidelines on its optimal value. In the following, when the value of $\alpha$ is not specified, it will always be set at $0.9$. The fact that its value is the same as the $90\%$ confidence objective is a coincidence and one should not suppose that $\alpha$ and the confidence objective should be the same. Indeed raising the $\alpha$ value would decrease the size of the set $D_\alpha$ and therefore decrease the proportion of sets containing the true disparity, \ie the global confidence rate of intervals. \Cref{fig:disparity_intervals_a,fig:disparity_intervals_b} graphically represent $D_\alpha$ for the cost curves of \Cref{fig:cost_curves_to_possibility}.

\begin{remark}
    We are modelling the epistemic uncertainty of the cost curves using possibility distributions. In the rest of this chapter, we use a possibility distribution for each cost curve because we think it is a correct model in itself for the uncertainty we encounter. It is however possible to have a probabilistic interpretation of possibilities, which we will share in this remark as it may be of interest for people thinking the underlying uncertainty can be modeled by probabilities.
    
    We saw in \cref{eq:credal_set_possibility} from \Cref{chap:representation_of_uncertainty} that one way of interpreting $\pi_{row,col}$ is that it defines a set of probability distributions, \ie a credal set $\M$. We can also define $D_\alpha$ using this set $\M$, as:
    \begin{align}
        D_\alpha=\{~ d ~|~ \exists P\in\M ~\st~ P(d)\geqslant \alpha\}
    \end{align}
    or in plain words, $D_\alpha$ is the set of all disparities $d$ for which there exists a probability in $\M$ whose value is greater than $\alpha$ for $d$. We will not rely on this interpretation in the rest of this chapter, and instead only reason in terms of possibilities.
\end{remark}

\todoroman{Figure montrant différentes courbes de coût transformées, CENSUS MCCNN}

\Cref{eq:set_of_possible_disparities} defines a set of disparities $D_\alpha$ which is not necessarily convex. We will rather consider confidence disparity intervals $I_\alpha$ deduced from $D_\alpha$ in the rest of this chapter:
\begin{align}
    I_\alpha = [\underline{I}_\alpha,~\overline{I}_\alpha]=[\min D_\alpha, ~\max D_\alpha]\label{eq:confidence_disparity_intervals}
\end{align}
\Cref{fig:disparity_intervals_c,fig:disparity_intervals_d} graphically represent how $I_\alpha$ is determined for the two possibility distributions of the cost curves in \Cref{fig:cost_curves_to_possibility}. In the rest of this section, we will not make any distinction between disparity confidence intervals, disparity intervals, confidence intervals or simply intervals. A disparity interval is the convex envelope of its disparity set $D_\alpha$, which is a conservative approach as observed in \Cref{fig:disparity_intervals_b,fig:disparity_intervals_d}.

Considering intervals thus presents the advantage of working with convex sets and only requiring two scalars to describe the set. Users of the \acrshort{dsm}s produces by stereo photogrammetry are also used to consider confidence intervals \cite{oksanen_digital_2006, wang_robust_2015, panagiotakis_validation_2018, deschamps-berger_apport_2021, hugonnet_uncertainty_2022}. This will also facilitate further processing in the rest of the stereo pipeline, as we only need to take into account $2$ bounds to characterize the possible disparities, instead of a set of arbitrary shape.

\begin{figure}
    \centering
    \begin{subfigure}[t]{0.47\linewidth}
        \centering
        \includegraphics[width=\linewidth]{Images/Chap_5/disparity_interval_1.png}
        \caption{$D_\alpha$ for the blue possibility of \Cref{fig:cost_curves_d}}
        \label{fig:disparity_intervals_a}
    \end{subfigure}\hfill
    \begin{subfigure}[t]{0.47\linewidth}
        \centering
        \includegraphics[width=\linewidth]{Images/Chap_5/disparity_interval_2.png}
        \caption{$D_\alpha$ for the orange possibility of \Cref{fig:cost_curves_d}}
        \label{fig:disparity_intervals_b}
    \end{subfigure}
    \begin{subfigure}[t]{0.47\linewidth}
        \centering
        \includegraphics[width=\linewidth]{Images/Chap_5/disparity_interval_3.png}
        \caption{$I_\alpha$ for the blue possibility of \Cref{fig:cost_curves_d}}
        \label{fig:disparity_intervals_c}
    \end{subfigure}\hfill
    \begin{subfigure}[t]{0.47\linewidth}
        \centering
        \includegraphics[width=\linewidth]{Images/Chap_5/disparity_interval_4.png}
        \caption{$I_\alpha$ for the orange possibility of \Cref{fig:cost_curves_d}}
        \label{fig:disparity_intervals_d}
    \end{subfigure}\hfill
    \caption{Set of possible disparities $D_\alpha$ and disparity intervals $I_\alpha$ with the same cost curves as in  \Cref{fig:cost_curves_to_possibility}, with $\alpha=0.9$. \Cref{fig:disparity_intervals_a,fig:disparity_intervals_b} represent the set of possible disparities $D_\alpha$ from \cref{eq:set_of_possible_disparities} in gray. \Cref{fig:disparity_intervals_c,fig:disparity_intervals_d} represent disparity intervals $I_\alpha$ from \cref{eq:confidence_disparity_intervals} in gray. There is no difference between $D_\alpha$ and $I_\alpha$ for the blue curve, contrary to the orange dashed curve.}
    \label{fig:disparity_sets_and_intervals}
\end{figure}

To qualitatively evaluate the behaviour of disparity confidence intervals $I_\alpha$, we will look at their values for different pixels of the same row. The rows considered are presented in \Cref{fig:cones_with_rows}: the upper row ($80$) is smaller so more details can be observed, while other rows provide broader estimations of the intervals. By construction, $I_\alpha$ will always contain the predicted disparity $\tilde{d}$, and we want to see if it can contain the true disparity $d_{true}$ even when the predicted disparity is far from it. \Cref{fig:intervals_ambiguous_row_80_1,fig:intervals_ambiguous_row_180_1,fig:intervals_ambiguous_row_240_1,fig:intervals_ambiguous_row_290_1} 

\begin{figure}
    \centering
    \includegraphics[width=0.5\linewidth]{Images/Chap_5/cones_with_rows.png}
    \caption{Left stereo image from Middlebury Cones. Disparity intervals $I_\alpha$ along orange lines are detailed in \Cref{fig:intervals_ambiguous_row_80,fig:intervals_ambiguous_row_180,fig:intervals_ambiguous_row_240,fig:intervals_ambiguous_row_290}}
    \label{fig:cones_with_rows}
\end{figure}

\begin{figure}
    \centering
    \begin{subfigure}[t]{0.8\linewidth}
        \centering
        \includegraphics[width=\linewidth]{Images/Chap_5/intervals_ambiguous_area_row_80_1.png}
        \caption{$I_\alpha$ without regularization in low confidence areas}
        \label{fig:intervals_ambiguous_row_80_1}
    \end{subfigure}\hfill
    \begin{subfigure}[t]{0.8\linewidth}
        \centering
        \includegraphics[width=\linewidth]{Images/Chap_5/intervals_ambiguous_area_row_80_2.png}
        \caption{$I_\alpha$ with regularization in low confidence areas}
        \label{fig:intervals_ambiguous_row_80_2}
    \end{subfigure}
    \caption{$I_\alpha$ with and without regularization in low confidence area, for row $80$ of the image of \Cref{fig:cones_with_rows}. Low confidence areas are indicated by the gray areas.}
    \label{fig:intervals_ambiguous_row_80}
\end{figure}
\begin{figure}
    \centering
    \begin{subfigure}[t]{\linewidth}
        \centering
        \includegraphics[width=\linewidth]{Images/Chap_5/intervals_ambiguous_area_row_180_1.png}
        \caption{$I_\alpha$ without regularization in low confidence areas}
        \label{fig:intervals_ambiguous_row_180_1}
    \end{subfigure}\hfill
    \begin{subfigure}[t]{\linewidth}
        \centering
        \includegraphics[width=\linewidth]{Images/Chap_5/intervals_ambiguous_area_row_180_2.png}
        \caption{$I_\alpha$ with regularization in low confidence areas}
        \label{fig:intervals_ambiguous_row_180_2}
    \end{subfigure}
    \caption{$I_\alpha$ with and without regularization in low confidence area, for row $180$ of the image of \Cref{fig:cones_with_rows}. Low confidence areas are indicated by the gray areas.}
    \label{fig:intervals_ambiguous_row_180}
\end{figure}

\begin{figure}
    \centering
    \begin{subfigure}[t]{\linewidth}
        \centering
        \includegraphics[width=\linewidth]{Images/Chap_5/intervals_ambiguous_area_row_240_1.png}
        \caption{$I_\alpha$ without regularization in low confidence areas}
        \label{fig:intervals_ambiguous_row_240_1}
    \end{subfigure}\hfill
    \begin{subfigure}[t]{\linewidth}
        \centering
        \includegraphics[width=\linewidth]{Images/Chap_5/intervals_ambiguous_area_row_240_2.png}
        \caption{$I_\alpha$ with regularization in low confidence areas}
        \label{fig:intervals_ambiguous_row_240_2}
    \end{subfigure}
    \caption{$I_\alpha$ with and without regularization in low confidence area, for row $240$ of the image of \Cref{fig:cones_with_rows}. Low confidence areas are indicated by the gray areas.}
    \label{fig:intervals_ambiguous_row_240}
\end{figure}

\begin{figure}
    \centering
    \begin{subfigure}[t]{\linewidth}
        \centering
        \includegraphics[width=\linewidth]{Images/Chap_5/intervals_ambiguous_area_row_290_1.png}
        \caption{$I_\alpha$ without regularization in low confidence areas}
        \label{fig:intervals_ambiguous_row_290_1}
    \end{subfigure}\hfill
    \begin{subfigure}[t]{\linewidth}
        \centering
        \includegraphics[width=\linewidth]{Images/Chap_5/intervals_ambiguous_area_row_290_2.png}
        \caption{$I_\alpha$ with regularization in low confidence areas}
        \label{fig:intervals_ambiguous_row_290_2}
    \end{subfigure}
    \caption{$I_\alpha$ with and without regularization in low confidence area, for row $290$ of the image of \Cref{fig:cones_with_rows}. Low confidence areas are indicated by the gray areas.}
    \label{fig:intervals_ambiguous_row_290}
\end{figure}

\todoroman{Attention Manue je m'en suis arrêté là :) }
Simple seuil mais avec une base théorique

\todoroman{Déplacer la phrase qui suit:} This model will however encounter some limitations in regions near depth discontinuities: as the \acrshort{sgm} regularization attempts to impose continuity on disparities, it results in cost curves that do not favor the correct disparity. Considering that the cost curves are equivalent to experts opinion in those areas can be over optimistic. We will see how we can adapt the model on those regions later in this section.

Ambiguity et régularization dans les zones de faibles confiance etc

\todoroman{Utilisation de l'ambiguité. Montrer differentes zones d'aggrégation avec différents paramètres. Dire qu'on pourrait utiliser toute bonne mesure de confiance, mais celle-ci à l'avantage d'être explicable, efficace, et déjà implémentée et utilisée dans la chaine CARS. Filtrage de la carte de disparité.}
\todoroman{Soit je mets les résultats des tests de disparité ici, soit je les mets après avoir fait tout le pipeline CARS.}
\todoroman{Expliquer aussi comment on peut remonter le LiDAR HD en géométrie Epi + rasterization.}

\subsection{Propagation in the CARS Pipeline}
Rasterization and small ideas with Gabriela?
Impact of tiling : cost volume min and max is not the same from one tile to another (statistically it should be close, provided that the tile is big enough and the disparity interval is also large enough)

\subsection{Unexplored leads}
Patatoides sur les lignes de visées. Un pixel est représenté par des coordonées précises mais est en fait l'aggrégation d'information radiométriques issues d'une surface au sol. On pourrait donc imaginer que les lignes de visées ne sont pas une ligne mais plus un cône/cylindre (approx), et donc qu'on pourrait donc représenter la position d'un point 3D comme l'intersection de deux volumes (cône + intervalles de disparité).   

\section{Coregistration}
DEM Compare 


\section{Evaluation of Disparity Intervals}
\section{Ablation Studies}
\section{Elevation Confidence Intervals}
\begin{table}[ht]
    \centering
    \begin{tabular}{|c||c|c|c|c|}
    \hline
        Area & Date for LiDAR & Date of Pléiades & Coupling mode & GT Size (0.5m)\\
        \hline\hline
        Bordeaux & 2023-09-15 & 2022-08-04 & Stereo & $6001\times 6001$\\\hline
        Grenoble & 2021-09-05 & 2020-09-17 & Stereo & $10 001\times 10 001$ \\\hline
        Jotunheinem & 2019-08-27 & 2019-08-27 & Stereo & \begin{tabular}{@{}c@{}} $7127\times 7298$ and\\ $3912\times2880$\end{tabular} \\\hline
        Langfjordjokelen & 2018-09-01 & 2018-09-01 & Stereo & $5841\times 3689$\\\hline
        Monaco & 2021-05-13 & 2020-08-30 & Stereo & $10 001\times 10 001$\\\hline
        Montpellier & 2021-05-28 & 2021-10-17 & Tri Stereo & $8001\times 8001$\\\hline
        Paris & 2023-03-03 & 2023-05-31 & Tri stereo & $10 001\times 10 001$\\\hline
        Peyto & 2016-09-13 & 2016-09-13 & Stereo & $13240\times 17874$\\\hline
        Pic du Midi & 2021-10-02 & 2021-10-16 & Stereo & $10001 \times 12001$ \\\hline
        Toulouse & 2022-05-29 & 2022-06-28 & Tri Stereo & $12001\times 8001$\\\hline
        Yosemite 1 &  \begin{tabular}{@{}c@{}} \tiny{2017-05-02 (snow) +} \\ \tiny{2014-08-27 (LiDAR)}\end{tabular} &  2017-05-01 & Tri stereo & \tiny{$17733\times 16800$ (3m)}\\\hline
        Yosemite 2 &  \begin{tabular}{@{}c@{}} \tiny{2017-05-02 (snow) +} \\ \tiny{2014-08-27 (LiDAR)}\end{tabular} &  2017-05-02 &  Tri stereo & \tiny{$17733\times 16800$ (3m)}\\ \hline
    \end{tabular}
    \caption{Acquisition date of Pléiade stereo or tri-stereo images, and the LiDAR ground truth}
    \label{tab:dates_pleiades_lidar_hd}
\end{table}

\begin{table}[ht]
    \centering
    \begin{tabular}{|c||c|c|c|c|}
    \hline
        Area & Date for LiDAR & Date of Pléiades & Coupling mode & GT Size (0.5m)\\
        \hline\hline
        Bordeaux & 2023-09-15 & 2022-08-04 & Stereo & $6001\times 6001$\\\hline
        Grenoble & 2021-09-05 & 2020-09-17 & Stereo & $10 001\times 10 001$ \\\hline
        Monaco & 2021-05-13 & 2020-08-30 & Stereo & $10 001\times 10 001$\\\hline
        Montpellier & 2021-05-28 & 2021-10-17 & Tri Stereo & $8001\times 8001$\\\hline
        Paris & 2023-03-03 & 2023-05-31 & Tri stereo & $10 001\times 10 001$\\\hline
        Pic du Midi & 2021-10-02 & 2021-10-16 & Stereo & $10001 \times 12001$ \\\hline
        Toulouse & 2022-05-29 & 2022-06-28 & Tri Stereo & $12001\times 8001$\\\hline
    \end{tabular}
    \caption{Acquisition date of Pléiade stereo or tri-stereo images, and the LiDAR ground truth}
    \label{tab:my_label}
\end{table}

\todoroman{Rasterizer le min max du lidar HD pour avoir une idée des intervalles.}
\pagebreak
\blankpage
    
    \setkeys{Gin}{draft=false}
    \chapter{Producing Elevation Confidence Intervals}\label{chap:elevation_intervals}
In \Cref{chap:epistemic_uncertainty}, we presented a method for computing disparity confidence intervals alongside each predicted disparity $\tilde{d}$. In the photogrammetry pipeline, the predicted disparity is triangulated, filtered (\Cref{sec:triangulation}) and then rasterized (\Cref{sec:rasterization}) to obtain the final \acrshort{dsm}. We will see in this chapter how we can apply those steps to disparity confidence intervals in order to propagate them into elevation confidence intervals, associated with the \acrshort{dsm} values. This chapter takes up work and data published in \cite{malinowski_robust_2024}.

In \Cref{sec:related_work} of \Cref{chap:stereophotogrammetry}, we presented related work where confidence intervals were computed based on a \acrshort{dsm}, mainly in \cite{oksanen_digital_2006,panagiotakis_validation_2018,deschamps-berger_apport_2021}. In those studies, a single global confidence interval associated with a \acrshort{dsm} is computed \textit{a posteriori}, \ie based solely on the \acrshort{dsm} (and reference points), regardless of the method used to obtain it (\ie, photogrammetry, \acrshort{lidar}, \acrshort{radar}). The methodology presented in this chapter estimates the uncertainty independently for each pixel, leading to small confidence intervals in confident areas, and bigger confidence intervals where the algorithms may have performed badly. We differ in this regard to previous work, and it does not seem relevant to compare our intervals to theirs, as their most similar characteristics is their name ``interval'', but neither share the same form (single \vs multiple intervals), nor are based on the same data.

\section{From Disparity Intervals to Elevation Intervals}\label{sec:elevation_intervals}
\subsection{Triangulation of Disparity Intervals}
As indicated by their name, disparity confidence intervals are expressed in pixel of disparities. Disparities are used to determine  pairs of intersecting lines of sight from different sensors, which are then triangulated to obtain a 3D point. With disparity confidence intervals, we now have 3 pairs of line of sight for each pixel instead of just one:
\begin{itemize}
    \item The pair composed using the predicted disparity $\tilde{d}$
    \item The two pairs composed using the upper and lower disparities from $I_\alpha=[\underline{I}_\alpha, ~\overline{I}_\alpha]$.
\end{itemize}
Intersecting each pair of line of sight yields $3$ 3D points, as presented in \Cref{fig:pairs_of_line_of_sight}. We deduce the first point $(x, ~y, ~z)$ from predicted disparity $\tilde{d}$, and the two other points $(\underline{x}, ~\underline{y}, ~\underline{z})$, $(\overline{x}, ~\overline{y}, ~\overline{z})$ are deduced from $\underline{I}_\alpha$ and $\overline{I}_\alpha$ respectively.
\begin{figure}
    \centering
    \includegraphics[width=\linewidth]{Images/Chap_6/Pairs_of_line_of_sight.png}
    \caption{Triangulation of the three pairs of lines of sight. The angle between lines of sight is exaggerated for the purpose of this illustration.}
    \label{fig:pairs_of_line_of_sight}
\end{figure}

Depending on the disposition of satellites, as well as which image is selected as the reference image, it is both possible that $\overline{z}\leqslant z \leqslant \underline{z}$ or $\underline{z}\leqslant z \leqslant \overline{z}$. In the following, and for simplicity, we will consider that $\underline{z}\leqslant z \leqslant \overline{z}$. This is not constraining, as we can just change the notations to unsure this inequality holds. We therefore have a 3D confidence ``interval'', defined as every point from the reference line of sight between $(\underline{x}, ~\underline{y}, ~\underline{z})$ and $(\overline{x}, ~\overline{y}, ~\overline{z})$. For instance in \Cref{fig:pairs_of_line_of_sight}, the reference line of sight is the one originating from the left satellite, and the confidence ``interval'' is the portion of this line of sight between the two orange crosses. It is not an interval \textit{per se}, as \acrshort{rpc} are polynomials and not straight lines, but they are approximated by straight lines in the computations, so the distinction is superfluous. 

The point cloud obtained from the triangulation of the disparity map is then filtered, as detailed in \Cref{sec:triangulation}. If a 3D point is filtered, then we naturally also filter its corresponding interval with it. 

\subsection{Rasterization of 3D ``Intervals''}
The final step of the pipeline is the rasterization, as our objective is to produce elevation confidence intervals associated with every value contained in the \acrshort{dsm}. However, rasterizing intervals along lines of sight as it stands raises an issue: we are not guaranteed that the rasterized intervals will remain coherent with the final \acrshort{dsm} when projected on the regular grid. A solution to circumvent this issue is to consider that the planimetric shift $\Delta XY=\sqrt{(x-\underline{x})^2+(y-\underline{y})^2}$ is small in comparison to the altimetric shift $\Delta Z=\sqrt{(z-\underline{z})^2}$ and that we can therefore neglect it. This hypothesis depends only on the incidence angle $\beta$ of the reference image. Indeed, it holds that $\dfrac{\Delta Z}{\Delta XY}=\dfrac{1}{\tan(\beta)}$, where $\beta$ is the incidence angle as depicted in \Cref{fig:incidence_angle}. Ideally, if the reference image has no incidence angle, then the $\Delta XY$ shift is null. \Cref{tab:angle_coupling_pleiades} details the different incidence angles encountered in our experiments. The incidence angles are typically around $10\degree$, \ie $\Delta Z$ is around $5.6$ times bigger than $\Delta XY$. For the scene in Paris where the incidence reaches $6.1\degree$, the ratio is around $10$, while the only scene with $21.3\degree$, near Peyto Lake, Canada, has a ratio around $2.6$. For this last acquisition, the hypothesis that $\Delta XY$ is small compared to $\Delta Z$ is debatable, but we will still neglect $\Delta XY$ in the rasterization in order to stay consistent across our experiments. This will also allow verifying if this hypothesis can impact the performance of our method.

The lower and upper bounds can thus be aligned directly above and below the predicted 3D point $(x, ~y, ~z)$, and they thus become:
\begin{align}
    (\overline{x}, ~\overline{y}, ~\overline{z}) ~\rightarrow~ (x, ~y, ~\overline{z})\\
    (\underline{x}, ~\underline{y}, ~\underline{z}) ~\rightarrow~ (x, ~y, ~\underline{z})
\end{align}
\Cref{fig:planimetric_shift} illustrates this modification, where the orange points along the line of sight are shifted in the $(X,~Y)$ plane to be vertically aligned with $(x, ~y, ~z)$. 

\begin{figure}
    \centering
    \includegraphics[width=0.8\linewidth]{Images/Chap_6/Incidence_angle.png}
    \caption{Acquisition angles of satellites}
    \label{fig:incidence_angle}
\end{figure}

\begin{figure}
    \centering
    \includegraphics[width=0.6\linewidth]{Images/Chap_6/Planimetric_shift.png}
    \caption{Aligning the confidence interval bounds along a line of sight}
    \label{fig:planimetric_shift}
\end{figure}

We remind here the general formulation of the rasterization step. Let $(x,~y)$ be a cell of the \acrshort{dsm}, and $PC$ be the point cloud considered in the rasterization. The value of the \acrshort{dsm} is computed using a weighted mean of $PC$:
\begin{align}
    \DSM(x,y) &= \dfrac{\sum\limits_{(x_i, y_i, z_i)\in PC}z_i\cdot w(x_i, y_i)}{\sum\limits_{(x_i,y_i,z_i)\in PC} w(x_i, y_i)}
\end{align}
where weights $w(x_i,y_i)$ are positive scalars. In CARS, the weights are computed using a Gaussian distribution, but other pipelines also use Inverse Distance Weightings which works similarly. In practice, only the points within a given distance of the center $(x,y)$ of the cell are considered in the mean. 

Because it holds that for any point: $\underline{z}\leqslant z\leqslant\overline{z}$, then computing the \acrshort{dsm}s independently using $(x, ~y, ~z)$, $(x, ~y, ~\overline{z})$ and $(x, ~y, ~\underline{z})$ will ensure the consistency of the resulting \acrshort{dsm}s:
\begin{eqnarray}
    \dfrac{\sum\limits_{(x_i, y_i, \underline{z}_i)\in PC}\underline{z}_i\cdot w(x_i, y_i)}{\sum\limits_{(x_i,y_i, \underline{z}_i)\in PC} w(x_i, y_i)}
    \leqslant&
    \dfrac{\sum\limits_{(x_i, y_i, z_i)\in PC}z_i\cdot w(x_i, y_i)}{\sum\limits_{(x_i,y_i,z_i)\in PC} w(x_i, y_i)}
    &\leqslant
    \dfrac{\sum\limits_{(x_i, y_i, \overline{z}_i)\in PC}\overline{z}_i\cdot w(x_i, y_i)}{\sum\limits_{(x_i,y_i,\overline{z}_i)\in PC} w(x_i, y_i)}\nonumber\\
    &&\nonumber\\
    \underline{\DSM}(x,y) \leqslant& \DSM(x,y) &\leqslant\overline{\DSM}(x,y)
\end{eqnarray}
where $\underline{\DSM}(x,y)$ is the \acrshort{dsm} computed using points $(x, ~y, ~\underline{z})$, and $\overline{\DSM}(x,y)$ is the \acrshort{dsm} computed using points $(x, ~y, ~\overline{z})$. For each cell $\DSM(x,y)$ we have now computed an elevation interval $[\underline{\DSM}(x,y),~\overline{\DSM}(x,y)]$.

As rasterization is the final step of the stereo pipeline, we now have propagated the confidence intervals all the way to the end of the pipeline while ensuring their coherency with the predicted \acrshort{dsm}. Additionally, The production of elevation confidence intervals did not influence the values of the final \acrshort{dsm}. We now need to evaluate the elevation confidence intervals on real data to verify if the potential errors occurring during the transformation of disparity intervals into elevation intervals do not question their accuracy.

\section{Acquiring and Processing Data for Evaluation}
\subsection{Satellite and DSM datasets}
This section will present the different images and ground truth \acrshort{dsm}s used to evaluate elevation intervals. We will use ground truth \acrshort{dsm}s that can be categorized into two categories.

The first category is composed of \acrshort{dsm}s acquired over mountainous regions containing glaciers. They have been kindly provided by Etienne Berthier from LEGOS, Liss Marie Andreassen from the Norwegian Water Resources and Energy Directorate (NVE) and Brian Menounos from the Natural Sciences and Engineering Research Council of Canada and the Tula Foundation (Hakai Institute). The data were acquired in the following regions:
\begin{itemize}
    \item Two different \acrshort{dsm}s in the mountainous region of Jotunheinem, Norway. We will refer to the two \acrshort{dsm}s as Hellmem (\Cref{fig:miniature_Hellmem}) and Graasubreen (\Cref{fig:miniature_Graasubreen}).
    \item Langfjordjøkelen glacier, Norway (\Cref{fig:miniature_Langfjordjokelen}).
    \item Mountains near Peyto lake, in the Alberta province of Canada. (\Cref{fig:miniature_Peyto})
\end{itemize}
The ground truth were acquired using a \acrshort{lidar} sensor, rasterized at 50 cm resolution. We did not apply the rasterization ourselves and only had access to the rasterized \acrshort{dsm}s.

The second category of ground truth data comes from the \acrshort{lidar} HD program (\cite{monnet_lidarhd_2023}, \url{https://geoservices.ign.fr/lidarhd}) which intends to cover the entirety of the French territory (except for French Guiana) by the end of 2026. It is a very important source of information, with around $10$ measured points per m$^2$ and a relative planimetric accuracy of 50 cm. As of the time we write this thesis, data over the French territory are not fully (publicly) available. From the data available at the time of this work, we selected different regions of interest in order to have a variety of landscapes (rural, urban, seaside \etc). Those landscapes all presented strong elevation variations in order to present a challenge for the stereo pipeline. We also want to determine if our method behaves differently depending on the nature of the scene. Here is the list of the considered regions:
\begin{itemize}
    \item The city of Bordeaux (\Cref{fig:miniature_Bordeaux})
    \item The city of Montpellier (\Cref{fig:miniature_Montpellier})
    \item The city of Paris (\Cref{fig:miniature_Paris})
    \item The city of Toulouse (\Cref{fig:miniature_Toulouse})
    \item Valleys and mountains near Grenoble, in the French Alps (\Cref{fig:miniature_Grenoble})
    \item Mediterranean coastline with the city of Monaco (\Cref{fig:miniature_Monaco})
    \item Valleys and mountains at Pic du Midi, french Pyrenees (\Cref{fig:miniature_pic_du_midi})
\end{itemize}
\acrshort{lidar} point clouds were then rasterized at 50 cm resolution with the same Gaussian rasterization method as in the stereo pipeline. \Cref{tab:dates_pleiades_lidar_hd} presents the different acquisition dates and shape of the ground truth \acrshort{dsm}s.

We used Pléiades images for producing both \acrshort{dsm}s and intervals that will be compared to the ground truth. We did not directly order the Pléiades acquisitions, as it is costly and hard to synchronize with the acquisition date of the ground truth. From our available catalog, we selected stereo pairs that were acquired as close as possible from the acquisition date of the ground truth, or if multiple months separated the two acquisitions, we rather selected a similar period of the previous/next year to minimize seasonal changes. \Cref{tab:dates_pleiades_lidar_hd} allows comparing dates of acquisition of the \acrshort{lidar} and Pléiades images.

We saw in \Cref{sec:elevation_intervals} that the planimetric and altimetric accuracy depend on the incidence and convergence angles of the lines of sights of the satellites. \Cref{tab:angle_coupling_pleiades} presents those angles for the different Pléiades images.

\begin{table}[ht]
    \centering
    \begin{tabular}{|c||c|c|c|}
    \hline
        Area & Date for LiDAR & Date of Pléiades &  GT Size (0.5m)\\
        \hline\hline
        Bordeaux & 2023-09-15 & 2022-08-04 & $6001\times 6001$\\\hline
        Grenoble & 2021-09-05 & 2020-09-17 & $10 001\times 10 001$ \\\hline
        Hellmem 1 & 2019-08-27 & 2019-08-27 & $7127\times 7298$ \\\hline
        Graasubreen 2 & 2019-08-27 & 2019-08-27 & $3912\times2880$ \\\hline
        Langfjordjøkelen & 2018-09-01 & 2018-09-01 & $5841\times 3689$\\\hline
        Monaco & 2021-05-13 & 2020-08-30 & $10 001\times 10 001$\\\hline
        Montpellier & 2021-05-28 & 2021-10-17 & $8001\times 8001$\\\hline
        Paris & 2023-03-03 & 2023-05-31 & $10 001\times 10 001$\\\hline
        Peyto & 2016-09-13 & 2016-09-13 & $13240\times 17874$\\\hline
        Pic du Midi & 2021-10-02 & 2021-10-16 & $10001 \times 12001$ \\\hline
        Toulouse & 2022-05-29 & 2022-06-28 & $12001\times 8001$\\\hline
    \end{tabular}
    \caption{Acquisition date of Pléiades stereo or tri-stereo images, and the \acrshort{lidar} ground truth}
    \label{tab:dates_pleiades_lidar_hd}
\end{table}

\begin{table}[ht]
    \centering
    \begin{tabular}{|c||c|c|c|c|}
        \hline
        Area & Convergence angle & Left incidence angle & Right incidence angle \\
        \hline\hline
        Bordeaux & 23.1\degree & 9.7\degree & 14.2\degree \\\hline
        Grenoble & 28.3\degree & 12.4\degree & 16.1\degree \\\hline
        Jotunheinem & 22.5\degree & 10.6\degree & 14.5\degree \\\hline
        Langfjordjøkelen & 21.4\degree & 10.2\degree & 13.8\degree \\\hline
        Monaco & 29.2\degree & 12.8\degree & 18.1\degree \\\hline
        Montpellier & 7.7\degree & 11.5\degree & 15\degree\\\hline
        Paris & 4.9\degree& 6.1\degree & 7.9\degree\\\hline
        Peyto & 17.2\degree & 21.3\degree & 21.6\degree \\\hline
        Pic du Midi & 29.0\degree & 13.7\degree & 15.3\degree \\\hline
        Toulouse & 11.4\degree & 12.5\degree & 18.7\degree\\\hline
    \end{tabular}
    \caption{Relevant angles of stereo pairs of Pléiades images. See \Cref{fig:incidence_angle} for a schematic representation of those angles.}
    \label{tab:angle_coupling_pleiades}
\end{table}

\begin{figure}
    \centering
    \begin{subfigure}[t]{0.48\linewidth}
        \flushleft
        \includegraphics[width=\linewidth]{Images/Chap_6/miniature_Bordeaux.png}
        \caption{\acrshort{rgb} image of Bordeaux (Pléiades © \acrshort{cnes} 2022, Distribution AIRBUS DS)}
        \label{fig:miniature_Bordeaux_rgb}
    \end{subfigure}\hfill
    \begin{subfigure}[t]{0.48\linewidth}
        \flushright
        \includegraphics[width=\linewidth]{Images/Chap_6/miniature_Bordeaux_gt.png}
        \caption{\acrshort{lidar} HD \acrshort{dsm}}
        \label{fig:miniature_Bordeaux_gt}
    \end{subfigure}
    \caption{\acrshort{rgb} image of Bordeaux and its associated ground truth \acrshort{dsm}.}
    \label{fig:miniature_Bordeaux}
\end{figure}

\begin{figure}
    \centering
    \begin{subfigure}[t]{0.48\linewidth}
        \flushleft
        \includegraphics[width=\linewidth]{Images/Chap_6/miniature_Grenoble.png}
        \caption{\acrshort{rgb} image of a region near Grenoble (Pléiades © \acrshort{cnes} 2020, Distribution AIRBUS DS)}
        \label{fig:miniature_Grenoble_rgb}
    \end{subfigure}\hfill
    \begin{subfigure}[t]{0.48\linewidth}
        \flushright
        \includegraphics[width=\linewidth]{Images/Chap_6/miniature_Grenoble_gt.png}
        \caption{\acrshort{lidar} HD \acrshort{dsm}}
        \label{fig:miniature_Grenoble_gt}
    \end{subfigure}
    \caption{\acrshort{rgb} image of Grenoble and its associated ground truth \acrshort{dsm}.}
    \label{fig:miniature_Grenoble}
\end{figure}

\begin{figure}
    \centering
    \begin{subfigure}[t]{0.48\linewidth}
        \flushleft
        \includegraphics[width=\linewidth]{Images/Chap_6/miniature_Graasubreen.png}
        \caption{\acrshort{rgb} image of Graasubreen (Pléiades © \acrshort{cnes} 2019, Distribution AIRBUS DS)}
        \label{fig:miniature_Graasubreen_rgb}
    \end{subfigure}\hfill
    \begin{subfigure}[t]{0.48\linewidth}
        \flushright
        \includegraphics[width=\linewidth]{Images/Chap_6/miniature_Graasubreen_gt.png}
        \caption{\acrshort{lidar} \acrshort{dsm}}
        \label{fig:miniature_Graasubreen_gt}
    \end{subfigure}
    \caption{\acrshort{rgb} image of Graasubreen and its associated ground truth \acrshort{dsm}.}
    \label{fig:miniature_Graasubreen}
\end{figure}

\begin{figure}
    \centering
    \begin{subfigure}[t]{0.48\linewidth}
        \flushleft
        \includegraphics[width=\linewidth]{Images/Chap_6/miniature_Hellmem.png}
        \caption{\acrshort{rgb} image of Hellmem (Pléiades © \acrshort{cnes} 2019, Distribution AIRBUS DS)}
        \label{fig:miniature_Hellmem_rgb}
    \end{subfigure}\hfill
    \begin{subfigure}[t]{0.48\linewidth}
        \flushright
        \includegraphics[width=\linewidth]{Images/Chap_6/miniature_Hellmem_gt.png}
        \caption{\acrshort{lidar} \acrshort{dsm}}
        \label{fig:miniature_Hellmem_gt}
    \end{subfigure}
    \caption{\acrshort{rgb} image of Hellmem and its associated ground truth \acrshort{dsm}.}
    \label{fig:miniature_Hellmem}
\end{figure}

\begin{figure}
    \centering
    \begin{subfigure}[t]{0.48\linewidth}
        \flushleft
        \includegraphics[width=\linewidth]{Images/Chap_6/miniature_Langfjordjokelen.png}
        \caption{\acrshort{rgb} image of Langfjordjøkelen (Pléiades © \acrshort{cnes} 2018, Distribution AIRBUS DS)}
        \label{fig:miniature_Langfjordjokelen_rgb}
    \end{subfigure}\hfill
    \begin{subfigure}[t]{0.48\linewidth}
        \flushright
        \includegraphics[width=\linewidth]{Images/Chap_6/miniature_Langfjordjokelen_gt.png}
        \caption{\acrshort{lidar} \acrshort{dsm}}
        \label{fig:miniature_Langfjordjokelen_gt}
    \end{subfigure}
    \caption{\acrshort{rgb} image of Langfjordjøkelen and its associated ground truth \acrshort{dsm}.}
    \label{fig:miniature_Langfjordjokelen}
\end{figure}

\begin{figure}
    \centering
    \begin{subfigure}[t]{0.48\linewidth}
        \flushleft
        \includegraphics[width=\linewidth]{Images/Chap_6/miniature_Monaco.png}
        \caption{\acrshort{rgb} image of Monaco (Pléiades © \acrshort{cnes} 2020, Distribution AIRBUS DS)}
        \label{fig:miniature_Monaco_rgb}
    \end{subfigure}\hfill
    \begin{subfigure}[t]{0.48\linewidth}
        \flushright
        \includegraphics[width=\linewidth]{Images/Chap_6/miniature_Monaco_gt.png}
        \caption{\acrshort{lidar} HD \acrshort{dsm}}
        \label{fig:miniature_Monaco_gt}
    \end{subfigure}
    \caption{\acrshort{rgb} image of Monaco and its associated ground truth \acrshort{dsm}.}
    \label{fig:miniature_Monaco}
\end{figure}

\begin{figure}
    \centering
    \begin{subfigure}[t]{0.48\linewidth}
        \flushleft
        \includegraphics[width=\linewidth]{Images/Chap_6/miniature_Montpellier.png}
        \caption{\acrshort{rgb} image of Montpellier (Pléiades © \acrshort{cnes} 2020, Distribution AIRBUS DS)}
        \label{fig:miniature_Montpellier_rgb}
    \end{subfigure}\hfill
    \begin{subfigure}[t]{0.48\linewidth}
        \flushright
        \includegraphics[width=\linewidth]{Images/Chap_6/miniature_Montpellier_gt.png}
        \caption{\acrshort{lidar} HD \acrshort{dsm}}
        \label{fig:miniature_Montpellier_gt}
    \end{subfigure}
    \caption{\acrshort{rgb} image of Montpellier and its associated ground truth \acrshort{dsm}.}
    \label{fig:miniature_Montpellier}
\end{figure}

\begin{figure}
    \centering
    \begin{subfigure}[t]{0.48\linewidth}
        \flushleft
        \includegraphics[width=\linewidth]{Images/Chap_6/miniature_Paris.png}
        \caption{\acrshort{rgb} image of Paris (Pléiades © \acrshort{cnes} 2023, Distribution AIRBUS DS)}
        \label{fig:miniature_Paris_rgb}
    \end{subfigure}\hfill
    \begin{subfigure}[t]{0.48\linewidth}
        \flushright
        \includegraphics[width=\linewidth]{Images/Chap_6/miniature_Paris_gt.png}
        \caption{\acrshort{lidar} HD \acrshort{dsm}}
        \label{fig:miniature_Paris_gt}
    \end{subfigure}
    \caption{\acrshort{rgb} image of  and its associated ground truth \acrshort{dsm}.}
    \label{fig:miniature_Paris}
\end{figure}

\begin{figure}
    \centering
    \begin{subfigure}[t]{0.48\linewidth}
        \flushleft
        \includegraphics[width=\linewidth]{Images/Chap_6/miniature_Peyto.png}
        \caption{\acrshort{rgb} image of Peyto (Pléiades © \acrshort{cnes} 2016, Distribution AIRBUS DS)}
        \label{fig:miniature_Peyto_rgb}
    \end{subfigure}\hfill
    \begin{subfigure}[t]{0.48\linewidth}
        \flushright
        \includegraphics[width=\linewidth]{Images/Chap_6/miniature_Peyto_gt.png}
        \caption{\acrshort{lidar} HD \acrshort{dsm}}
        \label{fig:miniature_Peyto_gt}
    \end{subfigure}
    \caption{\acrshort{rgb} image of Peyto and its associated ground truth \acrshort{dsm}.}
    \label{fig:miniature_Peyto}
\end{figure}

\begin{figure}
    \centering
    \begin{subfigure}[t]{0.48\linewidth}
        \flushleft
        \includegraphics[width=\linewidth]{Images/Chap_6/miniature_Pic_du_midi.png}
        \caption{\acrshort{rgb} image of Pic du Midi (Pléiades © \acrshort{cnes} 2021, Distribution AIRBUS DS)}
        \label{fig:miniature_pic_du_midi_rgb}
    \end{subfigure}\hfill
    \begin{subfigure}[t]{0.48\linewidth}
        \flushright
        \includegraphics[width=\linewidth]{Images/Chap_6/miniature_Pic_du_midi_gt.png}
        \caption{\acrshort{lidar} HD \acrshort{dsm}}
        \label{fig:miniature_pic_du_midi_gt}
    \end{subfigure}
    \caption{\acrshort{rgb} image of Pic du Midi and its associated ground truth \acrshort{dsm}.}
    \label{fig:miniature_pic_du_midi}
\end{figure}

\begin{figure}
    \centering
    \begin{subfigure}[t]{0.48\linewidth}
        \flushleft
        \includegraphics[width=\linewidth]{Images/Chap_6/miniature_Toulouse.png}
        \caption{\acrshort{rgb} image of Toulouse (Pléiades © \acrshort{cnes} 2022, Distribution AIRBUS DS)}
        \label{fig:miniature_Toulouse_rgb}
    \end{subfigure}\hfill
    \begin{subfigure}[t]{0.48\linewidth}
        \flushright
        \includegraphics[width=\linewidth]{Images/Chap_6/miniature_Toulouse_gt.png}
        \caption{\acrshort{lidar} HD \acrshort{dsm}}
        \label{fig:miniature_Toulouse_gt}
    \end{subfigure}
    \caption{\acrshort{rgb} image of Toulouse and its associated ground truth \acrshort{dsm}.}
    \label{fig:miniature_Toulouse}
\end{figure}

\subsection{Water Masks}
Both \acrshort{lidar} data and stereo correlation present poor results on water surfaces. For the \acrshort{lidar} data, wavelengths are often absorbed by still water and do not provide the sensor with any feedback (although turbid water can usually send back some signal). For moving waters, the provided signal is often noisy and needs post-processing to remove artifacts. Stereophotogrammetry has usually even poorer results, because open waters are texture-less or uniform surfaces on which the correlation does not perform well. Furthermore, for Pléiades acquisitions, the water may have moved between images, which prevents water pixels from being correctly triangulated. For those reasons, we chose to add a water mask on images of cities crossed by a river (Bordeaux, Monaco, Montpellier, Paris, and Toulouse) or near the seaside (Monaco). The water mask was obtained using an algorithm developed by \acrshort{cnes}, which uses a random forest trained on existing high resolution mapping of surface waters \cite{pekel_high-resolution_2016} and the Normalized Difference Water Index \cite{gao_ndwinormalized_1996}. \Cref{fig:paris_watermask_2} presents a water mask produced on the image of Paris.

\begin{figure}
    \centering
    \begin{subfigure}[t]{0.48\linewidth}
        \flushleft
        \includegraphics[width=\linewidth]{Images/Chap_6/miniature_Paris.png}
        \caption{\acrshort{rgb} image of Paris (Pléiades © \acrshort{cnes} 2023, Distribution AIRBUS DS)}
        \label{fig:paris_watermask_1}
    \end{subfigure}\hfill
    \begin{subfigure}[t]{0.48\linewidth}
        \flushright
        \includegraphics[width=\linewidth]{Images/Chap_6/watermask_Paris.png}
        \caption{Water mask}
        \label{fig:paris_watermask_2}
    \end{subfigure}
    \caption{\acrshort{rgb} image of Paris and its associated water mask. Water is indicated by blue pixels}
    \label{fig:paris_watermask}
\end{figure}

\subsection{Co-registration}
The \acrshort{dsm}s obtained from \acrshort{lidar} data and from stereophotogrammetry are not necessarily in the same projection reference system, and do not use the same altitude reference. Moreover, due to the limited accuracy of \acrshort{rpc} models and GPS measures of the \acrshort{lidar}, some planimetric and elevation biases may exist between the two \acrshort{dsm}s. Those biases do not allow their comparison as such. We therefore reproject the ground truth data in the same reference system as their corresponding \acrshort{dsm} produced by the CARS pipeline. We then rectify the planimetric and altimetric biases by applying the method presented in \cite{nuth_co-registration_2011}. This process is called co-registration. We quickly present the method for estimating the altimetric and planimetric biases, and refer to the original publication for additional details.

It is possible to observe that the measured differences between two shifted \acrshort{dsm}s vary with the slope angle and the orientation of the slope. The different parameters influencing the measured variations of elevation, presented in \Cref{fig:coregistration}, are the following:
\begin{itemize}
    \item We note $sl$ the angle of a slope. $sl=0\degree$ corresponds to a flat slope and $sl=90\degree$ to a vertical slope.
    \item We note $\psi$ the azimuth of the slope. If the slope faces north, then $\psi=0\degree$. If it faces west, $\psi=90\degree$ \etc
    \item The direction of the planimetric shift between \acrshort{dsm}s is given by an angle $\beta$, with the same conventions as the azimuth. The magnitude of the shift is $B$.
    \item $dh$ refers to measured local variations of elevation, and $\tilde{dh}$ is the global elevation shift between \acrshort{dsm}s.
\end{itemize}
The relation linking all those parameters is the following \cite{nuth_co-registration_2011}:
\begin{align}
    dh = B\cos(\psi-\beta)\tan(sl)+\tilde{dh}
\end{align}

\begin{figure}[ht!]
    \centering
    \includegraphics[width=\linewidth]{Images/Chap_6/Coregistration.png}
    \caption{Planimetric shift of magnitude $B$ and altimetric shift $\tilde{dh}$ between two \acrshort{dsm}s. We can see that local variations $\tilde{dh}$ vary depending on the slope $sl$. This diagram is in 2D, the angle of the planimetric shift and the azimuth of the slope are therefore not represented.}
    \label{fig:coregistration}
\end{figure}

The three unknowns are $B$, $\beta$ and $\tilde{dh}$, as the slope parameters can be computed by any \acrshort{gis} software from the \acrshort{dsm}s. For instance, the slope of $\DSM_{true}$ is computed as:
\begin{align}\label{eq:slope}
    sl(x,y) = \dfrac{1}{8}\sqrt{\left(\dfrac{\DSM_{true}\ast k_x(x,y)}{r_x}\right)^2 + \left(\dfrac{\DSM_{true}\ast k_y(x,y)}{r_y}\right)^2}
\end{align}
where $\ast$ denotes a convolution with two kernels $k_x=\begin{bmatrix}
-1 & 0 & 1\\
-2 & 0 & 2 \\
-1 & 0 & 1
\end{bmatrix}$, $k_y=\begin{bmatrix}
-1 & -2 & -1\\
0 & 0 & 0 \\
1 & 2 & 1
\end{bmatrix}$, and $r_x$, $r_y$ are the resolution in $x$ and $y$. We will use the slope in the evaluation of the different metrics.

The unknowns $B$, $\beta$ and $\tilde{dh}$ are determined by a least square optimization problem. Because the \acrshort{dsm} is not expressed analytically, the optimization is not guaranteed to be exact. Multiple iterations of the planimetric and altimetric shift estimation lead to a better final result. An illustration of the co-registration process is presented in \Cref{fig:coregistration_image}.

The problem we encounter is that the \acrshort{dsm} obtained from photogrammetry already possess some errors. In the least squared minimization problem, the residuals computed from the difference between the ground truth \acrshort{dsm} and the photogrammetry \acrshort{dsm} therefore also contain those errors. This deteriorates the quality of the co-registration. However, it remains the best solution for co-registering \acrshort{dsm}s. We apply this co-registration to our data, before computing any metric. The different metrics we consider will be presented in \Cref{sec:metrics_elevation}, while the following sections details the parameters used in our stereo pipeline.

\begin{figure}
    \centering
    \begin{subfigure}[t]{0.48\linewidth}
        \centering
        \includegraphics[width=\linewidth]{Images/Chap_6/coregisration_planimetric_shift_gt_toulouse.png}
        \caption{Ground Truth \acrshort{dsm}}
        \label{fig:coregistration_planimetric_gt}
    \end{subfigure}\hfill
    \begin{subfigure}[t]{0.48\linewidth}
        \flushright
        \includegraphics[width=\linewidth]{Images/Chap_6/coregisration_planimetric_shift_cars_toulouse.png}
        \caption{CARS \acrshort{dsm}}
        \label{fig:coregistration_planimetric_cars}
    \end{subfigure}\\
    \begin{subfigure}[t]{\linewidth}
        \centering
        \includegraphics[width=\linewidth]{Images/Chap_6/coregisration_altimetric_shift_toulouse.png}
        \caption{Elevation along a row of the \acrshort{dsm}}
        \label{fig:coregistration_altimetric}
    \end{subfigure}
    \caption{Planimetric shift and altimetric shift from the co-registration step over the city of Toulouse. In \subref{fig:coregistration_planimetric_gt} and \subref{fig:coregistration_planimetric_cars}, reference points in orange are located at the same row and columns in the two \acrshort{dsm}s. The planimetric shift is indicated with blue arrows. \subref{fig:coregistration_altimetric} presents the altimetric shift between the two \acrshort{dsm}s.}
    \label{fig:coregistration_image}
\end{figure}

\subsection{Configuration of the Photogrammetry Pipeline}
This section provides technical information about the configuration of the CARS pipeline used to process the satellite images.

We used the Copernicus DEM with a 30m resolution as the reference altitude, as the SRTM elevation models were not available for high latitudes, such as Norway for instance. We used the same range of considered disparity for every image: [-50pix, 50pix]. This range is quite high, especially for acquisition with a high altitude ratio $r_{alt}$, but ensures we are not limited to a restrictive range of possible elevations.

Although we developed a method for creating disparity intervals which works for both the CENSUS and MC-CNN cost functions, we will only present results using the CENSUS cost function. Two reasons motivate this choice. First, due to the quantity of scenes and the variety of analyses completed, it would quickly become quite overwhelming to present each result for both cost functions, especially when there is no major difference between the two. Secondly, we observed in the previous chapter that disparity intervals obtained using CENSUS were less accurate than those using MC-CNN. As we favor accuracy above interval size, we chose to use the CENSUS cost function throughout our results, as validating the accuracy requirements for CENSUS ensures that they are also validated using MC-CNN. We verified that there was no unexpected behavior for MC-CNN intervals across scenes, and since nothing notable appeared, we did not delve into so many details as we did for the CENSUS intervals.  

We detail here the list of the parameters used in the pipeline. The CENSUS cost function was computed using a 5$\times$5 window. The \acrshort{sgm} penalties used for regularizing the CENSUS cost volume were $P_1=8$ and $P_2=32$. Disparity intervals were produced using a possibility threshold of $0.9$. Intervals were regularized in low confidence areas for which the confidence from ambiguity was below $0.6$, after being minimized by a minitive kernel of shape $(1, 2\cdot k_{amb}+1)=(1,5)$. For each regularized pixel, we took the $90$\ith upper and $10$\ith lower quantiles over a low confidence area that extended $n_\N=3$ rows above and below. A V-fit refinement step was used to get sub-pixel disparities, and a median filter of shape $(3\times 3)$ was applied to the resulting disparity map. Pixels that did not validate the cross-checking criterion from \Cref{eq:cross-checking} were not considered for triangulation. Triangulated 3D points were filtered using \Cref{eq:statistical_outlier} with $k=5$ and $N=50$, and \Cref{eq:small_components} with $D_{max}=3$ m and $N_{min}=50$. The rasterization from \Cref{eq:rasterization} used $\sigma=0.3$ m and $r=3$ m.

\section{Metrics for Evaluating Elevation Intervals}\label{sec:metrics_elevation}
We introduce here the metrics used to evaluate the performance of elevation intervals. In order to stay consistent with \Cref{sec:metrics_disparity}, we consider the metrics introduced for disparity intervals and adapt them to elevation intervals. 

We will refer to the true elevation as $\DSM_{true}$, the predicted elevation as $\DSM$ and the elevation intervals as $[\underline{\DSM},~\overline{\DSM}]$.

\subsubsection{Elevation Accuracy Metric}
Similarly to \Cref{sec:metrics_disparity}, the first metric is the proportion of correct intervals, \ie the proportion of intervals containing the ground truth. We call this metric the elevation accuracy $Z_{acc}$:
\begin{align}\label{eq:relative_accuracy_elevation}
    Z_{acc} = \dfrac{\#\{\DSM_{true} ~|~ \st ~\DSM_{true}\in [\underline{\DSM},~\overline{\DSM}]\}}{\#\{\DSM_{true}\}}
\end{align}
\Cref{fig:Intervals_elevation_metrics_Z_acc} illustrates what this metric represents. $Z_{acc}$ is similar to $acc$ defined in \Cref{eq:accuracy} from \Cref{sec:metrics_disparity}, we thus want to maximize $Z_{acc}$. We also keep the objective of $90\%$ accuracy for our method, which was validated for the disparity intervals.

\begin{figure}
    \centering
    \includegraphics[width=0.7\linewidth]{Images/Chap_6/Intervals_elevation_metrics_Z_acc.png}
    \caption{Schematic representation of $Z_{acc}$}
    \label{fig:Intervals_elevation_metrics_Z_acc}
\end{figure}

\subsubsection{Residual Elevation Error Metric}
We are also interested in evaluating the magnitude of the errors. We thus define the residual elevation error $Z_{\varepsilon}$ for all intervals that do not contain the ground truth as:
\begin{align}
    Z_\varepsilon=\dfrac{1}{r_{alt}} \cdot\median \left( \min(|\DSM_{true}-\overline{\DSM}, ~ |\DSM_{true}-\underline{\DSM}|)\right)\label{eq:residual_error_elevation}
\end{align}
where $r_{alt}$ is the disparity to altitude ratio (or altimetric ratio), \ie the elevation difference resulting from a shift of one disparity. \Cref{fig:Intervals_elevation_metrics_Z_eps} illustrates what this metric represents. We want $Z_\varepsilon$ error to be as close to 0 as possible. $r_{alt}$ has been computed alongside the epipolar grids (\Cref{sec:epipolar_geometry}). Dividing by the ratio $r_{alt}$ effectively converts $\min(|\DSM_{true}-\overline{\DSM}|, |\DSM_{true}-\underline{\DSM}|)$ from meters to pixels. As an error of one disparity pixel in the dense matching step can result in an elevation error of 1m in a scene, and 2m in another, we divide by $r_{alt}$ in order to be able to compare stereo pairs with different convergence angles. \Cref{tab:elevation_metrics_global} contains the different altimetric ratios for the data we consider. $Z_{\varepsilon}$ is similar to $\varepsilon$ defined in \Cref{eq:residual_error} from \Cref{sec:metrics_disparity}.

\begin{figure}
    \centering
    \includegraphics[width=0.7\linewidth]{Images/Chap_6/Intervals_elevation_metrics_Z_eps.png}
    \caption{Schematic diagram of the quantity represented by $Z_{\varepsilon}$}
    \label{fig:Intervals_elevation_metrics_Z_eps}
\end{figure}

\subsubsection{Relative Elevation Size Metric}
Those two previous metrics allow measuring the accuracy and errors of elevation intervals. To evaluate the size of intervals, we use the relative elevation size $Z_{size}$ defined as :
\begin{align}\label{eq:relative_size_elevation}
    Z_{size} = \dfrac{1}{r_{alt}}\median\left(\overline{\DSM}-\underline{\DSM}\right)
\end{align}
We divide by the altimetric ratio for the same reasons as in $Z_\varepsilon$, \ie to compare stereo pairs with different convergence angles. \Cref{fig:Intervals_elevation_metrics_Z_size} illustrates what this metric represents. We want the relative elevation size to be as close to zero as possible. $Z_{size}$ is similar to $s_{rel}$ defined in \Cref{eq:relative_size_disparity} from \Cref{sec:metrics_disparity}.

\begin{figure}
    \centering
    \includegraphics[width=0.7\linewidth]{Images/Chap_6/Intervals_elevation_metrics_Z_size.png}
    \caption{Schematic diagram of the quantity represented by $Z_{size}$}
    \label{fig:Intervals_elevation_metrics_Z_size}
\end{figure}

\begin{remark}
    In the rasterization process, 3D points from low confidence area are mixed with points from high confidence areas. We therefore are unable to define interval regularization areas in the final \acrshort{dsm}. The relative over-estimation $o_{rel}$ has therefore no equivalent for elevation intervals.  
\end{remark}

\section{Elevation Intervals Results}\label{sec:results_elevation}
\begin{table}[ht]
    \centering
    \begin{tabular}{|c||c|c|c|c|c|c|}
        \hline
        Scene & $Z_{acc}$ & $Z_\varepsilon$ (pix) & $Z_{size}$ (pix) & $r_{alt}$ (m/pix) & invalid
        \\\hline\hline
        Bordeaux & 89.3$\%$ & 0.56 & 4.18 & 1.24  & 21.5$\%$\\\hline
        Graasubreen & 99.7$\%$ & 0.12 & 1.96 & 1.32  & 27.5$\%$\\\hline
        Hellmem & 98.8$\%$ & 2.76 & 2.02 & 1.32  & 34.7$\%$\\\hline
        Grenoble & 93.1$\%$ & 0.87 & 3.99 & 1.06  & 6.4$\%$\\\hline
        Langfjordjøkelen & 99.4$\%$ & 0.30 & 2.16 & 1.35  & 13.9$\%$\\\hline
        Monaco & 90.3$\%$ & 0.57 & 3.76 & 0.99  & 41.9$\%$\\\hline
        Montpellier & 89.1$\%$ & 0.28 & 2.05 & 3.64  & 1.9$\%$\\\hline
        Paris & 84.6$\%$ & 0.46 & 2.01 & 5.78  & 3.0$\%$\\\hline
        Peyto & 98.9$\%$ & 0.18 & 2.08 & 1.66  & 18.3$\%$\\\hline
        Pic du midi & 98.1$\%$ & 0.20 & 2.41 & 1.01  & 7.5$\%$\\\hline
        Toulouse & 92.0$\%$ & 0.38 & 2.04 & 2.43  & 6.9$\%$\\\hline
    \end{tabular}
    \caption{Elevation metrics for the different stereo pairs. The last column indicates the proportion of invalid pixels in the considered \acrshort{dsm}s.}
    \label{tab:elevation_metrics_global}
\end{table}

\begin{figure}
    \begin{subfigure}[t]{0.553\linewidth}
        \flushleft
        \includegraphics[width=\linewidth]{Images/Chap_6/Graasubreen_dsm.png}
        \caption{Photogrammetry \acrshort{dsm}. Orange area is detailed in \Cref{fig:Graasubreen_zoom}}
        \label{fig:Graasubreen_dsm}
    \end{subfigure}
    \begin{subfigure}[t]{0.447\linewidth}
        \flushright
        \includegraphics[width=\linewidth]{Images/Chap_6/Graasubreen_error.png}
        \caption{\acrshort{dsm} with wrong intervals in red.}
        \label{fig:Graasubreen_error}
    \end{subfigure}
    \caption{\acrshort{dsm} without and with wrong intervals over Graasubreen scene}
    \label{fig:Graasubreen_global}
\end{figure}

\begin{figure}
    \begin{subfigure}[t]{0.49\linewidth}
        \flushleft
        \includegraphics[width=\linewidth]{Images/Chap_6/Graasubreen_dsm_zoom.png}
        \caption{Photogrammetry \acrshort{dsm}. Orange line is detailed in \subref{fig:Graasubreen_zoom_row}}
        \label{fig:Graasubreen_dsm_zoom}
    \end{subfigure}\hfill
    \begin{subfigure}[t]{0.49\linewidth}
        \flushright
        \includegraphics[width=\linewidth]{Images/Chap_6/Graasubreen_error_zoom.png}
        \caption{\acrshort{dsm} with wrong intervals in red.}
        \label{fig:Graasubreen_error_zoom}
    \end{subfigure}
    \begin{subfigure}[t]{\linewidth}
        \centering
        \includegraphics[width=\linewidth]{Images/Chap_6/Graasubreen_row_1500.png}
        \caption{\acrshort{dsm}, ground truth and elevation intervals along the orange line of \subref{fig:Graasubreen_dsm_zoom}}
        \label{fig:Graasubreen_zoom_row}
    \end{subfigure}
    \caption{Details of the Graasubreen \acrshort{dsm}. This area corresponds to the orange square from \Cref{fig:Graasubreen_dsm}.}
    \label{fig:Graasubreen_zoom}
\end{figure}


\begin{figure}
    \begin{subfigure}[t]{0.549\linewidth}
        \flushleft
        \includegraphics[width=\linewidth]{Images/Chap_6/Toulouse_dsm.png}
        \caption{Photogrammetry \acrshort{dsm}. Orange area is detailed in \Cref{fig:toulouse_zoom}}
        \label{fig:toulouse_dsm}
    \end{subfigure}
    \begin{subfigure}[t]{0.451\linewidth}
        \flushright
        \includegraphics[width=\linewidth]{Images/Chap_6/Toulouse_error.png}
        \caption{\acrshort{dsm} with wrong intervals in red.}
        \label{fig:toulouse_error}
    \end{subfigure}
    \caption{\acrshort{dsm} without and with wrong intervals over Toulouse scene}
    \label{fig:toulouse_global}
\end{figure}

\begin{figure}
    \begin{subfigure}[t]{0.49\linewidth}
        \flushleft
        \includegraphics[width=\linewidth]{Images/Chap_6/Toulouse_dsm_zoom.png}
        \caption{Photogrammetry \acrshort{dsm}. Orange line is detailed in \Cref{fig:toulouse_zoom_row}}
        \label{fig:toulouse_dsm_zoom}
    \end{subfigure}\hfill
    \begin{subfigure}[t]{0.49\linewidth}
        \flushright
        \includegraphics[width=\linewidth]{Images/Chap_6/Toulouse_error_zoom.png}
        \caption{\acrshort{dsm} with wrong intervals in red.}
        \label{fig:toulouse_error_zoom}
    \end{subfigure}
    \begin{subfigure}[t]{\linewidth}
        \centering
        \includegraphics[width=\linewidth]{Images/Chap_6/Toulouse_row_4720.png}
        \caption{\acrshort{dsm}, ground truth and elevation intervals along the orange line of \Cref{fig:toulouse_dsm_zoom}}
        \label{fig:toulouse_zoom_row}
    \end{subfigure}
    \caption{Detailed \acrshort{dsm} over Wilson square, Toulouse. This area corresponds to the orange square from \Cref{fig:toulouse_dsm}.}
    \label{fig:toulouse_zoom}
\end{figure}

\begin{figure}
    \begin{subfigure}[t]{0.3\linewidth}
        \flushleft
        \includegraphics[width=\linewidth]{Images/Chap_6/Grenoble_dsm_full_3600-4000_4600-5000.png}
        \caption{Photogrammetry \acrshort{dsm}. Orange square is detailed in \subref{fig:grenoble_dsm_full}}
        \label{fig:grenoble_dsm_full}
    \end{subfigure}\hfill
    \begin{subfigure}[t]{0.3\linewidth}
        \centering
        \includegraphics[width=\linewidth]{Images/Chap_6/Grenoble_dsm_zoom_3600-4000_4600-5000.png}
        \caption{Photogrammetry \acrshort{dsm} zoom. Orange line is detailed in \subref{fig:grenoble_zoom_row}}
        \label{fig:grenoble_dsm_zoom}
    \end{subfigure}\hfill
    \begin{subfigure}[t]{0.3\linewidth}
        \flushright % k=h1/h2*w2/w1
        \includegraphics[width=0.807\linewidth]{Images/Chap_6/Grenoble_error_zoom_3600-4000_4600-5000.png}
        \caption{\acrshort{dsm} with wrong intervals in red.}
        \label{fig:grenoble_error_zoom}
    \end{subfigure}\\
    \begin{subfigure}[t]{\linewidth}
        \centering
        \includegraphics[width=\linewidth]{Images/Chap_6/Grenoble_row_3800.png}
        \caption{\acrshort{dsm}, ground truth and elevation intervals along the orange line of \Cref{fig:grenoble_dsm_zoom}}
        \label{fig:grenoble_zoom_row}
    \end{subfigure}
    \caption{Detailed \acrshort{dsm} over the mountainous region near Grenoble.}
    \label{fig:grenoble_zoom}
\end{figure}

\begin{figure}
    \begin{subfigure}[t]{0.3\linewidth}
        \flushleft
        \includegraphics[width=\linewidth]{Images/Chap_6/Monaco_dsm_full_6000-6500_4800-5300.png}
        \caption{Photogrammetry \acrshort{dsm}. Orange square is detailed in \subref{fig:monaco_dsm_full}}
        \label{fig:monaco_dsm_full}
    \end{subfigure}\hfill
    \begin{subfigure}[t]{0.3\linewidth}
        \centering
        \includegraphics[width=\linewidth]{Images/Chap_6/Monaco_dsm_zoom_6000-6500_4800-5300.png}
        \caption{Photogrammetry \acrshort{dsm} zoom. Orange line is detailed in \subref{fig:monaco_zoom_row}}
        \label{fig:monaco_dsm_zoom}
    \end{subfigure}\hfill
    \begin{subfigure}[t]{0.3\linewidth}
        \flushright
        \includegraphics[width=0.83\linewidth]{Images/Chap_6/Monaco_error_zoom_6000-6500_4800-5300.png}
        \caption{\acrshort{dsm} with wrong intervals in red.}
        \label{fig:monaco_error_zoom}
    \end{subfigure}\\
    \begin{subfigure}[t]{\linewidth}
        \centering
        \includegraphics[width=\linewidth]{Images/Chap_6/Monaco_row_6250.png}
        \caption{\acrshort{dsm}, ground truth and elevation intervals along the orange line of \subref{fig:monaco_dsm_zoom}}
        \label{fig:monaco_zoom_row}
    \end{subfigure}
    \caption{Detailed \acrshort{dsm} over the city of Monaco.}
    \label{fig:monaco_zoom}
\end{figure}

\begin{figure}
    \begin{subfigure}[t]{0.3\linewidth}
        \flushleft
        \includegraphics[width=\linewidth]{Images/Chap_6/Langfjordjokelen_dsm_full_1850-2250_2250-2750.png}
        \caption{Photogrammetry \acrshort{dsm}. Orange square is detailed in \subref{fig:langfjordjokelen_dsm_zoom}}
        \label{fig:langfjordjokelen_dsm_full}
    \end{subfigure}\hfill
    \begin{subfigure}[t]{0.3\linewidth}
        \centering
        \includegraphics[width=\linewidth]{Images/Chap_6/Langfjordjokelen_dsm_zoom_1850-2250_2250-2750.png}
        \caption{Photogrammetry \acrshort{dsm} zoom. Orange line is detailed in \subref{fig:langfjordjokelen_zoom_row}}
        \label{fig:langfjordjokelen_dsm_zoom}
    \end{subfigure}\hfill
    \begin{subfigure}[t]{0.3\linewidth}
        \flushright
        \includegraphics[width=0.809\linewidth]{Images/Chap_6/Langfjordjokelen_error_zoom_1850-2250_2250-2750.png}
        \caption{\acrshort{dsm} with wrong intervals in red.}
        \label{fig:langfjordjokelen_error_zoom}
    \end{subfigure}\\
    \begin{subfigure}[t]{\linewidth}
        \centering
        \includegraphics[width=\linewidth]{Images/Chap_6/Langfjordjokelen_row_2000.png}
        \caption{\acrshort{dsm}, ground truth and elevation intervals along the orange line of \subref{fig:langfjordjokelen_dsm_zoom}}
        \label{fig:langfjordjokelen_zoom_row}
    \end{subfigure}
    \caption{Detailed \acrshort{dsm} over the Langfjordjøkelen glacier.}
    \label{fig:langfjordjokelen_zoom}
\end{figure}

The different metrics have been evaluated between the photogrammetry \acrshort{dsm} and the associated ground truth \acrshort{dsm}. Results are presented in \Cref{tab:elevation_metrics_global}, and are commented in the following sections. We also displayed for informational purposes the proportion of invalid pixels of each scene, in the last column of the table. A pixel is invalid if the ground truth \acrshort{dsm} or predicted \acrshort{dsm} does not contain any data for this pixel, or if it was masked by a water mask. Invalid pixels are therefore pixels that were not considered when computing the metrics.

\Cref{fig:Graasubreen_global,fig:Graasubreen_zoom,fig:toulouse_global,fig:toulouse_zoom,fig:grenoble_zoom,fig:monaco_zoom,fig:langfjordjokelen_zoom} allow to qualitatively evaluate the performance of the elevation intervals in different scenes. We can see that intervals seem to have a better accuracy performance in rural scenes such as glaciers than in urban scenes. For instance, we can clearly see wrong intervals of the Toulouse scene in \Cref{fig:toulouse_error}. They are harder to distinguish in the Graasubreen scene in \Cref{fig:Graasubreen_error,fig:Graasubreen_error_zoom}. Another notable observation is that elevation intervals of the Graasubreen glacier in \Cref{fig:Graasubreen_zoom_row} are very small and do not present large size fluctuations along the considered cross-section. This is not the case for the Toulouse scene (\Cref{fig:toulouse_zoom_row}), where intervals present large variations in size. They are relatively small from columns 7000 to 7100, but their size increase from column 7400 to 7500. We can see on the right of \Cref{fig:toulouse_zoom_row}, near column 7450, that their size allow to correctly contain the ground truth despite the difference between $\DSM$ and $\DSM_{true}$. However, there are some very large intervals, sometimes reaching $80$ m in size, near columns 7150, or 7300. Those intervals can occur because the range of considered disparities in the dense matching step is $[-50, 50]$, for every scene. Converted into altitude, elevations intervals have a maximal size of $[-50\cdot r_{alt}, 50\cdot r_{alt}]$, centered around the reference altitude (which can vary across the scene). This means that we can have a maximal elevation interval of size 243m in the case of Toulouse, where $r_{alt}=2.43$. The previously observed intervals of size 80m in \Cref{fig:toulouse_zoom_row} thus originates from a disparity interval of 33 pixels in size, \ie a third of the disparity intervals. Other examples of large intervals include the Monaco scene where $r_{alt}=0.99$ and for which the largest interval has a size of $99$ m, or the Paris scene where  $r_{alt}=5.78$ and the largest interval has a size of 578m. We will see in section \ref{sec:unexplored_leads} that we can modify our method to filter out points leading to those intervals. For now, we will evaluate the method as it stands. The following sections are dedicated to discussing each metric individually.

\subsection{Elevation Accuracy}
The elevation accuracy $Z_{acc}$, defined in \Cref{eq:relative_accuracy_elevation}, measures the proportions of intervals which contains the ground truth elevation, which we want to be as high as possible. The first observation is that the $90\%$ accuracy objective is validated on $8$ of the $11$ scenes we considered. In particular, the different glaciers as well as the Pic du Midi elevation intervals present a great accuracy, between $98.1\%$ and $99.7\%$. Scenes that do not verify the $90\%$ objective are the urban stereo pairs over Bordeaux, Montpellier, and Paris. In general, it seems that urban \acrshort{dsm}s have a lower accuracy than rural ones, which was to be expected as there are more steep variations of elevation in cities due to the presence of buildings. The Montpellier \acrshort{dsm} and Bordeaux \acrshort{dsm} are close to the $90\%$ objective, but intervals over Paris present an elevation accuracy of only $84.6\%$ . We will see in \Cref{sec:other_errors} that some other sources of errors can explain why $Z_{acc}<90\%$ on those scenes, such as vibration of the satellite during images acquisition or ground truth non-synchronicity issues. Provided that those errors are not present, we can claim that elevation intervals are accurate enough. They correctly estimate the error committed during the dense matching step, and the propagation of disparity intervals to elevation intervals is properly carried out.

\subsection{Residual Elevation Error}
The residual elevation error $Z_\varepsilon$, defined in \Cref{eq:residual_error_elevation}, estimates the gap between intervals and the ground truth, when intervals do not contain the ground truth. We want it to be as small as possible. $Z_\varepsilon$ is less than, or almost equal to, half a pixel for the majority of scenes, which is also around $0.5\%$ of the disparity range. This is really low, especially when compared to disparity intervals of \Cref{tab:metric_average} where $\varepsilon$ was typically around $3\%$ of the disparity range. Because $Z_\varepsilon$ is the median of error, it means that half of wrong intervals are less than $Z_\varepsilon$ pixels away from one of the bounds of the elevation intervals. Consequently, extending intervals by $Z_\varepsilon$ would divide by two the number of wrong intervals. For instance, on the Monaco scene, $Z_\varepsilon=0.57$pix, $r_{alt}=0.99$ m/pix and $Z_{acc}=90.3\%$. Defining the extended intervals as:
\begin{align*}
    [\underline{\DSM}-Z_\varepsilon\cdot r_{alt}, \overline{\DSM}+Z_\varepsilon\cdot r_{alt}] \approx [\underline{\DSM}-0.56, \overline{\DSM}+0.56]
\end{align*}
As the accuracy is $90.3\%$, it means that $9.7\%$ of pixels are not accurate. Using $[\underline{\DSM}-0.56, \overline{\DSM}+0.56]$ would lead to a new accuracy of $90.3+\dfrac{9.7}{2}=95.15\%$. For every scene, validating the $90\%$ accuracy, $Z_\varepsilon$ provides an easy extension for obtaining intervals with $95\%$ accuracy.

\Cref{fig:histogram_elevation_error} displays histograms of the distribution of errors, from which $Z_\varepsilon$ is the median. For both urban and rural scenes, errors are mostly of a few pixels. The shape of distributions on other scenes are similar to the Toulouse or Pic du Midi scenes. There is one $Z_\varepsilon$ value that stands out, for the Hellmem scene, where it equals $2.76$ pixels. \Cref{fig:error_hist_hellmem} displays the distribution of errors on the Hellmem scene, where the majority of the distribution is still located on the first peak of the distribution. There is however the presence of a tail on the distributions, that was not present on other scenes. We have to keep in mind that the accuracy for this scene is $98.8\%$, so that $Z_\varepsilon$ is only computed on $1.2\%$ of valid intervals. $Z_\varepsilon=2.76$ is therefore not necessarily a concerning observation, as it concerns very few pixels. It is more surprising that it equals $0.12$ for the Graasubreen \acrshort{dsm} which already has the best accuracy ($99.7\%$) out of all scenes. This translates the great accuracy performance of the elevation intervals estimation on this area.

\begin{figure}
    \centering
    \begin{subfigure}[t]{0.49\linewidth}
        \flushleft
        \includegraphics[width=\linewidth]{Images/Chap_6/histogram_elevation_eps_Toulouse.png}
        \caption{Toulouse}
        \label{fig:error_hist}
    \end{subfigure}\hfill
    \begin{subfigure}[t]{0.49\linewidth}
        \flushright
        \includegraphics[width=\linewidth]{Images/Chap_6/histogram_elevation_eps_Pic_du_midi.png}
        \caption{Pic du Midi}
        \label{fig:error_hist_pic_du_midi}
    \end{subfigure}\vspace*{0.3cm}\\
    \centering
    \begin{subfigure}[t]{0.49\linewidth}
    \centering
        \includegraphics[width=\linewidth]{Images/Chap_6/histogram_elevation_eps_Hellmem.png}
        \caption{Hellmem}
        \label{fig:error_hist_hellmem}
    \end{subfigure}
    \caption{Histograms of the residual error over the Toulouse, Pic du Midi and Hellmem scenes. $Z_{\varepsilon}$ is defined as the median of those distributions.}
    \label{fig:histogram_elevation_error}
\end{figure}

\subsection{Relative Elevation Size}
The relative elevation size $Z_{size}$, defined in \Cref{eq:relative_size_elevation}, estimates the size of intervals, both correct and incorrect. $Z_{size}$ is usually around 2 pixels wide. The Bordeaux, Grenoble and Monaco datasets are the only exceptions, where $Z_{size}$ is closer to 4 pixels in size. It does not seem to be correlated neither to the accuracy, nor to the altimetric ratio or acquisition angles of the satellites. In any case, 4 pixels is not ideal, but it is not excessively large either, especially as the considered disparity range is $[-50,50]$ in pixels. We can deduce from those values of $Z_{size}$ that the predicted \acrshort{dsm} is close to the true elevation. 

\Cref{fig:histogram_elevation_size} presents the distributions of relative sizes over the Toulouse and Hellmem scenes. It provides a more complete overview of the distribution of sizes than simply the indicator $Z_{size}$ which is defined as the median of those distributions. The distributions on other scenes have a similar shape.

\begin{figure}
    \centering
    \begin{subfigure}[t]{0.49\linewidth}
        \flushleft
        \includegraphics[width=\linewidth]{Images/Chap_6/histogram_elevation_s_rel_Toulouse.png}
        \caption{Toulouse}
        \label{fig:size_hist}
    \end{subfigure}\hfill
    \begin{subfigure}[t]{0.49\linewidth}
        \flushright
        \includegraphics[width=\linewidth]{Images/Chap_6/histogram_elevation_s_rel_Pic_du_midi.png}
        \caption{Pic Du Midi}
        \label{fig:size_hist_pic_du_midi}
    \end{subfigure}
    \caption{Histograms of the relative size over the Toulouse and Pic du Midi scenes. $Z_{size}$ is defined as the median of those distributions.}
    \label{fig:histogram_elevation_size}
\end{figure}

\subsection{Altimetric Ratio and Invalid Pixels}
Looking at \Cref{tab:elevation_metrics_global}, it seems that high values of the altimetric ratio $r_{alt}$ are correlated to ``low'' values of accuracy. This correlation does not imply a causal relationship. Indeed, scenes with low accuracy are urban scenes, on which the 3D reconstruction is harder due to the strong variations of elevation. In order to limit the number of occlusions in urban scenes, we often reduce the convergence angle of satellites. A low convergence angle results in a high altimetric ratio, as seen in \Cref{fig:incidence_angle}. This explains why we observe a high altimetric ratio for scenes with low accuracy.

The proportion of invalid pixels that are not considered in the metrics can sometimes be quite high, as it is the case for Monaco (41$\%$) or Hellmem (34.7$\%$). This is due to the presence of water in large part of the image (\Cref{fig:miniature_Monaco_rgb}) or simply the absence of data from the provided ground truth (\Cref{fig:miniature_Hellmem_gt})

\section{Comparison with ``Naive'' Intervals}
The fact that most intervals have a relative size $Z_{size}$ of around 2 pixels may raise the following question: would ``naive'' intervals be accurate ? We define naive intervals $[\underline{\DSM}_{naive}, ~\overline{\DSM}_{naive}]$, as intervals centered on the predicted \acrshort{dsm} with a size of 2 times the altimetric ratio $r_{alt}$:
\begin{align}
    [\underline{\DSM}_{naive}, ~\overline{\DSM}_{naive}] = [\DSM-r_{alt}, \DSM+r_{alt}]
\end{align}
\Cref{tab:naive_accuracy} allows comparing the accuracy of the naive intervals with that of our interval method. The table also contains the mean and median errors of the \acrshort{dsm} defined as:
\begin{align}
    \varepsilon_{\mean} = \mean|\DSM - \DSM_{true}|\qquad \varepsilon_{\median} = \median|\DSM - \DSM_{true}|
\end{align}

\begin{table}[ht]
    \centering
    \begin{tabular}{|c||c|c|c|c|}
        \hline
        Scene & $Z_{acc}$ (pix) & ``Naive'' $Z_{acc}$ (pix) & $\varepsilon_{\median}$ (m) & $\varepsilon_{\mean}$ (m)
        \\\hline\hline
        Bordeaux & 89.3$\%$ & 62.9$\%$ & 0.86 & 1.93\\\hline
        Graasubreen & 99.7$\%$ & 98.8$\%$ & 0.19 & 0.28\\\hline
        Hellmem & 98.8$\%$ & 96.4$\%$ & 0.22 & 0.75\\\hline
        Grenoble & 93.1$\%$ & 65.7$\%$ & 0.61 & 2.08\\\hline
        Langfjordjøkelen & 99.4$\%$ & 88.1$\%$ & 0.26 & 1.31\\\hline
        Monaco & 90.3$\%$ & 61.2$\%$ & 0.69 & 1.75\\\hline
        Montpellier & 89.1$\%$ & 79.4$\%$ & 1.81 & 2.52\\\hline
        Paris & 84.6$\%$ & 82.4$\%$ & 2.35 & 3.59\\\hline
        Peyto & 98.9$\%$ & 95.6$\%$ & 0.32 & 0.56\\\hline
        Pic du midi & 98.1$\%$ & 86.1$\%$ & 0.31 & 1.08\\\hline
        Toulouse & 92.0$\%$ & 83.2$\%$ & 0.92 & 1.79\\\hline
    \end{tabular}
    \caption{Comparison of ``naive'' intervals accuracy with our method for elevation intervals. The median and mean error of the predicted $\DSM$ are also indicated for reference.}
    \label{tab:naive_accuracy}
\end{table}
\Cref{tab:naive_accuracy} indicates that all ``naive'' intervals have a lower accuracy than intervals defined using our method. The difference is sometimes quite substantial, as for the Bordeaux, Grenoble, or Monaco scenes where the accuracy drops from around $30\%$ between the two methods. It also corresponds to scenes with a relative elevation size $Z_{size}$ which is higher than on the others scenes. This confirms that our method correctly adapts the size of elevation intervals to each scene individually. It also proves that our method can provide good accuracy performance even when the predicted $\DSM$ is far from the ground truth (indicated by high values of $\varepsilon_{\median}$ and $\varepsilon_{\mean}$). 

\section{Influence of Slope on the Metrics}
Using \Cref{eq:slope}, it is possible to compute the slope of the scene from the ground truth \acrshort{dsm}. It would be interesting to take a deeper look at the behavior of elevation metrics depending on the slope. To do so, we compute the slope $sl$ for each pixel and divided the ranges of slopes in different sections. Those sections are delimited by the following values, from \cite{hugonnet_uncertainty_2022}:
$$
[0\degree, ~2.5\degree, ~5\degree, ~10\degree, ~15\degree, ~20\degree, ~30\degree, ~40\degree, ~50\degree, ~70\degree, ~90\degree]
$$
We then evaluated the metrics on each slope range separately. Results are presented in \Cref{fig:slope_size,fig:slope_eps}. \Cref{fig:slope_eps_all,fig:slope_eps_no_hellmem} are the same, except we removed one scene, Hellmem, in the right figure. As mentioned earlier, this scene has very few pixels in error, and the median is thus not a robust indicator on this scene. Those figures illustrate the fact that elevation intervals have better accuracy and smaller size on flat slopes than on steep slopes. 

\begin{figure}
    \begin{subfigure}[t]{0.495\linewidth}
        \flushleft
        \includegraphics[width=\linewidth]{Images/Chap_6/slope_acc.png}
        \caption{Accuracy $Z_{acc}$}
        \label{fig:slope_acc}
    \end{subfigure}
    \begin{subfigure}[t]{0.495\linewidth}
        \flushright
        \includegraphics[width=\linewidth]{Images/Chap_6/slope_s_rel.png}
        \caption{Relative size $Z_{size}$}
        \label{fig:slope_size}
    \end{subfigure}
    \caption{Accuracy $Z_{acc}$ and relative size $Z_{size}$ values depending on the slope (in degree). Each curve corresponds to a different scene.}
    \label{fig:slope_acc_size}
\end{figure}

\begin{figure}
    \begin{subfigure}[t]{0.495\linewidth}
        \flushleft
        \includegraphics[width=\linewidth]{Images/Chap_6/slope_eps.png}
        \caption{Residual error $Z_\varepsilon$}
        \label{fig:slope_eps_all}
    \end{subfigure}
    \begin{subfigure}[t]{0.495\linewidth}
        \flushright
        \includegraphics[width=\linewidth]{Images/Chap_6/slope_eps_no_Hellmem.png}
        \caption{Residual error $Z_\varepsilon$ without Hellmem yellow curve}
        \label{fig:slope_eps_no_hellmem}
    \end{subfigure}
    \caption{Residual error $Z_{\varepsilon}$ values depending on the slope (in degree). Each curve corresponds to a different scene. In \subref{fig:slope_eps_no_hellmem}, we removed the Hellmem yellow curve from \subref{fig:slope_eps_all} for more visibility.}
    \label{fig:slope_eps}
\end{figure}

Regarding the accuracy of intervals, it tends to drop with steeper slopes. This can be explained due to the fact that the dense matching algorithm has better accuracy performances on smooth surfaces, provided that they are not textureless. This is in part due to the \acrshort{sgm} regularization which penalizes changes in disparity, and therefore in elevation. 

The size of intervals is greater on steep slopes. This is due to the fact that we purposefully extended the disparity intervals near disparity fluctuations, which correspond to steeper slopes. Because disparity changes indicate a variation in elevation, elevation intervals obtained from disparity intervals are consequently larger on steeper slopes. 

\begin{figure}
    \centering
    \includegraphics[width=0.7\linewidth]{Images/Chap_6/Slope_effect_error.png}
    \caption{Influence of slope and planimetric error on the accuracy of intervals. The planimetric error increases the residual error on steeper slope}
    \label{fig:Slope_effect_error}
\end{figure}

Finally, the residual error also increases with steeper slope when a planimetric error is present between the ground truth and the predicted \acrshort{dsm}. This can also be understood using \Cref{fig:Slope_effect_error}, as a steeper slope on the right interval leads to a greater residual error.

The next section will investigate sources of errors independent of our method, that can affect the evaluation of metrics.

\section{Other Sources of Error}\label{sec:other_errors}
We saw in \Cref{sec:results_elevation} that elevation confidence intervals present high accuracy and low relative size. There were however three scenes that did not meet the accuracy objective of $90\%$: Bordeaux, Montpellier and Paris. We will see in this section that there are different factors which can explain why the accuracy objective was not verified on those scenes. Those reasons are: 
\begin{itemize}
    \item Asynchronicity of the ground truth and satellite acquisitions
    \item Rasterization of \acrshort{lidar} data over vegetation
    \item Vibrations of the satellite during the images acquisition
\end{itemize}

\subsection{Asynchronicity of Sources}
The dates of acquisition of \acrshort{lidar} data and Pléiades images were presented in \Cref{tab:dates_pleiades_lidar_hd}. In some scenes, large periods of time separate the \acrshort{lidar} data from the satellite image, partly because we also tried to minimize seasonality changes. Specifically, the time between acquisitions for Bordeaux, Grenoble, Monaco, and Montpellier ranges between 5 and 13 months. This can lead to major changes of elevation in those scenes. A remarkable example of this can be found in the Monaco scene. If we take a look at the errors on this scene in \Cref{fig:Monaco_errors_global}, we can see that there is a strong concentration of errors in the top left corner of the image. 

\begin{figure}
    \begin{subfigure}[t]{0.495\linewidth}
        \flushleft
        \includegraphics[height=5cm]{Images/Chap_6/Monaco_errors.png}
        \caption{Error $\DSM-\DSM_{true}$}
        \label{fig:Monaco_errors}
    \end{subfigure}
    \begin{subfigure}[t]{0.495\linewidth}
        \flushright
        \includegraphics[height=5cm]{Images/Chap_6/Monaco_wrong_intervals.png}
        \caption{Position of wrong intervals, in red.}
        \label{fig:Monaco_wrong_intervals}
    \end{subfigure}
    \caption{Errors and position of wrong intervals on the Monaco scene}
    \label{fig:Monaco_errors_global}
\end{figure}

\Cref{fig:Carriere} is a zoom on this area. By looking at the differences between the ground truth \acrshort{dsm} in \Cref{fig:Carriere_gt} and the \acrshort{rgb} image in \Cref{fig:Carriere_RGB}, we can see that there are some differences between the ground truth \acrshort{dsm} and the satellite image. Most notably, the bottom left of the quarry levels are more constricted on the ground truth than on the \acrshort{rgb} image. We can confidently state that the quarry was excavated between the images' acquisition and the \acrshort{lidar} acquisition that occurred a year later. This explains why so many intervals are wrong in the area, in comparison with the rest of the scene. This can also be observed on a cross-section of the \acrshort{dsm}s presented in \Cref{fig:Carriere_row}. Those intervals are thus false negatives, and lower the global accuracy of the scene, as they approximately represent around $1.5\%$ of the valid data on this scene.

\begin{figure}
    \begin{subfigure}[t]{0.33\linewidth}
        \flushleft
        \includegraphics[width=\linewidth]{Images/Chap_6/Carriere_gt_Monaco.png}
        \caption{Ground truth \acrshort{dsm}. Orange line is detailed in \subref{fig:Carriere_row}}
        \label{fig:Carriere_gt}
    \end{subfigure}\hfill
    \begin{subfigure}[t]{0.33\linewidth}
        \centering
        \includegraphics[width=\linewidth]{Images/Chap_6/Carriere_RGB_Monaco.png}
        \caption{\acrshort{rgb} image}
        \label{fig:Carriere_RGB}
    \end{subfigure}\hfill
    \begin{subfigure}[t]{0.33\linewidth}
        \flushright
        \includegraphics[width=\linewidth]{Images/Chap_6/Carriere_wrong_intervals_Monaco.png}
        \caption{Position of wrong intervals, in red.}
        \label{fig:Carriere_wrong_intervals}
    \end{subfigure}\\
    \begin{subfigure}[t]{\linewidth}
        \flushright
        \includegraphics[width=\linewidth]{Images/Chap_6/Carriere_row_800.png}
        \caption{\acrshort{dsm}, ground truth and elevation intervals along the orange line of \subref{fig:Carriere_gt}.}
        \label{fig:Carriere_row}
    \end{subfigure}
    \caption{Detail over the La Turbie quarry, in the top left corner of the Monaco scene.}
    \label{fig:Carriere}
\end{figure}

This type of false negative can be found on a smaller scale in most urban scenes due to asynchronicity between satellite and \acrshort{lidar} acquisitions. We did not manually detect and correct every building that was destroyed or built in-between acquisitions, as it was too time-consuming. The accuracy results could therefore be further improved in urban scenes. It is possible, although very unlikely, that asynchronicity causes false positives during the accuracy computation.

\subsection{LiDAR Data Over Vegetation}\label{eq:lidar_vegetation}
Intervals computed over the Paris scene have an accuracy of $84.6\%$, which is the lowest of all scenes. When trying to understand why intervals did not seem to perform correctly, we noticed that many pixels in error were representing vegetation. This can be observed in \Cref{fig:paris_error_tree}, where many intervals are wrong near trees. By looking at the intervals in \Cref{fig:paris_error_tree,fig:Bordeaux_zoom}, we can see that the photogrammetry \acrshort{dsm} over-estimates the height of the canopy: columns 4150 to 4170 and 4345 to 4450 of the Paris scene in \Cref{fig:paris_error_tree_line}, and columns 3400 to 3800 of the Bordeaux scene in \Cref{fig:Bordeaux_row}. Other buildings, or the ground, do not present such error. One probable hypothesis is that because the \acrshort{lidar} HD over Paris was acquired in early March, and the satellite images in the last day of May, tree foliage had major differences. Firstly, trees seem bigger on the \acrshort{rgb} image than on the \acrshort{lidar} \acrshort{dsm}, which is probably due to the increase of foliage in-between acquisitions. Also, the \acrshort{lidar} probably acquired points both on the top of trees and on the ground, through the sparse foliage. When rasterizing the \acrshort{lidar} point cloud using a Gaussian mean, we averaged both points on top of the trees and on the ground, resulting in a mix between the two elevations. Due to its resolution, the photogrammetry \acrshort{dsm} does not create multiple points at both the top and the base of the same tree. It therefore only predicts the elevation of the tree canopy, resulting in the observed error when comparing with the \acrshort{lidar} \acrshort{dsm}. We tried to mask the vegetation pixels, using a mask computed from the Normalized Difference Vegetation Index \cite{gao_ndwinormalized_1996}, but the quality of the masks were not sufficiently accurate, and tend to mask many pixels that were not vegetation. We therefore decided not to use them. Also, this effect was particularly present on the Paris scene, which is probably because the \acrshort{lidar} was acquired in early March, while the \acrshort{lidar} on other scenes was acquired from June to September, when the vegetation is denser. However, we will see in the next section that there is a greater source of error occurring on that scene, which can also be the source of errors leading to a low accuracy.

\begin{figure}
    \centering
    \begin{subfigure}[t]{\linewidth}
        \flushright
        \includegraphics[width=0.95\linewidth]{Images/Chap_6/paris_error_tree_gt.png}
        \caption{Ground truth \acrshort{dsm}. Orange line is detailed in \subref{fig:paris_error_tree_line}}
        \label{fig:paris_error_tree_gt}
    \end{subfigure}\\
    \begin{subfigure}[t]{\linewidth}
        \flushright
        \includegraphics[width=0.95\linewidth]{Images/Chap_6/paris_error_tree_clr.png}
        \caption{\acrshort{rgb} image}
        \label{fig:paris_error_tree_clr}
    \end{subfigure}\\
    \begin{subfigure}[t]{\linewidth}
        \flushright
        \includegraphics[width=0.95\linewidth]{Images/Chap_6/paris_error_tree_intervals.png}
        \caption{Position of wrong intervals, in red.}
        \label{fig:paris_error_tree_intervals}
    \end{subfigure}\\
    \begin{subfigure}[t]{\linewidth}
        \centering
        \includegraphics[width=\linewidth]{Images/Chap_6/paris_error_tree.png}
        \caption{\acrshort{dsm}, ground truth and elevation intervals along the orange line of \subref{fig:paris_error_tree_gt}}
        \label{fig:paris_error_tree_line}
    \end{subfigure}
    \caption{Zoom near Saint-Augustin, in the Paris scene}
    \label{fig:paris_error_tree}
\end{figure}

\begin{figure}
    \begin{subfigure}[t]{0.33\linewidth}
        \flushleft
        \includegraphics[width=\linewidth]{Images/Chap_6/Bordeaux_gt_zoom.png}
        \caption{Ground truth \acrshort{dsm}. Orange line is detailed in \subref{fig:Bordeaux_row}}
        \label{fig:Bordeaux_gt}
    \end{subfigure}\hfill
    \begin{subfigure}[t]{0.33\linewidth}
        \centering
        \includegraphics[width=\linewidth]{Images/Chap_6/Bordeaux_clr_zoom.png}
        \caption{\acrshort{rgb} image}
        \label{fig:Bordeaux_RGB}
    \end{subfigure}\hfill
    \begin{subfigure}[t]{0.33\linewidth}
        \flushright
        \includegraphics[width=\linewidth]{Images/Chap_6/Bordeaux_error_zoom.png}
        \caption{Position of wrong intervals, in red.}
        \label{fig:Bordeaux_wrong_intervals}
    \end{subfigure}\\
    \begin{subfigure}[t]{\linewidth}
        \flushright
        \includegraphics[width=\linewidth]{Images/Chap_6/Bordeaux_row_600.png}
        \caption{\acrshort{dsm}, ground truth and elevation intervals along the orange line of \subref{fig:Bordeaux_gt}.}
        \label{fig:Bordeaux_row}
    \end{subfigure}
    \caption{Detail over the Quinconces square, in the Monaco scene.}
    \label{fig:Bordeaux_zoom}
\end{figure}


\subsection{Vibration of the Satellite}\label{sec:vibrations}
The previous sources of errors concerned the quality of the ground truth \acrshort{dsm}. However, there is one source of errors that we did not take into consideration in our methods: the satellite vibration during acquisition. During the acquisition of images, push-broom captors may experience vibrations which are not modeled in the \acrshort{rpc} model. The vibrations occur on the pitch angle, whose axis is perpendicular to the direction of flight of the satellite, as depicted in \Cref{fig:Pleiade_angle_vibration}. This translates into low frequency along-track biases on the \acrshort{dsm}, sometimes called undulations \cite{hugonnet_uncertainty_2022}. This can be detected by looking at the difference between the photogrammetry and the ground truth \acrshort{dsm}s:
\begin{align}
    \DSM(\rowcol)-\DSM_{true}(\rowcol)
\end{align}
Without biases, this difference (or residual error), should not be correlated with the position $\rowcol$ in the image. \Cref{fig:vibrations} presents all scenes for which we detected a bias. The bias is small on \Cref{fig:vibrations_Hellmem} but is hard to miss on \Cref{fig:vibrations_Montpellier,fig:vibrations_Paris,fig:vibrations_Toulouse}. \Cref{fig:vibration_bias} indicates that those biases can reach 2.5m of amplitude.

\begin{figure}
    \centering
    \includegraphics[width=0.7\linewidth]{Images/Chap_6/Pleiade_angle_vibration.png}
    \caption{Angles of a satellite. Vibrations of the satellite occur on the pitch axis.}
    \label{fig:Pleiade_angle_vibration}
\end{figure}

\begin{figure}
    \begin{subfigure}[t]{0.48\linewidth}
        \flushleft
        \includegraphics[width=\linewidth]{Images/Chap_6/vibrations_Hellmem.png}
        \caption{Hellmem}
        \label{fig:vibrations_Hellmem}
    \end{subfigure}\hfill
    \begin{subfigure}[t]{0.48\linewidth}
        \flushright
        \includegraphics[width=\linewidth]{Images/Chap_6/vibrations_Montpellier.png}
        \caption{Montpellier}
        \label{fig:vibrations_Montpellier}
    \end{subfigure}\\
    \begin{subfigure}[t]{0.48\linewidth}
        \flushleft
        \includegraphics[width=\linewidth]{Images/Chap_6/vibrations_Paris.png}
        \caption{Paris}
        \label{fig:vibrations_Paris}
    \end{subfigure}\hfill
    \begin{subfigure}[t]{0.48\linewidth}
        \flushright
        \includegraphics[width=\linewidth]{Images/Chap_6/vibrations_Toulouse.png}
        \caption{Toulouse}
        \label{fig:vibrations_Toulouse}
    \end{subfigure}
    \caption{Difference $(\DSM-\DSM_{true})$ for different scenes. A bias along rows (along-track) appears, suggesting vibrations during the image acquisition. Biases on the Hellmem \acrshort{dsm} are smaller that for the other three scenes.}
    \label{fig:vibrations}
\end{figure}

When computing elevation confidence intervals, we did not consider potential errors of this magnitude for the \acrshort{rpc} models. Our method principally aimed to model and propagate the errors stemming from the dense matching step, which we considered to be the part of the pipeline where the largest errors could occur. Biases observed in \Cref{fig:vibrations} indicate that only considering the dense matching errors is not sufficient when producing \acrshort{dsm}s with Pléiades images. The CO3D satellites will not use a push-broom sensor, but rather a CCD Bayer matrix, which can potentially reduce vibration issues. However, the satellite will still need to stabilize itself during acquisition, therefore some vibrations can still be expected. 

In the scope of this thesis, we did not try to model the uncertainty associated with potential vibrations of the satellite. It would however be interesting to determine if Montpellier and Paris scenes could validate the $90\%$ accuracy objective without the effect of vibrations. We therefore propose a simple correction of the vibration effect, which is not meant to perfectly solve the issue, but rather to sufficiently reduce it for a better evaluation of the elevation interval metrics. In order to do so, we computed the residual error $(\DSM-\DSM_{true})$, and attempted to model the observed bias using a least square approach. We assume that the bias depends only on the row of the image and propose to simply model it by a cosine function carried by a linear function of the row:
\begin{align}\label{eq:bias_model}
    bias(row)=A_1\cos(A_2\cdot row+A_3) + A_4\cdot row+A_5
\end{align}
where $(A_1\enum A_5)$ are the parameters of our model. \Cref{fig:vibration_bias} shows scatter plots of the residual error $(\DSM-\DSM_{true})$ where the x-axis indicates image rows. We also computed the median of the residual error for each row as an indicator of the distribution of errors, and the optimized model from \Cref{eq:bias_model}.

\begin{figure}
    \begin{subfigure}[t]{\linewidth}
        \centering
        \includegraphics[width=\linewidth]{Images/Chap_6/vibration_bias_Montpellier.png}
        \caption{Montpellier scene}
        \label{fig:vibration_bias_Montpellier}
    \end{subfigure}
    \begin{subfigure}[t]{\linewidth}
        \centering
        \includegraphics[width=\linewidth]{Images/Chap_6/vibration_bias_Paris.png}
        \caption{Paris scene}
        \label{fig:vibration_bias_Paris}
    \end{subfigure}
    \caption{Scatter plot where gray points are the residual error $(\DSM-\DSM_{true})$ of a pixel. The x-axis represents the row of each pixel. As an indicator, we also computed the median residual error for each row, appearing in orange. The blue line is the optimized model from \Cref{eq:bias_model}.}
    \label{fig:vibration_bias}
\end{figure}

\begin{figure}
    \begin{subfigure}[t]{0.48\linewidth}
        \flushleft
        \includegraphics[width=\linewidth]{Images/Chap_6/vibrations_Montpellier.png}
        \caption{Residual error with bias}
        \label{fig:vibrations_Montpellier_with_bias}
    \end{subfigure}\hfill
    \begin{subfigure}[t]{0.48\linewidth}
        \flushright
        \includegraphics[width=\linewidth]{Images/Chap_6/vibration_corrected_Montpellier.png}
        \caption{Residual error without bias}
        \label{fig:vibrations_Montpellier_without_bias}
    \end{subfigure}
    \caption{Residual error $(\DSM-\DSM_{true})$ with and without the bias estimated in \Cref{fig:vibration_bias} for the Montpellier scene.}
    \label{fig:removed_bias_Montpellier}
\end{figure}

\begin{figure}
    \begin{subfigure}[t]{0.48\linewidth}
        \flushleft
        \includegraphics[width=\linewidth]{Images/Chap_6/vibrations_Paris.png}
        \caption{Residual error with bias}
        \label{fig:vibrations_Paris_with_bias}
    \end{subfigure}\hfill
    \begin{subfigure}[t]{0.48\linewidth}
        \flushright
        \includegraphics[width=\linewidth]{Images/Chap_6/vibration_corrected_Paris.png}
        \caption{Residual error without bias}
        \label{fig:vibrations_Paris_without_bias}
    \end{subfigure}
    \caption{Residual error $(\DSM-\DSM_{true})$ with and without the bias estimated in \Cref{fig:vibration_bias}, for the Paris scene.}
    \label{fig:removed_bias_Paris}
\end{figure}


Having estimated the bias, we can subtract it from the \acrshort{dsm} and its confidence intervals. This will not change the relative size of the intervals, as we apply the same elevation shift to both bounds. The accuracy and residual error are however impacted by this bias rectification. Here are the improvements:
\begin{itemize}
    \item For the Montpellier scene, the accuracy $Z_{acc}$ increases $89.1\%$ to $92.6\%$. The residual error $Z_\varepsilon$ also increases from $0.28$ pixels to $0.32$ pixels.
    \item For the Paris scene, the accuracy $Z_{acc}$ increases from $84.6\%$ to $88.1\%$. The residual error $Z_\varepsilon$ changes from $0.46$ pixels to $0.43$ pixels.  
\end{itemize}
The fact that the residual error slightly varies is due to the fact that as the accuracy increases, the set on which $Z_\varepsilon$ is computed changes, thus changing its median.

\begin{remark}
    In this section, we used the ground truth to correct the bias from the vibrations of the satellite. We used the ground truth because our objective was to verify if the biases were indeed the source of the missing accuracy of elevation intervals, so that we could safely assume our method was efficient, granted that no vibrations occurred during the acquisition. 
    
    Our aim was not to propose a general method for modeling or correcting the errors due to vibrations, in which case we could not have used the ground truth, as it is not available everywhere. However, if such were our intents, using a low resolution \acrshort{dsm} such as the Copernicus \acrshort{dsm} at 30m resolution could suffice to detect the presence of vibrations. This \acrshort{dsm} was already used in the \acrshort{cars} pipeline as our initial elevation. By plotting the differences between the low resolution \acrshort{dsm} and the predicted photogrammetry \acrshort{dsm}, we can indeed observe the presence of vibrations, as seen in the following figure:
    
    {\centering\includegraphics[width=7cm]{Images/Chap_6/copernicus_bias_Montpellier.png}
    \captionof{figure}{Difference $\DSM-\DSM_{low}$ between the photogrammetry model $\DSM$ and the low resolution model $\DSM_{low}$ over the Montpellier scene.}}
\end{remark}


The Paris scene still does not reach the $90\%$ accuracy objective, while being very close to it. Looking at \Cref{fig:vibrations_Paris_without_bias}, there still seems to be a bias, horizontally this time, as residuals are positive on the left side of the image and negative on the right side of the image. We can also distinguish that streets seem to have more positive residuals, which is once again probably due to the vegetation issue discussed in \Cref{eq:lidar_vegetation}. As the accuracy is very close to the $90\%$ objective, and we have good reasons to think there are false negative in our test, we consider that this scene confirms the aimed accuracy performance of our method for creating elevation confidence intervals. 

\section{Unexplored Leads}\label{sec:unexplored_leads}
In the previous sections, we evaluated the accuracy and size performance of elevation intervals. Those intervals were obtained by triangulation of disparity intervals, yielding an interval along a line of sight, which was then rasterized. As they stand, intervals are a reflection of the potential error committed during the dense matching step. During our investigations, we came up with other interesting ideas that we did not have time to get to the bottom of. In order to conclude this chapter, we will quickly go over them.

We saw in \Cref{fig:toulouse_zoom_row,fig:Carriere_row,fig:Bordeaux_row} that elevation intervals could sometimes reach very high values, like $80$ m in the case of Toulouse. The presence of those intervals reflect the fact that bounds of corresponding disparity intervals have a high possibility degree, and that the difference between choosing one bound or the other is not obvious. It could be wise to filter out the 3D points leading to such large intervals. After all, a large elevation interval suggests that we are not certain about the position of the point, so removing it seems natural. Additionally, we ignored the planimetric error $\Delta XY$ during the rasterization as presented in \Cref{fig:planimetric_shift}, but the hypothesis that $\Delta XY$ is small is probably not valid for large intervals. For acquisitions with incidence angles, around 11$\degree$, an elevation interval of 50m would mean that the planimetric uncertainty is around $50\tan(11\degree)\approx20$ m. For the Peyto scene, where the incidence is $21.3\degree$, this uncertainty reaches $46$ m. Not only are we unsure about the elevation of the point, but there is much uncertainty regarding which cell it falls into during the rasterization process. We can either filter those points directly after the dense matching step, based on the size of the intervals, or after they have been triangulated, based on the resulting planimetric/altimetric shift. The latter takes into account the acquisition angles of images, and is not limited to the information contained in the disparity map. 

When triangulating the disparity, we consider line of sights to be 3D lines. However, they represent the position of a pixel, which is not a point but a surface (determined by the resolution of images). Lines of sights are therefore not exactly lines, but rather a cone or a cylinder.  We could even go further with this model, by taking into consideration the accuracy of lines of sights into the definition of those cylinders. Intersecting ``cylinders'' of sights would then yield a 3D volume. We would then have to processed differently the obtained volumes during the rasterization step. 

The rasterization step uses a Gaussian mean based on the distance of 3D points to the center of a cell. We could additionally use the size of intervals to adapt the weights in the rasterization steps. Doing so, we would favor points with small intervals, \ie for which we are confident about their position, to improve the final \acrshort{dsm}. If we reason with 3D volumes, we could compute a density on each volume and project those densities on the final raster. 

The uncertainty information contained in the different intervals (disparity or elevation) could also be used to improve the final \acrshort{dsm}. As it stands, the computation of intervals does not modify the final \acrshort{dsm}.

\begin{conclusion}
    In this final chapter, we propagated the uncertainty associated with the disparity from \Cref{chap:epistemic_uncertainty} into elevation intervals at the end of the photogrammetry pipeline. Elevation intervals are a result of the uncertainty due to the dense matching step, which is the most complex part of the pipeline, and therefore the most important source of uncertainty. Intervals have an accuracy of $90\%$, and their size is proportionate to the potential errors between the predicted \acrshort{dsm} and the true elevation. Intervals have a better performance in terms of accuracy and size in natural landscape than in urban ones, as those scenes contain more elevation differences and steeper slopes. The methodology developed for computing elevation confidence intervals aims to satisfy requirements expressed by \acrshort{dsm} users. They can also provide a solution for the future CO3D mission requirements, as a performance map will need to be provided alongside the \acrshort{dsm}.
\end{conclusion}

\clearpage
    
    \pagestyle{conclusion}
    \setkeys{Gin}{draft=false}
    \addcontentsline{toc}{chapter}{Conclusion}
\chapter*{Conclusion}
\section*{Context}
Usage of elevation data for environmental applications, urban planning, risk assessment, \etc requires the massive production of \acrfull{dsm}s. \acrshort{dsm}s at a global scale can be produced using satellites orbiting the Earth, and using different techniques such as \acrshort{lidar} measures, radar interferometry or optical photogrammetry. In particular, optical sensors, which have become relatively low-cost, allow producing dense \acrshort{dsm}s with sub-meter resolution. In this context, \acrshort{cnes} and Airbus are preparing the launch of the \acrshort{co3d} mission, consisting in two pairs of low-cost optical satellites dedicated to the production of high-resolution \acrshort{dsm}s across the globe using stereophotogrammetry.

Photogrammetry is a complex task, consisting in multiple successive algorithms to process images and extract the elevation information contained through the parallax effect. It can be divided into the following main steps:
\begin{itemize}
    \item Resampling of stereo images into a convenient geometry.
    \item Dense matching of pixels between images.
    \item Triangulation of a 3D point for each match.
    \item Rasterization of the point cloud on a regular grid.
\end{itemize}
Along these steps, many uncertainties arise, which can lead to errors of varying magnitude. The objective of this thesis was to quantify and propagate the uncertainty alongside a photogrammetry pipeline, in preparation of the \acrshort{co3d} mission. In particular, we focused on the \acrshort{cars} pipeline which will be used to process the \acrshort{co3d} data. Furthermore, one of the requirements of the \acrshort{co3d} mission is to produce a performance map alongside \acrshort{dsm}s. Many \acrshort{dsm} users also seek to know the quality of the predicted \acrshort{dsm}, usually characterized by confidence intervals. The main contribution of this thesis is precisely the development of a method for computing confidence intervals, which can also be used as a performance map for the \acrshort{co3d} mission.

\section*{Contributions}
In this thesis, we used other uncertainty models than well-known probability distributions. We considered \textit{imprecise} probabilities and more specifically possibility distributions, more adapted to model epistemic uncertainty, \ie arising from a lack of knowledge, by opposition to the uncertainty due to a purely random process. Those models define credal sets, \ie convex sets of probability distributions. 

When propagating multiple sources of uncertainty, it is required to compute a multivariate model of uncertainty accounting for the different dependencies between uncertain sources. We proposed to use dependency models known as copulas to construct multivariate credal sets. We introduced three methods for aggregating marginal credal sets into multivariate credal sets using copulas, namely $\M_{robust}$, $\M_{mass}$ and $\M_{agg}$. Those models are not equivalent, and their computation presents varying degrees of complexity. We investigated the relations between those sets depending on the copula used to join them, as well as the specific models used to define marginal credal sets, such as the aforementioned possibility distributions.

We used the previous results concerning multivariate credal sets to propagate the uncertainty in a specific part of the photogrammetry pipeline: the dense matching step. More specifically, we consider the computation of the matching cost between every pixel of the stereo images. We showed that the models could correctly estimate the propagated uncertainty regarding the matching cost on a real pair of stereo images.

The cost volume is used to compute the disparity map, encoding the pairing of pixels to be triangulated. Computing a cost volume allowing to correctly estimate the disparity is the hardest part of the stereo pipeline. Correctly estimating the uncertainty on the predicted disparity is therefore crucial for the rest of the pipeline. However, only considering the propagated uncertainty on the cost volume, as we previously did, is not sufficient for the correct uncertainty estimation of the disparity map. We therefore proposed to use possibility distributions to model the epistemic uncertainty regarding the choice of each disparity. For each considered pixel, we used those possibility distributions to determine a disparity confidence interval. We evaluated the accuracy of the intervals using real stereo images, and observed that intervals contain the true disparity at least $90\%$ of the time. This method for creating intervals can be applied to a wide range of stereo algorithms and is not restricted to satellite photogrammetry. To the best of our knowledge, it is also the first time such disparity confidence intervals are computed. 

We then propagated those disparity confidence intervals all the way to the end of the pipeline, where they take the form of elevation confidence intervals associated with the predicted \acrshort{dsm}. We evaluated the performance of elevation intervals on real satellite images, for which we possess a reference high quality \acrshort{dsm}. Intervals are once again accurate $90\%$ of the time, validating the performances of our new method. We also implemented this method for estimating disparity and elevation confidence intervals into the \acrshort{cars} (\url{https://github.com/CNES/cars}) and Pandora (\url{https://github.com/CNES/Pandora}) software, developed by \acrshort{cnes}. They are already publicly available, and can be used to produce elevation confidence intervals for \acrshort{dsm}s.

\section*{Limitations and Perspectives}
We demonstrated that possibility distributions and copulas could be used to propagate uncertainty in a problem such as the evaluation of a cost volume. Implementing this propagation remains a difficult challenge. It requires using simple cost functions and other simplifying assumptions. It also requires large processing capacities to be carried out efficiently, as we joined the uncertainty of thousands of different pixels in our experiments.

We evaluated our method for producing elevation intervals using high resolution \acrshort{dsm}s obtained from \acrshort{lidar} data. However, we only had access to \acrshort{dsm}s provided by glaciologists, or high resolution point clouds from the \acrshort{lidar} HD program. This means that we only observed landscapes which either contained glaciers, or were located in France. Extending the evaluation of intervals to a broader diversity of locations would be valuable, such as deserts or American cities.

Another limitation of our method is that it does not take into consideration the potential errors occurring \textit{before} the dense matching step. Those errors could occur when defining the epipolar geometry for instance, or on the localization model itself, for instance caused by vibration of the satellite as seen at the end of \Cref{chap:elevation_intervals}. Combining uncertainty models of the sensor itself or on its geolocation model could lead to a better estimation of the overall uncertainty of the \acrshort{dsm}. 

Different future perspective regarding our work can be considered. The method we developed for computing confidence intervals is carried out alongside the processing of the main \acrshort{dsm}, without influencing the values it contains. However, we could imagine using the information contained in the confidence intervals to facilitate or improve the different disparity or elevation predictions. Here are a few interesting leads that could be explored:
\begin{itemize}
    \item Disparity intervals could be computed a first time before any \acrshort{sgm} regularization. Then, during the \acrshort{sgm} regularization, disparities that are not contained inside disparity intervals are ignored, which could greatly reduce the amount of computation necessary. 
    \item Another approach would be to use matching cost possibilities to apply a different strategy when computing the disparity map. Currently, disparities are determined using a \textit{winner-takes-all} strategy, meaning that for each pixel, the selected disparity is the one minimizing its cost curve. We then do the same for the cost volume of the right image, and remove matches that do not verify the cross-checking test of \Cref{eq:cross-checking}. However, we could see the choice of a disparity map as a stable marriage problem \cite{irving_matching_1998}, where possibilities are interpreted as degrees of preferences. This could improve the quality of the disparity map, as the choice of each disparity would consider more information than a single cost curve. It would also provide an alternative to the \textit{winner takes all} strategy, which accounts for the vast majority of strategies used in dense matching.  
    \item We currently extend intervals in low confidence areas using quantiles computed over a set of neighboring intervals. We could consider using instead other cost functions which present better performances near depth discontinuities, or using values of the cost volume before \acrshort{sgm} regularization to better process intervals in those low confidence areas.
    \item When triangulating disparity intervals, we could take into consideration the limited resolution of line of sights. We could for instance reason with uncertain 3D volumes when computing their intersection, which would be a more realistic model than the precise lines of sight currently considered.
    \item Before the rasterization step, we could discard points for which the interval size is too large, as the potential error committed by considering this point is too high. 
    \item During rasterization, we could use information from intervals to modify the weights of each point in the final value of the \acrshort{dsm}. Points with small elevation intervals would be granted more importance in the final product than those with large intervals. 
\end{itemize}

Ideas developed in this thesis for 1D matching could be extended to 2D matching, for instance in the Pandora2D tool (\url{https://github.com/CNES/Pandora2D}). Apart from photogrammetry, ideas developed in this thesis could be used to improve confidence criteria of the alignment of image bands for multi-spectral images, for instance in the TRISHNA mission \cite{lagouarde_indo-french_2019}, or alignment of multi-temporal images, such as Sentinel-2 \cite{yan_sentinel-2a_2018}.
Taking a step back from imagery, we showed in this thesis that using other models of uncertainty than well-known probabilities can lead to new methods for estimating and characterizing potential errors. Many methods using incomplete or imperfect data for diverse applications, such as clustering, classification or decision-making, can also benefit from using imprecise probabilities.

\clearpage
    
    \pagestyle{annex}
    \setkeys{Gin}{draft=false}
    % Redefining the chapter Heading
\makeatletter
\def\@makechapterhead#1{%
  \vspace*{50\p@}%
  {\parindent \z@ \raggedright \normalfont
        %\huge\bfseries \@chapapp\space \thechapter
        \Huge\bfseries \thechapter.\space Annex%
        %\par\nobreak
        %\vskip 20\p@
        \vskip 40\p@
  }}
\makeatother

\chapter{Annex}\label{chap:annex}
This annex contains results and experiments that we consider secondary in relation to the main content. \Cref{sec:dconvexity,sec:annex_dconvex_pbox} presents additional results concerning copulas presented in \Cref{sec:copulas} and used in \Cref{subsec:pboxes,subsec:pboxes}. \Cref{sec:median_filtering} presents results on the median filter used in \Cref{sec:coherence_disparity_intervals}. \Cref{sec:ablation_study} contains ablation studies of the method presented in \Cref{chap:epistemic_uncertainty}. 


\section{Directional Convexity/Concavity for Copulas}\label{sec:dconvexity}
This section will investigate a property shared by some copulas called directional convexity/concavity. This is a theoretical contribution, as we do not exploit them in the applications to stereophotogrammetry in \Cref{chap:propagating,chap:epistemic_uncertainty}. However, we will see that those properties are shared by many common families of copulas. The main result of this section is \cref{eq:convex_diff_hvol}, which was used to prove a specific relationship between multivariate uncertainty models in \Cref{subsec:pboxes,subsec:multiple_models}.

\begin{definition}[D-convexity, D-concavity]\label{def:convex}
    A copula $C$ is called directionally convex (D-convex) \cite{alvoni_dierent_2007} if for every $(u_1 \enum u_n)\in[0,1]^n$, $(v_1\enum v_n)\in[0,1]^n$, $i\in\opi1,n\cli$ and $t\in[0,1]$ it verifies:
    \begin{align}
        C(u_1\enum tu_i+(1-t)v_i\enum u_n) ~\leqslant~&t\cdot C(u_1\enum u_i\enum u_n)\nonumber\\
        &+ (1-t)\cdot C(u_1\enum v_i\enum u_n)\label{eq:convex_copula}
    \end{align}
    In other words, the copula is convex when fixing all but one of its variables. A copula is called directionally concave (D-concave) if the inequality is reversed.
\end{definition}

D-convexity/D-concavity is quite common in known families of copulas. The following paragraphs detail this property for copulas presented in \Cref{tab:family_of_copula}, in the case $n=2$.
As the copulas presented are symmetric regarding their variables, D-convexity/D-concavity is only detailed for $u_1$. We assume that \( \theta \) always belong to the domain of definition detailed in \Cref{tab:family_of_copula}, and that \(u_1, u_2\) are in \([0, 1]\). Regarding the Clayton and Gumbel families, the copula is defined by continuous extension in cases $u_1=0$ and $u_2=0$. Finally, if the restriction of a copula $C$ to one of its variable $u_i$ is twice differentiable, then proving its D-convexity for this variable can be done by proving that $\dfrac{\partial^2 C}{\partial {u_i}^2}\geqslant 0$.
\begin{description}
    \item[Ali-Mikhail-Haq copula] This copula is twice differentiable, and its second order partial derivative is
    $$\frac{\partial^2 C}{\partial {u_1}^2}=u_2(1-\theta(1-u_2))\frac{-2\theta(1-u_2)}{(1-\theta(1-u_1)(1-u_2))^3}$$
    Thus the Ali-Mikhail-Haq copula is D-convex for $\theta\in[-1,0]$ and D-concave for $\theta\in[0,1)$.
    \item[Clayton copula] This copula is not always differentiable on all of its domain, depending on which value is retained by the maximum function. It is however continuous as it is the maximum of two continuous function. For convenience, we work with $(u_1, u_2)$ in $\mathbb{I}^2$, where $\mathbb{I}$ is the open unit interval (the closed unit interval is then covered by continuity). Let $u_2\in\mathbb{I}$. We split the possible range $\mathbb{I}$ of $u_1$ in two:
\begin{itemize}
    \item the first domain $\mathcal{D}_1^{\theta,u_2}$ is where $u_1^{-\theta}+u_2^{-\theta}-1\leqslant0$, and thus $C(u_1, u_2)=0$. Here, $\frac{\partial^2 C}{\partial {u_1}^2}=0$ and the copula is both D-convex and D-concave.
    \item the second domain $\mathcal{D}_2^{\theta,u_2}$ is where $u_1^{-\theta}+u_2^{-\theta}-1>0$ and thus $C(u_1, u_2)\geqslant0$. Here it holds that
    $$\frac{\partial^2 C}{\partial {u_1}^2}=(1+\theta)(1-u_2^{-\theta})u_1^{-2-\theta}(u_1^{-\theta}+u_2^{-\theta}-1)^{-2-\frac{1}{\theta}}$$
    Because of the definition of $\mathcal{D}_2^{\theta,u_2}$, the sign of $\frac{\partial^2 C}{\partial {u_1}^2}$ on $\mathcal{D}_2^{\theta,u_2}$ is that of $1-u_2^{-\theta}$.
\end{itemize}

If $\theta>0$, then $D^\theta_2=\mathbb{I}$ and $\frac{\partial^2 C}{\partial {u_1}^2}\leqslant0$ which means that the copula is D-concave on all of its domain.

The case where $\theta<0$ is less straightforward. In that case, $\frac{\partial^2 C}{\partial {u_1}^2}\geqslant0$ on $\mathcal{D}_2^{\theta,u_2}$. The restrictions of the copula to $\mathcal{D}_1^{\theta,u_2}$ and $\mathcal{D}_2^{\theta,u_2}$ are both D-convex, but we need to prove that it is still true on their union. Let $u_1\in\mathcal{D}_1^{\theta,u_2}$, $v_1\in\mathcal{D}_2^{\theta,u_2}$ and $t\in[0,1]$. We note $w_1=(1-u_2^{-\theta})^{-\frac{1}{\theta}}$, such that $\mathcal{D}_1^{\theta,u_2}=]0,w_1]$ and $\mathcal{D}_2^{\theta,u_2}=]w_1, 1[$. By continuity, $C$ is D-convex on $\mathcal{D}_2^{\theta,u_2}\bigcup\{w_1\}$. Because $u_1,w_1\in \mathcal{D}_1^{\theta,u_2}$, it holds that:
    \begin{eqnarray*}
        tC(u_1,~u_2)+(1-t)C(v_1,~u_2) &=& tC(w_1,~u_2)+(1-t)C(v_1,~u_2)\\
        &\geqslant& C(tw_1+(1-t)v_1,~u_2)\\
        &&\text{by convexity of $C$ on $\mathcal{D}_2^{\theta,u_2}\bigcup\{w_1\}$}\\
        &\geqslant& C(tu_1+(1-t)v_1,~u_2)\\
        && \text{because $C$ is component-wise increasing}
    \end{eqnarray*}
which, by definition \ref{def:convex}, proves that $C$ is D-convex on $\mathcal{D}_1^{\theta,u_2}\bigcup \mathcal{D}_2^{\theta,u_2}$. By continuity, $C$ is D-convex on all of its domain.
\item[Frank copula] This copula is twice differentiable, and its second order partial derivative is
$$\frac{\partial^2 C}{\partial {u_1}^2}=\frac{(e^{-\theta u_2}-1)e^{-\theta u_1}\theta(e^{-\theta u_2}-e^{-\theta} )}{(e^{-\theta}-1+(e^{-\theta u_1}-1)(e^{-\theta u_2}-1))^2}$$
If $\theta\geqslant0$ then $\frac{\partial^2 C}{\partial {u_1}^2}\leqslant 0$ and $C$ is D-concave. If $\theta\leqslant0$ then $\frac{\partial^2 C}{\partial {u_1}^2}\geqslant 0$ and $C$ is D-convex.
\item[Gumbel copula] This copula is twice differentiable on $]0,1]^2$, and its second order partial derivative is
$$\frac{\partial^2 C}{\partial {u_1}^2}=-\theta\frac{u_2}{u_1}\ln(u_2)(1-\theta\ln(u_2))e^{-\theta\ln(u_1)\ln(u_2)}$$
It holds that for all $\theta\in(0,1]$, $\frac{\partial^2 C}{\partial {u_1}^2}\geqslant0$. By continuity, $C$ is always D-convex.
\end{description}

As there is no explicit formula for the family of multivariate Gaussian copulas, it is difficult to prove its D-concavity or D-convexity. However, numerical approximations in the case $n=2$ seem to indicate that a Gaussian copula would be D-convex if its marginals are positively correlated, and D-concave if they are negatively correlated. \Cref{fig:gaussian_copula_simu} present those observations, with solid lines representing positive correlation, and dashed lines representing negative correlations. In the case $n>2$, the copula can be neither D-convex nor D-concave, depending on the value of the correlation matrix. An example of this statement is provided in \Cref{fig:gaussian_copula_simu_n3}.

\begin{figure}
    \centering
    \begin{subfigure}{0.4\linewidth}
        \centering
        \includegraphics[width=\linewidth]{Images/Chap_2/Gaussian_copula/gaussian_copula_0.png}
    \end{subfigure}\hfill
    \begin{subfigure}{0.4\linewidth}
        \centering
        \includegraphics[width=\linewidth]{Images/Chap_2/Gaussian_copula/gaussian_copula_1.png}
    \end{subfigure}\hfill
    \begin{subfigure}{0.4\linewidth}
        \centering
        \includegraphics[width=\linewidth]{Images/Chap_2/Gaussian_copula/gaussian_copula_2.png}
    \end{subfigure}
    \caption{Gaussian 2-copulas in the direction $u_1$ for different $u_2$. Each figure present different plots for correlations $r$ between $u_1$ and $u_2$ ranging in $[-1,1]$. Solid lines represent negative correlation, while dashed lines represent positive correlations.}
    \label{fig:gaussian_copula_simu}
\end{figure}

The rest of this section will present different results regarding D-convex and D-concave copulas, leading to the main result of this section presented in proposition \ref{prop:convex_diff_hvol}.

\begin{figure}
    \centering
    \includegraphics[width=0.5\linewidth]{Images/Chap_2/Gaussian_copula/gaussian_copula_n3.png}
    \caption{Directional cut of a Gaussian 3-copula with in the direction $u_1$, with $R=\begin{bmatrix} 1 & -0.4 & 0.7\\ -0.4 & 1 & 0.3\\ 0.7 & 0.3 & 1 \end{bmatrix}$, $u_2=0.4$ and $u_3=0.6$. The copula is neither D-convex nor D-concave}
    \label{fig:gaussian_copula_simu_n3}
\end{figure}

\begin{remark}
    All D-convex copulas $C$ are dominated by the product copula. Similarly, all D-concave copulas dominate the product copula.

    Consider a D-convex copula $C$, and let $(u_1\enum u_n)\in[0,1]^n$.
    \begin{eqnarray*}
        C(u_1\enum u_n) &=& C((1-u_1)\cdot0+u_1\cdot 1, u_2\enum u_n)\\
        &\leqslant& (1-u_1)C(0, u_2\enum u_n)+u_1C(1, u_2\enum u_n)\\
        &\leqslant& u_1C(1, u_2\enum u_n)
    \end{eqnarray*}
    Repeating this for $u_2\enum u_n$ yields:
    \begin{eqnarray*}
        C(u_1\enum u_n) &\leqslant& u_1\dots u_nC(1\enum1) = u_1\dots u_n = C_\Pi(u_1\enum u_n)
    \end{eqnarray*}
    The proof for D-convexity is similar.
\end{remark}

\begin{proposition}\label{prop:sup_additivity}
    If $C$ is a D-convex copula, then it verifies for all $(u_1\enum u_n)\in[0,1]^n$, $(v_1\enum v_n)\in[0,1]^n$ \st $\forall i\in\opi1,n\cli$, $u_i+v_i\leqslant 1$:
    \begin{eqnarray}
        C(u_1\enum u_i+v_i\enum u_n)\geqslant& &C(u_1\enum u_i\enum u_n)\nonumber\\
        &+ &C(u_1\enum v_i\enum u_n)\label{eq:convex_sum}
    \end{eqnarray}
    Similarly, if $\forall i\in\opi1,n\cli$, $u_i-v_i\geqslant 0$, it verifies:
    \begin{eqnarray}
        C(u_1\enum u_i-v_i\enum u_n)\leqslant& &C(u_1\enum u_i\enum u_n)\nonumber\\
        &- &C(u_1\enum v_i\enum u_n)\label{eq:convex_diff}
    \end{eqnarray}
    The inequalities are reversed for D-concave copulas.
\end{proposition}

\begin{proof}
    Let $(u_1\enum u_n)\in[0,1]^n$, $(v_1\enum v_n)\in[0,1]^n$ \st $\forall i\in\opi1,n\cli$, $u_i+v_i\leqslant 1$. Let $i\in\opi1,n\cli$. Applying the definition of convexity \eqref{eq:convex_copula} with $v_i=0$ yields:
    \begin{eqnarray*}
        \forall t\in[0,1],~C(u_1\enum tu_i\enum u_n) \leqslant tC(u_1\enum u_n)
    \end{eqnarray*}
    
    Let $w_i=u_i+v_i\in ]0,1]$ (the case where $u_i=0$ or $v_i=0$ is trivial). It is possible to write $u_i=tw_i$ and $v_i=(1-t)w_i$, with $t=\frac{u_i}{w_i}\in[0,1]$. Then it holds that:
    \begin{eqnarray*}
            C(u_1\enum u_i\enum u_n) =& C(u_1\enum t\cdot w_i\enum u_n)\\
            \leqslant& t\cdot C(u_1\enum w_i\enum u_n)\\
            C(u_1\enum v_i\enum u_n) =& C(u_1\enum (1-t)\cdot w_i\enum u_n)\\
            \leqslant& (1-t)\cdot C(u_1\enum w_i\enum u_n)
    \end{eqnarray*}
    Summing the above equations yields:
    \begin{eqnarray*}
        C(u_1\enum u_i\enum u_n) + C(u_1\enum v_i\enum u_n) &\leqslant& C(u_1\enum w_i\enum u_n)\\
        &\leqslant& C(u_1\enum u_i+v_i\enum u_n)
    \end{eqnarray*}
    which proves \eqref{eq:convex_sum}.

    Let $w_i=u_i-v_i\in [0,1]$, clearly $w_i+v_i\leqslant1$. Using \cref{eq:convex_sum}, it holds that:
    \begin{align*}
        &C(u_1\enum w_i\enum u_n) + &&\\
        &C(u_1\enum v_i\enum u_n) &\leqslant~&C(u_1\enum w_i+v_i\enum u_n)\\
        \Leftrightarrow~ &C(u_1\enum w_i\enum u_n) &\leqslant~&C(u_1\enum w_i+v_i\enum u_n)\\
        &&&-C(u_1\enum v_i\enum u_n)\\
        \Leftrightarrow~ &C(u_1\enum u_i-v_i\enum u_n) &\leqslant~&C(u_1\enum u_i\enum u_n)\\
        &&&-C(u_1\enum v_i\enum u_n)
    \end{align*}
    which proves \eqref{eq:convex_diff}.
\end{proof} 

\begin{proposition}\label{prop:convex_diff_hvol}
    If $C$ is a D-convex copula, then it verifies for all $(u_1\enum u_n)\in[0,1]^n$, $(v_1\enum v_n)\in[0,1]^n$ \st $\forall i\in\opi1,n\cli$, $u_i-v_i\geqslant 0$:
    \begin{eqnarray}
        C(u_1-v_1\enum u_i-v_i\enum u_n-v_n)\leqslant& &H^{u_1\enum u_i\enum u_n}_{v_1\enum v_i\enum v_n}\label{eq:convex_diff_hvol}
    \end{eqnarray}
    where $H$ is the H-volume of $C$. The inequality is reversed for D-concave copulas.
\end{proposition}

\begin{proof}
    The result is straightforward by induction using \cref{eq:convex_diff}.
\end{proof}

\section{Joining P-boxes Using D-Convex/Concave Copulas}\label{sec:annex_dconvex_pbox}
This section contains the proof of \Cref{prop:convexity_pbox} from \Cref{chap:joining_credal_sets}. We remind here the proposition:
\begin{proposition}
    When joining marginals represented by p-boxes using the natural ordering from \eqref{eq:order_pbox} with a copula $C$, it holds that:
    \begin{itemize}
        \item if $C$ is D-convex, then $\M_{mass}\subseteq \M_{agg}$.
        \item if $C$ is D-concave, then $\M_{agg}\subseteq\M_{mass}$.
    \end{itemize}
\end{proposition}

\begin{proof}
    Consider $n$ p-boxes $[\underline{F}_1,~\overline{F}_1] \enum [\underline{F}_n,~\overline{F}_n]$, a D-convex copula $C$ and the natural order on focal sets $(a^i_k)_{1\leqslant k \leqslant N_i}$ of each marginal p-box $[\underline{F}_i,~\overline{F}_i]$. We will denote $m_C$ the joint mass functions obtained using \Cref{eq:joint_mass} and $\Bel_C$ its associated belief function. We will also refer to $\underline{P}$ as the lower probability associated with $\M_{agg}$ from \Cref{eq:copula_on_lower_proba}.

    When considering the natural order on focal sets of a p-box \eqref{eq:order_pbox}, it holds that for every focal set $a^i_p$, the set $\{k~|~a^i_k\subseteq a^i_p\}$ is composed of consecutive integers. In the following, we denote by $\underline{p}$ and $\overline{p}$ the lowest and highest indices of the focal sets included in $a^i_p$. This means that $\{a^i_{\underline{p}} \enum a^i_{\overline{p}}\}$ is the set of all focal sets included in $a^i_p$. 

    Let $a^1_{p_1} \enum a^n_{p_n}$ be focal sets of $m_1 \enum m_n$. We note $u^i_p=\sum_{k=1}^p m_i(a^i_k)$. It then holds that:
    \begin{align*}
        \Bel_C(a^1_{p_1} \enum a^n_{p_n}) &= \sum_{a^1_{k_1}\subseteq a^1_{p_1}}\ldots\sum_{a^n_{k_n}\subseteq a^n_{p_n}}m_C(a^1_{k_1} \enum  a^n_{k_n})\\
        &= \sum_{k_1=\underline{p}_1}^{\overline{p}_1}\ldots\sum_{k_n=\underline{p}_n}^{\overline{p}_n}m_C(a^1_{k_1} \enum  a^n_{k_n})\\
        &= \sum_{k_1=\underline{p}_1}^{\overline{p}_1}\ldots\sum_{k_n=\underline{p}_n}^{\overline{p}_n}H_{u^1_{k_1-1} \enum u^n_{k_n-1}}^{u^1_{k_1} \enum u^n_{k_n}}
    \end{align*}
    As the H-volume is computed over a partitioning of $[u^1_{\underline{p}_1-1}, u^1_{\overline{p}_1}]\tdt[u^n_{\underline{p}_n-1} , u^n_{\overline{p}_n}]$, it is possible to greatly simplify the sums. The proof is the same as the proof of \eqref{eq:joint_mass} except that the \acrshort{cdf} is computed over $[u^1_{\underline{p}_1-1}, u^1_{\overline{p}_1}]\tdt[u^n_{\underline{p}_n-1} , u^n_{\overline{p}_n}]$ and not $[0,1]^n$. This yields:
    \begin{align*}
        \Bel_C(a^1_{p_1} \enum a^n_{p_n}) = H_{u^1_{\underline{p}_1-1} \enum u^n_{\underline{p}_n-1}}^{u^1_{\overline{p}_1} \enum u^n_{\overline{p}_n}}
    \end{align*}
    On the other hand, it holds that:
    \begin{align*}
        \underline{P}(a^1_{p_1} \enum a^n_{p_n}) &= C(\Bel_1(a^1_{p_1}) \enum \Bel_n(a^n_{p_n}))\\
        &= C(\sum_{k_1=\underline{p}_1}^{\overline{p}_1} m_1(a^1_{k_1}) \enum \sum_{k_n=\underline{p}_n}^{\overline{p}_n} m_1(a^n_{k_n}))\\
        &= C(u^1_{\overline{p}_1} - u^1_{\underline{p}_1-1} \enum u^n_{\overline{p}_n} - u^n_{\underline{p}_n-1})
    \end{align*}
    Using \cref{eq:convex_diff_hvol} yields:
    \begin{align*}
        \Bel_C(a^1_{p_1} \enum a^n_{p_n}) \geqslant \underline{P}(a^1_{p_1} \enum a^n_{p_n})
    \end{align*}
    The inequality is reversed if $C$ is D-concave, which concludes the proof.
\end{proof}

\Cref{fig:copula_convex} illustrates the difference between $\Bel_C$ and $\underline{P}$ in the case $n=2$.

\begin{figure}[!ht]
    \centering
    \begin{tikzpicture}
        \begin{axis}[%
            xlabel=$u_1$,
            ylabel=$u_2$,
            xmin=0, xmax=1,
            ymin=0, ymax=1,
            xtick={0.35, 0.55,  0.9, 1},
            xticklabels={$u^1_{\overline{p}_1} - u^1_{\underline{p}_1-1}~~$, $u^1_{\underline{p}_1-1}$, $u^1_{\overline{p}_1}$, $1$},
            ytick={0, 0.3, 0.5,  0.8, 1},
            yticklabels={$0$, $u^2_{\overline{p}_2} - u^2_{\underline{p}_2-1}$, $u^2_{\underline{p}_2-1}$, $u^2_{\overline{p}_2}$, $1$},
            xticklabel style={rotate=30},
            xtick pos=left,
            ytick pos=left,],
            
            \addplot [domain=0:1,samples=40,draw opacity=0.3,color=gray]({x},{1-x});
            \addplot [domain=0:1,samples=40,draw opacity=0.3,color=gray]({x},{1.2-x}); 
            \addplot [domain=0:1,samples=40,draw opacity=0.3,color=gray]({x},{1.4-x}); 
            \addplot [domain=0:1,samples=40,draw opacity=0.3,color=gray]({x},{1.6-x}); 
            \addplot [domain=0:1,samples=40,draw opacity=0.3,color=gray]({x},{1.8-x});

            \node[rotate=-45, fill=white, rounded corners=2pt, inner sep=1pt] (x) at (0.1, 0.9) {\color{gray}0};
            \node[rotate=-45, fill=white, rounded corners=2pt, inner sep=1pt] (x) at (0.3, 0.9) {\color{gray}0.2};
            \node[rotate=-45, fill=white, rounded corners=2pt, inner sep=1pt] (x) at (0.5, 0.9) {\color{gray}0.4};
            \node[rotate=-45, fill=white, rounded corners=2pt, inner sep=1pt] (x) at (0.7, 0.9) {\color{gray}0.6};
            \node[rotate=-45, fill=white, rounded corners=2pt, inner sep=1pt] (x) at (0.9, 0.9) {\color{gray}0.8};
            
            
            \node (a) at (0.9, 0.8) {};
            \node (b) at (0.55, 0.8) {};
            \node (c) at (0.55, 0.5) {};
            \node (d) at (0.9, 0.5) {};
            \draw[ultra thick, black] (c.center) rectangle (a.center) node[pos=0.5]{$\Bel_C$};
            
            \node (e) at (0, 0.8) {};
            \node (f) at (0, 0.5) {};
            \node (g) at (0.9, 0) {};
            \node (h) at (0.55, 0) {};
            \draw[thick, dashed, black] (e.center) -- (b.center);
            \draw[thick, dashed, black] (f.center) -- (c.center);
            \draw[thick, dashed, black] (g.center) -- (d.center);
            \draw[thick, dashed, black] (h.center) -- (c.center);
            
            \node (i) at (0.35, 0.3) {};
            \node (j) at (0, 0) {};
            \draw[ultra thick, black] (i.center) rectangle (j.center) node[pos=0.5]{$\underline{P}$};
            
            \node (m) at (0.55, 0.45) {};
            \node (n) at (0.9, 0.45) {};
            \node (o) at (0.5, 0.8) {};
            \node (p) at (0.5, 0.5) {};
            \draw [stealth-stealth, thick, black] (m.center) -- (n.center) node[pos=0.5, below]{$u^1_{\overline{p}_1} - u^1_{\underline{p}_1-1}$};
            \draw [stealth-stealth, thick, black] (o.center) -- (p.center) node[pos=0.5, left]{$u^2_{\overline{p}_2} - u^2_{\underline{p}_2-1}$};
        
        \end{axis}
    \end{tikzpicture}
    \caption{Bird's-eye view of the \L ukaciewicz 2-copula $C_L$, where the gray lines are the isolines of the copula. $\Bel_C$ and $\underline{P}$ are represented in the case where the marginals are p-boxes. The thick rectangles represent the bounds on which to compute the H-volume. Numbers $u^i_k$ use the notation of the proof of \Cref{prop:convexity_pbox}.}
    \label{fig:copula_convex}
 
\end{figure}
\section{Consistency of the Median Filtering}\label{sec:median_filtering}
This section demonstrates a result used in \Cref{sec:coherence_disparity_intervals}. We define the median of a set of $n$ sorted values $X=\{x_1\enum x_n\}$ is defined as
    \begin{align}
        &\text{if }n=2l+1,\qquad \median X =& x_{l+1}\\
        &\text{if }n=2l,\qquad \median X =& \frac{x_{l}+x_{l+1}}{2}
    \end{align}
    where $l$ is an integer.

\begin{proposition}\label{prop:median_consistency}
    Let $n\in\mathbb{N}^*$. Let $X=\{x_1\enum x_n\}$ and $Y=\{y_1\enum y_n\}$ be two sets of integers such that for all $i$ in $\opi1,n\cli$, $x_i\leqslant y_i$. Then:
    \begin{align}
        \median X \leqslant \median Y
    \end{align}
\end{proposition}

\begin{proof}
    Let $\sigma_X:\opi1,n\cli\rightarrow\opi1,n\cli$ be a bijection sorting $X$, \ie:
    \begin{align}
        x_{\sigma_X(1)}\leqslant\dots\leqslant x_{\sigma_X(i)}\leqslant\dots\leqslant x_{\sigma_X(n)}
    \end{align}
    We define $\sigma_Y:\opi1,n\cli\rightarrow\opi1,n\cli$ similarly, this time sorting $Y$. Notice that it does not necessarily holds that $x_{\sigma_X(i)}\leqslant y_{\sigma_Y(i)}$, but only that $x_{\sigma_Y(i)}\leqslant y_{\sigma_Y(i)}$ 
    
    Suppose that $n=2l+1$ where $l$ is an integer. Then the median of $X$ equals $x_{\sigma_X(l+1)}$ and the median of $Y$ equals $y_{\sigma_Y(l+1)}$ 
    It holds that for all $i\in\opi1,l+1\cli$:
    \begin{align}
        y_{\sigma_Y(l+1)}\geqslant y_{\sigma_Y(i)}\geqslant x_{\sigma_Y(i)}
    \end{align}
    The median of $Y$ is thus greater than at least $l+1$ elements of $X$. Because the $(l+1)$\ith smallest element of $X$ is the median of $X$, then the median of $Y$ is necessarily greater than the median of $X$.
    
    The case $n=2l$, is somehow similar. With the same arguments, we can say that $y_{\sigma_Y(l+1)}$ is greater than $l+1$ elements of $X$, thus it is greater than its $(l+1)$\ith smallest element $x_{\sigma_X(l+1)}$. Similarly, $y_{\sigma_Y(l)}$ is greater than $l$ elements of $X$, thus it is greater than its $l$\ith smallest element $x_{\sigma_X(l)}$. Therefore, it holds that:
    \begin{align}
        \frac{y_{\sigma_Y(l)}+y_{\sigma_Y(l+1)}}{2}\geqslant\frac{x_{\sigma_X(l)}+x_{\sigma_X(l+1)}}{2}
    \end{align}
    Which also means that the median of $Y$ is necessarily greater than the median of $X$.
\end{proof}

\section{Ablation Studies for Disparity Confidence Intervals}\label{sec:ablation_study}
This section will present ablation studies regarding the different parameters used to create disparity confidence intervals in \Cref{chap:epistemic_uncertainty}. When not specified, the values of the different parameters are fixed to the same values used in our experiments, \ie:\newline\medskip
\begin{minipage}{0.5\linewidth}
\begin{align*}
    &\alpha = 0.9\\
    &\tau_{amb} = 0.6\\
    &n_\N = 2
\end{align*}
\end{minipage}
\begin{minipage}{0.5\linewidth}
\begin{align*}
    &q = 0.9\\
    &k_{amb} = 2\\
\end{align*}
\end{minipage}\bigskip\newline
Those values were determined by evaluating different metrics on specific scenes of the Middlebury dataset. We present some results on the Middlebury Cones stereo images. In our experiments, we used the same values of parameters for every scene and for the two considered cost functions: CENSUS and MC-CNN. More in-depth analyses could lead to sets of parameter values specialized for each cost function. This is left as future work, considering that only the CENSUS cost function is used in \Cref{chap:elevation_intervals}.

\subsection{Possibility Threshold}
In \Cref{eq:set_of_possible_disparities}, we considered a parameter $\alpha$ used as a possibility threshold to compute a set of most possible disparities. We recall here the equations used in this section to compute confidence intervals:
\begin{align*}
    D_\alpha&=\{~ d ~|~ \pi_{\rowcol}(d)\geqslant\alpha\}\\
    I_\alpha &= [\underline{I}_\alpha,~\overline{I}_\alpha]=[\min D_\alpha, ~\max D_\alpha]
\end{align*}

Higher values of $\alpha$ will lead to smaller disparity sets, and thus to smaller intervals. In \Cref{chap:epistemic_uncertainty,chap:elevation_intervals}, we chose to use $\alpha=0.9$, but other values could be considered. \Cref{fig:ablation_study_alpha_1,fig:ablation_study_alpha_2} presents the evolution of metrics for different values of $\alpha$. As expected, the accuracy decreases with higher values of $\alpha$, as displayed in \Cref{fig:ablation_study_alpha_acc}.

\Cref{fig:ablation_study_alpha_eps} displays the influence of $\alpha$ over the residual error $\varepsilon$. Variations of $\varepsilon$ have a magnitude of around $2\%$, which indicates that $\alpha$ has very little influence over this metric. $\varepsilon$ is not monotone with regard to $\alpha$. This can be explained as follows: as $\alpha$ increases, the size of intervals decreases. There are thus more intervals that do not contain the ground truth, effectively modifying the set for which $\varepsilon$ is the median. There is therefore no guarantee that $\epsilon$ is monotone with regard to $\alpha$. 

\Cref{fig:ablation_study_alpha_s_rel} displays the influence of $\alpha$ over the relative size $s_{rel}$ of intervals in high confidence areas. $\alpha$ does not have an influence over $s_{rel}$ for the CENSUS cost function, and small influence for the MC-CNN cost function (fluctuations of $3\%$ only). 

In \Cref{fig:ablation_study_alpha_o_rel}, we can observe that the relative over-estimation also decreases with higher values of $\alpha$. For the CENSUS cost function, the relative over-estimation $o_{rel}$ decreases less rapidly for values $\alpha\geqslant0.9$. $\alpha=0.9$ ensures that the relative over-estimation is less than $50\%$ for both cost functions, which we consider a satisfying size for intervals in low confidence areas. 
\begin{figure}
    \centering
    \begin{subfigure}[t]{0.7\linewidth}
        \centering
        \includegraphics[width=\linewidth]{Images/X_Annex/ablation_study_cones_alpha_acc.png}
        \caption{Accuracy $acc$ for different values of the parameter $\alpha$ for constructing intervals.}
        \label{fig:ablation_study_alpha_acc}
    \end{subfigure}\vspace*{0.3cm}\\
    \begin{subfigure}[t]{0.7\linewidth}
        \centering
        \includegraphics[width=\linewidth]{Images/X_Annex/ablation_study_cones_alpha_eps.png}
        \caption{Residual error $\varepsilon$ for different values of the parameter $\alpha$ for constructing intervals.}
        \label{fig:ablation_study_alpha_eps}
    \end{subfigure}
    \caption{Influence of the parameter $\alpha$ used for constructing intervals in \Cref{eq:set_of_possible_disparities} for Middlebury Cones stereo images. Intervals were regularized in low confidence areas.}
    \label{fig:ablation_study_alpha_1}
\end{figure}
\begin{figure}
    \centering
    \begin{subfigure}[t]{0.7\linewidth}
        \centering
        \includegraphics[width=\linewidth]{Images/X_Annex/ablation_study_cones_alpha_s_rel.png}
        \caption{Relative size $s_{rel}$ for different values of the parameter $\alpha$ for constructing intervals.}
        \label{fig:ablation_study_alpha_s_rel}
    \end{subfigure}\vspace*{0.3cm}\\
    \begin{subfigure}[t]{0.7\linewidth}
        \centering
        \includegraphics[width=\linewidth]{Images/X_Annex/ablation_study_cones_alpha_o_rel.png}
        \caption{Relative over-estimation $o_rel$ for different values of the parameter $\alpha$ for constructing intervals.}
        \label{fig:ablation_study_alpha_o_rel}
    \end{subfigure}
    \caption{Influence of the parameter $\alpha$ used for constructing intervals in \Cref{eq:set_of_possible_disparities} for Middlebury Cones stereo images. Intervals were regularized in low confidence areas.}
    \label{fig:ablation_study_alpha_2}
\end{figure}

\subsection{Low Confidence Areas}
In \Cref{chap:epistemic_uncertainty}, we noticed that intervals that did not contain the ground truth were usually located in low confidence areas. We thus proposed to detect low confidence areas in order to process the intervals differently there. We recall \Cref{eq:low_confidence_area,eq:n_N} used to define low confidence areas. A pixel $(\rowcol)$ is considered to be in a low confidence area if it verifies:
\begin{align*}
    \min c_{amb}(\rowcol+k)=\min_{-k_{amb}~\leqslant~ k~\leqslant ~k_{amb}}c_{amb}(\rowcol+k)\leqslant \tau_{amb}
\end{align*}
This equation translates the fact that we use a minitive kernel to filter the confidence from ambiguity map, of size $(1, 2\times k_{amb}+1)$. High values of $k_{amb}$ lead to lower values of the confidence from ambiguity, and consequently to more low confidence pixels. Similarly, higher values of the threshold $\tau_{amb}$ lead to more low confidence pixels. In \Cref{sec:regularization_of_intervals}, we also considered low confidence neighboring $\N(\rowcol)$ of a low confidence pixel, defined as:
\begin{align*}
    \N(row,~col) = \{ p\in &\bigcup_{-n_\N\leqslant k\leqslant n_\N}S(row+k, ~col_k)\\
    &\st~S(row + (k+1), ~col_{k+1})\text{ is adjacent to }S(row+k, ~col_{k})\}
\end{align*}
where $S$ is defined as:
\begin{align*}
    S(\rowcol) = \{~ (row,~col') \text{ \st }~\forall c\in\opi col, col'\cli,~\min c_{amb}(row, ~c+k)\leqslant \tau_{amb}~\}
\end{align*}
Given a low confidence pixel $(\rowcol)$, the parameter $n_\N$ defines the number of rows above and below a pixel that will be explored when defining its neighboring. High values of $n_\N$ will lead to neighboring with more low confidence pixels, but does not influence the total amount of low confidence pixels present in an image.

In \Cref{chap:epistemic_uncertainty,chap:elevation_intervals}, we chose to use $\tau_{amb}=0.6$, $k_{amb}=2$ and $n_\N=2$. We detail here the influence of those parameters.

We want low confidence areas to include as many wrong intervals as possible, but they should not be covering the entire image either. We also want to consider true intervals in the low confidence areas so that our quantile regularization is based on a sufficient pool of correct intervals. For this reason, computing the F1-score is not relevant. We therefore consider two metrics separately: the coverage and the proportion of low confidence pixels. The coverage is the proportion of wrong intervals that are also low confidence pixels:
\begin{align*}
    \mathrm{Coverage} = \dfrac{\#\{(\rowcol)~\st~d_{true}\not\in I_\alpha\text{ and }\min c_{amb}(\rowcol+k)\leqslant \tau_{amb} \}}{\#\{\rowcol)~\st~d_{true}\not\in I_\alpha\}}
\end{align*}
The proportion of low confidence pixels is simply the proportion of low confidence pixels in the entire left image.
\Cref{fig:ablation_study_tau_coverage,fig:ablation_study_tau_size} respectively display the evolution of the coverage and of the proportion of low confidence pixels for different values of $\tau_{amb}$. \Cref{fig:ablation_study_k_coverage,fig:ablation_study_k_size} display the same metrics, but for different values of $k_{amb}$. We can see that the value $\tau_{amb}=0.6$ guarantees a coverage superior to $60\%$ while maintaining a proportion of low confidence pixels smaller than $20\%$ in the case of Middlebury Cones, for both considered cost functions. Similarly, $k_{amb}=2$ guarantees a coverage superior to $60\%$ while maintaining a proportion of low confidence pixels smaller than $20\%$ in the case of Middlebury Cones, for both considered cost functions. In general, the coverage and proportion of low confidence pixels are more sensitive to a variation of $0.1$ of $\tau_{amb}$ than a variation of $1$ of the kernel size parameter $k_{amb}$.

\begin{figure}
    \centering
    \begin{subfigure}[t]{0.7\linewidth}
        \centering
        \includegraphics[width=\linewidth]{Images/X_Annex/ablation_study_cones_tau_amb.png}
        \caption{Coverage for different values of the parameter $\tau_{amb}$ for detecting wrong intervals in low confidence areas.}
        \label{fig:ablation_study_tau_coverage}
    \end{subfigure}\vspace*{0.3cm}\\
    \begin{subfigure}[t]{0.7\linewidth}
        \centering
        \includegraphics[width=\linewidth]{Images/X_Annex/ablation_study_cones_tau_amb_2.png}
        \caption{Proportion of pixels detected as low confidence pixels depending on the value $\tau_{amb}$}
        \label{fig:ablation_study_tau_size}
    \end{subfigure}
    \caption{Influence of the parameter $\tau_{amb}$ used for interval regularization in \Cref{eq:low_confidence_area} for Middlebury Cones stereo images.}
    \label{fig:ablation_study_tau}
\end{figure}

\begin{figure}
    \centering
    \begin{subfigure}[t]{0.7\linewidth}
        \centering
        \includegraphics[width=\linewidth]{Images/X_Annex/ablation_study_cones_k_amb.png}
        \caption{Coverage for different values of the parameter $k_{amb}$ for detecting wrong intervals in low confidence areas.}
        \label{fig:ablation_study_k_coverage}
    \end{subfigure}\vspace*{0.3cm}\\
    \begin{subfigure}[t]{0.7\linewidth}
        \centering
        \includegraphics[width=\linewidth]{Images/X_Annex/ablation_study_cones_k_amb_2.png}
        \caption{Proportion of pixels detected as low confidence pixels depending on the value $k_{amb}$}
        \label{fig:ablation_study_k_size}
    \end{subfigure}
    \caption{Influence of the parameter $k_{amb}$ used for interval regularization in \Cref{eq:low_confidence_area} for Middlebury Cones stereo images.}
    \label{fig:ablation_study_k}
\end{figure}

\Cref{fig:ablation_study_n_N} displays the influence of $n_\N$ over the accuracy, and over the number of pixels in low confidence neighboring $n_\N$. \Cref{fig:ablation_study_n_N_acc} shows that $n_\N$ increases the global accuracy, but it is less sensitive to variations of $n_\N$ than it is to variation of $\alpha$ from \Cref{fig:ablation_study_alpha_acc}. \Cref{fig:ablation_study_n_N_hist_1,fig:ablation_study_n_N_hist_2} shows that $n_\N\geqslant2$ strongly reduces the number of low confidence neighboring with less than $50$ pixels. However, high values of $n_\N$ increases the computation complexity for a small gain of accuracy. We therefore choose $n_\N=2$ as a trade-off.

\begin{figure}
    \centering
    \begin{subfigure}[t]{0.7\linewidth}
        \centering
        \includegraphics[width=\linewidth]{Images/X_Annex/ablation_study_cones_n_N_1.png}
        \caption{Accuracy $acc$ for different values of $n_\N$}
        \label{fig:ablation_study_n_N_acc}
    \end{subfigure}\vspace*{0.3cm}\\
    \begin{subfigure}[t]{0.7\linewidth}
        \centering
        \includegraphics[width=\linewidth]{Images/X_Annex/ablation_study_cones_n_N_2.png}
        \caption{Histograms of the number of pixels in low confidence areas $\N(\rowcol)$ for $n_\N\in\{1,2,3\}$.}
        \label{fig:ablation_study_n_N_hist_1}
    \end{subfigure}\vspace*{0.3cm}\\
    \begin{subfigure}[t]{0.7\linewidth}
        \centering
        \includegraphics[width=\linewidth]{Images/X_Annex/ablation_study_cones_n_N_3.png}
        \caption{Histograms of the number of pixels in low confidence areas $\N(\rowcol)$ for $n_\N\in\{4,5,6\}$.} 
        \label{fig:ablation_study_n_N_hist_2}
    \end{subfigure}
    \caption{Influence of the parameter $n_\N$ used for interval regularization in \Cref{eq:n_N} using the CENSUS cost function on Middlebury cones.}
    \label{fig:ablation_study_n_N}
\end{figure}

\subsection{Quantile Regularization}
Once the neighboring $\N$ of a low confidence pixel has been computed, we update the value of its confidence interval based on the distributions of intervals of its neighbors. We recall here \Cref{eq:quantile_reg} used for the regularization of intervals:
\begin{align*}
    I^{reg}_\alpha=[&\mathcal{Q}_{1-q}(\{\underline{I}_\alpha(r,~c)~|~(r,~c)\in\N(\rowcol)\}),\nonumber\\
    &\mathcal{Q}_q(\{\overline{I}_\alpha(r,~c)~|~(r,~c)\in\N(\rowcol)\})]
\end{align*}
where $\mathcal{Q}_q$ refers to the $q$\ith quantile of a set. In our experiments, we used $q=0.9$.

\begin{figure}
    \centering
    \begin{subfigure}[t]{0.7\linewidth}
        \centering
        \includegraphics[width=\linewidth]{Images/X_Annex/ablation_study_cones_q_acc.png}
        \caption{Accuracy $acc$ with regularization for different quantiles $q$}
        \label{fig:ablation_study_q_no_reg}
    \end{subfigure}\vspace*{0.3cm}\\
    \begin{subfigure}[t]{0.7\linewidth}
        \centering
        \includegraphics[width=\linewidth]{Images/X_Annex/ablation_study_cones_q_o_rel.png}
        \caption{Relative over-estimation $o_{rel}$ with regularization for different quantiles $q$}
        \label{fig:ablation_study_q_reg}
    \end{subfigure}
    \caption{Influence of the parameter $q$ used for interval regularization in \Cref{eq:quantile_reg} for Middlebury Cones stereo images.}
    \label{fig:ablation_study_q}
\end{figure}

\Cref{fig:ablation_study_q} displays the influence of the parameter $q$ over the accuracy and relative over-estimation for the Middlebury Cones stereo images. As expected, both the accuracy and relative over-estimation increase with the value of $q$.

\begin{figure}
    \begin{subfigure}[t]{0.5\linewidth}
        \flushleft
        \includegraphics[width=\linewidth]{Images/X_Annex/ablation_study_acc_reg_no_reg_CENSUS.png}
        \caption{CENSUS cost function}
        \label{fig:acc_no_reg_CENSUS}
    \end{subfigure}
    \begin{subfigure}[t]{0.5\linewidth}
        \flushright
        \includegraphics[width=\linewidth]{Images/X_Annex/ablation_study_acc_reg_no_reg_MCCNN.png}
        \caption{MC-CNN cost function}
        \label{fig:acc_no_reg_MCCNN}
    \end{subfigure}
    \caption{Accuracy with and without regularization of intervals in low confidence areas for different scenes of the Middlebury dataset.}
    \label{fig:acc_reg_no_reg}
\end{figure}

\Cref{fig:acc_reg_no_reg} displays the accuracy with and without regularization for scenes of the Middlebury dataset. Without the regularization, many scenes do not reach an accuracy of $90\%$. This is especially true for the CENSUS cost function. This justifies the use of the regularization step from \Cref{sec:regularization_of_intervals}.

\clearpage
    
    \pagestyle{plain}
    \pagenumbering{roman}
    \setcounter{page}{1}  % Resetting the page number
    \singlespace
    \bibliography{thesis}

\end{document}
